% -*- coding: utf-8; mode: latex -*-
% 自動生成: html2tex.py  (52410_51061.html → LuaLaTeX)
% ─────────────────────────────────────────────
\documentclass[
  paper=a5,
  fontsize=10pt,
  jafontsize=9.5pt,
  tate,                  % 縦組み (横組みなら削除)
  open_bracket_pos=nibu, % 括弧字取り
]{jlreq}

% ── フォント ──────────────────────────────────
\usepackage{fontspec}
\usepackage{luatexja-fontspec}
\usepackage{luatexja-ruby}

% 本文日本語フォント (Yu Mincho / IPA系)
\setmainjfont{YuMincho}[
  BoldFont         = {YuMincho Demibold},
  ItalicFont       = {YuMincho},
]
\setmonojfont{YuGothic}

% ── 外字フォールバック ───────────────────────
\newcommand{\AozoraTooltip}[2]{%
  #1%
}
\newcommand{\AozoraGaijiUnicode}[1]{{\footnotesize #1}}
\newcommand{\AozoraGaijiUnknown}[1]{{\AozoraTooltip{\textbf{[GAIJI]}}{#1}}}

% ── jigmo 外字フォント定義 ────────────────────
% fontspec でフォントファミリを宣言
% LuaLaTeX は kpsewhich 経由で otf ファイルを発見する
\newfontface\JigmoFont  {jigmo}  [BoldFont={jigmo},ItalicFont={jigmo}]
\newfontface\JigmoTwoFont{jigmo2} [BoldFont={jigmo2},ItalicFont={jigmo2}]
\newfontface\JigmoThreeFont{jigmo3}[BoldFont={jigmo3},ItalicFont={jigmo3}]

% ── ページスタイル ────────────────────────────
\usepackage{geometry}   % jlreq が制御するが、追加調整用
\usepackage{microtype}  % LuaLaTeX でのマイクロタイポグラフィ

% ── タイトル情報 ──────────────────────────────
\title{三国志　桃園の巻}
\author{吉川英治}
\date{}

% ─────────────────────────────────────────────
\begin{document}
\maketitle

% ── 本文 ──────────────────────────────────────
\\

\subsection{\ruby{黄巾賊}{こうきんぞく}}

\\

\subsubsection{一}

\\


　\ruby{後漢}{ごかん}の\ruby{建寧}{けんねい}元年のころ。\\


　今から約千七百八十年ほど前のことである。\\


　一人の旅人があった。\\


　腰に、一剣を\ruby{佩}{は}いているほか、身なりはいたって見すぼらしいが、\ruby{眉}{まゆ}は\ruby{秀}{ひい}で、\ruby{唇}{くち}は\ruby{紅}{あか}く、とりわけ\ruby{聡明}{そうめい}そうな\ruby{眸}{ひとみ}や、\ruby{豊}{ゆた}かな頬をしていて、つねにどこかに微笑をふくみ、総じて\ruby{賤}{いや}しげな\ruby{容子}{ようす}がなかった。\\


　年の頃は二十四、五。\\


　草むらの中に、ぽつねんと坐って、膝をかかえこんでいた。\\


　\ruby{悠久}{ゆうきゅう}と水は行く――\\


　微風は\ruby{爽}{さわ}やかに\ruby{鬢}{びん}をなでる。\\


　涼秋の八月だ。\\


　そしてそこは、黄河の\ruby{畔}{ほとり}の――黄土層の低い\ruby{断}{き}り\ruby{岸}{ぎし}であった。\\


「おーい」\\


　誰か河でよんだ。\\


「――そこの若い者ウ。なにを見ているんだい。いくら待っていても、そこは渡し舟の着く所じゃないぞ」\\


　小さな漁船から\ruby{漁夫}{りょうし}がいうのだった。\\


　青年は\ruby{笑}{え}くぼを送って、\\


「ありがとう」と、少し頭を下げた。\\


　漁船は、下流へ流れ去った。けれど青年は、同じ所に、同じ姿をしていた。膝をかかえて坐ったまま遠心的な眼をうごかさなかった。\\


「おい、おい、旅の者」\\


　こんどは、後ろを通った人間が呼びかけた。近村の百姓であろう。ひとりは鶏の足をつかんでさげ、ひとりは農具をかついでいた。\\


「――そんな所で、今朝からなにを待っているんだね。このごろは、\ruby{黄巾賊}{こうきんぞく}とかいう悪徒が立ち廻るからな。役人衆に\ruby{怪}{あや}しまれるぞよ」\\


　青年は、振りかえって、\\


「はい、どうも」\\


　おとなしい\ruby{会釈}{えしゃく}をかえした。\\


　けれどなお、腰を上げようとはしなかった。\\


　そして、幾千万年も、こうして流れているのかと思われる黄河の水を、\ruby{飽}{あ}かずに眺めていた。\\


（――どうしてこの河の水は、こんなに黄色いのか？）\\


　\ruby{汀}{みぎわ}の水を、仔細に見ると、それは水その物が黄色いのではなく、\ruby{砥石}{といし}を粉にくだいたような黄色い\ruby{沙}{すな}の\ruby{微粒}{びりゅう}が、水に\ruby{混}{ま}じっていちめんにおどっているため、\ruby{濁}{にご}って見えるのであった。\\


「ああ……、この土も」\\


　青年は、大地の土を、一つかみ\ruby{掌}{て}に\ruby{掬}{すく}った。そして眼を――はるか西北の空へじっと放った。\\


　支那の大地を作ったのも、黄河の水を黄色くしたのも、みなこの沙の微粒である。そしてこの沙は中央\ruby{亜細亜}{アジア}の沙漠から吹いてきた物である。まだ人類の生活も始まらなかった何万年も前の大昔から――不断に吹き送られて、積り積った大地である。この広い\ruby{黄土}{こうど}と黄河の流れであった。\\


「わたしのご先祖も、この河を\ruby{下}{くだ}って……」\\


　彼は、自分の体に今、脈うっている血液がどこからきたか、その遠い根元までを想像していた。\\


　支那を\ruby{拓}{ひら}いた漢民族も、その沙の来る亜細亜の山岳を越えてきた。そして黄河の流れに添いつつ次第にふえ、\ruby{苗族}{びょうぞく}という未開人を追って、農業を\ruby{拓}{ひら}き、産業を\ruby{興}{おこ}し、ここに何千年の文化を植えてきたものだった。\\


「ご先祖さま、みていて下さいまし。いやこの\ruby{劉備}{りゅうび}を、\ruby{鞭}{むち}打って下さい。劉備はきっと、漢の民を興します。漢民族の血と平和を守ります」\\


　天へ向って誓うように、劉備青年は、空を拝していた。\\


　するとすぐ後ろへ、誰か突っ立って、彼の頭からどなった。\\


「\jruby{うさん}{・・・・・・}な\ruby{奴}{やつ}だ。やいっ、汝は、\ruby{黄巾賊}{こうきんぞく}の仲間だろう？」\\

\\

\subsubsection{二}

\\


　劉備は、おどろいて、何者かと振りかえった。\\


　\ruby{咎}{とが}めた者は、\\


「どこから来たっ」と、彼の襟がみをもう\ruby{用捨}{ようしゃ}なくつかんでいた。\\


「……？」\\


　見ると、役人であろう、胸に県の\ruby{吏章}{りしょう}をつけている。近頃は\ruby{物騒}{ぶっそう}な世の中なので、地方の小役人までが、平常でもみな武装していた。二人のうち一名は\ruby{鉄弓}{てっきゅう}を持ち、一名は\ruby{半月槍}{はんげつそう}をかかえていた。\\


「\ruby{{\fontspec{jigmo}\char"6DBF}県}{たくけん}の者です」\\


　劉備青年が答えると、\\


「{\fontspec{jigmo}\char"6DBF}県はどこか」と、たたみかけていう。\\


「はい、{\fontspec{jigmo}\char"6DBF}県の\ruby{楼桑村}{ろうそうそん}（現在・京広線の北京―保定間）の生れで、今でも母と共に、楼桑村に住んでおります」\\


「商売は」\\


「\ruby{蓆}{むしろ}を\ruby{織}{お}ったり\ruby{簾}{すだれ}をつくって、売っておりますが」\\


「なんだ、\ruby{行商人}{ぎょうしょうにん}か」\\


「そんなものです」\\


「だが……」\\


　と、役人は急にむさい物からのくように襟がみを放して、劉備の腰の剣をのぞきこんだ。\\


「この剣には、黄金の\ruby{佩環}{はいかん}に、\ruby{琅{\fontspec{jigmo}\char"7395}}{ろうかん}の\ruby{緒珠}{おだま}がさがっているのではないか、\ruby{蓆売}{むしろう}りには過ぎた刀だ。どこで盗んだ？」\\


「これだけは、父の\ruby{遺物}{かたみ}で持っているのです。盗んだ物などではありません」\\


　素直ではあるが、\ruby{凛}{りん}とした答えである。役人は、劉備青年の眼を見ると、急に眼をそらして、\\


「しかしだな、こんなところに、半日も坐りこんで、いったい何を見ておるのか。怪しまれても仕方があるまい。――折も折、ゆうべもこの近村へ、黄巾賊の群れが\ruby{襲}{よ}せて、\ruby{掠奪}{りゃくだつ}を働いて逃げた所だ。――見るところ大人しそうだし、賊徒とは思われぬが、一応疑ってみねばならん」\\


「ごもっともです。……実は私が待っているのは、今日あたり\ruby{江}{こう}を下ってくると聞いている\ruby{洛陽船}{らくようぶね}でございます」\\


「ははあ、誰か身寄りの者でもそれへ便乗して来るのか」\\


「いいえ、茶を求めたいと思って。――待っているのです」\\


「茶を」\\


　役人は眼をみはった。\\


　彼らはまだ茶の味を知らなかった。茶という物は、\ruby{瀕死}{ひんし}の病人に与えるか、よほどな貴人でなければのまないからだった。それほど高価でもあり貴重に思われていた。\\


「誰にのませるのだ。重病人でもかかえているのか」\\


「病人ではございませんが、生来、私の母の大好物は茶でございます。貧乏なので、めったに買ってやることもできませんが、一両年\ruby{稼}{かせ}いでためた\ruby{小費}{こづかい}もあるので、こんどの旅の土産には、買って戻ろうと考えたものですから」\\


「ふーむ。……それは感心なものだな。おれにも息子があるが、親に茶をのませてくれるどころか――あの通りだわえ」\\


　二人の役人は、顔を見合せてそういうと、もう劉備の疑いも解けた\ruby{容子}{ようす}で、何か語らいながら立ち去ってしまった。\\


　\ruby{陽}{ひ}は西に傾きかけた。\\


　\ruby{茜}{あかね}ざした夕空を、赤い黄河の流れに対したまま、劉備はまた、黙想していた。\\


　と、やがて、\\


「おお、船旗が見えた。洛陽船にちがいない」\\


　彼は初めて草むらを起った。そして\ruby{眉}{まゆ}に手をかざしながら、上流のほうを眺めた。\\

\\

\subsubsection{三}

\\


　ゆるやかに、江を下ってくる船の影は、\ruby{舂}{うすず}く\ruby{陽}{ひ}を負って黒く、徐々と眼の前に近づいてきた。ふつうの客船や貨船とちがい、洛陽船はひと目でわかる。無数の紅い\ruby{龍舌旗}{りゅうぜつき}を帆ばしらにひるがえし、\ruby{船楼}{せんろう}は五\ruby{彩}{さい}に塗ってあった。\\


「おうーい」\\


　\ruby{劉備}{りゅうび}は手を振った。\\


　しかし船は一個の彼に見向きもしなかった。\\


　おもむろに\ruby{舵}{かじ}を曲げ、スルスルと帆をおろしながら、黄河の流れにまかせて、そこからずっと下流の岸へ着いた。\\


　百戸ばかりの\ruby{水村}{すいそん}がある。\\


　今日、洛陽船を待っていたのは、劉備ひとりではない。岸にはがやがやと沢山な人影がかたまっていた。\ruby{驢}{ろ}をひいた仲買人の群れだの、\ruby{鶏車}{チイチャー}と呼ぶ手押し車に、土地の糸や綿を積んだ百姓だの、獣の肉や果物を\ruby{籠}{かご}に入れて待つ物売りだの――すでにそこには、洛陽船を迎えて、\ruby{市}{いち}が立とうとしていた。\\


　なにしろ、黄河の上流、洛陽の都には今、\ruby{後漢}{ごかん}の第十二代の帝王、\ruby{霊帝}{れいてい}の\ruby{居城}{きょじょう}があるし、珍しい物産や、文化の\ruby{粋}{すい}は、ほとんどそこでつくられ、そこから全支那へ行きわたるのである。\\


　幾月かに一度ずつ、文明の製品を積んだ洛陽船が、この地方へも\ruby{下江}{かこう}してきた。そして沿岸の小都市、村、部落など、市の立つところに船を寄せて、\ruby{交易}{こうえき}した。\\


　ここでも。\\


　夕方にかけて、おそろしく騒がしくまたあわただしい取引が始まった。\\


　劉備は、そのやかましい人声と人影の中に立ちまじって、まごついていた。彼は、自分の求めようとしている茶が、仲買人の手にはいることを心配していた。一度、商人の手に移ると、\ruby{莫大}{ばくだい}な値になって、とても自分の貧しい\ruby{嚢中}{のうちゅう}では\ruby{購}{あがな}えなくなるからであった。\\


　またたく間に、市の取引は終った。仲買人も百姓も物売りたちも、三々五々、夕闇へ散ってゆく。\\


　劉備は、船の商人らしい男を見かけてあわててそばへ寄って行った。\\


「茶を売って下さい、茶が欲しいんですが」\\


「え、茶だって？」\\


　\ruby{洛陽}{らくよう}の商人は、\ruby{鷹揚}{おうよう}に彼を振向いた。\\


「あいにくと、お前さんに\ruby{頒}{わ}けてやるような安茶は持たないよ。\ruby{一葉}{ひとは}いくらというような佳品しか船にはないよ」\\


「結構です。たくさんは\ruby{要}{い}りませんが」\\


「おまえ茶をのんだことがあるのかね。地方の衆が何か葉を煮てのんでいるが、あれは茶ではないよ」\\


「はい。その、ほんとの茶を\ruby{頒}{わ}けていただきたいのです」\\


　彼の声は、懸命だった。\\


　茶がいかに貴重か、高価か、また地方にもまだない物かは、彼もよくわきまえていた。\\


　その\ruby{種子}{たね}は、遠い熱帯の異国からわずかにもたらされて、\ruby{周}{しゅう}の代にようやく宮廷の秘用にたしなまれ、漢帝の\ruby{代々}{よよ}になっても、\ruby{後宮}{こうきゅう}の茶園に少し\ruby{摘}{つ}まれる物と、民間のごく貴人の所有地にまれに栽培されたくらいなものだとも聞いている。\\


　また別な説には、一日に百\ruby{草}{そう}を\ruby{嘗}{な}めつつ人間に食物を教えた\ruby{神農}{しんのう}はたびたび毒草にあたったが、茶を得てからこれを噛むとたちまち毒をけしたので、以来、秘愛せられたとも伝えられている。\\


　いずれにしろ、劉備の身分でそれを求めることの無謀は、よく知っていた。\\


　――だが、彼の懸命な\ruby{面}{おも}もちと、\ruby{真面目}{まじめ}に、欲するわけを話す態度を見ると、洛陽の商人も、やや心を動かされたとみえて、\\


「では少し頒けてあげてもよいが、お前さん、失礼だが、その代価をお持ちかね？」と訊いた。\\

\\

\subsubsection{四}

\\


「持っております」\\


　彼は、\ruby{懐中}{ふところ}の\ruby{革嚢}{かわぶくろ}を取出し、銀や砂金を取りまぜて、相手の\ruby{両掌}{りょうて}へ、惜しげもなくそれを皆あけた。\\


「ほ……」\\


　洛陽の商人は、\ruby{掌}{て}の上の\ruby{目量}{めかた}を計りながら、\\


「あるねえ。しかし、\ruby{銀}{ぎん}があらかたじゃないか。これでは、よい茶はいくらも上げられないが」\\


「何ほどでも」\\


「そんなに欲しいのかい」\\


「母が眼を細めて、よろこぶ顔が見たいので――」\\


「お前さん、商売は？」\\


「\ruby{蓆}{むしろ}や\ruby{簾}{すだれ}を作っています」\\


「じゃあ、失礼だが、これだけの\ruby{銀}{かね}をためるにはたいへんだろ」\\


「二年かかりました。自分の食べたい物も、着たい物も、節約して」\\


「そう聞くと、断われないな。けれどとても、これだけの銀と替えたんじゃ引合わない。なにかほかにないかね」\\


「これも添えます」\\


　\ruby{劉備}{りゅうび}は、剣の\ruby{緒}{お}にさげている\ruby{琅{\fontspec{jigmo}\char"7395}}{ろうかん}の珠を解いて出した。洛陽の商人は琅{\fontspec{jigmo}\char"7395}などは珍しくない顔つきをして見ていたが、\\


「よろしい。おまえさんの孝心に免じて、茶と交易してやろう」\\


　と、やがて船室の中から、\ruby{錫}{すず}の小さい\ruby{壺}{つぼ}を一つ持ってきて、劉備に与えた。\\


　黄河は暗くなりかけていた。西南方に、\ruby{妖猫}{ようびょう}の眼みたいな大きな星がまたたいていた。その星の光をよく見ていると虹色の\ruby{暈}{かさ}がぼっとさしていた。\\


　――世の中がいよいよ乱れる\ruby{凶兆}{きょうちょう}だ。\\


　と、近頃しきりと、世間の者が\ruby{怖}{こわ}がっている星である。\\


「ありがとうございました」\\


　劉備青年は、錫の小壺を、\ruby{両掌}{りょうて}に持って、やがて岸を離れてゆく船の影を拝んでいた。もう\ruby{瞼}{まぶた}に、母のよろこぶ顔がちらちらする。\\


　しかし、ここから故郷の\ruby{{\fontspec{jigmo}\char"6DBF}県楼桑村}{たくけんろうそうそん}までは、百里の余もあった。幾夜の泊りを重ねなければ帰れないのである。\\


「今夜は寝て――」と、考えた。\\


　\ruby{彼方}{かなた}を見ると、\ruby{水村}{すいそん}の\ruby{灯}{ひ}が二つ三つまたたいている。彼は村の\ruby{木賃}{きちん}へ眠った。\\


　すると夜半頃。\\


　木賃の亭主が、あわただしく起しにきた。眼をさますと、\ruby{戸外}{おもて}は真っ赤だった。むうっと蒸されるような熱さの中にどこかでパチパチと、火の燃える物音もする。\\


「あっ、火事ですか」\\


「\ruby{黄巾賊}{こうきんぞく}がやってきたのですよ旦那、\ruby{洛陽船}{らくようぶね}と交易した仲買人たちが、今夜ここに泊ったのを\ruby{狙}{ねら}って――」\\


「えっ。……賊？」\\


「旦那も、交易した一人でしょう。奴らが、まっ先に狙うのは、今夜泊った仲買たちです。次にはわしらの番だが、はやく裏口からお逃げなさい」\\


　劉備はすぐ剣を\ruby{佩}{は}いた。\\


　裏口へ出てみるともう近所は焼けていた。家畜は、異様なうめきを放ち、女子どもは、焔の下に悲鳴をあげて逃げまどっていた。\\


　昼のように大地は明るい。\\


　見れば、\ruby{夜叉}{やしゃ}のような人影が、\ruby{矛}{ほこ}や\ruby{槍}{やり}や\ruby{鉄杖}{てつじょう}をふるって、逃げ散る旅人や村の者らを見あたり次第にそこここで\ruby{殺戮}{さつりく}していた。――眼をおおうような地獄がえがかれているではないか。\\


　昼ならば眼にも見えよう。それらの悪鬼は皆、\ruby{結髪}{けっぱつ}のうしろに、黄色の\ruby{巾}{きれ}をかりているのだ。黄巾賊の名は、そこから起ったものである。本来は支那の――この国のもっとも尊い色であるはずの黄土の国色も、今は、善良な民の眼をふるえ上がらせる、悪鬼の\ruby{象徴}{しるし}になっていた。\\

\\

\subsubsection{五}

\\


「ああ、\ruby{酸鼻}{さんび}な――」\\


　\ruby{劉備}{りゅうび}は、つぶやいて、\\


「ここへ自分が泊り合せたのは、天が、天に代って、この\ruby{憐}{あわ}れな民を救えとの、\ruby{思}{おぼ}し\ruby{召}{めし}かも知れぬ。……おのれ、鬼畜どもめ」\\


　と、剣に手をかけながら、家の\ruby{扉}{と}を蹴って、躍りだそうとしたが、いや待て――と思い直した。\\


　母がある。――自分には自分を頼みに生きているただ一人の母がある。\\


　黄巾の乱賊はこの地方にだけいるわけではない。\ruby{蝗}{いなご}のように天下いたるところに\ruby{群}{むれ}をなして\ruby{跳梁}{ちょうりょう}しているのだ。\\


　一剣の勇では、百人の賊を斬ることもむずかしい。百人の賊を斬っても、天下は救われはしないのだ。\\


　母を悲しませ、百人の賊の\ruby{生命}{いのち}を自分の一命と取換えたとて何になろう。\\


「そうだ。……わしは今日も黄河の\ruby{畔}{ほとり}で天に誓ったではないか」\\


　劉備は、眼をおおって、裏口からのがれた。\\


　彼は、闇夜を駈けつづけ、ようやく村をはなれた山道までかかった。\\


「もうよかろう」\\


　汗をぬぐって振りかえると、焼きはらわれた水村は、\ruby{曠野}{こうや}の果ての\ruby{焚火}{たきび}よりも小さい火にしか見えなかった。\\


　空を仰いで、\ruby{白虹}{はっこう}のような星雲をかけた宇宙と見くらべると、この世の山岳の大も、黄河の長さも、支那大陸の\ruby{偉}{い}なる広さも、むしろ\ruby{愍}{あわ}れむべき小さい存在でしかない。\\


　まして人間の小ささ――一個の自己のごときは――と劉備は、我というものの無力を\ruby{嘆}{なげ}いたが、\\


「\ruby{否}{いな}！　否！　人間あっての宇宙だ。人間がない宇宙はただの\ruby{空虚}{うつろ}ではないか。人間は宇宙より偉大だ」と、われを忘れて、天へ向ってどなった。すると後ろのほうで、\\


　――\ruby{然}{しか}なり。然なり。\\


　と、誰かいったような気がしたが、振りかえって見たが、人影なども見あたらなかった。\\


　ただ、樹木の蔭に、一宇の古い\ruby{孔子廟}{こうしびょう}があった。\\


　劉備は、近づいて、廟にぬかずきながら、\\


「そうだ、孔子、今から七百年前に、\ruby{魯}{ろ}の国（山東省）に生れて、世の乱れを正し、今に至るまで、こうして人の心に生き、人の魂を救っている。人間の偉大を証拠だてたお方だ。その孔子は文を以て、世に立ったが、わしは武を以て、民を救おう――。今のように\ruby{黄魔鬼畜}{こうまきちく}の\ruby{跳梁}{ちょうりょう}にまかせている暗黒な世には、文を布く前に、武を以て、地上に平和をたてるしかない」\\


　多感な劉備青年は、あたりに人がいないとのみ思っていたので、孔子廟へ向って、誓いを立てるように、思わず情熱的な声を放って云った。\\


　――と、廟の中で、\\


「わはははは」\\


「あははは」\\


　大声で笑った者がある。\\


　びっくりして、劉備がたちかけると、廟の\ruby{扉}{と}を\ruby{蹴}{け}って、突然、\ruby{豹}{ひょう}のように躍りだしてきた男があって、\\


「こら、待て」\\


　劉備の\ruby{襟首}{えりくび}を抑えた。\\


　同時に、もう一人の大男は、廟の内から劉備の眼の前へと、孔子の木像を蹴とばして、\\


「ばか野郎、こんな物が貴様ありがたいのか。どこが偉大だ」と、\ruby{罵}{ののし}った。\\


　孔子の木像は首が折れて、わかれわかれに転がった。\\

\\

\subsubsection{六}

\\


　\ruby{劉備}{りゅうび}は怖れた。これは悪い者に出合ったと思った。\\


　二人の\ruby{巨男}{おおおとこ}を見るに、結髪を黄色の布で包んでいるし、胴には鉄甲を\ruby{鎧}{よろ}い、脚には獣皮の靴をはき、腰には大剣を横たえている。\\


　問うまでもなく、黄巾賊の仲間である。しかも、その\ruby{頭分}{かしらぶん}の者であることは、\ruby{面構}{つらがま}えや服装でもすぐ分った。\\


「\ruby{大方}{だいほう}。こいつを、どうするんですか」\\


　劉備の襟がみをつかんだのが、もう一人のほうに向って訊くと、孔子の木像を蹴とばした男は、\\


「離してもいい。逃げればすぐ叩っ斬ってしまうまでのことだ。おれが睨んでいる前からなんで逃げられるものか」と、いった。\\


　そして廟の前の\ruby{玉石}{たまいし}に腰を悠然とおろした。\\


　\ruby{大方}{だいほう}、\ruby{中方}{ちゅうほう}、\ruby{小方}{しょうほう}などというのは、\ruby{方師}{ほうし}（\ruby{術者}{じゅつしゃ}・祈祷師）の称号で、その位階をも現わしていた。黄巾賊の仲間では、部将をさして、みなそう呼ぶのであった。\\


　けれど、総大将の\ruby{張角}{ちょうかく}のことは、そうよばない。張角と、その二人の弟に向ってだけは、特に、\\


　\ruby{大賢良師}{だいけんりょうし}、張角\\


　\ruby{天公将軍}{てんこうしょうぐん}、\ruby{張梁}{ちょうりょう}\\


　\ruby{地公将軍}{ちこうしょうぐん}、\ruby{張宝}{ちょうほう}\\


　というように尊称していた。\\


　その下に、大方、中方などとよぶ部将をもって組織しているのであった――で今、劉備の前に腰かけている男は、張角の配下の\ruby{馬元義}{ばげんぎ}という黄巾賊の一頭目であった。\\


「おい、\ruby{甘洪}{かんこう}」と、馬元義は手下の甘洪が、まだ危ぶんでいる様子に、\ruby{顎}{あご}で大きくいった。\\


「そいつを、もっと前へ引きずってこい――そうだ俺の前へ」\\


　劉備は、襟がみを持たれたまま、馬元義の足もとへ引き据えられた。\\


「やい、百姓」\\


　\ruby{馬}{ば}はねめつけて、\\


「\ruby{汝}{われ}は今、孔子廟へ向って、大それた\ruby{誓願}{せいがん}を立てていたが、一体うぬは、正気か狂人か」\\


「はい」\\


「はいではすまねえ。\ruby{黄魔鬼畜}{こうまきちく}を討ってどうとかぬかしていたが、黄魔とは、誰のことだ、鬼畜とは、何をさしていったのだ」\\


「べつに意味はありません」\\


「意味のないことを独りでいう\jruby{たわけ}{・・・・・・}があるか」\\


「あまり山道が淋しいので、怖ろしさをまぎらすために出たらめに、声を放って歩いてきたものですから」\\


「相違ないか」\\


「はい」\\


「――で、何処まで行くのだ。この真夜中に」\\


「\ruby{{\fontspec{jigmo}\char"6DBF}県}{たくけん}まで帰ります」\\


「じゃあまだ道は遠いな。俺たちも夜が明けたら、北のほうの町まで行くが、てめえのために眼をさましてしまった。もう二度寝もできまい。ちょうど荷物があって困っていた所だから、俺の荷をかついで、供をしてこい――おい、\ruby{甘洪}{かんこう}」\\


「へい」\\


「荷物はこいつにかつがせて、\ruby{汝}{われ}は俺の半月槍を持て」\\


「もう出かけるんですか」\\


「峠を降りると夜が明けるだろう。その間に奴らも、今夜の仕事をすまして、後から追いついてくるにちげえねえ」\\


「では、歩き歩き、通ったしるしを残して行きましょう」と、甘洪は、\ruby{廟}{びょう}の壁に何か書き残したが、半里も歩くとまた、道ばたの木の枝に、黄色の\ruby{巾}{きれ}を結びつけて行く――\\


　\ruby{大方}{だいほう}の\ruby{馬元義}{ばげんぎ}は、悠々と、\ruby{驢}{ろ}に乗って先へ進んで行くのであった。\\

\\

\subsection{\ruby{流行}{はや}る\ruby{童歌}{どうか}}

\\

\subsubsection{一}

\\


　驢は、北へ向いて歩いた。\\


　鞍上の馬元義は、ときどき南を振り向いて、\\


「奴らはまだ追いついてこないがどうしたのだろう」と、つぶやいた。\\


　彼の半月槍をかついで、驢の後からついてゆく手下の\ruby{甘洪}{かんこう}は、\\


「どこかで道を取っ違えたのかも知れませんぜ。いずれ\ruby{冀州}{きしゅう}（河北省保定の南方）へ行けば落ち合いましょうが」と、いった。\\


　いずれ賊の仲間のことをいっているのであろう――と\ruby{劉備}{りゅうび}は察した。とすれば、自分がのがれてきた黄河の水村を襲ったあの連中を待っているのかも知れない、と思った。\\


（何しろ、従順をよそおっているに\ruby{如}{し}くはない。そのうちには、逃げる機会があるだろう）\\


　劉備は、賊の荷物を負って、黙々と、驢と半月槍のあいだに挟まれながら歩いた。丘陵と河と平原ばかりの道を、四日も歩きつづけた。\\


　幸い雨のない日が続いた。十方\ruby{碧落}{へきらく}、一\ruby{朶}{だ}の雲もない秋だった。\ruby{黍}{きび}のひょろ長い穂に、時折、驢も人の\ruby{背丈}{せたけ}もつつまれる。\\


「ああ――」\\


　旅に\ruby{倦}{う}んで、馬元義は大きなあくびを見せたりした。甘も\ruby{気}{け}だるそうに居眠り半分、足だけを動かしていた。\\


　そんな時、劉備はふと、\\


　――今だっ。\\


　という衝動にかられて、幾度か剣に手をやろうとしたが、もし仕損じたらと、母を想い、身の大望を考えて、じっと辛抱していた。\\


「おう、甘洪」\\


「へえ」\\


「飯が食えるぞ。冷たい水にありつけるぞ――見ろ、むこうに寺があら」\\


「寺が」\\


　黍の間から伸び上がって、\\


「ありがてえ。\ruby{大方}{だいほう}、きっと酒もありますぜ。坊主は酒が好きですからね」\\


　夜は冷え渡るが、昼間は焦げつくばかりな炎熱であった。――水と聞くと、劉備も思わず伸び上がった。\\


　低い丘陵が彼方に見える。\\


　丘陵に抱かれている\ruby{一叢}{ひとむら}の木立と沼があった。沼には紅白の\ruby{蓮花}{はちす}がいっぱい咲いていた。\\


　そこの石橋を渡って、荒れはてた寺門の前で、馬元義は驢をおりた。門の扉は、一枚はこわれ、一枚は形だけ残っていた。それに黄色の紙が貼ってあって、次のような文が書いてあった。\\

\ruby{蒼天已死}{そうてんすでにしす}\\

\ruby{黄夫当レ立}{こうふまさにたつべし}\\

\ruby{歳在二甲子一}{としこうしにありて}\\

\ruby{天下大吉}{てんかだいきち}\\


　　○\\

\ruby{大賢良師張角}{だいけんりょうしちょうかく}\\



「大方ご覧なさい。ここにもわが党の\ruby{盟符}{めいふ}が貼ってありまさ。この寺も黄巾の仲間に入っている奴ですぜ」\\


「誰かいるか」\\


「ところが、いくら呼んでも誰も出てきませんが」\\


「もう一度、どなってみろ」\\


「おうい、誰かいねえのか」\\


　――薄暗い堂の中を、どなりながら覗いてみた。何もない堂の真ん中に、\ruby{曲{\fontspec{jigmo}\char"5F54}}{きょくろく}に腰かけている骨と皮ばかりな老僧がいた。しかし老僧は眠っているのか、死んでいるのか、\ruby{木乃伊}{ミイラ}のように、\ruby{空虚}{うつろ}な眼を\ruby{梁}{うつばり}へ向けたまま、\ruby{寂然}{じゃくねん}と――答えもしない。\\

\\

\subsubsection{二}

\\


「やい、老いぼれ」\\


　\ruby{甘洪}{かんこう}は、半月槍の柄で、老僧の\ruby{脛}{すね}をなぐった。\\


　老僧は、やっとにぶい眼をあいて、眼の前にいる甘と、馬と、\ruby{劉}{りゅう}青年を見まわした。\\


「食物があるだろう。おれたちはここで腹支度をするのだ、はやく支度をしろ」\\


「……ない」\\


　老僧は、\ruby{蝋}{ろう}のような青白い顔を、力なく振った。\\


「ない？　――これだけの寺に食物がないはずはねえ。俺たちをなんだと思う。\ruby{頭髪}{あたま}の\ruby{黄巾}{きれ}を見ろ。大賢良師張角様の\ruby{方将}{ほうしょう}、馬元義というものだ。家探しして、もし食物があったら、素ッ首をはね落すがいいか」\\


「……どうぞ」\\


　老僧は、うなずいた。\\


　馬は甘をかえりみて、\\


「ほんとにないのかもしれねえな。あまり落着いていやがる」\\


　すると老僧は、\ruby{曲{\fontspec{jigmo}\char"5F54}}{きょくろく}にかけていた枯木のような\ruby{肘}{ひじ}を上げて、後ろの祭壇や、壁や四方をいちいちさして、\\


「ない！　ない！　ない！　……仏陀の像さえない！　一物もここにはないっ」と、いった。\\


　泣くがような声である。\\


　そしてにぶい\ruby{眸}{ひとみ}に、怨みの光をこめてまたいった。\\


「みんな、お前方の仲間が持って行ってしまったのだ。\ruby{蝗}{いなご}の群れが通ったあとの田みたいだよここは……」\\


「でも、何かあるだろう。何か喰える物が」\\


「ない」\\


「じゃあ、冷たい水でも汲んでこい」\\


「井戸には、毒が投げこんである。飲めば死ぬ」\\


「誰がそんなことをした」\\


「それも、\ruby{黄巾}{こうきん}をつけたお前方の仲間だ。前の\ruby{地頭}{じとう}と戦った時、残党が隠れぬようにと、みな毒を投げこんで行った」\\


「しからば、泉があるだろう。あんな美麗な\ruby{蓮花}{はちす}が咲いている池があるのだから、どこぞに、冷水が\ruby{湧}{わ}いているにちがいない」\\


「――あの蓮花が、なんで美しかろう。わしの眼には、\ruby{紅蓮}{ぐれん}も\ruby{白蓮}{びゃくれん}も、無数の民の\ruby{幽魂}{ゆうこん}に見えてならない。一花、一花\ruby{呪}{のろ}い、恨み、\ruby{哭}{な}き\ruby{戦}{おのの}きふるえているような」\\


「こいつめが、妙な世まい\ruby{言}{ごと}を……」\\


「嘘と思うなら池をのぞいてみるがよい。紅蓮の下にも、白蓮の根元にも、\ruby{腐爛}{ふらん}した人間の死骸がいっぱいだよ。お前方の仲間に殺された善良な農民や女子供の死骸だの、また、黄巾の党に入らないので、\ruby{縊}{くび}り殺された地頭やら、その\ruby{夫人}{おくさん}やら、戦って死んだ役人衆やら――何百という死骸がのう」\\


「あたり前だ。大賢良師張角様に\ruby{反}{そむ}くやつらは、みな天罰でそうなるのだ」\\


「…………」\\


「いや。よけいなことは、どうでもいい。食べ物もなく水もなく、一体それでは、てめえは何を喰って生きているのか」\\


「わしの喰ってる物なら」と、老僧は、自分の\ruby{沓}{くつ}のまわりを指さした。\\


「……そこらにある」\\


　馬元義は、何気なく、床を見まわした。根を\ruby{噛}{か}んだ\ruby{生草}{なまぐさ}だの、虫の足だの、鼠の骨などが散らかっていた。\\


「こいつは参った。ご\ruby{饗応}{きょうおう}はおあずけとしておこう。おい\ruby{劉}{りゅう}、\ruby{甘洪}{かんこう}、行こうぜ」\\


　と出て行きかけた。\\


　すると、その時はじめて、賊の供をしている劉備の存在に気づいた老僧は、穴のあくほど、劉青年の顔を見つめていたが、突然、\\


「あっ？」と、打たれたような\ruby{愕}{おどろ}きを声に放って、曲{\fontspec{jigmo}\char"5F54}から突っ立った。\\

\\

\subsubsection{三}

\\


　老僧の落ちくぼんでいる眼は大きく\ruby{驚異}{きょうい}にみはったまま\ruby{劉備}{りゅうび}の\ruby{面}{おもて}をじいと見すえたきり、\ruby{眼}{ま}ばたきもしなかった。\\


　やがて、独りで、うーむと\ruby{唸}{うな}っていたが、なに思ったか、\\


「あ、あ！　あなただっ」\\


　膝を折って、床に坐り、あたかも現世の\ruby{文殊}{もんじゅ}\ruby{弥勒}{みろく}でも見たように、何度も礼拝して止まなかった。\\


　劉備は、迷惑がって、\\


「老僧、何をなさいます」と、手を取った。\\


　老僧は、彼の手にふれると、なおさら、\ruby{随喜}{ずいき}の涙を流さぬばかりふるえて、額に押しいただきながら、\\


「青年。――わしは長いこと待っていたよ。まさしく、わしの待っていたのはあなただ。――あなたこそ\ruby{魔魅跳梁}{まみちょうりょう}を退けて、暗黒の国に楽土を\ruby{創}{た}て、\ruby{乱麻}{らんま}の世に道を示し、\ruby{塗炭}{とたん}の底から大民を救ってくれるお方にちがいない」と、いった。\\


「とんでもない。私は\ruby{{\fontspec{jigmo}\char"6DBF}県}{たくけん}から迷ってきた貧しい\ruby{蓆売}{むしろう}りです。老僧はなしてください」\\


「いいや、あなたの人相骨がらに現われておるよ。青年、聞かしておくれ。あなたの祖先は、帝系の流れか、王侯の血をひいていたろう」\\


「ちがう」\\


　劉備は、首を振って、「父も、祖父も、楼桑村の百姓でした」\\


「もっと先は……」\\


「わかりません」\\


「分らなければ、わしの言を信じたがよい。あなたが\ruby{佩}{は}いている剣は誰にもらったのか」\\


「\ruby{亡父}{ちち}の\ruby{遺物}{かたみ}」\\


「もっと前から、家におありじゃったろう。古びて見る面影もないがそれは\ruby{凡人}{ただびと}の\ruby{佩}{は}く剣ではない。\ruby{琅{\fontspec{jigmo}\char"7395}}{ろうかん}の\ruby{珠}{たま}がついていたはず、\ruby{戛玉}{かつぎょく}とよぶ珠だよ。\ruby{剣帯}{けんたい}に革か\ruby{錦}{にしき}の\ruby{腰帛}{ようはく}もついていたのだよ。王者の\ruby{佩}{はい}とそれを呼ぶ。何しろ、\ruby{刀身}{なかみ}も\ruby{無双}{むそう}な名剣にまちがいない。試してみたことがおありかの」\\


「……？」\\


　堂の外へ先に出たが、後から劉備が出てこないので、足を止めていた賊の馬元義と甘洪は、老僧のぶつぶついっていることばを、聞きすましながら振向いていた。が、――しびれをきらして、\\


「やいっ劉。いつまで何をしているんだ。荷物を持って早くこいっ」と、どなった。\\


　老僧は、まだ何か、いいつづけていたが、馬の大声に\ruby{恟}{すく}んで、急に口をつぐんだ。劉備はその\ruby{機}{しお}に、堂の外へ出てきた。\\


　\ruby{驢}{ろ}をつないでいる以前の門を踏みだすと、馬元義は、驢の手綱をときかける手下の甘を止めて、\\


「劉、そこへ掛けろ」と、木の根を指さし、自分も石段に腰かけて、大きく構えた。\\


「今、聞いていると、てめえは行く末、偉い者になる人相を備えているそうだな。まさか、王侯や将軍になれっこはあるめえが、俺も実は、てめえは見込みのある野郎だと見ているんだ――どうだ、俺の部下になって、黄巾党の仲間へ加盟しないか」\\


「はい。有難うございますが」と、劉備はあくまで、\ruby{素直}{すなお}をよそおって、\\


「私には、\ruby{故郷}{くに}に一人の母がいますので、せっかくですが、お仲間には入れません」\\


「おふくろなぞは、あってもいいじゃねえか。\ruby{喰}{く}い\ruby{扶持}{ぶち}さえ送ってやれば」\\


「けれど、こうして、私が旅に出ている間も、痩せるほど子の心配ばかりしている、至って\ruby{子煩悩}{こぼんのう}な母ですから」\\


「そりゃそうだろう。貧乏ばかりさせておくからだ。黄巾党に入って、腹さえふくらせておけば、なに、\ruby{嬰児}{あかご}じゃあるめえし、子の心配などしているものか」\\

\\

\subsubsection{四}

\\


　馬元義は、功名に燃えやすい青年の心をそそるように、それから黄巾党の勢力やら、世の中の将来やらを、談義しはじめた。\\


「狭い目で見ている奴は、俺たちが良民いじめばかりしていると思っているが、俺たちの総大将\ruby{張角}{ちょうかく}様を、神の如く\ruby{崇}{あが}めている地方だってかなりある――」\\


　と、\ruby{前提}{まえおき}して、まず、黄巾党の起りから説きだすのだった。\\


　今から十年ほど前。\\


　\ruby{鉅鹿郡}{きょろくぐん}（河北省）の人で、張角という無名の士があった。\\


　張角はしかし\ruby{稀世}{きせい}の秀才と、郷土でいわれていた。その張角が、あるとき、山中へ薬をとりに入って、道で異相の\ruby{道士}{どうし}に出会った。道士は手に\ruby{藜}{あかざ}の杖をもち、\\


（お前を待っていること久しかった）と、さしまねくので、ついて行ってみると、白雲の\ruby{裡}{うち}の洞窟へ\ruby{誘}{いざな}い、張角に三巻の書物を授けて、（これは、太平要術という書物である。この書をよく体して、天下の\ruby{塗炭}{とたん}を救い、道を興し、善を施すがよい。もし自身の\ruby{我意栄耀}{がいえいよう}に酔うて、悪心を起す時は、天罰たちどころに身を亡ぼすであろう）と、いった。\\


　張角は、再拝して、\ruby{翁}{おきな}の名を問うと、\\


（我は\ruby{南華老仙}{なんかろうせん}なり）と答え、姿は、一\ruby{颯}{さつ}の白雲となって飛去ってしまったというのである。\\


　張角は、そのことを、山を降りてから、里の人々へ自分から話した。\\


　正直な、里の人々は、（わしらの郷土の秀才に、神仙が宿った）と\ruby{真}{ま}にうけて、たちまち張角を、救世の\ruby{方師}{ほうし}と\ruby{崇}{あが}めて、触れまわった。\\


　張角は、門を閉ざし、\ruby{道衣}{どうい}を着て、\ruby{潔斎}{けっさい}をし、常に南華老仙の書を帯びて、昼夜行いすましていたが、或る年\ruby{悪疫}{あくえき}が流行して、村にも毎日おびただしい死人が出たので、\\


（今は、神が我をして、出でよと命じ給う日である）\\


　と、おごそかに、\ruby{草門}{そうもん}を開いて、病人を救いに出たが、その時もう、彼の門前には、五百人の者が、弟子にしてくれといって、\ruby{蝟集}{いしゅう}してぬかずいていたということである。\\


　五百人の弟子は、彼の命に依って、\ruby{金仙丹}{きんせんたん}、\ruby{銀仙丹}{ぎんせんたん}、\ruby{赤神丹}{せきしんたん}の秘薬をたずさえ、おのおの、悪疫の地を視て廻った。そして、張角方師の\ruby{功徳}{くどく}を語り聞かせ、男子には金仙丹を、女子には銀仙丹を、幼児には赤神丹を与えると、神薬のききめはいちじるしく、皆、数日を出でずして\ruby{癒}{なお}った。\\


　それでも、癒らぬ者は、張角自身が行って、\ruby{大喝}{だいかつ}の\ruby{呪}{じゅ}を\ruby{唱}{とな}え、病魔を家から追うと称して、\ruby{符水}{ふすい}の法を施した。それで起きない病人はほとんどなかった。\\


　体の病人ばかりでなく、次には心に病のある者も集まってきて、張角の前に\ruby{懺悔}{ざんげ}した。貧者も来た。富者も来た。美人も来た。力士や武術者も来た。それらの人々は皆、張角の\ruby{帷幕}{いばく}に参じたり、\ruby{厨房}{ちゅうぼう}で働いたり、彼のそば近く\ruby{侍}{じ}したり、また多くの弟子の中に交じって、弟子となったことを誇ったりした。\\


　たちまち、諸州にわたって、彼の勢力はひろまった。\\


　張角は、その弟子たちを、三十六の方を立たせ、階級を作り、大小に分かち、頭立つ者には\ruby{軍帥}{ぐんすい}の称を許し、また方帥の称呼を授けた。\\


　\ruby{大方}{だいほう}を行う者、一万余人。小方を行う者六、七千人。その部の内に、部将あり\ruby{方兵}{ほうへい}あり、そして張角の兄弟、\ruby{張梁}{ちょうりょう}、張宝のふたりを、\ruby{天公将軍}{てんこうしょうぐん}、\ruby{地公将軍}{ちこうしょうぐん}とよばせて、最大の権威をにぎらせ、自身はその上に君臨して、\ruby{大賢良師}{だいけんりょうし}\ruby{張角}{ちょうかく}と、\ruby{称}{とな}えていた。\\


　これがそもそもの、黄巾党の起りだとある。初め張角が、常に、結髪を黄色い\ruby{巾}{きれ}でつつんでいたので、その\ruby{風}{ふう}が全軍にひろまって、いつか党員の\ruby{徽章}{きしょう}となったものである。\\

\\

\subsubsection{五}

\\


　また、黄巾軍の徒党は、全軍の旗もすべて黄色を用い、その\ruby{大旆}{おおはた}には、\\

\ruby{蒼天已死}{そうてんすでにしす}\\

\ruby{黄夫当レ立}{こうふまさにたつべし}\\

\ruby{歳在二甲子一}{としこうしにありて}\\

\ruby{天下大吉}{てんかだいきち}\\



　という宣文を書き、党の楽謡部は、その宣文に、童歌風のやさしい作曲をつけて、党兵に唄わせ、部落や村々の地方から郡、県、市、都へと熱病のようにうたい\ruby{流行}{はや}らせた。\\



大賢良師張角！\\


大賢良師張角！\\



　今は、三歳の児童も、その名を知らぬはなく、\\


（――蒼天スデニ死ス。黄夫マサニ立ツベシ）\\


　と唄った後では、張角の名を\ruby{囃}{はや}して、今にも、天上の楽園が地上に実現するような感を民衆に抱かせた。\\


　けれど、黄巾党が\ruby{跋扈}{ばっこ}すればするほど、\ruby{楽土}{らくど}はおろか、一日の\ruby{安穏}{あんのん}も土民の中にはなかった。\\


　張角は自己の勢力に服従してくる愚民どもへは、\\


（太平を楽しめ）と、\ruby{逸楽}{いつらく}を許し、\\


（わが世を謳歌せよ）と、\ruby{暗}{あん}に掠奪を奨励した。\\


　その代りに、\ruby{逆}{さか}らう者は、\ruby{仮借}{かしゃく}なく罰し、人間を殺し、財宝を\ruby{掠}{かす}めとることが、党の日課だった。\\


　\ruby{地頭}{じとう}や地方の官吏も、防ぎようはなく、中央の\ruby{洛陽}{らくよう}の王城へ、急を告げることもひんぴんであったが、現下、漢帝の宮中は、\ruby{頽廃}{たいはい}と内争で乱脈をきわめていて、地方へ兵をやるどころではなかった。\\


　天下一統の大業を完成して、後漢の代を興した光武帝から、今は二百余年を経、宮府の内外にはまた、ようやく\ruby{腐爛}{ふらん}と\ruby{崩壊}{ほうかい}の\ruby{兆}{ちょう}があらわれてきた。\\


　十一代の帝、\ruby{桓帝}{かんてい}が\ruby{逝}{ゆ}いて、十二代の帝位についた霊帝は、まだ十二、三歳の幼少であるし、輔佐の重臣は、幼帝をあざむき合い、\ruby{朝綱}{ちょうこう}を\ruby{猥}{みだ}りにし、\ruby{佞智}{ねいち}の者が勢いを得て、真実のある人材は、みな野に追われてしまうという状態であった。\\


　心ある者は、ひそかに、\\


（どうなり行く世か？）と、憂えているところへ、地方に\ruby{蜂起}{ほうき}した黄巾賊の口々から、\\


　――\ruby{蒼天已死}{そうてんすでにしす}\\


　の童歌が流行ってきて、後漢の末世を暗示する声は、\ruby{洛陽}{らくよう}の城下にまで、満ちていた。\\


　そうした折にまた、こんなこともひどく人心を不安にさせた。\\


　ある年。\\


　幼帝が\ruby{温徳殿}{うんとくでん}に出御なされると、にわかに、狂風がふいて、\ruby{長}{たけ}二丈余の青蛇が、\ruby{梁}{はり}から帝の椅子のそばに落ちてきた。帝はきゃっと、床に\ruby{仆}{たお}れて気を失われてしまった。殿中の騒動はいうまでもなく、\ruby{弓箭}{きゅうせん}や\ruby{鳳尾槍}{ほうびそう}をもった禁門の武士がかけつけて、青蛇を\ruby{刺止}{しと}めんとしたところが、突如、\ruby{雹}{ひょう}まじりの大風が王城をゆるがして、青蛇は雲となって飛び、その日から三日三夜、大雨は底のぬけるほど降りつづいて、洛陽の民家の\ruby{浸水}{みずつ}くもの二万戸、崩壊したもの千何百戸、溺死怪我人算なし――というような大災害を生じた。\\


　また、つい近年には。\\


　赤色の\ruby{彗星}{すいせい}が現れたり、風もない真昼、\ruby{黒旋風}{こくせんぷう}が突然ふいて、王城の屋根望楼を飛ばしたり、五\ruby{原山}{げんざん}の山つなみに、部落数十が、一夜に地底へ埋没してしまったり――凶兆ばかり年ごとに起った。\\

\\

\subsubsection{六}

\\


　そんな凶兆のあるたびに、黄巾賊の「蒼天スデニ死ス――」の歌は、盲目的にうたわれて行き、賊党に加盟して、掠奪、横行、\ruby{殺戮}{さつりく}――の自由にできる「我党の太平を楽しめ」とする者が、ふえるばかりだった。\\


　思想の悪化、組織の混乱、道徳の\ruby{頽廃}{たいはい}。――これをどうしようもない後漢の末期だった。\\


　\ruby{燎原}{りょうげん}の火とばかり、魔の手をひろげて行った黄巾賊の勢力は、今では\ruby{青州}{せいしゅう}、\ruby{幽州}{ゆうしゅう}、\ruby{徐州}{じょしゅう}、\ruby{冀州}{きしゅう}、\ruby{荊州}{けいしゅう}、\ruby{揚州}{ようしゅう}、\ruby{{\fontspec{jigmo}\char"5157}州}{えんしゅう}、\ruby{予州}{よしゅう}等の諸地方に及んでいた。\\


　州の諸侯をはじめ、郡県市部の\ruby{長}{おさ}や官吏は、逃げ散るもあり、\ruby{降}{くだ}って賊となるもあり、\ruby{屍}{かばね}を積んで、\ruby{焚}{や}き殺された者も数知れなかった。\\


　富豪は皆、財を捧げて、生命を乞い、寺院や民家は戸ごとに、大賢良師張角――と書いた例の\ruby{黄符}{こうふ}を\ruby{門}{かど}に貼って、絶対服従を誓い、まるで鬼神をまつるように、\ruby{崇}{あが}め恐れた。そうした現状にあった。\\


　さて。……\\


　長々と、そうした現状や、黄巾党の\ruby{勃興}{ぼっこう}などを、自慢そうに語りきたって、\\


「\ruby{劉}{りゅう}――」と、\ruby{大方}{だいほう}\ruby{馬元義}{ばげんぎ}は、腰かけている石段から、寺の門を、\ruby{顎}{あご}でさした。\\


「そこでも、黄色い貼紙を見たろう。書いてある文句も読んだろう。この地方もずっと、俺たち黄巾党の勢力範囲なのだ」\\


「…………」\\


　劉備は、終始黙然と聞いているのみだった。\\


「いや、この地方や、十州や二十州はおろかなこと、今に天下は黄巾党のものになる。後漢の代は亡び、次の新しい代になる」\\


　劉備は、そこで初めて、こう訊ねた。\\


「では、張角良師は、後漢を亡ぼした後で、自分が帝位につく\ruby{肚}{はら}なんですか」\\


「いやいや。張角良師には、そんなお考えはない」\\


「では、誰が、次の帝王になるのでしょう」\\


「それはいえない。……だが劉備、てめえが俺の部下になると約束するなら聞かせてやるが」\\


「なりましょう」\\


「きっとか」\\


「母が許せばです」\\


「――では打明けてやるが、帝王の問題は、今の漢帝を亡ぼしてから後の重大な評議になるんだ。\ruby{匈奴}{きょうど}（\ruby{蒙古族}{もうこぞく}）のほうとも相談しなければならないから」\\


「へえ？　……なぜです。どうして支那の帝王を決めるのに、昔から\ruby{秦}{しん}や\ruby{趙}{ちょう}や\ruby{燕}{えん}などの\ruby{国境}{さかい}を侵して、われわれ漢民族を\ruby{脅}{おびや}かしてきた異国の匈奴などと相談する必要があるのですか」\\


「それは大いにあるさ」と、馬は当然のように――\\


「いくら俺たちが暴れ廻ろうたって、俺たちの\ruby{背後}{うしろ}から、軍費や兵器をどしどし廻してくれる黒幕がなくっちゃ、こんな短い年月に、後漢の天下を\ruby{攪乱}{かくらん}することはできまいじゃねえか」\\


「えっ。では黄巾賊のうしろには、異国の匈奴がついているわけですか」\\


「だから絶対に、俺たちは\ruby{敗}{ま}けるはずはないさ。どうだ劉、俺がすすめるのは、貴様の出世のためだ。部下になれ、すぐここで、黄巾賊に加盟せぬか」\\


「結構なお話です。母も聞いたら\ruby{歓}{よろこ}びましょう。……けれど、親子の中にも礼儀ですから、一応、母にも告げた上でご返辞を……」\\


　云いかけているのに、馬元義は不意に起ち上がって、\\


「やっ、来たな」と、彼方の平原へ向って、眉に手をかざした。\\

\\

\subsection{\ruby{白芙蓉}{びゃくふよう}}

\\

\subsubsection{一}

\\


　それは約五十名ほどの賊の小隊であった。中に\ruby{驢}{ろ}に乗っている二、三の賊将が\ruby{鉄鞭}{てつべん}を\ruby{指}{さ}して、何かいっていたように見えたが、やがて、馬元義の姿を見かけたか、寺のほうへ向って、一散に近づいてきた。\\


「やあ、\ruby{李朱氾}{りしゅはん}。遅かったじゃないか」\\


　こなたの馬元義も、石段から伸び上がっていうと、\\


「おう\ruby{大方}{だいほう}、これにいたか」と、李と呼ばれた男も、そのほかの仲間も、つづいて驢の鞍から降りながら、\\


「峠の\ruby{孔子廟}{こうしびょう}で待っているというから、あれへ行った所、姿が見えないので、俺たちこそ、大まごつきだ。遅いどころじゃない」と、汗をふきふき、かえって馬元義に向って、不平を並べたが、同類の冗談半分とみえて、責められた\ruby{馬}{ば}のほうも、げらげら笑うのみだった。\\


「ところで、ゆうべの\ruby{収穫}{みいり}はどうだな。洛陽船を\ruby{的}{あて}に、だいぶ諸方の\ruby{商人}{あきんど}が泊っていた筈だが」\\


「大していう程の収穫もなかったが、一村焼き払っただけの物はあった。その財物は皆、荷駄にして、例の通りわれわれの営倉へ送っておいたが」\\


「近頃は人民どもも、金は\ruby{埋}{い}けて隠しておく方法をおぼえたり、商人なども、隊伍を組んで、俺たちが襲うまえに、うまく逃げ散ってしまうので、だんだん以前のようにうまいわけには行かなくなったなあ」\\


「ウム、そういえば、先夜も一人惜しいやつを取逃がしたよ」\\


「惜しい奴？　――それは何か高価な財宝でも持っていたのか」\\


「なあに、砂金や宝石じゃないが、洛陽船から、茶を交易した男があるんだ。知っての通り、盟主張角様には、茶ときては、眼のない好物。これはぜひ\ruby{掠}{かす}めとって、\ruby{大賢良師}{だいけんりょうし}へご献納もうそうと、そいつの泊った\ruby{旅籠}{はたご}も目ぼしをつけておき、その近所から焼き払って踏みこんだところ、いつの間にか、逃げ\ruby{失}{う}せてしまって、とうとう見つからない。――こいつあ近頃の失策だったよ」\\


　賊の\ruby{李朱氾}{りしゅはん}は、劉備のすぐそばで、それを大声で話しているのだった。\\


　劉備は、驚いた。\\


　そして思わず、\ruby{懐中}{ふところ}に秘していた\ruby{錫}{すず}の小さい\ruby{茶壺}{ちゃつぼ}をそっとさわってみた。\\


　すると、馬元義は、\\


「ふーむ」と、うめきながら、改めて後ろにいる劉青年を振向いてから、さらに、李へ向って、\\


「それは、\ruby{幾歳}{いくつ}ぐらいな男か」\\


「そうさな。俺も見たわけでないが、\ruby{嗅}{か}ぎつけた部下のはなしによると、まだ若いみすぼらしい\ruby{風態}{ふうてい}の男だが、どこか\ruby{凛然}{りんぜん}としているから、油断のならない人間かも知れないといっていたが」\\


「じゃあ、この男ではないのか」\\


　馬元義は、すぐ傍らにいる劉備を指さして、いった。\\


「え？」\\


　李は、意外な顔をしたが、馬元義から\ruby{仔細}{しさい}を聞くとにわかに怪しみ疑って、\\


「そいつかもしれない。――おういっ、\ruby{丁峰}{ていほう}、丁峰」\\


　と、池畔に\ruby{屯}{たむろ}させてある部下の群れへ向ってどなった。\\


　手下の丁峰は、呼ばれて、屯の中から馳けてきた。李は、黄河で茶を交易した若者は、この男ではないかと、劉の顔を指さして、質問した。\\


　丁は、劉青年を見ると、惑うこともなくすぐ答えた。\\


「あ。この男です。この若い男に違いありません」\\


「よし」\\


　李は、そういって、丁峰を退けると、馬元義と共に、いきなり劉備の両手を左右からねじあげた。\\

\\

\subsubsection{二}

\\


「こら、貴様は茶をかくしているというじゃないか。その茶壺をこれへ出してしまえ」\\


　馬元義も責め、\ruby{李朱氾}{りしゅはん}も共に、劉備のきき腕を、ねじ抑えながら\ruby{脅}{おど}した。\\


「出さぬと、ぶった斬るぞ。今もいった通り、張角良師のご好物だが、良師のご威勢でさえ、めったに手にはいらぬ程の物だ。貴様のような\ruby{下民}{げみん}などが、茶を持ったところで、何となるものか。われわれの手を経て、良師へ献納してしまえ」\\


　劉備は、云いのがれのきかないことを、はやくも観念した。しかし、\ruby{故郷}{くに}の母が、いかにそれを楽しみに待っているかを思うと、自分の\ruby{生命}{いのち}を求められたより辛かった。\\


（何とか、ここをのがれる工夫はないものか）\\


　となお、未練をもって、両手の痛みをこらえていると、李朱氾の靴は、気早に劉備の腰を蹴とばして、「\ruby{唖}{おし}か、つんぼか、おのれは」と、\ruby{罵}{ののし}った。\\


　そして、よろめく劉備の襟がみを、つかみもどして、\\


「あれに、血に飢えている五十の部下がこちらを見て、\ruby{餌}{え}を欲しがっているのが、眼に見えないか。返辞をしろ」と、\ruby{威猛高}{いたけだか}にいった。\\


　劉備は二人の土足の前へ、そうしてひれ伏したまま、まだ、母の歓びを売って、この場を助かる気持になれないでいたが、ふと、眼を上げると、寺門の陰にたたずんで、こちらを覗いていた最前の老僧が、\\


（物など惜しむことはない。求める物は、何でも与えてしまえ、与えてしまえ）\\


　と、手真似をもって、しきりと彼の善処をうながしている。\\


　劉備もすぐ、（そうだ。この身体を傷つけたら、母にも大不孝となる）と思って、心をきめたが、それでもまだ\ruby{懐中}{ふところ}の茶壺は出さなかった。腰に\ruby{佩}{は}いている剣の帯革を解いて、\\


「これこそは、父の\ruby{遺物}{かたみ}ですから、自分の\ruby{生命}{いのち}の次の物ですが、これを献上します。ですから、茶だけは見のがして下さい」と哀願した。\\


　すると、馬元義は、\\


「おう、その剣は、俺がさっきから眼をつけていたのだ。貰っておいてやる」と\ruby{奪}{と}り上げて、「茶のことは、俺は知らん」と、空うそぶいた。\\


　\ruby{李朱氾}{りしゅはん}は、前にもまして怒りだして、一方へ剣を渡して、俺になぜ茶壺を渡さないかと責めた。\\


　劉備は、やむなく、肌深く持っていた\ruby{錫}{すず}の小壺まで出してしまった。李は、\ruby{宝珠}{ほうしゅ}をえたように、\ruby{両掌}{りょうて}を捧げて、\\


「これだ、これだ。洛陽の\ruby{銘葉}{めいよう}に違いない。さだめし良師がおよろこびになるだろう」と、いった。\\


　賊の小隊はすぐ先へ出発する予定らしかったが、ひとりの物見が来て、ここから十里ほどの先の河べりに、県の吏軍が約五百ほど野陣を張り、われわれを捜索しているらしいという報告をもたらした。で、にわかに、「では、今夜はここへ泊れ」となって、約五十の黄巾賊は、そのまま寺を宿舎にして、携帯の\ruby{糧嚢}{りょうのう}を解きはじめた。\\


　夕方の炊事の混雑をうかがって、劉備は今こそ逃げるによい\ruby{機}{しお}と、薄暮の門を、そっと外へ踏みだしかけた。\\


「おい。どこへ行く」\\


　賊の\ruby{哨兵}{しょうへい}は、見つけるとたちまち、大勢して彼を包囲し、奥にいる馬元義と李朱氾へすぐ知らせた。\\

\\

\subsubsection{三}

\\


　劉備は\ruby{縛}{いまし}められて、\ruby{斎堂}{さいどう}の丸柱にくくりつけられた。\\


　そこは床に瓦を敷き詰め、太い丸柱と、小さい窓しかない石室だった。\\


「やい劉。貴様は、おれの眼をかすめて、逃げようとしたそうだな。察するところ、てめえは官の密偵だろう。いいや違えねえ。きっと県軍のまわし者だ。――今夜、十里ほど先まで、県軍がきて野陣を張っているそうだから、それへ連絡を取るために、\ruby{脱}{ぬ}け出そうとしたのだろう」\\


　馬元義と李朱氾は、かわるがわるに来て、彼を\ruby{拷問}{ごうもん}した。\\


「――道理で、貴様の面がまえは、\ruby{凡者}{ただもの}でないはずだ。県軍のまわし者でなければ、洛陽の直属の隠密か。いずれにしても、官人だろうてめえは。――さ、泥を吐け。いわねば、痛い思いをするだけだぞ」\\


　しまいには、馬と李と、二人がかりで、劉を蹴って\ruby{罵}{ののし}った。\\


　劉は一口も物をいわなかった。こうなったからには、天命にまかせようと観念しているふうだった。\\


「こりゃひと筋縄では口をあかんぞ」\\


　李は、持てあまし気味に、\ruby{馬}{ば}へ向ってこう提議した。\\


「いずれ明日の早暁、俺はここを出発して、張角良師の総督府へ参り、例の茶壺を献上かたがた良師のご機嫌伺いに出るつもりだが、その折、こいつも引っ立てて行って、大方軍本部の軍法会議にさし廻してみたらどうだろう。思いがけない拾いものになるかもしれぬぜ」\\


　よかろうと、馬も同意だ。\\


　斎堂の扉は、かたく閉められてしまった。夜が更けると、ただ一つの高い窓から、今夜も銀河の秋天が冴えて見える。けれどとうてい、そこからのがれ出る工夫はない。\\


　どこかで、馬のいななきがする。官の県軍が攻めてきたのならよいが――と劉備は、望みをつないだが、それは物見から帰ってきた二、三の賊兵らしく、後は\ruby{寂}{せき}として、物音もなかった。\\


「母へ孝養を努めようとして、かえって大不孝の子となってしまった。死ぬる身は惜しくもないが、老母の余生を悲しませ、不孝の\ruby{屍}{むくろ}を野にさらすのは悲しいことだ」\\


　劉備は、星を仰いで\ruby{嘆}{なげ}いた。そして、孝行するにも、身に不相応な望みを持ったのが悪かったと悔いた。\\


　賊府へひかれて、人中で生恥さらして殺されるよりは、いっそ、ここで、ひと思いに死なんか――と考えた。\\


　死ぬにも、身に剣はなかった。柱に頭を打ちつけて憤死するか。舌を噛んで星夜をねめつけながら\ruby{呪死}{じゅし}せんか。\\


　劉備は、悶々と、迷った。\\


　――すると彼の\ruby{眸}{ひとみ}の前に一筋の縄が下がってきた。それは神の意志によって下がってくるように、高い切窓の口から石の壁に伝わってスルスルと垂れてきたのである。\\


「……あ？」\\


　人影もなにも見えない、ただ四角な星空があるだけだった。\\


　劉備は、身を起しかけた。しかしすぐ無益であることを知った。身は\ruby{縛}{いまし}めにかかっている、この縄目の解けない以上、救い手がそこまで来ていても、すがりつく\ruby{術}{すべ}はない。\\


「……ああ、誰だろう？」\\


　誰か、窓の下へ、救いに来ている。外で自分を待っていてくれる者がある。劉備は、なおさらもがいた。\\


　と、――彼の行動が遅いので、早くしろとうながすように、外の者は\ruby{焦}{じ}れているのであろう。高窓から垂れている縄が左右に動いた。そして縄の端に\ruby{結}{ゆ}いつけてあった短剣が、白い魚のように、コトコトと\ruby{瓦}{かわら}の\ruby{床}{ゆか}を打って躍った。\\

\\

\subsubsection{四}

\\


　足の先で、短剣を寄せた。そしてようやく、それを手にして、自身の縄目を断ち切ると、劉備は、窓の下に立った。\\


（早く。早く）といわんばかりに、無言の縄は外から意志を伝えて、ゆれうごいている。\\


　劉備は、それにつかまった。石壁に足をかけて、窓から外を見た。\\


「……オオ」\\


　外にたたずんでいたのは、昼間、ただひとりで\ruby{曲{\fontspec{jigmo}\char"5F54}}{きょくろく}に腰かけていたあの老僧だ。骨と皮ばかりのような彼の細い影であった。\\


「――今だよ」\\


　その手がさしまねく。\\


　劉備はすぐ地上へ跳びおりた。待っていた老僧は、彼の身を抱えるようにして、物もいわず馳けだした。\\


　寺の裏に、\ruby{疎林}{そりん}があった。樹の間の細道さえ、銀河の秋はほの明るい。\\


「老僧、老僧。いったいどっちへ逃げるんですか」\\


「まだ、逃げるのじゃない」\\


「では、どうするんです」\\


「あの\ruby{塔}{とう}まで行ってもらうのじゃよ」\\


　走りながら、老僧は指さした。\\


　見るとなるほど、疎林の奥に、疎林の\ruby{梢}{こずえ}よりも、高くそびえている古い塔がある。老僧は、あわただしく古塔の\ruby{扉}{と}をひらいて中へ隠れた。そしてあんなに急いだのに、なかなか出てこなかった。\\


「どうしたのだろう？」\\


　劉備は気を揉んでいる。そして賊兵が追ってきはしまいかと、あちこち見まわしているとやがて、\\


「青年、青年」\\


　小声で呼びながら、塔の中から老僧は何かひきながら出てきた。\\


「おや？」\\


　劉備は眼をみはった。老僧が引っぱっているのは駒の手綱だった。銀毛のように美しい白馬がひかれだしたのである。\\


　いや、いや、白馬の毛並の見事さや、背の鞍の華麗などはまだいうも\ruby{愚}{おろ}かであった。その駒に続いて、後ろから歩みも\ruby{嫋}{たおや}かに、世間の風にも怖れるもののように、\ruby{楚々}{そそ}と姿をあらわした美人がある。\ruby{眉}{まゆ}の麗しさ、耳の白さ、また、眼にふくむ\ruby{愁}{うれ}いの悩ましいばかりなど、思いがけぬ場合ではあり、星夜の光に見るせいか、この世の人とも思えぬのであった。\\


「青年。わしがお前を助けて上げたことを、恩としてくれるなら、逃げるついでに、このお\ruby{嬢}{じょう}さまを連れて、ここから十里ほど北へ向った所の河べりに陣している県軍の隊まで、届けて上げてくれぬか。わずか十里じゃ、この白馬に鞭打てば――」\\


　老僧のことばに、劉備は、\ruby{否}{いな}やもなく、はいと答えるべきであるが、その任務よりも、届ける人のあまりに美し過ぎるので、なんとなくためらわれた。\\


　老僧は、彼のためらいを、どう解釈したか。\\


「そうだ、\ruby{氏素性}{うじすじょう}も知れない婦人をと、疑ぐっておるのじゃろうが、心配するな。このお方は、つい先頃までの、この地方県城を預かっておられた領主のお\ruby{嬢}{じょう}さまじゃ。黄巾賊の乱入にあって、県城は焼かれ、ご領主は殺され、家来は四散し、ここらの寺院さえ、あの通りに成り果てたが、その乱軍の中から迷うてござったお嬢さまを、実はわしが、ここの塔へそっと\ruby{匿}{かくも}うて――」\\


　と、老僧の眼がふと、古塔の\ruby{頂}{いただき}を見上げた時、疎林を渡る秋風の外に、にわかに、人の跫音や馬のいななきが聞えだした。\\

\\

\subsubsection{五}

\\


　\ruby{劉備}{りゅうび}が、眼をくばると、\\


「いや、動かぬがよい。しばらくは、かえってここに、じっとしていたほうが……」\\


　と、老僧が彼の袖をとらえ、そんな危急の中になお、語りつづけた。\\


　県の城長の娘は、名を\ruby{芙蓉}{ふよう}といい姓は\ruby{鴻}{こう}ということ。また、今夜近くの河畔にきて宿陣している県軍は、きっと先に四散した城長の家臣が、残兵を集めて、黄巾賊へ報復を計っているに違いないということ。\\


　だから、芙蓉の身を、そこまで届けてくれさえすれば、後は以前の家来たちが守護してくれる――白馬の背へ二人してのって、抜け道から一気に逃げのびて行くように――と、\ruby{祷}{いの}るようにいうのだった。\\


「承知しました」\\


　劉備は、勇気を示して答えた。\\


「けれど\ruby{和上}{わじょう}、あなたはどうしますか」\\


「わしかの」\\


「そうです。私たちを逃がしたと賊に知られたら、和上の身は、ただでは済まないでしょう」\\


「案じることはない。生きていたとて、このさき幾年生きていられよう。ましてこの十数日は、草の根や虫など食うて、露命をつないでいたはかない身じゃ。それも\ruby{鴻家}{こうけ}の\ruby{阿嬢}{おむすめ}を助けて上げたい一心だけで生きていたが――今は、そのことも、頼む者に頼み果てたし、あなたという者をこの世に見出したので、思い残りは少しもない」\\


　老僧はそう云い終ると、風の如く、塔の中へ影をかくした。\\


　あれよと、芙蓉は、老僧を\ruby{慕}{した}って追いすがったが、とたんに、塔の口もとの扉は内から閉じられていた。\\


「和上さま。和上さま！」\\


　芙蓉は慈父を失ったように、扉をたたいて泣いていたが、その時、高い塔の\ruby{頂}{いただき}で、再び老僧の声がした。\\


「青年。わしの指をご覧。わしの指さすほうをご覧。――ここの疎林から西北だよ。\ruby{北斗星}{ほくとせい}がかがやいておる。それを\ruby{的}{あて}にどこまでも逃げてゆくがよい。南も東も\ruby{蓮池}{はすいけ}の\ruby{畔}{ほとり}も、寺の近くにも、賊兵の影が道をふさいでいる。逃げる道は、西北しかない。それも今のうちじゃ。はやく白馬に\ruby{鞭}{むち}打たんか」\\


「はいっ」\\


　答えながら仰ぐと、老僧の影は、塔上の\ruby{石欄}{せきらん}に立って、一方を指さしているのだった。\\


「\ruby{佳人}{かじん}。はやくおのりなさい。泣いているところではない」\\


　劉備は、彼女の\ruby{細腰}{さいよう}を抱き上げて、白馬の鞍にすがらせた。\\


　芙蓉の体はいと軽かった。柔軟で高貴な\ruby{薫}{かお}りがあった。そして彼女の手は、劉備の肩にまとい、劉の頬は、彼女の黒髪にふれた。\\


　劉備も木石ではない。かつて知らない\ruby{動悸}{ときめき}に、血が熱くなった。けれどそれは、地上から鞍の上まで、彼女の身を移すわずかな間でしかなかった。\\


「ご免」といいながら、劉備ものって一つ鞍へまたがった。そして片手に彼女をささえ、片手に白馬の手綱をとって、老僧の指さした方角へ馬首を向けた。\\


　塔上の老僧は、それを見おろすと、わが事おわれり――と思ったか、突然、歓喜の声をあげて、\\


「見よ、見よ。\ruby{凶雲}{きょううん}\ruby{没}{ぼっ}して、\ruby{明星}{みょうじょう}出づ。\ruby{白馬}{はくば}\ruby{翔}{か}けて、\ruby{黄塵}{こうじん}\ruby{滅}{めっ}す。――ここ数年を出でないうちじゃろう。青年よ、はや行け。おさらば」\\


　云い終ると、みずから舌を噛んで、塔上の石欄から百尺下の大地へ、身を躍らして、五体の骨を自分でくだいてしまった。\\

\\

\subsection{\ruby{張飛卒}{ちょうひそつ}}

\\

\subsubsection{一}

\\


　白馬は\ruby{疎林}{そりん}の細道を西北へ向ってまっしぐらに駆けて行った。秋風に舞う木の葉は、鞍上の\ruby{劉備}{りゅうび}と\ruby{芙蓉}{ふよう}の影を、\ruby{征箭}{そや}のようにかすめた。\\


　やがて\ruby{曠}{ひろ}い野に出た。\\


　野に出ても、二人の身をなお、\ruby{箭}{や}うなりがかすめた。今度のは木の葉のそれではなく、鋭い\ruby{鏃}{やじり}をもった鉄弓の矢であった。\\


「オ。あれへ行くぞ」\\


「女をのせて――」\\


「では違うのか」\\


「いや、やはり劉備だ」\\


「どっちでもいい。逃がすな。女も逃がすな」\\


　賊兵の声々であった。\\


　疎林の陰を出たとたんに、黄巾賊の一隊は早くも見つけてしまったのである。\\


　獣群の声が、\ruby{鬨}{とき}をつくって、白馬の影を追いつめて来た。\\


　劉備は、振り向いて、\\


「しまった！」\\


　思わずつぶやいたので、彼と白馬の脚とを唯一の頼みにしがみついていた芙蓉は、\\


「ああ、もう……」\\


　消え入るようにおののいた。\\


　万が一つも、助からぬものとは観念しながらも、劉備は励まして、\\


「大丈夫、大丈夫。ただ、振り落されないように、駒の\ruby{鬣}{たてがみ}と、私の帯に、必死でつかまっておいでなさい」と、いって、\ruby{鞭}{むち}打った。\\


　芙蓉はもう返事もしない。ぐったりと鬣に顔をうつ伏せている。その\ruby{容貌}{かんばせ}の白さはおののく\ruby{白芙蓉}{びゃくふよう}の花そのままだった。\\


「河まで行けば。県軍のいる河まで行けば！　……」\\


　劉備の打ちつづけていた\ruby{生木}{なまき}の鞭は、皮がはげて白木になっていた。\\


　低い\ruby{土坡}{どは}のうねりを躍り越えた。遠くに帯のように流れが見えてきた。しめたと、劉備は勇気をもり返したが、河畔まで来てもそこには何物の影もなかった。宵に\ruby{屯}{たむろ}していたという県軍も、賊の勢力に怖れをなしたか、陣を払って何処かへ去ってしまったらしいのである。\\


「待てッ」\\


　\ruby{驢}{ろ}にのった\ruby{精悍}{せいかん}な影は、その時もう五騎六騎と、彼の前後を包囲してきた。いうまでもなく黄巾賊の\ruby{小方}{しょうほう}（\ruby{小頭目}{しょうとうもく}）らである。\\


　驢を持たない徒歩の卒どもは、駒の足に続ききれないで、途中であえいでしまったらしいが、\ruby{李朱氾}{りしゅはん}をはじめとして、騎馬の小方たち七、八騎はたちまち追いついて、\\


「止れッ」\\


「射るぞ」と、どなった。\\


　鉄弓の\ruby{弦}{つる}をはなれた一\ruby{矢}{し}は、白馬の\ruby{環囲}{かんい}に突きささった。\\


　\ruby{喉}{のど}に矢を立てた白馬は、\ruby{棹立}{さおだ}ちに躍り上がって、一\ruby{声}{せい}いななくと、どうと横ざまに仆れた。\ruby{芙蓉}{ふよう}の身も、劉備の体も、共に大地へほうり捨てられていた。\\


　そのまま芙蓉は身動きもしなかったが、劉備は起ち上がって、\\


「何かっ！」と、さけんだ。彼は今日まで、自分にそんな大きな声量があろうとは知らなかった。百獣も為に\ruby{怯}{ひる}み、曠野を\ruby{野彦}{のびこ}して渡るような\ruby{大喝}{だいかつ}が、\ruby{唇}{くち}から無意識に出ていたのである。\\


　賊は、ぎょっとし、劉備の大きな眼の光におどろき、驢は彼の大喝に、\ruby{蹄}{ひづめ}をすくめて止った。\\


　だが、それは一瞬、\\


「何を、青二才」\\


「手むかう気か」\\


　驢を跳びおりた賊は、鉄弓を捨てて大剣を抜くもあり、槍を舞わして、劉備へいきなり突っかけてくるもあった。\\

\\

\subsubsection{二}

\\


　どういう悪日と\ruby{凶}{わる}い方位をたどってきたものだろうか。\\


　黄河の\ruby{畔}{ほとり}から、ここまでの間というものは、劉備は、幾たび死線を\ruby{彷徨}{ほうこう}したことか知れない。これでもかこれでもかと、彼を試さんとする百難が、次々に形を変えて待ちかまえているようだった。\\


「もうこれまで」\\


　劉備もついに観念した。避けようもない賊の包囲だ。斬り\ruby{死}{じに}せんものと覚悟をきめた。\\


　けれど身には寸鉄も帯びていない。少年時代から片時もはなさず持っていた父の\ruby{遺物}{かたみ}の剣も、先に賊将の馬元義に\ruby{奪}{と}られてしまった。\\


　劉備は、しかし、\\


「ただは死なぬ」と思い、石ころをつかむが早いか、近づく者の顔へ投げつけた。\\


　見くびっていた賊の一名は、不意を喰らって、\\


「あッ」と、鼻ばしらをおさえた。\\


　劉備は、飛びついて、その槍を奪った。そして大音に、\\


「四民を悩ます害虫ども、もはや\ruby{免}{ゆる}しはおかぬ。\ruby{{\fontspec{jigmo}\char"6DBF}県}{たくけん}の\ruby{劉備玄徳}{りゅうびげんとく}が腕のほどを見よや」\\


　といって、捨身になった。\\


　賊の小方、\ruby{李朱氾}{りしゅはん}は笑って、\\


「この百姓めが」と半月槍をふるってきた。\\


　もとより劉備はさして武術の達人ではない。田舎の\ruby{楼桑村}{ろうそうそん}で、多少の武技の稽古はしたこともあるが、それとて程の知れたものだ。武技を磨いて身を立てることよりも、\ruby{蓆}{むしろ}を織って母を養うことのほうが常に彼の急務であった。\\


　でも、必死になって、七人の賊を相手に、ややしばらくは、一命をささえていたが、そのうちに、槍を打落され、よろめいて倒れたところを、李朱氾に馬のりに組み敷かれて、李の大剣は、ついに、彼の胸いたに突きつけられた。\\


　――おおういっ。\\


　すると、……いやさっきからその声は遠くでしたのだが、\ruby{剣戟}{けんげき}のひびきで、誰の耳にも入らなかったのである。\\


　遥か彼方の野末から、\\


「――おおういっ。待ってくれい」\\


　呼ばわる声が近づいてくる。\\


　野彦のように凄い声は、思わず賊の頭を振り向かせた。\\


　両手を振りながら\ruby{韋駄天}{いだてん}と、こなたへ馳けてくる人影が見える。その迅いことは、まるで疾風に一葉の木の葉が舞ってくるようだった。\\


　だがまたたく間に近づいてきたのを見ると、木の葉どころか身の\ruby{丈}{たけ}七尺もある\ruby{巨漢}{おおおとこ}だった。\\


「やっ、\ruby{張卒}{ちょうそつ}じゃないか」\\


「そうだ。近頃、卒の中に入った下ッ端の\ruby{張飛}{ちょうひ}だ」\\


　賊は、不審そうに、顔見合せて云い合った。自分らの部下の中にいる張飛という一卒だからである。他の大勢の歩卒は、騎馬に追いつけず皆、途中で遅れてしまったのに、張卒だけが、たとえひと足遅れたにせよ、このくらいの差で追いついてきたのだから、その脚力にも、賊将たちは\ruby{愕}{おどろ}いたに違いなかった。\\


「なんだ、張卒」\\


　\ruby{李朱氾}{りしゅはん}は、膝の下に、劉備の体を抑えつけ、\ruby{右手}{めて}に大剣を持って、その胸いたに\ruby{擬}{ぎ}しながら振り向いていった。\\


「小方、小方。殺してはいけません。その人間は、わしに渡して下さい」\\


「何？　……誰の命令で貴様はそんなことをいうのか」\\


「卒の張飛の命令です」\\


「ばかっ。張飛は、貴様自身じゃないか。卒の分際で」\\


　と、いう言葉も終らぬ間に、そう\ruby{罵}{ののし}っていた李朱氾の体は、二丈もうえの空へ飛んで行った。\\

\\

\subsubsection{三}

\\


　卒の張飛が、いきなり李朱氾をつまみ上げて、宙へ投げ飛ばしたので、\\


「やっ、こいつが」と、賊の小方たちは、劉備もそっちのけにして、彼へ総掛りになった。\\


「やい張卒、なんで貴様は、味方の李小方を投げおったか。また、おれ達のすることを邪魔だてするかっ」\\


「ゆるさんぞ。ふざけた真似すると」\\


「党の軍律に照らして、成敗してくれる。それへ直れ」\\


　ひしめき寄ると、張は、\\


「わははははは。吠えろ吠えろ。\ruby{胆}{きも}をつぶした野良犬めらが」\\


「なに、野良犬だと」\\


「そうだ。その中に一匹でも、人間らしいのがおるつもりか」\\


「うぬ。\ruby{新米}{しんまい}の卒の分際で」\\


　\ruby{喚}{おめ}いた一人が、槍もろとも、躍りかかると、張飛は、\ruby{団扇}{うちわ}のような大きな手で、その横顔をはりつけるや否や、槍を引ッたくって、よろめく尻をしたたかに打ちのめした。\\


　槍の柄は折れ、打たれた賊は、腰骨がくだけたように、ぎゃっと\jruby{もんどり}{・・・・・・}打った。\\


　思わぬ裏切者が出て、賊は狼狽したが、日頃から図抜けた\ruby{巨漢}{おおおとこ}の鈍物と、小馬鹿にしていた卒なので、その怪力を眼に見ても、まだ張飛の真価を信じられなかった。\\


　張飛は、さながら岩壁のような胸いたをそらして、\\


「まだ来るか。\jruby{むだ}{・・・・・・}な\ruby{生命}{いのち}を捨てるより、おとなしく逃げ帰って、\ruby{鴻家}{こうけ}の姫と\ruby{劉備}{りゅうび}の身は、先頃、県城を焼かれて鴻家の亡びた時、降参と\ruby{偽}{いつわ}って、黄巾賊の卒にはいっていた張飛という者の手に渡しましたと、\ruby{有態}{ありてい}に報告しておけ」\\


「あっ！　……では汝は、鴻家の旧臣だな」\\


「いま気がついたか。此方は県城の\ruby{南門衛少督}{なんもんえいしょうとく}を勤めていた鴻家の武士で名は\ruby{張飛}{ちょうひ}、\ruby{字}{あざな}は\ruby{翼徳}{よくとく}と申すものだが無念や此方が他県へ公用で留守の間に、黄巾賊の\ruby{輩}{やから}のために、県城は焼かれ、主君は殺され、領民は苦しめられ、一夜に城地は焦土と化してしまった。――その無念さ、いかにもして\ruby{怨}{うら}みをはらしてくれんものと、身を偽り、敗走の兵と化けて、一時、其方どもの賊の中に、卒となって隠れていたのだ。――大方馬元義にも、また、総大将の兇賊張角にも、よく申しておけ。いずれいつかはきっと、張飛翼徳が思い知らせてくるるぞと」\\


　\ruby{雷}{いかずち}のような声だった。\\


　\ruby{豹頭環眼}{ひょうとうかんがん}、張飛がそういってくわっと\ruby{睨}{ね}めつけると、賊の小方らは、足もすくんでしまったらしいが、まだ衆をたのんで、\\


「さては、鴻家の残兵だったか。そう聞けばなおのこと、生かしてはおけぬ」と、一度に打ってかかった。\\


　張飛は、腰の剣も抜かず、寄りつく者をとっては投げた。投げられた者は皆、\ruby{脳骨}{のうこつ}をくだき、\ruby{眼窩}{がんか}は飛びだし、またたくうちに\ruby{碧血}{へきけつ}の大地、惨として、二度と起き上がる者はなかった。\\


　劉備は、茫然と、張飛の働きをながめていた。\ruby{燕飛龍{\fontspec{jigmo}\char"9B02}}{えんぴりゅうびん}、蹴れば雲を生じ、\ruby{吠}{ほ}ゆれば風が起るようだった。\\


「なんという豪傑だろう？」\\


　残る二、三人は、\ruby{驢}{ろ}に飛びついて逃げうせたが、張飛は笑って追いもしなかった。そして\ruby{踵}{きびす}をめぐらすと、劉備のほうへ大股に近づいてきて、\\


「いや旅の人。えらい目に遭いましたなあ」\\


　と、何事もなかったような顔して話しかけた。そして直ぐ、腰に帯びていた二剣のうちの一つをはずし、また、\ruby{懐中}{ふところ}から見おぼえのある茶の小壺を取出して、「これはあなたの物でしょう。賊に奪り上げられたあなたの剣と茶壺です。さあ取っておきなさい」と、劉の手へ渡した。\\

\\

\subsubsection{四}

\\


「あ。私のです」\\


　劉備は、失くした珠が返ってきたように、剣と茶壺の二品を、張飛の手から受取ると、幾度も感謝をあらわして、「すでに\ruby{生命}{いのち}もないところを救っていただいた上に、この大事な二品まで、自分の手に戻るとは、なんだか、夢のような心地がします。\ruby{大人}{たいじん}のお名前は、さきほど聞きました。心に銘記しておいて、ご恩は生涯忘れません」と、いった。\\


　張飛は、\ruby{首}{こうべ}を振って、\\


「いやいや徳は\ruby{孤}{こ}ならずで、貴公がそれがしの旧主、\ruby{鴻家}{こうけ}の姫を助けだしてくれた義心に対して、自分も義をもってお答え申したのみです。ちょうど最前、古塔のあたりから白馬にのって逃げた者があると、哨兵の知らせに、こよい黄巾賊の将兵が泊っていたかの寺が、すわと一度に、混雑におちた\ruby{隙}{すき}をうかがい、夕刻見ておいた貴公のその二品を、馬元義と李朱氾の眠っていた内陣の壇からすばやく奪い返し、追手の卒と共にこれまで馳けてきたものでござる。貴公の孝心と、誠実を天もよみし賜うて、自然お手に戻ったものでしょう」\\


　と、\ruby{理由}{わけ}をはなした。張飛が武勇に誇らない謙遜なことばに、劉備はいよいよ感じて、感銘のあまり二品のうちの剣のほうを差しだして、\\


「大人、失礼ですが、これはお礼として、あなたに差上げましょう。茶は、\ruby{故郷}{くに}に待っている母の土産なので、\ruby{頒}{わか}つことはできませんが、剣は、あなたのような\ruby{義胆}{ぎたん}の豪傑に持っていただけば、むしろ剣そのものも本望でしょうから」と、再び、張飛の手へ授けて云った。\\


　張飛は、眼をみはって、\\


「えっ、この品をそれがしに、賜ると仰っしゃるのですか」\\


「劉備の寸志です。どうか納めておいて下さい」\\


「自分は根からの武人ですから、実をいえば、この剣の世に稀な名刀だということは知っていますから、欲しくてならなかったところです。けれど、同時に貴公とこの剣との来歴も聞いていましたから、望むに望めないでおりましたが」\\


「いや、\ruby{生命}{いのち}の恩人へ酬いるには、これをもってしても、まだ足りません。しかも剣の真価を、そこまで、分っていて下されば、なおさら、差上げても張合いがあり、自分としても満足です」\\


「そうですか。しからば、ほかならぬ品ですから、頂戴しておこう」\\


　と、張飛は、自身の剣をすぐ解き捨て、\ruby{渇望}{かつぼう}の名剣を身に\ruby{佩}{は}いていかにもうれしそうであった。\\


「じゃあ早速ですが、まだ賊が押し返してくるにきまっている。それがしは鴻家のご息女を立てて、旧主の残兵を集め事を\ruby{謀}{はか}る考えですが――貴公も一刻もはやく、郷里へさしてお帰りなさい」\\


　張飛のことばに、\\


「おお、それでは」\\


　と、劉備は、\ruby{芙蓉}{ふよう}の身を\ruby{扶}{たす}けて、張飛に託し、自分は、賊の捨てた\ruby{驢}{ろ}をひろってまたがった。\\


　張飛は、先に自分が解き捨てた剣を劉備の腰に\ruby{佩}{は}かせてやりながら、\\


「こんな剣でも帯びておいでなされ。まだ、\ruby{{\fontspec{jigmo}\char"6DBF}県}{たくけん}までは、数百里もありますから」といった。\\


　そして張飛自身も、芙蓉の身を抱いて、白馬の上に移り、名残り惜しげに、\\


「いつかまた、再会の日もありましょうが、ではご機嫌よく」\\


「おお、きっとまた、会う日を待とう。あなたも武運めでたく、鴻家の再興を成しとげらるるように」\\


「ありがとう。では」\\


「おさらば――」\\


　劉備の驢と、芙蓉を抱えた張飛の白馬とは、\ruby{相顧}{あいかえ}りみながら、西と東に別れ去った。\\

\\

\subsection{桑の家}

\\

\subsubsection{一}

\\


　\ruby{{\fontspec{jigmo}\char"6DBF}県}{たくけん}の\ruby{楼桑村}{ろうそうそん}は、戸数二、三百の小駅であったが、春秋は北から南へ、南から北へと流れる旅人の多くが、この宿場で\ruby{驢}{ろ}をつなぐので、酒を売る\ruby{旗亭}{きてい}もあれば、\ruby{胡弓}{こきゅう}を\ruby{弾}{ひ}くひなびた\ruby{妓}{おんな}などもいて相当に賑わっていた。\\


　この地はまた、\ruby{太守}{たいしゅ}\ruby{劉焉}{りゅうえん}の領内で、\ruby{校尉}{こうい}\ruby{鄒靖}{すうせい}という代官が役所をおいて支配していたが、なにぶん、近年の物情騒然たる\ruby{黄匪}{こうひ}の\ruby{跳梁}{ちょうりょう}に\ruby{脅}{おびや}かされているので、楼桑村も例にもれず、夕方になると明るいうちから村はずれの城門をかたく閉めて、旅人も居住者も、いっさいの往来は止めてしまった。\\


　城門の\ruby{鉄扉}{てっぴ}が閉まる時刻は、大陸の\ruby{西{\fontspec{jigmo}\char"5393}}{さいがい}にまっ赤な太陽が沈みかける頃で、望楼の役人が、六つの\ruby{鼓}{こ}を叩くのが合図だった。\\


　だからこの辺の住民は、そこの門のことを、六\ruby{鼓門}{こもん}と呼んでいたが、今日もまた、赤い夕陽が鉄の\ruby{扉}{と}にさしかける頃、望楼の鼓が、もう二つ三つ四つ……と鳴りかけていた。\\


「待って下さい。待って下さいっ」\\


　彼方から驢を飛ばしてきたひとりの旅人は、危うく一足ちがいで、一夜を城門の外に明かさなければならない間ぎわだったので、手をあげながら馳けてきた。\\


　最後の鼓の一つが鳴ろうとした時、からくも旅人は、城門へ着いて、\\


「おねがい致します。通行をおゆるし下さいまし」\\


　と、驢をそこで降りて、型のごとく\ruby{関門}{かんもん}調べを受けた。\\


　役人は、旅人の顔を見ると、「やあ、お前は\ruby{劉備}{りゅうび}じゃないか」と、いった。\\


　劉備は、ここ楼桑村の住民なので、誰とも顔見知りだった。\\


「そうです。今、旅先から帰って参ったところです」\\


「お前なら、顔が\ruby{手形}{てがた}だ、何も調べはいらないが、いったい何処へ行ったのだ。こんどの旅はまた、ばかに長かったじゃないか」\\


「はい、いつもの商用ですが、なにぶん、どこへ行っても近頃は、\ruby{黄匪}{こうひ}の横行で、思うように\ruby{商}{あきな}いもできなかったものですから」\\


「そうだろう。関門を通る旅人も、毎日へるばかりだ。さあ、早く通れ」\\


「ありがとう存じます」\\


　再び驢にのりかけると、\\


「そうそう、お前の母親だろう、よく関門まで来ては、きょうもまだ息子は帰りませぬか、今日も劉備は通りませぬかと、夕方になると訊ねにきたのが、この頃すがたが見えぬと思ったらわずらって寝ているのだぞ。はやく帰って顔を見せてやるがよい」\\


「えっ。では母は、留守中に、病気で寝ておりますか」\\


　劉備はにわかに胸さわぎを覚え、驢を急がせて、関門から城内へ馳けた。\\


　久しく見ない町の暮色にも、眼もくれないで彼は驢を家路へ向けた。道幅の狭い、そして短い宿場町はすぐとぎれて、道はふたたび悠長な田園へかかる。\\


　ゆるい小川がある。水田がある。秋なのでもう村の人々は刈入れにかかっていた。そして所々に見える農家のほうへと、田の人影も水牛の影も戻って行く。\\


「ああ、わが家が見える」\\


　劉備は、驢の上から手をかざした。\ruby{舂}{うすず}く\ruby{陽}{ひ}のなかに黒くぽつんと見える一つの屋根と、そして遠方から見ると、まるで大きな\ruby{車蓋}{しゃがい}のように見える桑の木。劉備の生れた家なのである。\\


「どんなに自分をお待ちなされておることやら。……思えば、わしは孝養を励むつもりで、実は不孝ばかり重ねているようなもの。母上、済みません」\\


　彼の心を知るか、驢も足を早めて、やがて\ruby{懐}{なつ}かしい桑の大樹の下までたどりついた。\\

\\

\subsubsection{二}

\\


　この桑の大木は、何百年を経たものか、村の古老でも知る者はない。\\


　\ruby{沓}{くつ}や\ruby{蓆}{むしろ}をつくる劉備の家――と訊けば、あああの桑の樹の家さと指さすほど、それは村の何処からでも見えた。\\


　古老がいうには、\\


「楼桑村という地名も、この桑の木が茂る時は、まるで緑の楼台のように見えるから、この樹から起った村の名かもしれない」とのことであった。\\


　それはともかく、劉備は今、ようやく帰り着いたわが家の裏に驢をつなぐとすぐ、\\


「おっ母さん、今帰りました。\ruby{玄徳}{げんとく}です。玄徳ですよ」\\


　と、広い家の中へ馳けこむようにはいって行った。\\


　旧家なので、家は大きいが、何一つあるではなく、中庭は、\ruby{沓}{くつ}を\ruby{編}{あ}んだり\ruby{蓆}{むしろ}を織る仕事場になっており、そこも劉備の留守中は職人も通っていないので、荒れたままになっていた。\\


「オヤ、どうしたのだろう。\ruby{燈火}{あかり}もついていないじゃないか」\\


　彼は召使いの老婆と、\ruby{下僕}{しもべ}の名を呼びたてた。\\


　ふたりとも、返辞もない。\\


　劉備は、舌打ちしながら、\\


「おっ母さん」\\


　母の部屋をたたいた。\\


　\ruby{阿備}{あび}か――と飛びつくように迎えてくれるであろうと思っていた母の姿も見えなかった。いや母の部屋だけにたった一つあった\ruby{箪笥}{たんす}も寝台も見えなかった。\\


「や？　……どうしたのだろう」\\


　茫然、胸さわぎを抱いて、たたずんでいると、暗い中庭のほうで、かたん、かたん――と\ruby{蓆}{むしろ}を織る音がするのであった。\\


「おや」\\


　廊へ出てみると、そこの仕事場にだけ、うす暗い灯影がたった一つかかげてあった。その灯の下に、白髪の母の影が後ろ向きに腰かけていた。ただ一人で、星の下に、蓆を織っているのだった。\\


　母は、彼が帰ってきたのも気がついていないらしかった。劉備がすがりつかんばかり馳け寄って、\\


「今、帰りました」\\


　と顔を見せると、母は、びっくりしたように起ってよろめきながら、\\


「おお、阿備か、阿備か」\\


　乳呑み児でも抱きしめるようにして、何を問うよりも先に、うれし涙を眼にいっぱいためたまま、しばしは、母は子の肌を、子は母親のふところを、相擁して\ruby{温}{ぬく}め合うのみであった。\\


「城門の番人に、おまえの母親は病気らしいぞといわれて、気もそぞろに帰ってきたのですが、おっ母さん、どうしてこんな夜露の冷える外で、今頃、\ruby{蓆}{むしろ}など織っていらっしゃるのですか」\\


「病気？　……ああ城門の番人さんは、そういったかも知れないね。毎日のように関門までおまえの帰りを見に行っていたわたしが、この十日ばかりは行かないでいたから」\\


「では、ご病気ではないんですか」\\


「病気などはしていられないよ、おまえ」と、母はいった。\\


「寝台も箪笥もありませんが……」\\


　劉備が問うと、\\


「税吏が来て、持って行ってしまった。\ruby{黄匪}{こうひ}を討伐するために、年々軍費がかさむというので、ことしはとほうもなく税が上がり、おまえが用意しておいただけでは間に合わない程になったんだよ」\\


「婆やが見えませんが、婆やはどうしましたか」\\


「息子が、黄匪の仲間にはいっているという疑いで、縛られて行った」\\


「若い下僕は」\\


「兵隊にとられて行ったよ」\\


「――ああ！　すみませんでしたおっ母さん」\\


　劉備は、母の足もとに、ひれ伏して\ruby{詫}{わ}びた。\\

\\

\subsubsection{三}

\\


　詫びても詫びても詫び足らないほど、劉備は母に対して済まない心地であった。けれども母は、久しぶりに旅から帰ってきた我が子が、そんな自責に泣きかなしむことは、かえって\ruby{不愍}{ふびん}やら気の毒やらで、自分の胸も\ruby{傷}{いた}むらしく、\\


「阿備や、泣いておくれでない。何を詫びることがあるものかね。お前のせいではありはしない。世の中が悪いのだよ。……どれ\ruby{粟}{あわ}でも煮て、久しぶりに、ふたりして晩のお膳をかこもうね。さだめし疲れているだろうに、今、湯を沸かしてあげるから、汗でも拭いたがよい」\\


　と、\ruby{蓆機}{むしろばた}の前から立ちかけた。\\


　子の機嫌をとって、子の罪を責めない母のあまりなやさしさに、劉備はなおさら大愛の姿にぬかずいて、\\


「もったいない。私が戻りましたからには、そんなことは、玄徳がいたします。もうご不自由はさせません」\\


「いいえお前はまた、あしたから働いておくれ。稼ぎ人だからね、婆やも下僕もいなくなったのだから、台所のことぐらいは、わたしがしましょうよ」\\


「留守中、そんなことがあろうとは、少しも知らず、つい旅先で長くなって、思わぬご苦労をかけました。さあ、こんな大きな息子がいるんですから、おっ母さんは部屋へはいって、安楽に寝台で寝ていて下さい」と、いって劉備はむりに母の手を\ruby{誘}{いざな}ったが、考えて見ると、その寝台も税吏に税の代わりに持って行かれてしまったので、母の部屋には、身を横たえる物もなかった。\\


　いや、寝台や\ruby{箪笥}{たんす}だけではない。それから彼が\ruby{灯}{あか}りを持って、台所へ行って見ると、鍋もなかった。四、五羽の鶏と一匹の牛もいたのであるが、そうした家畜類まで、すべて領主の軍需と税に徴発されて、目ぼしい物は何も残っていなかった。\\


「こんなにまで、領主の軍費も詰まってきたのか」\\


　劉備は、身の生活を考えるよりも、もっと大きな意味で、\ruby{暗澹}{あんたん}となった。\\


　そして直ぐ、\\


「これも、黄匪の害の一つのあらわれだ。ああどうなるのだろう？」\\


　世の行く末を思いやると、彼はいよいよ暗い心にとざされた。\\


　物置をあけて、彼は\ruby{夕餉}{ゆうげ}にする\ruby{粟}{あわ}や豆の俵を見まわした。驚いたことには、多少その中に蓄えておいた穀物も\ruby{干}{ほ}し肉も、天井につるしておいた\ruby{乾菜}{かんさい}まできれいに失くなっているのだった。――もう母に訊くまでもないことと、彼はまた、そこで\ruby{茫失}{ぼうしつ}していた。\\


　すると、むりに部屋へ入れて休ませておいた母が部屋の中で、何か小さい物音をさせていた。行って見ると、床板を上げて、土中の\ruby{瓶}{かめ}の中から、わずかな粟と食物を取出している。\\


「……ア。そんな所に」\\


　劉備の声に、彼女はふり向いて、浅ましい自分を笑うように、\\


「すこし隠しておいたのだよ。生きてゆくだけの物はないと困るからね」\\


「…………」\\


　世の中は急転しているのだ。これはもうただ事ではない。何億の人間が、生きながら餓鬼となりかけているのだ。反対に、一部の黄巾賊が、その血をすすり肉をくらって、不当な\ruby{富貴}{ふっき}と\ruby{悪辣}{あくらつ}な\ruby{栄華}{えいが}をほしいままにしているのだ。\\


「阿備や……。灯りを持っておいで、粟が煮えたよ。何もないけれど、二人して喰べれば、おいしかろう」\\


　やがて、老いたる母は、貧しい卓から子を呼んでいた。\\

\\

\subsubsection{四}

\\


　貧しいながら、母子は久しぶりで共にする晩の食事を楽しんだ。\\


「おっ母さん、あしたの朝は、きっと歓んでいただけると思います。こんどの旅から、私はすばらしいお\ruby{土産}{みやげ}を持って帰ってきましたから」\\


「お土産を」\\


「ええ。おっ母さんの、大好きな物です」\\


「ま。何だろうね？」\\


「生きているうちに、もいちど味わってみたいと、いつか仰っしゃったことがありましたろう。それですよ」\\


　母を楽しませるために、劉備も、それが\ruby{洛陽}{らくよう}の銘茶であるということを、しばらく明かさなかった。\\


　母は、わが子のその気持だけでも、もう眼を細くして歓んでいるのである。\ruby{焦}{じ}らされていると知りながら、\\


「織物かえ」と訊いた。\\


「いいえ。今もいったとおり、味わう物ですよ」\\


「じゃあ、食べ物？」\\


「――に、近いものです」\\


「何じゃろ。わからないよ、阿備や。わたしにそんな\ruby{好物}{こうぶつ}があるかしら」\\


「望んでも、望めない物と、\ruby{諦}{あきら}めの中に忘れておしまいになったんでしょう。一生に一度は、とおっ母さんが何年か前に云ったことがあるので、私も、一生に一度はと、おっ母さんにその望みをかなえてあげたいと、今日まで願望に抱いておりました」\\


「まあ、そんなに長年、心にかけてかえ？　……なおさら、分らなくなってしもうたよ阿備。……いったいなんだねそれは」\\


「おっ母さん、実は、これですよ」\\


　\ruby{錫}{すず}の小さい\ruby{茶壺}{ちゃつぼ}を取出して、劉備は、卓の上に置いた。\\


「洛陽の銘茶です。……おっ母さんの大好きなお茶です。……あしたの朝は、うんと早起きしましょう。そしておっ母さんは、裏の桃園に\ruby{莚}{むしろ}をお敷きなさい。私は\ruby{驢}{ろ}に乗って、ここから四里ほど先の\ruby{鶏村}{けいそん}まで行くと、とてもいい清水の\ruby{湧}{わ}いている所がありますから、番人に頼んでひと\ruby{桶}{おけ}清水を\ruby{汲}{く}んできます」\\


「…………」\\


　母は眼をまるくしたまま錫の小壺を見つめて、物もいえなかった。ややしばらくしてから怖い物でもさわるように、そっと\ruby{掌}{て}に乗せて、壺の横に貼ってある\ruby{詩箋}{しせん}のような文字などを見ていた。そして大きな溜息をつきながら、眼を息子の顔へあげて、\\


「阿備や。……お前、いったいこれは、どうしたのだえ」\\


　声までひそめて訊ねるのだった。\\


　劉備は、母が疑いの余り案じてはならないと考えて、自分の気持や、それを手に入れたことなど、噛んでふくめるようにして話して聞かせた。民間ではほとんど手に入れがたい品にはちがいないが、自分が求めたのは、正当な手続きで\ruby{購}{あがな}ったのだから少しも懸念をする必要はありません――とつけ加えていった。\\


「ああ、お前は！　……なんてやさしい子だろう」\\


　母は、茶壺を置いて、わが子の劉備に\ruby{掌}{て}をあわせた。\\


　劉備は、あわてて、\\


「おっ母さん、滅相もない。そんなもったいない真似はよして下さい。ただ歓んでさえいただければ」\\


　と、手を取った。そうして相擁したまま、劉備は自分の気もちの\ruby{酬}{むく}いられたうれしさに泣き、母は子の孝心に感動の余り涙にくれていた。\\


　翌る朝――\\


　まだ夜も白まぬうちに起きて、劉備は驢の背に水桶を結いつけ、自分ものって、\ruby{鶏村}{けいそん}まで水を汲みに行った。\\

\\

\subsubsection{五}

\\


　もちろん劉備が出かけた頃、彼の母も\ruby{夙}{と}く起きていた。\\


　母はその間に、\ruby{竈}{かまど}の下に\ruby{豆莢}{まめ}\jruby{がら}{・・・・・・}を\ruby{焚}{た}いて、朝の\ruby{炊}{かし}ぎをしておき、やがて家の裏のほうへ出て行った。\\


　桑の大木の下を通って、裏へ出ると、牛のいない牛小屋があり、鶏のいない\ruby{鶏}{とり}小屋があり、何もかも荒れ果てて、いちめんに秋草がのびている。\\


　だが、そこから百歩ほど歩くと這うような姿をした果樹が、背を並べて、何千坪かいちめんに揃っていた。それはみんな桃の樹であった。秋は葉も落ちて淋しいが、春の花のさかりには、この先の\ruby{蟠桃河}{ばんとうが}が落花で紅くなるほどだったし、桃の実は\ruby{市}{まち}に売り出して、村の家何軒かで分け合って、それが一年の生計の重要なものになった。\\


「……オオ」\\


　彼女は、ひとりでに出たような声をもらした。桃園の彼方から陽が昇りかけたのだ。金色の日輪は、密雲を噛み破るように、端だけ見えていた。今や何か尊いものがこの世に生れかけているような感銘を彼女もうけた。\\


「…………」\\


　彼女は、ひざまずいて、三礼を\ruby{施}{ほどこ}した。子どものことを\ruby{祷}{いの}っているらしかった。\\


　それから、\ruby{箒}{ほうき}を持った。\\


　たくさんの落葉がちらかっている。桃園は村の共有なので、日ごろ誰も掃除などはしない。彼女も一部を掃いただけであった。\\


　新しい\ruby{莚}{むしろ}をそこへ敷いた。そして一箇の\ruby{土炉}{どろ}と茶碗など運んだ。彼女はもともと\ruby{氏素性}{うじすじょう}の\ruby{賤}{いや}しくない人の娘であったし、劉家も元来正しい家柄なので、そういう品もどこかに何十年も使用せずにしまってあった。\\


　清掃した桃園に坐って、彼女は水を汲みに行った息子が、やがて鶏村から帰るのを、心静かに待っていた。\\


　桃園の梢の\ruby{湖}{うみ}を、秋の\ruby{小禽}{ことり}が来てさまざまな音いろを\ruby{転}{まろ}ばした。陽はうらうらと雲を越えて、朝霧はまだ紫ばんだまま大陸によどんでいた。\\


「わしは\ruby{倖}{しあわ}せ者よ」\\


　彼女は、この一朝の満足をもって、死んでもいいような気がした。いやいや、そうでないとも思う。独り強くそう思う。\\


「あの子の\ruby{将来}{ゆくすえ}を見とどけねば……」\\


　ふと彼方を見ると、その劉備の姿が近づいてきた。水を汲んで帰ってきたのである。驢にのって、驢の鞍に小さい桶を結いつけて。\\


「おお。おっ母さん」\\


　桃園の小道をぬって、劉備は間もなくそこへきた。そして水桶をおろした。\\


「鶏村の水は、とてもいい水ですね。さだめし、これで茶を煮たらおいしいでしょう」\\


「ま。ご苦労だったね。鶏村の水のことはよく聞いているけれど、あそこはとても恐い谷間だというじゃないか。後でわたしはそれを心配していたよ」\\


「なあに、道なんかいくら\ruby{嶮}{けわ}しくても何でもありませんがね、清水には水番がいまして、なかなかただはくれません。少しばかり金をやってもらって来ました」\\


「黄金の水、洛陽のお茶、それにお前の孝心。王侯の母に生れてもこんないい思いにはめぐり会えないだろうよ」\\


「おっ母さん、お茶はどこへ置きましたか」\\


「そうそう、私だけがいただいてはすまないと思い、ご先祖のお仏壇へ上げておいたが」\\


「そうですか、盗まれたらたいへんです。すぐ取って参りましょう」\\


　劉備は、家のほうへ馳けて、\ruby{宝珠}{ほうしゅ}を抱くように、茶壺を捧げてきた。\\


　母は、\ruby{土炉}{どろ}へ、火をおこしていた。その前にひざまずいて劉備が茶壺を差出すと、その時、何が母の眼に映ったのであろうか、母は手を出そうともしないで、劉備の身のまわりを改まった\ruby{眸}{ひとみ}でじっと見つめた。\\

\\

\subsubsection{六}

\\


　劉備は、母がにわかに改まって自分の身なりを見ているので、\\


「どうしたのですかおっ母さん」\\


　いぶかしげに訊いた。\\


　母は、いつになく厳粛な\ruby{容子}{ようす}をつくって、\\


「\ruby{阿備}{あび}」と、声まで、常とはちがって呼んだ。\\


「はい。何ですか」\\


「お前の\ruby{佩}{は}いている剣は、それは誰の剣ですか」\\


「わたくしのですが」\\


「嘘をおいい。旅に出る前の物とはちがっている。お前の剣は、お父さんから\ruby{遺物}{かたみ}にいただいた――ご先祖から伝わっている剣のはずです。それを、何処へやってしまったのです？」\\


「……はい」\\


「はいではありません。片時でも肌身から離してはなりませぬぞと、わしからもくれぐれいってあるはずです。どうしたのだえ、あの大事な剣は」\\


「実は、その……」\\


　劉備はさしうつ向いてしまった。\\


　母の顔は、いよいよ峻厳に変っていた。劉備が口ごもっていると、なお追及して、\\


「まさか手放してしまったのではあるまいね」と念を押した。\\


　劉備は、両手をつかえて、\\


「申しわけありません。実は旅から帰る途中、ある者に礼として与えてしまいましたので」\\


　いうと、母は、「えっ、人に与えてしまったッて。――ま！　あの剣を」と、顔いろを変えた。\\


　劉備はそこで、黄巾賊の一群につかまって、人質になったことや、茶壺も剣も奪り上げられてしまったことや、それからようやく救われて、賊の群れから脱出してきたが、再び追いつかれて黄匪の重囲に陥ち、すでに斬り死しようとした時、\ruby{卒}{そつ}の張飛という者が、一命を助けてくれたので、うれしさの余り、何か礼を与えようと思ったが、身に持っている物は、剣と茶壺しかなかったので、やむなく、剣を彼に与えたのです――とつぶさに話して、\\


「賊に捕まった時も、張卒に助けられた時も、その折はもう何も要らないという気持になっていたんです。……けれど、この銘茶だけは、生命がけでも持って帰って、おっ母さんに上げたいと思っていました。剣を手放したのは申しわけありませんが、そんなわけで、この銘茶を、生命から二番目の物として、持ち帰ったのでございます」\\


「…………」\\


「剣は、先祖伝来の物で、大事な物には違いありませんが、\ruby{沓}{くつ}や\ruby{蓆}{むしろ}をつくって\ruby{生活}{くら}しているあいだは、張卒から貰ったこれでも決して間にあわないこともありませんから……」\\


　母の惜しがる気持をなだめるつもりで彼がそういうと、何思ったか劉備の母は、\\


「ああ――わしは、お前のお父様に申しわけがない。亡き良人に顔向けがなりません。――わたしは、子の育て方を\ruby{過}{あやま}った！」と、\ruby{慟哭}{どうこく}して叫んだ。\\


「何を仰っしゃるんです。おっ母さん！　どうしてそんなことを」\\


　母の心を酌みかねて、劉備がおろおろというと、母はやにわに、眼の前にあった\ruby{錫}{すず}の小さい茶壺を取上げ、\\


「\ruby{阿備}{あび}、おいで！」と、きつい顔して、彼の腕を片手で引っ張った。\\


「何処へです。おっ母さん。……ど、どこへいらっしゃろうというんですか」\\


「…………」\\


　彼の母は、答えもせず、劉備の腕くびを固くつかんだまま、桃園の果てへ馳けだして行った。そしてそこの\ruby{蟠桃河}{ばんとうが}の岸までくると、持っていた\ruby{錫}{すず}の茶壺を、河の中ほど目がけてほうり捨ててしまった。\\

\\

\subsubsection{七}

\\


「あッ。何で？」\\


　びっくりした劉備は、われを忘れて、母の\ruby{手頸}{てくび}をとらえたが、母の手から投げられた茶の壺は、小さいしぶきを見せて、もう河の底に沈んでいた。\\


「おっ母さん！　……おっ母さん！　……一体、なにがお気にさわったのですか。なんで折角の茶を、河へ捨てておしまいになったんですか」\\


　劉備の声は、ふるえていた。母によろこばれたいばかりに、百難の中を、生命がけで持ってきた茶であった。\\


　母は、歓びの余りに、気が\ruby{狂}{ふ}れたのではあるまいか？\\


「……何をいうのです。\ruby{譟}{さわ}がしい！」\\


　母は、劉備の手を払った。\\


　そして\ruby{亡父}{ちち}のような顔をした。\\


「…………」\\


　劉備は、きびしい母の眉に、思わず後ろへ退がった。生れてから初めて、母にも怖い姿があることを知った。\\


「阿備。お坐りなさい」\\


「……はい」\\


「お前が、わしを歓ばせるつもりで、はるばる苦労して持っておいでた茶を、河へ捨ててしもうた母の心がわかりますか」\\


「……わかりません。おっ母さん、\ruby{玄徳}{げんとく}は愚鈍です。どこが悪い、なにが気にいらぬと、叱って下さい。仰っしゃって下さい」\\


「いいえ！」\\


　母は、つよく頭を振り、\\


「勘ちがいをおしでない。母は自分の気ままから叱るのではありません。――大事な剣を人手に渡すようなお前を育ててきたことを、わたしは母として、ご先祖にも、死んだお父さんにも、済まなく思うたからです」\\


「私が悪うございました」\\


「お黙りっ！　……そんな簡単に聞かれては、母の\ruby{叱言}{こごと}がおまえに分っているとはいえません。――私が怒っているのは、お前の心根がいつのまにやら\ruby{萎}{な}えしぼんで、楼桑村の水呑百姓になりきってしもうたかと――それが口惜しいのです。残念でならないのです」\\


　母は、子を叱るために励ましているわれとわが声に泣いてしまって、\ruby{袍}{ほう}の袖を、老いの眼に当てた。\\


「……お忘れかえ、阿備。おまえのお父様も、お\ruby{祖父}{じい}様も、おまえのように\ruby{沓}{くつ}を作り\ruby{蓆}{むしろ}を織り、土民の中に埋もれたままお果てなされてはいるけれど、もっともっと先のご先祖をたずねれば、漢の\ruby{中山靖王}{ちゅうざんせいおう}\ruby{劉勝}{りゅうしょう}の正しい血すじなのですよ。おまえはまぎれもなく\ruby{景帝}{けいてい}の\ruby{玄孫}{げんそん}なのです。この支那をひとたびは統一した帝王の血がおまえの体にながれているのです。あの剣は、その\ruby{印綬}{いんじゅ}というてもよい物です」\\


「…………」\\


「だが、こんなことは、めったに口に出すことではない。なぜならば、今の\ruby{後漢}{ごかん}の帝室は、わたし達のご先祖を亡ぼして立った帝王だからです。景帝の玄孫とわかれば、とうに私たちの家すじは\ruby{断}{た}ちきられておるでしょう。……だからというて、お前までが、土民になりきってしまってよいものか」\\


「…………」\\


「わたしは、そんな教育を、お前にした覚えはない。\ruby{揺籃}{ゆりかご}に入れて、子守唄をうとうて聞かせた頃から――また、この母が膝に抱いて眠らせた頃から――おまえの耳へ母はご先祖のお心を血の中へおしえこんだつもりです。――時の来ぬうちはぜひもないが、時節が来たら、世のために、また、漢の正統を再興するために、剣をとって、\ruby{草廬}{そうろ}から起たねばならぬぞと」\\


「……はい」\\


「阿備。――その剣を人手に渡して、そなたは、生涯、蓆を織っている気か。剣よりも茶を大事にお思いか。……そんな\ruby{性根}{しょうね}の子が求めてきた茶などを、歓んで飲む母とお思いか。……わたしは腹が立つ。わたしはそれが悲しい」\\


　と、母は\ruby{慟哭}{どうこく}しながら、劉備の襟をつかまえて、\ruby{嬰児}{あかご}を\ruby{懲}{こ}らすように\ruby{折檻}{せっかん}した。\\

\\

\subsubsection{八}

\\


　母に打ちすえられたまま、劉備は身うごきもしなかった。\\


　\ruby{打々}{ちょうちょう}と、母が打つたびに、母の大きな愛が、骨身にしみ、さんさんと涙がとまらなかった。\\


「すみません」\\


　母の手をいたわるように、劉備はやがて、打つ手を抑えて、自分の額に、押しいただいた。\\


「わたくしの考え違いでございました。まったく玄徳の愚かがいたした落度でございます。仰っしゃる通り、玄徳もいつか、土民の中に貧窮しているため、心まで土民になりかけておりました」\\


「分かりましたか。阿備、そこへ気がつきましたか」\\


「ご\ruby{打擲}{ちょうちゃく}をうけて、幼少のご訓言が、骨身からよび起されて参りました。――大事な剣を失いましたことは、ご先祖へも、申しわけありませんが……ご安心下さいお母さん……玄徳の魂はまだ\ruby{此身}{ここ}にございます」\\


　――するとそれまで、老いの手が\ruby{痺}{しび}れるほど子を打っていた母の手は、やにわに阿備のからだをひしと抱きしめて、\\


「おお！　阿備や！　……ではお前にも、一生土民で朽ち果てまいと思う気持はおありかえ。まだ忘れないで、わたしの言葉を、魂のなかにお持ちかえ」\\


「なんで忘れましょう。わたくしが忘れても景帝の玄孫であるこの血液が忘れるわけはありません」\\


「よう云いなすった。……阿備よ。それを聞いて母は安心しました。ゆるしておくれ、……ゆるしておくれよ」\\


「何をなさるんです。わが子へ手をついたりして、もったいない」\\


「いいえ。心まで落ちぶれ果てたかと、悲しみと怒りの余り、お前を打擲したりして」\\


「ご恩です。大愛です。今のご打擲は、わたくしにとって、真の勇気をふるいたたせる\ruby{神軍}{しんぐん}の\ruby{鼓}{つづみ}でございました。\ruby{仏陀}{ぶっだ}の\ruby{杖}{つえ}でございました。――もしきょうのお怒りを見せて下さらなければ、玄徳は何を胸に考えていても、おっ母さんが世にあるうちはと、\ruby{卑怯}{ひきょう}な土民をよそおっていたかも知れません。いいえそのうちについ年月を過して、ほんとの土民になって朽ちてしまったかもしれません」\\


「――ではお前は、何を思っても、この母が心配するのを怖れて、母が生きているうちはただ無事に暮していることばかり願っていたのだね。……ああ、そう聞けば、なおさらわたしのほうが済まない気がします」\\


「もう私も、肚がきまりました。――でなくても、今度の旅で、諸州の乱れやら、黄匪の惨害やら、地上の民の苦しみを、眼の痛むほど見てきたのです。おっ母さん、玄徳が今の世に生れ出たのは、天上の諸帝から、何か使命をうけて世に出たような気がされます」\\


　彼が、真実の心を吐くと彼の母は、天地に黙祷をささげて、いつまでも、両の\ruby{腕}{かいな}の中に額をうめていた。\\


　しかし、この日の朝のことは、どこまでも、\ruby{母子}{おやこ}ふたりだけの\ruby{秘}{ひそ}かごとだった。\\


　劉備の家には、相変らず\ruby{蓆機}{むしろばた}を織る音が、何事もなげに、毎日、外へもれていた。\\


　土民の手あらの者が、職人として雇われてきて、日ごとに中庭の作業場で、\ruby{沓}{くつ}を編み、蓆を荷造りして、それが溜ると、城内の\ruby{市}{いち}へ持って行って、穀物や布や、母の持薬などと交易してきた。\\


　変ったことといえば、それくらいなもので、家の\ruby{東南}{たつみ}にある高さ五丈余の桑の大樹に、春は\ruby{禽}{とり}が歌い、秋は落葉して、いつかここに三、四年の星霜は過ぎた。\\


　すると、浅春の\ruby{一日}{あるひ}。\\


　白い\ruby{山羊}{やぎ}の背に、二箇の\ruby{酒瓶}{さかがめ}を乗せて、それをひいてきた旅の老人が、桑の下に立って、独りで何やら感嘆していた。\\

\\

\subsubsection{九}

\\


　誰か、のっそりと、無断で家の横から中庭へはいってきた。\\


　劉備は、母と二人で、蓆を織っていた。無断といっても、土塀は崩れたままだし、門はないし、通り抜けられても、\ruby{咎}{とが}めるわけにもゆかないほどな家ではあったが――\\


「……おや？」\\


　振向いた\ruby{母子}{おやこ}は目をみはった。そこに立った旅の老人よりも、酒瓶を背にのせている山羊の毛の雪白な美しさに、すぐ気をとられたのである。\\


「ご精が出るのう」\\


　老人は、なれなれしい。\\


　\ruby{蓆機}{むしろばた}のそばに腰をおろし、なにか話しかけたい顔だった。\\


「お爺さん、どこから来なすったね。たいそう毛のいい山羊だな」\\


　いつまでも黙っているので、かえって劉備から口を切ってやると、老人はさもさも何か感じたように、独りで首を振りながら云った。\\


「息子さんかの。このお方は」\\


「はい」と、母が答えると、\\


「よい子を生みなすったな、わしの山羊も自慢だが、この息子にはかなわない」\\


「お爺さんは、この山羊をひいて、城内の\ruby{市}{いち}へ売りに来なすったのかね」\\


「なあに、この山羊は、売れない。誰にだって、売れないさ。わしの息子だものな。わしの売物は酒じゃよ。だが道中で\ruby{悪漢}{わるもの}に\ruby{脅}{おど}されて、酒は呑まれてしもうたから、瓶は二つとも空っぽじゃ。何もない。はははは」\\


「では、せっかく遠くから来て、おかねにも換えられずに帰るんですか」\\


「帰ろうと思って、ここまで来たら、偉い物を見たよわしは」\\


「なんですか」\\


「お宅の桑の樹さ」\\


「ああ、あれですか」\\


「今まで、何千人、いや何万人となく、村を通る人々が、あの樹を見たろうが、誰もなんともいった者はいないかね」\\


「べつに……」\\


「そうかなあ」\\


「珍しい樹だ、桑でこんな大木はないとは、誰もみないいますが」\\


「じゃあ、わしが告げよう。あの樹は、\ruby{霊木}{れいぼく}じゃ。この家から必ず貴人が生れる。\ruby{重々}{ちょうちょう}、\ruby{車蓋}{しゃがい}のような枝が皆、そういってわしへ囁いた。……遠くない、この春。桑の葉が青々とつく頃になると、いい友達が訪ねてくるよ。\ruby{蛟龍}{こうりょう}が雲をえたように、それからここの\ruby{主}{あるじ}はおそろしく身の上が変ってくる」\\


「お爺さんは、易者かね」\\


「わしは、\ruby{魯}{ろ}の\ruby{李定}{りてい}という者さ。というて年中\ruby{飄々}{ひょうひょう}としておるから、故郷にいたためしはない。山羊をひっぱって、酒に酔うて、時々、市へ行くので、皆が\ruby{羊仙}{ようせん}といったりする」\\


「羊仙さま。じゃあ世間の人は、あなたを仙人と思っているので？」\\


「はははは。迷惑なはなしさ。何しろきょうはよろこばしい人とはなし、珍しい霊木を見た。この子のおっ母さん」\\


「はい」\\


「この山羊を、お祝いに献上しよう」\\


「えっ？」\\


「おそらく、この子は、自分の誕生日も、祝われたことはあるまい。だが、今度は祝ってやんなさい。この\ruby{瓶}{かめ}に酒を買い、この山羊を\ruby{屠}{ほふ}って、血は神壇に捧げ、肉は\ruby{羹}{あつもの}に煮て」\\


　初めは、戯れであろうと、半ば笑いながら聞いていたところ、羊仙はほんとに山羊を置いて、立ち去ってしまった。\\


　驚いて、桑の下まで馳けだし、往来を見まわしたが、もう姿は見えなかった。\\

\\

\subsection{\ruby{橋畔風談}{きょうはんふうだん}}

\\

\subsubsection{一}

\\


　\ruby{蟠桃河}{ばんとうが}の水は紅くなった。両岸の桃園は\ruby{紅霞}{こうか}をひき、夜は眉のような月が香った。\\


　けれど、その水にも、詩を\ruby{詠}{よ}む人を乗せた一艘の舟もないし、杖をひいて逍遥する\ruby{雅人}{がじん}の影もなかった。\\


「おっ母さん、行ってきますよ」\\


「ああ、行っておいで」\\


「なにか城内からおいしい物でも買ってきましょうかね」\\


　劉備は、家を出た。\\


　\ruby{沓}{くつ}や\ruby{蓆}{むしろ}をだいぶ納めてある城内の問屋へ行って、\ruby{価}{あたい}を取ってくる日だった。\\


　\ruby{午}{ひる}から出ても、用達をすまして陽のあるうちに、らくに帰れる道のりなので、劉備は\ruby{驢}{ろ}にものらなかった。\\


　いつか羊仙のおいて行った山羊がよく馴れて、劉備の後についてくるのを、母が後ろで呼び返していた。\\


　城内は、\ruby{埃}{ほこり}ッぽい。\\


　雨が久しくなかったので、\ruby{沓}{くつ}の裏がぽくぽくする。劉備は、問屋から銭を受け取って、\ruby{脂}{あぶら}光りのしている\ruby{市}{いち}の軒なみを見て歩いていた。\\


　\ruby{蓮根}{れんこん}の菓子があった。劉備はそれを少し買い求めた。――けれど少し歩いてから、\\


「蓮根は、母の持病に悪いのじゃないか」と、取換えに戻ろうかと迷っていた。\\


　がやがやと沢山な人が辻に集まっている。いつもそこは、\ruby{野鴨}{のがも}の丸揚げや餅など売っている場所なので、その混雑かと思うていたが、ふと見ると、大勢の頭の上に、高々と、立札が見えている。\\


「何だろ？」\\


　彼も、好奇にかられて、人々のあいだから\ruby{高札}{こうさつ}を仰いだ。\\


　見ると――\\

\ruby{遍}{あまね}く天下に義勇の士を\ruby{募}{つの}る\\



　という布告の文であった。\\

\ruby{黄巾}{こうきん}の\ruby{匪}{ひ}、諸州に蜂起してより、年々の害、鬼畜の毒、惨として\ruby{蒼生}{そうせい}に\ruby{青田}{せいでん}なし。\\


今にして、鬼賊を\ruby{誅}{ちゅう}せずんば、天下知るべきのみ。\\

\ruby{太守}{たいしゅ}\ruby{劉焉}{りゅうえん}、遂に、子民の\ruby{泣哭}{きゅうこく}に奮って討伐の天鼓を鳴らさんとす。故に、隠れたる\ruby{草廬}{そうろ}の君子、野に\ruby{潜}{ひそ}むの義人、旗下に参ぜよ。\\


欣然、各子の武勇に依って、府に迎えん。\\

\ruby{{\fontspec{jigmo}\char"6DBF}郡}{たくぐん}\ruby{校尉}{こうい}\ruby{鄒靖}{すうせい}


「なんだね、これは」\\


「兵隊を\ruby{募}{つの}っているのさ」\\


「ああ、兵隊か」\\


「どうだ、志願して行って、ひと働きしては」\\


「おれなどはだめだ。武勇もなにもない。ほかの能もないし」\\


「誰だって、そう能のある者ばかり集まるものか。こう書かなくては、勇ましくないからだよ」\\


「なるほど」\\


「\ruby{憎}{にく}い\ruby{黄匪}{こうひ}めを討つんだ、槍の持ち方が分らないうちは、馬の\ruby{飼糧}{かいば}を刈っても\ruby{軍}{いくさ}の手伝いになる。おれは行く」\\


　ひとりがつぶやいて去ると、そのつぶやきに決心を固めたように、二人去り、三人去り、皆、城門の役所のほうへ力のある足で急いで行った。\\


「…………」\\


　劉備は、時勢の\ruby{跫音}{あしおと}を聞いた。民心のおもむく\ruby{潮}{うしお}を見た。\\


　――が。蓮根の菓子を手に持ったまま、いつまでも、考えていた。誰もいなくなるまで、高札と睨み合って考えていた。\\


「……ああ」\\


　気がついて、間がわるそうに、そこから離れかけた。すると、誰か、\ruby{楊柳}{ようりゅう}のうしろから、\\


「若人。待ち給え」\\


　と、呼んだ者があった。\\

\\

\subsubsection{二}

\\


　さっきから楊柳の下に腰かけて、\ruby{路傍}{みちばた}の酒売りを相手に、声高に話していた男のあったことは、劉備も知っていた。\\


　自分の容子を、横目ででも見ていたのだろうか、二、三歩、高札から足を退けると、\\


「貴公、それを読んだか」\\


　片手に、\ruby{酒杯}{さかずき}を持ち、片手に剣の\ruby{把}{つか}を握って不意に起ってきたのである。\\


　楊柳の幹より大きな肩幅を、後ろ向きに見ていただけであったが、立上がったのを見ると、実に見上げるばかりの偉丈夫であった。突然、山が立ったように見えた。\\


「……私ですか」\\


　劉備はさらに改めて、その人を見なおした。\\


「うむ。貴公よりほかに、もう誰もいないじゃないか」\\


　\ruby{黒漆}{こくしつ}の\ruby{髯}{ひげ}の中で、\ruby{牡丹}{ぼたん}のような口を開いて笑った。\\


　声も年頃も、劉備と幾つも違うまいと思われたが、偉丈夫は、髪から\ruby{腮}{あご}まで、隙間もないように艶々しい髯をたくわえていた。\\


「――読みました」\\


　劉備の答えは\ruby{寡言}{かごん}だった。\\


「どう読んだな、貴公は」と、彼の問いは深刻で、その眼は、\ruby{烱々}{けいけい}として鋭い。\\


「さあ？」\\


「まだ考えておるのか。あんなに長い間、高札と睨み合っていながら」\\


「ここで語るのを好みません」\\


「おもしろい」\\


　偉丈夫は、酒売りへ、銭と\ruby{酒杯}{さかずき}を渡して、ずかずかと、劉備のそばへ寄ってきた。そして劉備の口真似しながら、\\


「ここで語るのを好みません……いや愉快だ。その言葉に、おれは真実を聴く。さ、何処かへ行こう」\\


　劉備は困ったが、「とにかく歩きましょう。ここは人目の多い市ですから」\\


「よし歩こう」\\


　偉丈夫は、\ruby{闊歩}{かっぽ}した。劉備は並行してゆくのに骨が折れた。\\


「あの\ruby{虹橋}{こうきょう}の辺はどうだ」\\


「よいでしょう」\\


　偉丈夫の指さすところは町はずれの楊柳の多い池のほとりだった。虹をかけたような石橋がある。そこから先は\ruby{廃苑}{はいえん}であった。何とかいう学者が池をほって、聖賢の学校を建てたが、時勢は聖賢の道と逆行するばかりで、真面目に通ってくる生徒はなかった。\\


　学者は、それでも根気よく、石橋に立って道を説いたが、市の住民や童は、（気狂いだ）と、耳もかさない。それのみか、\ruby{小賢}{こざか}しい奴だと、石を投げる者もあったりした。\\


　学者は、いつのまにか、ほんとの狂人になってしまったとみえ、ついには、あらぬ事を絶叫して、学苑の中をさまよっていたが、そのうちに蓮池の中に、あわれ死体となって浮び上がった。\\


　そういう遺蹟であった。\\


「ここはいい。掛け給え」\\


　偉丈夫は、虹橋の\ruby{石欄}{せきらん}へ腰をかけ、劉備にもすすめた。\\


　劉備は、ここまで来る間に、偉丈夫の人物をほぼ観ていた。そして、（この人間は\ruby{偽}{にせ}ものでない）と思ったので、ここへ来た時は、彼もかなりな落着きと本気を示していた。\\


「時に、失礼ですが、尊名から先に承りたいものです。私はここからほど遠くない楼桑村の住人で、劉備玄徳という者ですが」\\


　すると偉丈夫は、いきなり劉備の肩を打って、「好漢。それはもう聞いておるじゃないか。この方の名だって、よくご承知のはずだが」といった。\\

\\

\subsubsection{三}

\\


「えっ？　……私も以前からご存じの方ですって」\\


「お忘れかな。ははは」\\


　偉丈夫は、肩をゆすぶって、\ruby{腮}{あご}の黒い\ruby{髯}{ひげ}をしごいた。\\


「――無理もない。頬の刀傷で、容貌も少し変った。それにここ三、四年はつぶさに浪人の辛酸をなめたからなあ。貴公とお目にかかった頃には、まだこの黒髯もたくわえてなかった時じゃ」\\


　そういわれても、劉備はまだ思い出せなかったが、ふと、偉丈夫の腰に\ruby{佩}{は}いている剣を見て、思わずあっと口をすべらせた。\\


「おお、恩人！　思い出しました。あなたは数年前、私が黄河から\ruby{{\fontspec{jigmo}\char"6DBF}}{たく}県のほうへ帰ってくる途中、\ruby{黄匪}{こうひ}に囲まれてすでに危うかった所を助けてくれた\ruby{鴻家}{こうけ}の浪士、\ruby{張飛翼徳}{ちょうひよくとく}と仰っしゃったお方ではありませんか」\\


「そうだ」\\


　張飛はいきなり腕をのばして、劉備の手を握りしめた。その手は鉄のようで、劉備の\ruby{掌}{て}を握ってなお、五指が余っていた。\\


「よく覚えていて下された。いかにもその折の張飛でござる。かくの如く、髯をたくわえ、容貌を変えているのも、以来、志を得ずに、世の裏に潜んでおるがためです。――で実は、貴公に分るかどうか試してみたわけで、最前からの無礼はどうかゆるされい。」\\


　偉丈夫に似あわず、礼には\ruby{篤}{あつ}かった。\\


　すると劉備は、より以上、\ruby{慇懃}{いんぎん}にいった。\\


「豪傑。失礼はむしろ私のほうこそ咎めらるべきです。恩人のあなたを見忘れるなどということは、たとえいかに当時とお変りになっているにせよ、相済まないことです。どうか、劉備の罪はおゆるし下さい」\\


「やあ、ご鄭重で恐れいる。ではまあ、お互いとしておこう」\\


「時に、豪傑。あなたは今、この県城の\ruby{市}{まち}に住んでおるのですか」\\


「いや、話せば長いことになるが、いつかも打明け申した通り、どうかして黄巾賊に奪われた主家の県城を取返さんものと、民間にかくれては兵を\ruby{興}{おこ}し、また、敗れては民間に隠れ、幾度も幾度も事を\ruby{謀}{はか}ったが、黄匪の勢力はさかんになるばかりで、近頃はもう矢も尽き刀も折れたという恰好です。……で先頃から、この\ruby{{\fontspec{jigmo}\char"6DBF}県}{たくけん}に流れてきて、山野の\ruby{猪}{いのこ}を狩って肉を\ruby{屠}{ほふ}り、それを市にひさいで露命をつないでおるような状態です。おわらい下さい。ここのところ、張飛も\ruby{尾羽}{おは}\ruby{打枯}{うちか}らした\ruby{態}{てい}たらくなので」\\


「そうですか。少しも知りませんでした。そんなことなら、なぜ楼桑村の私の家を訪ねてくれなかったのですか」\\


「いや、いつかは一度、お目にかかりに参る心ではいたが、その折には、ぜひ尊公に、うんと承知してもらいたいことがあるので――その準備がまだこっちにできていないからだ」\\


「この劉備に、お頼みとは、いったい何事ですか」\\


「劉君」\\


　張飛は、鏡のような眼をした。らんらんとそのなかに胸中の\ruby{炬火}{きょか}が燃えているのを劉備は認めた。\\


「尊公は今日、市で県城の布令を読まれたであろう」\\


「うむ。あの高札ですか」\\


「あれを見て、どう思われましたか。黄匪討伐の兵を募るという文を見て――」\\


「べつに、どうといって、なんの感じもありません」\\


「ない？」\\


　張飛は、斬りこむような語気でいった。明らかに、激怒の血を、顔にうごかしてである。\\


　けれど劉備は、\\


「はい。何も思いません。なぜなら、私には、ひとりの母がありますから。――従って、兵隊に出ようとは思いませんから」\\


　水のように冷静にいった。\\

\\

\subsubsection{四}

\\


　秋かぜが橋の下を吹く。\\


　虹橋の下には、\ruby{枯蓮}{かれはす}の葉がからから鳴っていた。\\


　びらっと、色羽の\ruby{征箭}{そや}が飛んだと見えたのは、水を離れた\ruby{翡翠}{かわせみ}だった。\\


「嘘だっ」\\


　張飛は、静かな話し相手へ、いきなり呶鳴って、腰かけていた橋の石欄から突っ立った。\\


「劉君。貴公は、本心を人に秘して、この張飛へも、深くつつんでおられるな。いや、そうだ。張飛をご信用なさらぬのだ」\\


「本心？　……私の本心は今いった通りです。なにを、あなたにつつむものか」\\


「しからば貴公は、今の天下を眺めて、なんの感じも抱かれないのか」\\


「黄匪の害は見ていますが、小さい\ruby{貧屋}{ひんおく}に、ひとりの母さえ養いかねている身には」\\


「人は知らず、張飛にそんなことを仰っしゃっても、張飛はあなたを、ただの土民と見ることはできぬ。打明けて下さい。張飛も武士です。他言は断じて致さぬ\ruby{漢}{おとこ}です」\\


「困りましたな」\\


「どうしても」\\


「お答えのしようがありません」\\


「ああ――」\\


　\ruby{憮然}{ぶぜん}として、張飛は、\ruby{黒漆}{こくしつ}の髯を秋かぜに吹かせていたが、何か、思い出したように、突然、\ruby{佩}{は}いている剣帯を解いて、\\


「お覚えがあるでしょう」と、\ruby{鞘}{さや}を握って、劉備の\ruby{面}{おもて}へ、横ざまに突きつけていった。\\


「これはいつか、貴公から礼にと手前へ賜わった剣です。また、私から所望した剣であった。――だが不肖は、いつか尊公に再び巡り合ったら、この品は、お手もとへ返そうと思っていた。なぜならば、これは張飛の如き\ruby{匹夫}{ひっぷ}が持つ剣ではないからだ」\\


「…………」\\


「血しぶく戦場で、――また、\ruby{戦}{いくさ}に敗れて落ち行く草枕の寝覚めに――幾たびとなく拙者はこの剣を抜き払ってみた。そして、そのたびに、拙者は剣の声を聞いた」\\


「…………」\\


「劉君、\ruby{其許}{そこもと}は聞いたことがあるか、この剣の声を！」\\


「…………」\\


「一\ruby{揮}{き}して、風を断てば、剣は\ruby{啾々}{しゅうしゅう}と泣くのだ。星\ruby{衝}{つ}いて、\ruby{剣把}{けんぱ}から\ruby{鋩子}{ぼうし}までを\ruby{俯仰}{ふぎょう}すれば、\ruby{朧夜}{おぼろよ}の雲とまがう光の\ruby{斑}{ふ}は、みな剣の涙として拙者には見える」\\


「…………」\\


「いや、剣は、剣を持つ者へ訴えていうのだ。いつまで、わが身を、\ruby{為}{な}すなく室中に閉じこめておくぞと。――劉備どの、嘘と思わば、その耳に、剣の声を聞かそうか、剣の涙を見せようか」\\


「……あっ」\\


　劉備も思わず石欄から腰を立てた。――止める間はなかった。張飛は、剣を払って、ぴゅっと、秋風を斬った。正しく、剣の声が走った。しかもその声は、劉備の\ruby{腸}{はらわた}を断つばかり胸をうった。\\


「君聞かずや！」\\


　張飛は、いいながら、またも一振り二振りと、虚空に剣光を描いて、\\


「何の声か。そも」と、叫んだ。\\


　そしてなおも、答えのない劉備を見ると、もどかしく思ったのか、橋の石欄へ片足を踏みかけて、枯蓮の池を望みながら独り云った。\\


「\ruby{可惜}{あたら}、治国愛民の宝剣も、いかにせん持つ人もなき\ruby{末世}{まっせ}とあってはぜひもない。霊あらば剣も\ruby{恕}{じょ}せ。\ruby{猪肉売}{いのこう}りの浪人の腰にあるよりは、むしろ池中に葬って――」\\


　あなや、剣は、虹橋の下に投げ捨てられようとした。劉備は驚いて、走り寄るなり彼の腕を支え、「豪傑待ち給え」と、叫んだ。\\

\\

\subsubsection{五}

\\


　張飛はもとより折角の名剣を泥池に捨ててしまうのは本意ではないから、止められたのを幸いに、\\


「何か？」と、わざと身を\ruby{退}{ひ}いて、劉備の言を待つもののように見まもった。\\


「まず、お待ちなさい」\\


　劉備は言葉しずかに、張飛の悲壮な顔いろをなだめて、\\


「真の勇者は\ruby{慷慨}{こうがい}せずといいます。また、大事は\ruby{蟻}{あり}の穴より漏るというたとえもある。ゆるゆるはなすとしましょう。しかし、\ruby{足下}{そっか}が偽ものでないことはよく認めました。偉丈夫の心事を一時でも疑った罪はゆるして下さい」\\


「おっ。……では」\\


「風にも耳、水にも眼、大事は路傍では語れません。けれど自分は何をつつもう、漢の\ruby{中山靖王}{ちゅうざんせいおう}\ruby{劉勝}{りゅうしょう}の\ruby{後胤}{こういん}で、景帝の玄孫にあたるものです。……なにをか好んで、\ruby{沓}{くつ}を作り\ruby{蓆}{むしろ}を織って、\ruby{黄荒}{こうこう}の\ruby{末季}{まっき}を心なしに見ておりましょうや」と、声は小さく\ruby{語韻}{ごいん}はささやく如くであったが、\ruby{凛}{りん}たるものをうちに\ruby{潜}{ひそ}めていい、そしてにこと笑ってみせた。\\


「豪傑。これ以上、もう多言は吐く必要はないでしょう。折を見てまた会いましょう。きょうは市へきた出先で、遅くなると母も案じますから――」\\


　張飛は獅子首を突きだして、噛みつきそうな眼をしたまま、いつまでも無言だった。これは感きわまった時にやるかれの癖なのである。それからやがて唸るような息を吐いて、大きな胸をそらしたと思うと、\\


「そうだったのか！　やはりこの張飛の眼には誤りはなかった！　いやいつか古塔の上から跳び降りて死んだかの老僧のいったことが、今思いあたる。……ウウム、あなたは景帝の\ruby{裔孫}{えいそん}だったのか。治乱興亡の長い星霜のあいだに、名門名族は\ruby{泡沫}{うたかた}のように消えてゆくが、血は一滴でも残されればどこかに伝わってゆく。ああ有難い。生きていたかいがあった。今月今日、張飛は会うべきお人に会った」\\


　独りしてそう\ruby{呻}{うめ}いていたかと思うと、彼はにわかに、石橋の石の上にひざまずき、剣を奉じて、劉備へいった。\\


「謹んで、剣は、尊手へおかえしします。これはもともと、やつがれなどの身に\ruby{佩}{は}くものではない。――が、ただしです。あなたはこの剣を受け取らるるや否や、この剣を佩くからには、この剣と共にある使命もあわせて佩かねばならぬが」\\


　劉備は、手を伸ばした。\\


　何か、おごそかな姿だった。\\


「\ruby{享}{う}けましょう」\\


　剣は、彼の手にかえった。\\


　張飛は、いく度も、拝姿の礼を、くり返して、\\


「では、そのうちに、きっと楼桑村へ、お訪ねして参るぞ」\\


「おお、いつでも」\\


　劉備は、今まで佩いていた剣と佩きかえて、前の物は、張飛へ戻した。それは張飛に救われた数年前に、取換えた物だったからである。\\


「日が暮れかけてきましたな。じゃあ、いずれまた」\\


　夕闇の中を、劉備は先に、足を早めて別れ去った。風にふかれて行く水色の服は汚れていたが、剣は眼に見える\ruby{黄昏}{たそがれ}の万象の中で、なによりも異彩を放って見えた。\\


「体に持っている気品というものは争えぬものだ。どこか貴公子の風がある」\\


　張飛は見送りながら、独り虹橋の上に立ち暮れていたが、やがてわれにかえった顔をして、\\


「そうだ、\ruby{雲長}{うんちょう}にも聞かせて、早く歓ばしてやろう」と、何処ともなく馳けだしたが、劉備とちがってこれはまた、一陣の風が黒い物となって飛んで行くようだった。\\

\\

\subsection{\ruby{童学草舎}{どうがくそうしゃ}}

\\

\subsubsection{一}

\\


　城壁の望楼で、今しがた、\ruby{鼓}{こ}が鳴った。\\


　\ruby{市}{いち}は宵の燈となった。\\


　張飛は一度、市の辻へ帰った。そして昼間ひろげていた\ruby{猪}{いのこ}の露店をしまい、猪の股や肉切り庖丁などを\ruby{苞}{つと}にくくって持つとまた馳けだした。\\


「やあ、遅かったか」\\


　城内の街から城外へ通じるそこの関門は、もう閉まっていた。\\


「おうい、開けてくれっ」\\


　張飛は、望楼を仰いで、駄々っ子のようにどなった。\\


　関門のかたわらの小さい兵舎から五、六人ぞろぞろ出てきた。とほうもない馬鹿者に訪れられたように、からかい半分に叱りとばした。\\


「こらっ。なにをわめいておるか。関門が閉まったからには、\ruby{霹靂}{へきれき}が墜ちても、開けることはできない。なんだ貴様は一体」\\


「毎日、城内の市へ、猪の肉を売りに出ておる者だが」\\


「なるほど、こやつは肉売りだ。なんで今頃、寝ぼけて関門へやってきたのか」\\


「用が遅れて、閉門の時刻までに、帰りそびれてしまったのだ。開けてくれ」\\


「正気か」\\


「酔うてはいない」\\


「ははは。こいつ酔っぱらっているに違いない。三べんまわってお辞儀をしろ」\\


「なに」\\


「三度ぐるりと廻って、俺たちを三拝したら通してやる」\\


「そんなことはできぬが、このとおりお辞儀はする。さあ、開けてくれ」\\


「帰れ帰れ。何百ぺん頭を下げても、通すわけにはゆかん。市の軒下へでも寝て、あした通れ」\\


「あした通っていいくらいなら頼みはせん。通さぬとあれば、汝らをふみつぶして、城壁を躍り越えてゆくがいいか」\\


「こいつが……」と、呆れて、\\


「いくら酒の上にいたせ、よいほどに引っ込まぬと、\ruby{素}{そ}ッ\ruby{首}{くび}を\ruby{刎}{は}ね落すぞ」\\


「では、どうしても、通さぬというか。おれに頭を下げさせておきながら」\\


　張飛は、そこらを見廻した。酔いどれとは思いながら、雲つくような\ruby{巨漢}{おおおとこ}だし、無気味な眼の光にかまわずにいると、ずかずかと歩みだして、城壁の下に立ち、役人以外は登ることを厳禁している\ruby{鉄梯子}{かなばしご}へ片足をかけた。\\


「こらっ。どこへ行く」\\


　ひとりは、張飛の腰の\ruby{紐帯}{ちゅうたい}をつかんだ。他の関門兵は、槍をそろえて向けた。\\


　張飛は、髯の中から、白い歯を見せて、人なつこい笑い方をした。\\


「いいじゃないか。\ruby{野暮}{やぼ}をいわんでも……」\\


　そしてたずさえている猪の肉の\ruby{片股}{かたもも}と、肉切り庖丁とを、彼らの目のまえに突き出した。\\


「これをやろう。貴公らの身分では、めったに肉も喰らえまい。これで寝酒でもやったほうが、俺になぐり殺されるより遥かに\jruby{まし}{・・・・・・}じゃろうが」\\


「こいつが、いわしておけば――」\\


　また一人、組みついた。\\


　張飛は、猪の股を振り上げて、突きだしてくる槍を束にして払い落した。そして自分の腰と首に組みついている二人の兵は、蠅でもたかっているように、そのまま振りのけもせず、二丈余の鉄梯子を馳け登って行った。\\


「や、やっ」\\


「\ruby{狼藉者}{ろうぜきもの}っ」\\


「関門破りだっ」\\


「出合え。出合えっ」\\


　狼狽して、わめき合う人影のうえに、城壁の上から、二箇の人間が飛んできた。もちろん、投げ落された人間も\ruby{血漿}{けっしょう}の粉になり、下になった人間も、\ruby{肉餅}{にくぺい}のように圧しつぶされた。\\

\\

\subsubsection{二}

\\


　物音に、望楼の守兵と、役人らが出て見た時は、張飛はもう、二丈余の城壁から、関外の大地へとび降りていた。\\


「\ruby{黄匪}{こうひ}だっ」\\


「間諜だ」\\


　\ruby{警鼓}{けいこ}を鳴らして、関門の上下では騒いでいたが、張飛はふりむきもせず、疾風のように馳けて行った。\\


　五、六里も来ると、一条の河があった。\ruby{蟠桃河}{ばんとうが}の支流である。河向うに約五百戸ほどの村が墨のような\ruby{夜靄}{よもや}のなかに沈んでいる。村へはいってみるとまだそう夜も\ruby{更}{ふ}けていないので、所々の家の灯皿に薄暗い明りがゆらいでいる。\\


　楊柳に囲まれた寺院がある。塀にそって張飛は大股に曲がって行った。すると大きな\ruby{棗}{なつめ}の木が五、六本あって、隠士の住居とも見える閑寂な庭があった。門柱はあるが扉はない。そしてそこの入口に、\\

\ruby{童学草舎}{どうがくそうしゃ}\\



　という看板がかかっていた。\\


「おういっ。もう寝たのか。\ruby{雲長}{うんちょう}、雲長」\\


　張飛は、烈しく、奥の家の扉をたたいた。すると横の窓に、うすい灯がさした。\ruby{帳}{とばり}を揚げて誰か窓から首を出したようであった。\\


「だれだ」\\


「それがしだ」\\


「張飛か」\\


「おう、雲長」\\


　窓の灯が、中の人影といっしょに消えた。間もなく、たたずんでいる張飛の前の扉がひらかれた。\\


「何用だ。今頃――」\\


　手燭に照らされてその人の\ruby{面}{おもて}が昼みるよりもはっきり見えた。まず驚くべきことは、張飛にも劣らない背丈と広い胸幅であった。その胸にはまた、張飛よりも長い\ruby{腮髯}{あごひげ}がふっさりと垂れていた。毛の\ruby{硬}{こわ}い者は粗暴で神経もあらいということがほんとなら、雲長というその者の髯のほうが、彼のものよりは軟かで素直でそして長いから、同時に張飛よりもこの人のほうが智的にすぐれているといえよう。\\


　智的といえば、額もひろい。眼は\ruby{鳳眼}{ほうがん}であり、\ruby{耳朶}{じだ}は豊かで、総じて、体の\ruby{巨}{おお}きいわりに\ruby{肌目}{きめ}こまやかで、音声もおっとりしていた。\\


「いや、夜中とは思ったが、一刻もはやく、尊公にも聞かせたいと思って――よろこびを\ruby{齎}{もたら}してきたのだ」\\


　張飛のことばに、\\


「また、それを\ruby{肴}{さかな}に、飲もうというのじゃないかな」\\


「ばかをいえ。それがしを、そう飲んだくれとばかり思うているから困る。平常の酒は、\ruby{鬱懐}{うっかい}をはらすために飲むのだ。今夜はその鬱懐もいっぺんに散じて、愉快でならない吉報をたずさえて来たのだ。酒がなくても、ずいぶん話せることだ。あればなおいいが」\\


「ははははは。まあ入れ」\\


　暗い廊を歩いて、一室へ二人はかくれた。その部屋の壁には、孔子やその弟子たちの聖賢の図がかかっていた。また、たくさんな机が置いてあった。門柱に見えるとおり、童学草舎は村の寺子屋であり、\ruby{主}{あるじ}は村童の先生であった。\\


「雲長――いつも話の上でばかり語っていたことだが、俺たちの夢がどうやらだんだん夢ではなく、現実になってきたらしいぞ。実はきょう、前からも心がけていたが――かねて尊公にもはなしていた\ruby{劉備}{りゅうび}という\ruby{漢}{おとこ}――それに偶然市で出会ったのだ。突っこんだ話をしてみたところ、果たして、ただの土民ではなく、漢室の\ruby{宗族}{そうぞく}景帝の\ruby{裔孫}{えいそん}ということが分った。しかも\ruby{英邁}{えいまい}な青年だ。さあ、これから楼桑村の彼の家を訪れよう。雲長、支度はそれでよいか」\\

\\

\subsubsection{三}

\\


「相かわらずだのう」\\


　\ruby{雲長}{うんちょう}は笑ってばかりいる。張飛がせきたてても、なかなか腰を上げそうもないので、張飛は、「何が相かわらずだ」と、やや突っかかるような言葉で反問した。\\


「だって」と、雲長はまた笑い、「これから楼桑村へゆけば、\ruby{真夜中}{まよなか}を過ぎてしまう。初めての家を訪問するのに、あまり礼を知らぬことに当ろう。なにも、明日でも明後日でもよいではないか。\jruby{さあ}{・・・・・・}といえば、\jruby{それ}{・・・・・・}というのが、貴公の性質だが、偉丈夫たる者はよろしくもっと沈重な態度であって欲しいなあ」\\


　せっかく、一刻もはやくよろこんでもらおうと思ってきたのに、案外、雲長が気のない返辞なので、\\


「ははあ。雲長。尊公はまだそれがしの話を、半信半疑で聞いておるんじゃないか。それで、渋ッたい\ruby{面}{おもて}をしておるのだろう。おれのことを、いつも短気というが、尊公の性質は、むしろ優柔不断というやつだ。\ruby{壮図}{そうと}を抱く勇者たる者は、もっと事に当って、果断であって欲しいものだ」\\


「ははははは。やり返したな。しかしおれは考えるな。なんといわれても、もっと熟慮してみなければ、うかつに、景帝の玄孫などという男には会えんよ。――世間に、よくあるやつだから」\\


「そら、その通り、拙者の言を疑っておるのではないか」\\


「疑ぐるのが常識で、疑わない貴公が元来、生一本のばか正直というものじゃ」\\


「聞き捨てにならんことをいう。おれがどうしてばか正直か」\\


「ふだんの生活でも、のべつ人に\ruby{騙}{だま}されておるではないか」\\


「おれはそんなに人に騙されたおぼえはない」\\


「騙されても、騙されたと覚らぬほど、尊公はお人が好いのだ。それだけの武勇をもちながら、いつも生活に困って、窮迫したり流浪したり、皆、尊公の浅慮がいたすところである。その上、短気ときているので、怒ると、途方もない暴をやる。だから張飛は悪いやつだと反対な誤解をまねいたりする。すこし反省せねばいかん」\\


「おい雲長。拙者は今夜、なにも貴公の\ruby{叱言}{こごと}を聞こうと思って、こんな夜中、やって来たわけではないぜ」\\


「だが、貴公とわしとは、かねて、お互いの大志を打明け、義兄弟の約束をし、わしは兄、貴公は弟と、固く心を結び合った仲だ。――だから弟の短所を見ると、兄たるわしは、憂えずにはいられない。まして、秘密の上にも秘密にすべき大事は、世間へ出て、二度や三度会ったばかりの\ruby{漢}{おとこ}へ、軽率に話したりなどするのはよろしくないことだ。そのうえ人の言をすぐ信じて、真夜中もかまわず直ぐ訪れようなんて……どうもそういう\ruby{浅慮}{あさはか}では案じられてならん」\\


　雲長は、劉備の家を訪問するなどもってのほかだといわぬばかりなのである。彼は、張飛にとって、いわゆる義兄弟の義兄ではあるし、物分りもすぐれているので、話が、理になってくると、いつも頭は上がらないのであった。\\


　出\jruby{ばな}{・・・・・・}をくじかれたので、張飛はすっかり\ruby{悄気}{しょげ}てしまった。雲長は気の毒になって、彼の好きな酒を出して与えたが、\\


「いや、今夜は飲まん」\\


　と、張飛はすっかり無口になって、その晩は、雲長の家で寝てしまった。\\


　夜が明けると、学舎に通う村童が、わいわいと集まってきた。雲長は、よく子供らにも\ruby{馴}{な}じまれていた。彼は、子どもらに\ruby{孔孟}{こうもう}の書を読んで聞かせ、文字を教えなどして、もう他念なき\ruby{村夫子}{そんぷうし}になりすましていた。\\


「また、そのうちに来るよ」\\


　学舎の窓から雲長へいって、張飛は黙々とどこかへ出て行った。\\

\\

\subsubsection{四}

\\


　むっとして、張飛は、雲長の家の門を出た。門を出ると、振向いて、\\


「ちぇっ。なんていう煮え切らない\ruby{漢}{おとこ}だろう」と、門へ\ruby{罵}{ののし}った。\\


　楽しまない顔色は、それでも\ruby{癒}{い}えなかった。村の居酒屋へくると、ゆうべから\ruby{渇}{かわ}いていたように、すぐ呶鳴った。\\


「おいっ、酒をくれい」\\


　朝の空き腹に、\ruby{斗酒}{としゅ}をいれて、張飛はすこし、眼のふちを赤黒く染めた。\\


　やや気色が晴れてきたとみえて居酒屋の亭主に、\ruby{冗戯}{じょうだん}などいいだした。\\


「おやじ、お前んとこの\ruby{鶏}{とり}は、おれに喰われたがって、おれの足もとにばかりまとってきやがる。喰ってもいいか」\\


「旦那、召しあがるなら、毛をむしって、丸揚げにしましょう」\\


「そうか。そうしてくれればなおいいな。あまり鶏めが慕ってくるから、\ruby{生}{なま}で\ruby{喰}{や}ろうと思っていたんだが」\\


「生肉をやると腹に虫がわきますよ、旦那」\\


「ばかをいえ。\ruby{鶏}{とり}の肉と馬の肉には寄生虫は\ruby{棲}{す}んでおらん」\\


「ヘエ。そうですか」\\


「体熱が高いからだ。すべて低温動物ほど寄生虫の巣だ。国にしてもそうだろう」\\


「へい」\\


「おや、鶏がいなくなった。おやじもう釜へ入れたのか」\\


「いえ。お代さえいただけば、揚げてあるやつを直ぐお出しいたしますが」\\


「銭はない」\\


「ごじょうだんを」\\


「ほんとだよ」\\


「では、お酒のお代のほうは」\\


「この先の寺の横丁を曲がると、童学草舎という寺子屋があるだろう。あの雲長のとこへ行って貰ってこい」\\


「弱りましたなあ」\\


「なにが弱る。雲長という\ruby{漢}{おとこ}は、武人のくせに、金に困らぬやつだ。雲長はおれの\ruby{兄哥}{あにき}だ。弟の張飛が飲んで行ったといえば、払わぬわけにはゆくまい。――おいっ、もう一杯ついでこい」\\


　亭主は、如才なく、彼をなだめておいて、その間に、女房を裏口からどこかへ走らせた。雲長の家へ問合せにやったものとみえる。間もなく、帰ってきて何かささやくと、\\


「そうかい。じゃあ飲ませても間違いあるまい」\\


　おやじはにわかに、態度を変えて、張飛の飲みたい放題に、酒をつぎ鶏の丸揚げも出した。\\


　張飛は、丸揚げを見ると、\\


「こんな、鶏の\ruby{乾物}{ひもの}など、おれの口には合わん。おれは動いている奴を喰いたいのだ」\\


　と、そこらにいる鶏をとらえようとして、往来まで追って行った。\\


　鶏は羽ばたきして、彼の肩を跳び越えたり、彼の危うげな股をくぐって、逃げ廻ったりした。\\


　すると、しきりに、村の軒並を物色してきた捕吏が、張飛のすがたを認めると、\ruby{率}{ひ}きつれている十名ほどの兵へにわかに命令した。\\


「あいつだ。ゆうべ関門を破った上、衛兵を殺して逃げた賊は。――要心してかかれ」\\


　張飛は、その声に、\\


「何だろ？」と、いぶかるように、あたりを酔眼で見まわした。一羽の若鶏が彼の手に脚をつかまえられて、けたたましく啼いたり羽ばたきをうっていた。\\


「賊っ」\\


「\ruby{遁}{のが}さん」\\


「神妙に縄にかかれ」\\


　捕吏と兵隊に取囲まれて、張飛ははじめて、おれのことかと気づいたような面もちだった。\\


「何か用か」\\


　まわりの槍を見まわしながら、張飛は、若鶏の脚を引っ裂いて、その股の肉を横にくわえた。\\

\\

\subsubsection{五}

\\


　酔うと酒くせのよくない張飛であった。それといたずらに\ruby{殺伐}{さつばつ}を好む癖は、二つの欠点であるとは常々、雲長からもよくいわれていることだった。\\


　\ruby{鶏}{とり}を裂いて、股を喰らうぐらいな酒の上は、彼としては、いと穏当な芸である。――だが、捕吏や兵隊は驚いた。鶏の血は張飛の唇のまわりを染め、その\ruby{炯々}{けいけい}たる眼は怖ろしく不気味であった。\\


「なに？　……おれを捕まえにきたと。……わははははは。あべこべに取っつかまって、この通りになるなよ」\\


　裂いた鶏を、眼の高さに、上げて示しながら、張飛は取囲む捕吏と兵隊を\ruby{揶揄}{やゆ}した。\\


　捕吏は怒って、\\


「それっ、酔どれに、愚図愚図いわすな。突き殺してもかまわん。かかれっ」と、呶号した。\\


　だが、兵隊たちは、近寄れなかった。槍ぶすまを並べたまま、彼の周囲を巡りまわったのみだった。\\


　張飛は、変な腰つきをして、犬みたいにつく\ruby{這}{ば}った。それがよけいに捕吏や兵隊を恐怖させた。彼の眼が向ったほうへ飛びかかってくる支度だろうと思ったからである。\\


「さあ、大きな鶏どもめ、一羽一羽、ひねりつぶすから逃げるなよ」\\


　張飛はいった。\\


　彼の頭にはまだ鶏を追いかけ廻している\ruby{戯}{たわむ}れが連続していて、捕吏の頭にも、兵隊の頭にも、\ruby{鶏冠}{とさか}が生えているように見えているらしかった。\\


　大きな鶏どもは呆れかつ怒り\ruby{心頭}{しんとう}に発して、\\


「野郎っ」と、\ruby{喚}{わめ}きながら一人が槍でなぐった。槍は正確に、張飛の肩へ当ったが、それは猛虎の髯にふれたも同じで、張飛の酔いをして\ruby{勃然}{ぼつぜん}と遊戯から殺伐へと転向させた。\\


「やったな」\\


　槍を引ったくると、張飛はそれで、\ruby{莚}{むしろ}の\ruby{豆幹}{まめがら}でも叩くように、まわりの人間を叩き出した。\\


　叩かれた捕吏や兵隊も、はじめて死にもの狂いになり始めた。張飛は、面倒といいながら槍を虚空へ投げた。虚空へ飛んだ槍は、唸りを起したままどこまで飛んで行ったか、なにしろその附近には落ちてこなかった。\\


　鶏の悲鳴以上な\ruby{叫喚}{きょうかん}が、一瞬のまに起って、一瞬の間にやんでしまった。\\


　居酒屋のおやじ、居合せた客、それから往来の者や、附近の人たちは皆、家の中や木蔭にひそんで、どうなることかと、息をころしていたが、余りにそこが、急に墓場のような\ruby{寂寞}{しじま}になったので、そっと首を出して往来をながめると、ああ――と誰も\ruby{呻}{うめ}いたままで口もきけなかった。\\


　首を払われた死骸、血へどを吐いた死骸、眼のとびだしている死骸などが、惨として、太陽の\ruby{下}{もと}にさらされている。\\


　半分は、逃げたのだろう。捕吏も兵隊も、誰もいない。\\


　張飛は？\\


　と見ると、これはまた、悠長なのだ。村はずれのほうへ、後ろ姿を見せて、\ruby{寛々}{かんかん}と歩いてゆく。\\


　その\ruby{袂}{たもと}に、春風はのどかに動いていた。酒のにおいが、遠くにまで、\ruby{漂}{ただよ}ってくるように――。\\


「たいへんだ。おい、はやくこのことを、雲長先生の家へ知らせてこい。あの\ruby{漢}{おとこ}が、ほんとに、先生の舎弟なら、これはあの先生も、ただでは済まないぞ」\\


　居酒屋のおやじは、自分のおかみさんへ\ruby{喚}{わめ}いた。だが、彼の妻はふるえているばかりで役に立たないので、ついに自分であたふたと、童学草舎の横丁へ、馳けよろめいて行った。\\

\\

\subsection{三\ruby{花}{か}一\ruby{瓶}{ぺい}}

\\

\subsubsection{一}

\\


　母と子は、仕事の庭に、きょうも他念なく、\ruby{蓆機}{むしろばた}に向って、蓆を織っていた。\\


　がたん……\\


　ことん\\


　がたん\\


　水車の\ruby{回}{まわ}るような単調な音がくり返されていた。\\


　だが、その音にも、きょうはなんとなく活気があり、歓喜の\ruby{譜}{ふ}があった。\\


　黙々、仕事に精だしてはいるが、母の胸にも、\ruby{劉備}{りゅうび}の心にも、今日この頃の大地のように、希望の芽が生々と息づいていた。\\


　ゆうべ。\\


　劉備は、城内の市から帰ってくると、まっ先に、二つの吉事を告げた。\\


　一人の良き友に出会った事と、かねて手放した家宝の剣が、計らず再び、自分の手へ返ってきた事と。\\


　そう二つの歓びを告げると、彼の母は、\\


「一\ruby{陽来復}{ようらいふく}。おまえにも時節が来たらしいね。劉備や……心の支度もよいかえ」\\


　と、かえって静かに声を低め、劉備の覚悟を\ruby{糺}{ただ}すようにいった。\\


　時節。……そうだ。\\


　長い長い冬を経て、桃園の花もようやく\ruby{蕾}{つぼみ}を破っている。土からも草の芽、木々の枝からも緑の芽、生命のあるもので、\ruby{萌}{も}え出ない物はなに一つない。\\


　がたん……\\


　ことん……\\


　\ruby{蓆機}{むしろばた}は単調な音をくりかえしているが、劉備の胸は単調でない。こんな春らしい春をおぼえたことはない。\\


　――我は青年なり。\\


　空へ向って言いたいような気持である。いやいや、老いたる母の肩にさえ、どこからか舞ってきた桃花の\ruby{一片}{ひとひら}が、\ruby{紅}{あか}く点じているではないか。\\


　すると、どこかで、歌う者があった。十二、三歳の少女の声だった。\\

\ruby{妾}{ショウ}ガ髪初メテ\ruby{額}{ヒタイ}ヲ\ruby{覆}{オオ}ウ\\


花ヲ折ッテ門前ニ\ruby{戯}{タワム}レ\\

\ruby{郎}{ロウ}ハ竹馬ニ騎シテ\ruby{来}{キタ}リ\\

\ruby{床}{ショウ}ヲ\ruby{繞}{メグ}ッテ青梅ヲ\ruby{弄}{ロウ}ス\\



　劉備は、耳を澄ました。\\


　少女の美音は、近づいてきた。\\



……十四\ruby{君}{キミ}ノ婦ト\ruby{為}{ナ}ッテ\\

\ruby{羞顔}{シュウガン}\ruby{未}{イマ}ダ\ruby{嘗}{カツ}テ開カズ\\


十五初メテ\ruby{眉}{マユ}を\ruby{展}{ノ}ベ\\


願ワクバ\ruby{塵}{チリ}ト灰トヲ共ニセン\\


常ニ\ruby{抱柱}{ホウチュウ}ノ信ヲ\ruby{存}{ソン}シ\\

\ruby{豈}{アニ}\ruby{上}{ノボ}ランヤ\ruby{望夫台}{ボウフダイ}\\


十六\ruby{君}{キミ}遠クヘ行ク\\



　近所に住む少女であった。早熟な彼女はまだ青い\ruby{棗}{なつめ}みたいに小粒であったが、劉備の家のすぐ\ruby{墻隣}{かきどなり}の息子に恋しているらしく、星の晩だの、人気ない折の真昼などうかがっては、墻の外へきて、よく歌をうたっていた。\\


「…………」\\


　劉備は、木蓮の花に\ruby{黄金}{きん}の\ruby{耳環}{みみわ}を通したような、少女の\ruby{貌}{かお}を眼にえがいて、隣の息子を、なんとなく\ruby{羨}{うらや}ましく思った。\\


　そしてふと、自分の心の底からも一人の麗人を思い出していた。それは、三年前の旅行中、古塔の下であの折の老僧にひき合わされた\ruby{鴻家}{こうけ}の息女、\ruby{鴻芙蓉}{こうふよう}のその後の消息であった。\\


　――どうしたろう。あれから先。\\


　張飛に訊けば、知っている筈である。こんど張飛に会ったら――など独り考えていた。\\


　すると、\ruby{墻}{かき}の外で、しきりに歌をうたっていた少女が、犬にでも噛まれたのか、突然、きゃっと悲鳴をあげて、どこかへ逃げて行った。\\

\\

\subsubsection{二}

\\


　少女は、犬に\ruby{咬}{か}まれたわけではなかった。\\


　自分のうしろに、この辺で見たこともない、剣を\ruby{佩}{は}いた巨きな\ruby{髯漢}{ひげおとこ}が、いつのまにか来ていて、\\


「おい、小娘、劉備の家はどこだな」と、訊ねたのだった。\\


　けれど、少女は、振向いてその\ruby{漢}{おとこ}を仰ぐと、姿を見ただけで、\ruby{胆}{きも}をつぶし、きゃっといって、逃げ走ってしまったのであった。\\


「あははは。わははは」\\


　髯漢は、小娘の驚きを、滑稽に感じたのか、独りして笑っていた。\\


　その笑い声が止むと一緒に、後ろの\ruby{墻}{かき}の内でも、はたと、\ruby{蓆機}{むしろばた}の音が止んでいた。\\


　墻といっても\ruby{匪賊}{ひぞく}に備えるためこの辺では、すべてといってよい程、土民の家でも、土の塀か、石で組上げた物でできていたが、劉家だけは、泰平の頃に建てた旧家の慣わしで、高い樹木と灌木に、細竹を渡して結ってある生垣だった。\\


　だから、背の高い張飛は、首から上が、生垣の上に出ていた。劉備の庭からもそれが見えた。\\


　ふたりは顔を見合って、\\


「おう」\\


「やあ」\\


　と、十年の知己のように呼び合った。\\


「なんだ、ここか」\\


　張飛は、外から木戸口を見つけてはいって来た。ずしずしと地が鳴った。劉家はじまって以来、こんな大きな跫音が、この家の庭を踏んだのは初めてだろう。\\


「きのうは失礼しました。君に会ったことや、剣のことを、母に話したところ、母もゆうべは歓んで、夜もすがら希望に\ruby{耽}{ふけ}って、語り明かしたくらいです」\\


「あ。こちらが貴公の\ruby{母者人}{ははじゃひと}か」\\


「そうです。――母上、このお方です。きのうお目にかかった\ruby{翼徳}{よくとく}張飛という豪傑は」\\


「オオ」\\


　劉備の母は、\ruby{機}{はた}の前からすっと立って張飛の礼をうけた。どういうものか、張飛は、その母公の姿から、劉備以上、気高い威圧をうけた。\\


　また、実際、劉備の母にはおのずから備わっている名門の気品があったのであろう。世の常の甘い母親のように、息子の友達だからといって、やたらに小腰をかがめたりチヤホヤはしなかった。\\


「劉備からおはなしは聞きました。失礼ですが、お見うけ申すからに頼もしい偉丈夫。どうか、柔弱なわたしの一子を、これから\ruby{叱咤}{しった}して下さい。おたがいに\ruby{鞭撻}{べんたつ}し合って、大事をなしとげて下さい」と、いった。\\


「はっ」\\


　張飛は、自然どうしても、頭を下げずにはいられなかった。\ruby{長上}{ちょうじょう}に対する礼儀のみからではなかった。\\


「母公。安心して下さい。きっと男児の素志をつらぬいて見せます。――けれどここに、遺憾なことが一つ起りました。で、実はご子息に相談に来たわけですが」\\


「では、男同士のはなし、わたくしは部屋へ行っていましょう。ゆるりとおはなしなさい」\\


　母は、奥へかくれた。\\


　張飛は、その後の\ruby{床几}{しょうぎ}へ腰かけて、実は――と、自分の盟友、いや義兄とも仰いでいる、雲長のことを話しだした。\\


　雲長も、自分が見込んだ\ruby{漢}{おとこ}で、何事も打明け合っている仲なので、早速、ゆうべ訪れて、仔細を話したところ、意外にも、彼は少しも歓んでくれない。\\


　のみならず、景帝の\ruby{裔孫}{えいそん}などとは、むしろ怪しむべき者だ。そんな路傍の\jruby{まやかし}{・・・・・・}者と、大事を語るなどは、もってのほかであると叱られた。\\


「残念でたまらない。雲長めは、そういって疑うのだ。……ご足労だが、貴公、これから拙者と共に、彼の住居まで行ってくれまいか。貴公という人間を見せたら、彼も恐らくこの張飛の言を信じるだろうと思うから――」\\

\\

\subsubsection{三}

\\


　張飛は、疑いが嫌いだ。疑われることはなお嫌いだ。雲長が、自分の言を信じてくれないのが、心外でならないのである。\\


　だから劉備を連れて行って、その人物を実際に示してやろう――こう考えたのも張飛らしい考えであった。\\


　しかし、劉備は、「……さあ？」と、いって、考えこんだ。\\


　信じない者へ、\ruby{強}{し}いて、自己を押しつけて、信じろというのも、好ましくないとする風だった。\\


　すると、廊のほうから、\\


「劉備。行っておいでなさい」\\


　彼の母がいった。\\


　母は、やはり心配になるとみえて、\ruby{彼方}{かなた}で張飛のはなしを聞いていたものとみえる。\\


　もっとも、張飛の声は、この家の中なら、どこにいても聞えるほど大きかった。\\


「やあ、お許し下さるか。母公のおゆるしが出たからには、劉君、何もためらうことはあるまい」\\


　促すと、母も共に、「時機というものは、その時をのがしたら、またいつ\ruby{巡}{めぐ}ってくるか知れないものです。――何やら、今はその天機が巡ってきているような気がするのです。\ruby{些細}{ささい}な気持などにとらわれずに、お誘いをうけたものなら、張飛どのにまかせて、行ってごらんなさい」\\


　劉備は、母のことばに、\\


「では、参ろう」と決心の腰を上げた。\\


　二人は並んで、廊のほうへ、\\


「では、行ってきます」\\


　礼をして、\ruby{墻}{かき}の外へ出て行った。\\


　すると、道の彼方から、約百人ほどの軍隊が、まっしぐらに馳けてきた。騎馬もあり徒歩の兵もあった。\ruby{埃}{ほこり}の中に、青龍刀の白い光がつつまれて見えた。\\


「あ……、また来た」\\


　張飛のつぶやきに、劉備はいぶかって、\\


「なんです、あれは」\\


「城内の兵だろう」\\


「関門の兵らしいですね。何事があったのでしょう」\\


「たぶん、この張飛を、召捕らえにきたのかも知れん」\\


「え？」\\


　劉備は、驚きを\ruby{喫}{きっ}して、\\


「では、こっちへむかって来る軍隊ですか」\\


「そうだ。もう疑いない。劉君、あれをちょっと片づける間、貴公はどこかに休んで見物していてくれないか」\\


「弱りましたな」\\


「なに、大したことはない」\\


「でも、州郡の兵隊を\ruby{殺戮}{さつりく}したら、とてもこの土地にはおられませんぞ」\\


　云っている間に、もう百余名の州郡の兵は張飛と劉備を包囲してわいわい騒ぎだした。\\


　だが、容易に手は下してはこなかった。張飛の武力を二度まで知っているからであろう。けれど二人は一歩もあるくことはできなかった。\\


「邪魔すると、蹴殺すぞ」\\


　張飛は、一方へこう呶鳴って歩きかけた。わっと兵は\ruby{退}{ひ}いたが、背後から矢や鉄槍が飛んできた。\\


「面倒っ」\\


　またしても、張飛は持ち前の短気を出して、すぐ剣の\ruby{柄}{つか}へ手をかけた。\\


　――すると、彼方から一頭の\ruby{逞}{たくま}しい\ruby{鹿毛}{かげ}を飛ばして、\\


「待てっ、待てえ」\\


　と呼ばわりながら馳けてくる者があった。州郡の兵も、張飛も、何気なく眼をそれへはせて振向くと、胸まである\ruby{黒髯}{こくぜん}を春風になぶらせ、腰に\ruby{偃月刀}{えんげつとう}の\ruby{佩環}{はいかん}を\ruby{戛々}{かつかつ}とひびかせながら、手には\ruby{緋総}{ひぶさ}のついた\ruby{鯨鞭}{げいべん}を持った偉丈夫が、その鞭を上げつつ近づいてくるのであった。\\

\\

\subsubsection{四}

\\


　それは、雲長であった。\\


　\ruby{童学草舎}{どうがくそうしゃ}の\ruby{村夫子}{そんぷうし}も、武装すれば、こんなにも威風堂々と見えるものかと、眼をみはらせるばかりな雲長の風貌であった。\\


「待て諸君」\\


　乗りつけてきた鹿毛の鞍から跳び降りると、雲長は、兵の中へ割って入り、そこに囲まれている張飛と劉備を後ろにして、大手をひろげながらいった。\\


「貴公らは、関門を守備する領主の兵と見うけるが、五十や百の小人数をもって、一体なにをなさろうとするのか。――この\ruby{漢}{おとこ}を召捕ろうとするならば」と、背後にいる張飛へ、顎を振向けて、\\


「まず五百か千の人数をそろえてきて、半分以上の\ruby{屍}{しかばね}はつくる覚悟がなければからめ捕ることはできまい。諸君は、この\ruby{翼徳張飛}{よくとくちょうひ}という人間が、どんな力量の漢か知るまいが、かつて、幽州の\ruby{鴻家}{こうけ}に仕えていた頃、重さ九十\ruby{斤}{きん}、長さ一丈八尺の\ruby{蛇矛}{じゃぼこ}をふるって、\ruby{黄巾賊}{こうきんぞく}の大軍中へ馳けこみ、\ruby{屍山血河}{しざんけつが}をつくって、半日の合戦に八百八\ruby{屍}{し}の死骸を積み、張飛のことを、八百八屍将軍と\ruby{綽名}{あだな}して、\ruby{黄匪}{こうひ}を戦慄させたという勇名のある漢だ。――それを、\ruby{素手}{すで}にもひとしい小人数で、からめ捕ろうなどは、檻へ入って、虎と組むようなもの、\ruby{各{\fontspec{jigmo}\char"303B}}{おのおの}が皆、死にたいという願いで、この漢へかまうなら知らぬこと、命知らずな真似はやめたらどうだ。生命の欲しい者は足もとの明るいうちに帰れ。ここは、かくいう雲長にまかせて、ひとまず引揚げろ」\\


　雲長は、実に雄弁だった。一息にここまで演説して、まったく相手の気をのんでしまい、さらに語をついでいった。\\


「――こういったら諸公は、わしを何者ぞと疑い、また、巧みに張飛を逃がすのではないかと、疑心を抱くであろうが、さに\ruby{非}{あら}ず、不肖はかりそめにも、童学草舎を営み子弟の\ruby{薫陶}{くんとう}を任とし、常に聖賢の道を本義とし、国主を尊び、法令を\ruby{遵守}{じゅんしゅ}すべきことを、身にも守り、子弟に教えている\ruby{雲長関羽}{うんちょうかんう}という者である。そして、これにいる翼徳張飛は、何をかくそう自身の義弟にあたる人間でもある。――だが、昨夜から今朝にかけて、張飛が官の吏兵を殺害し、関門を破り、酒の上で暴行したことを聞き及んで、ゆるしがたく思い、この上多くの\ruby{犠牲}{いけにえ}を出さんよりは、義兄たるわが手に召捕りくれんものと、かくは身固め致して、官へ願い出で、宙を馳せてこれへ駆けつけてきたわけでござる――。張飛はこの雲長が召捕って、後刻、太守の県城へまで送り届けん。諸公は、ここの事実を見とどけて、その由、先へご報告おきねがう」\\


　雲長は、\ruby{沓}{くつ}をめぐらして、きっと張飛のほうへ今度は向きなおった。\\


　そして、\ruby{大喝}{だいかつ}一声、\\


「ここな不届き者っ」\\


　と、鯨の\ruby{鞭}{むち}で、張飛の肩を打ちすえた。\\


　張飛は、むかっとしたような眼をしたが、雲長はさらに、\\


「\ruby{縛}{ばく}につけ」と、跳びかかって、張飛の両手を後ろへまわした。\\


　張飛は、雲長の心を疑いかけたが、より以上、雲長の人物を信じる心のほうが強かった。\\


　で――何か考えがあることだろうと、神妙に縄を受けて、大地へ坐ってしまった。\\


「見たか、諸公」\\


　雲長は再び、呆っ気にとられている捕吏や兵の顔を見まわして、\\


「張飛は、後刻、それがしが県城へ直接参って渡すから、諸公は先へここを引揚げられい。それともなお、この雲長を怪しみ、それがしの言葉を疑うならば、ぜひもない、縄を解いて、この猛虎を、諸公の中へ放つが、どうだ」\\


　いうと、捕吏も兵も、逃げ足早く、物もいわず皆、退却してしまった。\\

\\

\subsubsection{五}

\\


　誰もいなくなると、雲長はすぐ張飛の縄を解いて、\\


「よく俺を信じて、神妙にしていてくれた。事なく助ける策謀のためとはいえ、貴様を手にかけた罪はゆるしてくれ」\\


　詫びると、張飛も、\\


「それどころではない。また無益の\ruby{殺生}{せっしょう}を重ねるところを、尊兄のお蔭で助かった」と、今朝のむかっ腹もわすれて、いつになく、素直に謝った。そして、「――だが雲長。その身なりは一体、どうしたことか。俺を助けにくるためにしては、余りに物々しい装いではないか」\\


　怪しんで問うと、\\


「張飛。なにをとぼけたことをいう。それでは昨夜、あんなに熱をこめて、時節到来だ、良き盟友をえた、いざ、かねての約束を、実行にかかろうといったのは、嘘だったのか」\\


「嘘ではないが、大体、尊兄が不賛成だったろう。俺のいうこと何ひとつ、信じてくれなかったじゃないか」\\


「それは、あの場のことだ。召使いもいる、女どももいる。貴様のはなしは、秘密秘密といいながら、あの大声だ。洩れてはならない――そう考えたから一応冷淡に聞いていたのだ」\\


「なんだ、それなら、尊兄もわしの言葉を信じ、かねての計画へ乗りだす\ruby{肚}{はら}を固めてくれたのか」\\


「おぬしの言葉よりも、実は、相手が楼桑村の劉備どのと聞いたので、即座に心はきめていたのだ。かねがね、わしの村まで孝子という噂の聞えている劉備どの、それによそながら、ご素姓や平常のことなども、ひそかに調べていたので」\\


「人が悪いな。どうも尊兄は、智謀を\ruby{弄}{ろう}すので、\ruby{交際}{つきあ}いにくいよ」\\


「ははは。貴様から交際いにくいといわれようとは思わなかった。人を殺し、酒屋を飲みたおし、その\ruby{尻尾}{しっぽ}は童学草舎へ持って行けなどという乱暴者から、そういわれてはたまらない」\\


「もう行ったか」\\


「酒屋の勘定ぐらいならよいが、官の\ruby{捕手}{とりて}を殺したのは、雲長の\ruby{義弟}{おとうと}だと分ったひには、童学草舎へも子供を通わせる親はあるまい。いずれ官からこの雲長へも、やかましく出頭を命じてくるにきまっている」\\


「なるほど」\\


「\ruby{他人事}{ひとごと}のように聞くな」\\


「いや、済まん」\\


「しかし、これはむしろ、よい\ruby{機}{しお}だ。天意の命じるものである。こう考えたから、今朝、召使いや女どもへ、みな暇を出した上、通学してくる子供たちの親も呼んで、都合によって学舎を閉鎖するといい渡し、心おきなく、身一つになって、かくは貴様の後を追って来たわけだ。――さ。これから改めて、劉備どのの家へお目にかかりに行こう」\\


「いや。劉備どのなら、そこにいる」\\


「え？　……」\\


　雲長は、張飛の指さす所へ、眼を振り向けた。\\


　劉備は最前から、少し離れた所に立っていた。そして、張飛と雲長との二人の仲の\ruby{睦}{むつ}まじさと、その信義に篤い様子を見て、感にたえている面もちだった。\\


「あなたが劉備様ですか」\\


　雲長は、近づいて行くと、彼の足もとへ最初から膝を折って、\\


「初めてお目にかかります。自分は\ruby{河東解良}{かとうかいりょう}（\ruby{山西省}{さんせいしょう}・\ruby{解県}{かいけん}）の産で、\ruby{関羽}{かんう}\ruby{字}{あざな}は雲長と申し、長らく\ruby{江湖}{こうこ}を\ruby{流寓}{りゅうぐう}のすえ、四、五年前よりこの近村に住んで、村夫子となって草裡にむなしく月日を送っていた者です。かねてひそかに心にありましたが、計らずも今日、拝姿の栄に会い、こんな歓ばしいことはありません。どうかお見知りおき下さい」\\


　と、最高な礼儀をとって、\ruby{慇懃}{いんぎん}にいった。\\

\\

\subsubsection{六}

\\


　劉備はあえて、\ruby{卑下}{ひげ}しなかったが、べつに尊大に構えもしなかった。雲長関羽の礼に対して、当り前に礼を返しながら、\\


「ご丁寧に。……どうも申し遅れました。私は、楼桑村に永らく住む百姓の劉玄徳という者ですが、かねて、\ruby{蟠桃河}{ばんとうが}の\ruby{上流}{かみ}の村に、醇風良俗の桃源があると聞きました。おそらく先生の高風に化されたものでありましょう。なにをいうにも、ここは路傍ですから、すぐそこの\ruby{茅屋}{あばらや}までお越しください」\\


　と、誘えば、\\


「おお、お供しよう」\\


　関羽も歩み、張飛も肩を並べ、共にそこからほど近い劉備の家まで行った。\\


　劉備の母は、また新しい客がふえたので、不審がったが、張飛から紹介されて、関羽の人物を見、よろこびを現して、\\


「ようぞ、\ruby{茅屋}{あばらや}へ」と心から歓待した。\\


　その晩は、母もまじって、夜更けまで語った。劉備の母は、劉家の古い歴史を、覚えている限り話した。\\


　生れてからまだ劉備さえ聞いていない話もあった。\\


（いよいよ漢室のながれを汲んだ\ruby{景帝}{けいてい}の\ruby{裔孫}{えいそん}にちがいない）\\


　張飛も、関羽も、今は少しの疑いも抱かなかった。\\


　同時に、この人こそ、義挙の盟主になすべきであると肚にきめていた。\\


　しかし、劉玄徳の母親思いのことは知っているので、この母親が、\\


（そんな危ない\ruby{企}{たくら}みに息子を加えることはできない）\\


　と、断られたらそれまでになる。関羽は、それを考えて、ぼつぼつと母の胸をたずねてみた。\\


　すると劉備の母は、みなまで聞かないうちにいった。\\


「ねえ劉や、今夜はもうおそいから、おまえも寝み、お客様にも\ruby{臥床}{ふしど}を作っておあげなさい。――そして明日はいずれまた、お三名して将来の相談もあろうし、大事の門出でもありますし、母が一生一度の馳走をこしらえてあげますからね」\\


　それを聞いて、関羽は、この母親の胸を問うなど\ruby{愚}{ぐ}であることを知った。張飛も共に、頭を下げて、「ありがとうござる」と、心服した。\\


　劉備は、\\


「では、お言葉に甘えて、明日はおっ母さんに、一世一代の祝いを\ruby{奢}{おご}っていただきましょう。けれどそのご馳走は、吾々ばかりでなく、祭壇を設けて、先祖にも上げていただきたいものです」\\


「では、ちょうど今は、桃園の花が真盛りだから、桃園の中に\ruby{蓆}{むしろ}を敷こうかね」\\


　張飛は手を打って、\\


「それはいい。では吾々も、あしたは朝から桃園を\ruby{浄}{きよ}めて、せめて祭壇を作る手助けでもしよう」\\


　と、いった。\\


　客の二人に\ruby{床}{しょう}を与えて、眠りをすすめ、劉備と母のふたりは、暗い\ruby{厨}{くりや}の片隅で、藁をかぶって寝た。劉が眼をさましてみると、母はもういなかった。夜は明けていたのである。どこかでしきりに、山羊の啼く声がしていた。\\


　厨の\ruby{窯}{かまど}の下には、どかどかと\ruby{薪}{まき}がくべられていた。こんなに景気よく窯に薪の焚かれた\ruby{例}{ためし}は、劉備が少年の頃から覚えのないことであった。春は桃園ばかりでなく、貧しい劉家の台所に訪れてきたように思われた。\\

\\

\subsection{\ruby{義盟}{ぎめい}}

\\

\subsubsection{一}

\\


　桃園へ行ってみると、関羽と張飛のふたりは、近所の男を雇ってきて、園内の中央に、もう祭壇を作っていた。\\


　壇の四方には、\ruby{笹竹}{ささだけ}を建て、\ruby{清縄}{せいじょう}をめぐらして\ruby{金紙}{きんし}\ruby{銀箋}{ぎんせん}の\ruby{華}{はな}をつらね、土製の白馬を\ruby{贄}{いけにえ}にして天を祭り、烏牛を\ruby{屠}{ほふ}ったことにして、地神を\ruby{祠}{まつ}った。\\


「やあ、おはよう」\\


　\ruby{劉備}{りゅうび}が声をかけると、\\


「おお、お目ざめか」\\


　張飛、関羽は、振向いた。\\


「見事に祭壇ができましたなあ。寝る間はなかったでしょう」\\


「いや、張飛が、興奮して、寝てから話しかけるので、ちっとも眠る間はありませんでしたよ」\\


　と、関羽は笑った。\\


　張飛は劉備のそばへきて、\\


「祭壇だけは立派にできたが、酒はあるだろうか」\\


　心配して訊ねた。\\


「いや、母が何とかしてくれるそうです。今日は、一生一度の祝いだといっていますから」\\


「そうか、それで安心した。しかし劉兄、いいおっ母さんだな。ゆうべからそばで見ていても、\ruby{羨}{うらやま}しくてならない」\\


「そうです。自分で自分の母を褒めるのもへんですが、子に優しく世に強い母です」\\


「気品がある、どこか」\\


「失礼だが、劉兄には、まだ\ruby{夫人}{おくさん}はないようだな」\\


「ありません」\\


「はやくひとり\ruby{娶}{めと}らないと、母上がなんでもやっている様子だが、あのお年で、お気の毒ではないか」\\


「…………」\\


　劉備は、そんなことを訊かれたので、またふと、忘れていた\ruby{鴻芙蓉}{こうふよう}の佳麗なすがたを思い出してしまった。\\


　で、つい答えを忘れて、何となく眼をあげると、眼の前へ、白桃の花びらが、\ruby{霏々}{ひひ}と情あるもののように散ってきた。\\


「劉備や。皆さんも、もうお支度はよろしいのですか」\\


　厨に見えなかった母が、いつの間にか、三名の後ろにきて告げた。\\


　三名が、いつでもと答えると、母はまた、いそいそと\ruby{厨房}{ちゅうぼう}のほうへ去った。\\


　近隣の人手を借りてきたのであろう。きのう張飛の姿を見て、きゃっと\ruby{魂消}{たまげ}て逃げた娘も、その娘の恋人の隣家の息子も、ほかの家族も、大勢して手伝いにきた。\\


　やがて、まず一人では持てないような\ruby{酒瓶}{さかがめ}が祭壇の\ruby{莚}{むしろ}へ運ばれてきた。\\


　それから豚の仔を丸ごと油で煮たのや、山羊の吸物の鍋や、\ruby{干菜}{かんさい}を\ruby{牛酪}{ぎゅうらく}で煮つけた物だの、年数のかかった漬物だの――運ばれてくるごとに、三名は、その豪華な珍味の鉢や大皿に眼を奪われた。\\


　劉備さえ、心のうちで、\\


「これは一体、どうしたことだろう」と、母の算段を心配していた。\\


　そのうちにまた、村長の家から、\ruby{花梨}{かりん}の立派な卓と\ruby{椅子}{いす}がかつがれてきた。\\


「大饗宴だな」\\


　張飛は、子どものように、歓喜した。\\


　準備ができると、手伝いの者は皆、母屋へ退がってしまった。\\


　三名は、\\


「では」\\


　と、眼を見合せて、祭壇の前の\ruby{蓆}{むしろ}へ坐った。そして天地の神へ、\\


「われらの大望を成就させ給え」\\


　と、祈念しかけると、関羽が、\\


「ご両所。少し待ってくれ」\\


　と、なにか改まっていった。\\

\\

\subsubsection{二}

\\


「ここの祭壇の前に坐ると同時に、自分はふと、こんな考えを呼び起されたが、両公の所存はどんなものだろうか」\\


　関羽は、そう云いだして、劉備と張飛へ、こう相談した。\\


　すべて物事は、\ruby{体}{たい}を\ruby{基}{もと}とする。体形を整えていないことに成功はあり得ない。\\


　偶然、自分たち三人は、その精神において、\ruby{合致}{がっち}を見、きょうを出発として大事をなそうとするものであるが、三つの者が寄り合っただけでは、体をなしていない。\\


　今は、小なる三人ではあるが、理想は遠大である。三体一心の体を整えおくべきではあるまいか。\\


　事の中途で、仲間割れなど、よくある例である。そういう結果へ到達させてはならない。神のみ\ruby{祷}{いの}り、神のみ\ruby{祀}{まつ}っても、人事を尽さずして、大望の成就はあり得べくもあるまい。\\


　関羽の説くところは、道理であったが、さてどういう体を備えるかとなると、張飛にも劉備にもさし当ってなんの考えもなかった。\\


　関羽は、語をつづけて、\\


「まだ兵はおろか、兵器も金も一頭の馬すら持たないが、三名でも、ここで義盟を結べば、即座に一つの軍である。軍には将がなければならず、武士には主君がなければならぬ。行動の中心に正義と報国を奉じ、個々の中心に、主君を持たないでは、それは徒党の乱に終り、\ruby{烏合}{うごう}の\ruby{衆}{しゅう}と化してしまう。――張飛もこの関羽も、今日まで、\ruby{草田}{そうでん}に隠れて時を待っていたのは、実に、その中心たるお人が容易にないためだった。折ふし劉備玄徳という、しかも血統の正しいお方に会ったのが、急速に、今日の義盟の会となったのであるから、今日ただいま、ここで劉備玄徳どのを、自分らの主君と仰ぎたいと思うが、張飛、おまえの考えはどうだ」\\


　訊くと、張飛も、手を打って、\\


「いや、それは拙者も考えていたところだ。いかにも、\ruby{兄}{けい}のいう通り、きめるならば、今ここで、神に\ruby{祷}{いの}るまえに、神へ誓ったほうがよい」\\


「玄徳様、ふたりの熱望です。ご承知くださるまいか」\\


　左右から詰めよられて、劉備玄徳は、黙然と考えていたが、\\


「待って下さい」\\


　と、二人の意気ごみを\ruby{抑}{おさ}え、なおややしばらく沈思してから、身を正していった。\\


「なるほど、自分は漢の宗室のゆかりの者で、そうした系図からいえば、主たる位置に坐るべきでしょうが、生来鈍愚、久しく\ruby{田舎}{でんしゃ}の\ruby{裡}{うち}にひそみ、まだなにも\ruby{各{\fontspec{jigmo}\char"303B}}{おのおの}の上に立って主君たるの修養も徳も積んでおりませぬ。どうか今しばらく待って下さい」\\


「待ってくれと仰っしゃるのは」\\


「実際に当って、徳を積み、身を修め、果たして主君となるの資才がありや否や、それを自身もあなたたちも見届けてから約束しても、遅くないと思われますから」\\


「いや。それはもう、われわれが見届けてあるところです」\\


「\ruby{左}{さ}はいえ、私はなお、\ruby{憚}{はばか}られます。――ではこうしましょう。君臣の誓いは、われわれが一国一城を持った上として、ここでは、三人義兄弟の約束を結んでおくことにして下さい。君臣となって後も、なお三人は、末永く義兄弟であるという約束をむしろ私はしておきたいのですが」\\


「うむ」\\


　関羽は、長い髯を持って、自分の顔を引っぱるように大きくうなずいた。\\


「結構だ。張飛、おぬしは」\\


「異論はない」\\


　改めて三名は、祭壇へ向って牛血と酒をそそぎ、ぬかずいて、天地の\ruby{神祇}{しんぎ}に\ruby{黙祷}{もくとう}をささげた。\\

\\

\subsubsection{三}

\\


　年齢からいえば、関羽がいちばん年上であり、次が劉備、その次が張飛という順になるのであるが、義約のうえの義兄弟だから年順をふむ必要はないとあって、「長兄には、どうか、あなたがなって下さい。それでないと、張飛の我ままにも、おさえが利きませんから」と、関羽がいった。\\


　張飛も、ともども、\\


「それは是非、そうありたい。いやだといっても、二人して、長兄長兄と\ruby{崇}{あが}めてしまうからいい」\\


　劉備は強いて\ruby{拒}{いな}まなかった。そこで三名は、\ruby{鼎座}{ていざ}して、将来の理想をのべ、\ruby{刎頸}{ふんけい}の\ruby{誓}{ちかい}い［＃「\ruby{誓}{ちかい}い」はママ］をかため、やがて壇をさがって桃下の卓を囲んだ。\\


「では、永く」\\


「変るまいぞ」\\


「変らじ」\\


　と、兄弟の杯を交わし、そして、三人一体、協力して国家に報じ、下万民の\ruby{塗炭}{とたん}の\ruby{苦}{く}を救うをもって、大丈夫の生涯とせんと申し合った。\\


　張飛は、すこし酔うてきたとみえて、声を大にし、杯を高く挙げて、\\


「ああ、こんな吉日はない。実に愉快だ。再び天にいう。われらここにあるの三名。同年同月同日に生まるるを\ruby{希}{ねが}わず、願わくば同年同月同日に死なん」\\


　と、呶鳴った。そして、\\


「飲もう。大いに、きょうは飲もう――ではありませんか」\\


　などと、劉備の杯へも、やたらに酒をついだ。そうかと思うと、自分の頭を、ひとりで叩きながら、「愉快だ。実に愉快だ」と、子供みたいにさけんだ。\\


　あまり彼の酒が、上機嫌に発しすぎる傾きが見えたので、関羽は、\\


「おいおい、張飛。今日のことを、そんなに歓喜してしまっては、先の歓びは、どうするのだ。今日は、われら三名の義盟ができただけで、大事の成功不成功は、これから後のことじゃないか。少し有頂天になるのが早すぎるぞ」と、たしなめた。\\


　だが、一たん上機嫌に昇ってしまうと、張飛の機嫌は、なかなか水をかけても\ruby{醒}{さ}めない。関羽の生真面目を、手を打って笑いながら、\\


「わはははは、今日かぎり、もう村夫子は廃業したはずじゃないか。お互いに軍人だ。これからは\ruby{天空海闊}{てんくうかいかつ}に、\ruby{豪放磊落}{ごうほうらいらく}に、武人らしく\ruby{交際}{つきあ}おうぜ。なあ長兄」\\


　と、劉備へも、すぐ\ruby{馴々}{なれなれ}といって、肩を抱いたりした。\\


「そうだ。そうだ」と、劉備玄徳は、にこにこ笑って、張飛のなすがままになっていた。\\


　張飛は、牛の如く飲み、馬のごとく喰ってから、\\


「そうそう。ここの席に、\ruby{劉母公}{りゅうぼこう}がいないという法はない。われわれ三人、兄弟の杯をしたからには、俺にとっても、尊敬すべきおっ母さんだ。――ひとつおっ母さんをこれへ連れてきて、乾杯しなおそう」\\


　急に、そんなことを云いだすと、張飛はふらふら母屋のほうへ馳けて行った。そしてやがて、劉母公を、無理に、自分の背中へ負って、ひょろひょろ戻ってきた。\\


「さあ、おっ母さんを、連れてきたぞ。どうだ、俺は親孝行だろう――さあおっ母さん、大いに祝って下さい。われわれ孝行息子が三人も揃いましたからね――いやこれは、独りおっ母さんにとって祝すべき孝行息子であるのみではない。支那の――国家にとってもだ、われわれこう三名は得がたい忠良息子ではあるまいか――そうだ、おっ母さんの孝行息子万歳、国家の忠良息子万歳っ」\\


　そしてやがて、こう三人の中では、酒に対しても一番の誠実息子たるその張飛が、まっ先に酔いつぶれて、桃花の下に大いびきで寝てしまい、夜露の降りるころまで、眼を醒まさなかった。\\

\\

\subsubsection{四}

\\


　大丈夫の誓いは結ばれた。しかし\ruby{徒手空拳}{としゅくうけん}とはまったくこの三人のことだった。しかも志は天下にある。\\


「さて、どうしたものか」\\


　翌日はもう酒を飲んでただ\ruby{快哉}{かいさい}をいっている日ではない。理想から実行へ、第一歩を踏みだす日である。\\


　朝飯を食べると、すぐその卓の上で、いかに実行へかかるかの問題がでた。\\


「どうかなるよ。男児が、しかも三人一体で、やろうとすれば」\\


　張飛は、理論家でない。また計画家でもない。\ruby{遮}{しゃ}二\ruby{無}{む}二、実行力に燃える\ruby{猪突邁進家}{ちょとつまいしんか}なのである。\\


「どうかなるって、ただ貴公のように、\ruby{力}{りき}んでばかりいたってどうもならん。まず、一郡の士を持たんとするには、一旗の兵がいる。一旗の兵を持つには、すくなくとも相当の軍費と、兵器と、馬とが必要だな」\\


　が、関羽は、常識家であった。二人のことばを\ruby{飽和}{ほうわ}すると、そこにちょうどよい情熱と常理との推進力が\ruby{醸}{かも}されてくる。\\


　劉備は、そのいずれへも、うなずきを与えて、\\


「そうです。こう三人の一念をもってすれば、必ず大事を成しうることは目に見えていますが、さし当って、兵隊です。――これをひとつ\ruby{募}{つの}りましょう」\\


「馬も、兵器も、金もなく、募りに応じてくれる者がありましょうか」\\


　関羽の憂いを、劉備はかろく微笑をもって打消し、\\


「いささか、自信があります。――というのは、実はこの楼桑村の内にも、平常からそれとなく、私が目にかけた、同憂の志を持っている青年たちが少々あります。――また近郷にわたって、\ruby{檄}{げき}を飛ばせば、おそらく今の時勢に、\ruby{鬱念}{うつねん}を感じている者もすくなくはありませんから、きっと、三十人や四十人の兵はすぐできるかと思います」\\


「なるほど」\\


「ですから、恐れいるが、関羽どのの筆で、ひとつ\ruby{檄文}{げきぶん}を起草して下さい。それを配るのは、私の知っている村の青年にやらせますから」\\


「いや、手前は、生来悪文の\ruby{質}{たち}ですから、ひとつそれは、劉長兄に起草していただこう」\\


「いいや、あなたは多年塾を持って、子弟を教育していたから、そういう子弟の気持を打つことは、よくお心得のはずだ。どうか書いて下さい」\\


　すると張飛がそばからいった。\\


「こら関羽、\ruby{怪}{け}しからんぞ」\\


「なにが怪しからん」\\


「長兄劉玄徳のことば、主命の如く\ruby{反}{そむ}くまいぞ、昨日、約束したばかりじゃないか」\\


「やあ、これは一本、張飛にやられたな、よし早速書こう」\\


　\ruby{飛檄}{ひげき}はでき上がった。\\


　なかなか名文である。荘重なる\ruby{慷慨}{こうがい}の気と、憂国の文字は、読む者を打たずにおかなかった。\\


　それが近郷へ飛ばされると、やがてのこと、劉玄徳の破れ家の門前には、毎日、七名十名ずつとわれこそ天下の豪傑たらんとする熱血の壮士が集まってきた。\\


　張飛は、門前へ出て、\\


「お前達は、われわれの檄を見て、兵隊になろうと望んできたのか」\\


　と、採用係の試験官になって、いちいち姓名や生国や、また、その志を質問した。\\


「そうです、\ruby{大人}{たいじん}がたのお名前と、義挙の趣旨に賛同して、旗下に馳せ参じてきた者どもです」\\


　壮士らは異口同音にいった。\\


「そうか、どれを見ても、たのもしい\ruby{面魂}{つらだましい}、早速、われわれの旗挙げに、加盟をゆるすが、しかしわれらの志は、黄巾賊の輩の如く、野盗掠奪を旨とするのとは違うぞ。天下の\ruby{塗炭}{とたん}を救い、害賊を討ち、国土に即した公権を確立し、やがては永遠の平和と民福を計るにある。分っておるかそこのところは！」\\


　張飛は、一場の訓示を垂れて、それからまた、次のように誓わせた。\\

\\

\subsubsection{五}

\\


「われわれの旗下に加盟するからには、即ち、われわれの奉じる軍律に服さねばならん。今、それを読み聞かすゆえ、謹んで\ruby{承}{うけたまわ}れ」\\


　張飛は、志願してきた壮士たちへいって、うやうやしく、\ruby{懐中}{ふところ}から一通を取出して、声高く読んだ。\\



一　\ruby{卒}{そつ}たる者は、将たる者に、絶大の服従と礼節を守る。\\


一　目前の利に惑わず、大志を遠大に備う。\\


一　一身を浅く思い、一世を深く思う。\\


一　\ruby{掠奪断首}{りゃくだつだんしゅ}。\\


一　\ruby{虐民極刑}{ぎゃくみんきょっけい}。\\


一　軍紀を\ruby{紊}{みだ}る行為一切死罪。\\



「わかったかっ」\\


　あまり厳粛なので、壮士たちも、しばらく黙っていたが、やがて、\\


「分りました」と、異口同音にいった。\\


「よし、しからば、今よりそれがしの部下として用いてやる。ただし、当分の間は給料もつかわさんぞ。また、食物その他も、お互いにある物を分けて喰い、いっさい不平を申すことならん」\\


　それでも、募りに応じてきた若者\ruby{輩}{ばら}は、元気に兵隊となって、劉備、関羽らの命に服した。\\


　四、五日のうちに、約七、八十人も集まった。望外な成功だと、関羽はいった。\\


　けれど、すぐ困りだしたのは食糧であった。ゆえに、一刻もはやく、戦争をしなければならない。\\


　\ruby{黄匪}{こうひ}の害に泣いている地方はたくさんある。まずその地方へ行って、黄巾賊を追っぱらうことだ。その後には、正しい税と食物とが収穫される。それは掠奪でない。\ruby{天禄}{てんろく}だ。\\


　するとある日。\\


「張将軍、張将軍。馬がたくさん通りますぞ、馬が」\\


　と、一人の部下が、ここの本陣へ馳せてきて注進した。\\


　何者か知らないが、何十頭という馬を\ruby{数珠}{じゅず}つなぎにひいて、この先の峠を越えてくる者があるという報告なのだ。\\


　馬と聞くと、張飛は、「そいつは何とか欲しいものだなあ」と正直にうなった。\\


　実際いま、\ruby{喉}{のど}から手の出るほど欲しい物は馬と金と兵器だった。だが、義挙の軍律というものを立てて部下にも示してあるので、「掠奪して来い」とは、命じられなかった。\\


　張飛は、奥へ行って、\\


「関羽、こういう報告があるが、なんとか、手に入れる\ruby{工夫}{くふう}はあるまいか。実に天の与えだと思うのだが」と、相談した。\\


　関羽は聞くと、\\


「よし、それでは、自分が行って、掛合ってみよう」と、部下数名をつれて、峠へ急いで行った。麓の近くで、その一行とぶつかった。物見の兵の注進に\ruby{過}{あやま}りなく、成程、四、五十頭もの馬匹をひいて、一隊の者がこっちへ下ってくる。近づいて見ると皆、商人ていの男なので、これならなんとか、話合いがつくと、関羽は得意の雄弁をふるうつもりで待ち構えていた。\\


　ここへきた馬\ruby{商人}{あきんど}の一隊の\ruby{頭}{かしら}は、\ruby{中山}{ちゅうざん}の豪商でひとりは\ruby{蘇双}{そそう}、ひとりは\ruby{張世平}{ちょうせいへい}という者だった。\\


　関羽は、それに着くと、自分ら三人が義軍を興すに至った、愛国の衷情をもって、切々訴えた。今にして、誰か、この覇業を建て、人天の正明をたださなければ、この世は永遠の闇黒であろうといった。支那大陸は、ついに、\ruby{胡北}{こほく}の武民に征服され終るであろうと嘆いた。\\


　張世平と蘇双の両人は、なにか小声で相談していたが、やがて、\\


「よく分りました。この五十頭の馬が、そういうことでお役に立てば満足です。差上げますからどうぞ曳いて行って下さい」と、意外にも、いさぎよく云った。\\

\\

\subsubsection{六}

\\


　いずれ易々とは承知しまい。最悪な場合までを関羽は考えていたのである。それが案外な返辞に、\\


「ほ。……いや\ruby{忝}{かたじ}けない。早速の快諾に、申しては失礼だが、利に\ruby{敏}{さと}い商人たるお身らが、どうしてそう一言のもとに、多くの\ruby{馬匹}{ばひつ}を無料でそれがしへ引渡すといわれたか」\\


　掛合いにきた目的は達しているのに、こう先方へ\ruby{要}{い}らざる念を押すのも妙なはなしだと思ったが、あまり不審なので、関羽はこう訊ねてみた。\\


　すると、張世平はいった。\\


「はははは。あまりさっぱりお渡しするといったので、かえってお疑いとみえますな。いやごもっともです。けれど手前は、第一にまず\ruby{大人}{たいじん}が悪人でないことを認めました。第二に、ご計画の義兵を挙げることは、すこぶる\ruby{時宜}{じぎ}をえておると存じます。第三は、あなた方のお力をもって、自分らの恨みをはらしていただきたいと思ったからです」\\


「恨みとは」\\


「黄巾賊の大将張角一門の暴政に対する恨みでございます。手前も以前は中山で一といって二と下らない豪商といわれた者ですが、かの地方もご承知の通り黄匪の\ruby{蹂躙}{じゅうりん}にあって秩序は破壊され、財産は掠奪され、町に少女の影を見ず、\ruby{家苑}{かえん}の\ruby{小禽}{ことり}すら\ruby{啼}{な}かなくなってしまいました。――手前の店なども一物もなく没収され、あげくの果てに、妻も娘も、暴兵にさらわれてしまったのです」\\


「むむ。なるほど」\\


「で、甥の\ruby{蘇双}{そそう}と二人して、馬商人に身を落し、市から馬匹を購入して、北国へ売りに行こうとしたのですが、途中まで参ると、北辺の山岳にも、黄賊が道をふさいで、旅人の持物を奪い、虐殺をほしいままにしておるとのことに、むなしくまた、この群馬をひいて立ち帰ってきたわけです。南へ行くも賊国、北へおもむくも賊国、こうして馬とともに漂泊しているうちには、ついに賊に生命まで共に奪われてしまうのは知れきっています。恨みのある賊の手に武力となる馬匹を与えるよりも、貴下の如きお志を抱く人に、進上申したほうが、はるかに意味のあることなんです。よろこんで手前がお渡しする気持というのは、そんなわけでございます」\\


「やあ、そうか」\\


　関羽の疑問も氷解して、\\


「では、楼桑村まで、馬をひいて一緒に来てくれないか。われわれの盟主と仰ぐ劉玄徳と仰っしゃる人にひきあわせよう」\\


「おねがい致します。手前も根からの商人ですから、以上申上げたような理由でもって、無料で馬匹を進上しましても、やはりそこはまだ正直、利益のことを考えておりますからな」\\


「いや、玄徳様へお目にかかっても、ただ今のところ、代金はお下げになるわけにはゆかぬぞ」\\


「そんな目先のことではありません。遠い将来でよろしいので。……はい。もしあなたがたが大事を成しとげて、一国を取り、十州二十州を平らげ、あわよくば天下に号令なさろうという筋書きのとおりに行ったらば、私へも充分に、利をつけて、今日の馬代金を払っていただきたいのでございます。私は、あなたの計画を聞いて、これがあなたがたの夢ではなく、わたしども民衆が待っていたものであるという点から、きっと成功するものと信じております。ですから、今日この処分に困っている馬を使っていただくのは、商人として、手前にも遠大な利殖の方法を見つけたわけで、まったくこんなよろこばしいことはありません」\\


　張世平は、そういって、甥の蘇双と共に、関羽に案内されてついて行ったが、その途中でも、関羽へ対して、こう意見を述べた。\\


「事を計るうえは、人物はお揃いでございましょうし、馬もこれで整いました。これで一体、あなた方のご計画の内輪には、よく経済を切りまわして糧食兵費の内助の役目をする算数の達識が控えているのでございますか。\ruby{算盤}{そろばん}というものも、充分お考えのうえでこのお仕事にかかっておいででございますか？」\\

\\

\subsubsection{七}

\\


　張世平に、そう指摘されてみると、関羽は、自分らの仲間に、大きな\ruby{欠陥}{けっかん}のあるのを見いだした。\\


　経営ということであった。\\


　自分はもとより、張飛にも、劉玄徳にも、経済的な観念は至ってない。武人\ruby{銭}{ぜに}を愛さずといったような思想がはなはだ古くから頭の隅にある。経済といえばむしろ卑しみ、銭といえば横を向くをもって清廉の士とする風が高い。一個の人格にはそれも高風と仰ぎうるが、国家の大計となればそれでは不具を意味する。\\


　一軍を持てばすでに経営を思わねばならぬ。武力ばかりでふくらもうとする軍は暴軍に化しやすい。古来、理想はあっても、そのため、暴軍と\ruby{堕}{だ}し、乱賊と終った者、史上決してすくなくない。\\


「いや、いいことを聞かしてくれた。劉玄徳様にも、大いにそのへんのことをはなして貰いたいものだ」\\


　関羽は、正直、教えられた気がしたのである。一商人のことばといえども、これは将来の大切な問題だと考えついた。\\


　やがて、楼桑村に着く。\\


　関羽はすぐ張世平と蘇双のふたりを、劉玄徳の前につれて来た。もちろん、玄徳も張飛も、張の好意を聞いて非常によろこんだ。\\


　張は五十頭の馬匹を、無償で提供するばかりでなく、玄徳に会ってから玄徳の人物をさらに見込んで、それに加うるに、駿馬に積んでいた鉄一千\ruby{斤}{きん}と、百\ruby{反}{たん}の獣皮織物と、金銀五百両を挙げてみな、「どうか、軍用の費に」と、献上した。\\


　その際も、張はいった。\\


「最前も、みちみち、申しました通り、手前はどこまでも、利を道とする商人です。武人に武道あり、聖賢に文道あるごとく、商人にも利道があります。ご献納申しても、手前はこれをもって、義心とは誇りません。その代り、今日さし上げた馬匹金銀が、十年後、三十年後には、莫大な利を生むことを望みます。――ただその利は、自分一個で\ruby{飽慾}{ほうよく}しようとは致しません。困苦の底にいる万民にお\ruby{頒}{わか}ちください。それが私の希望であり、また私の商魂と申すものでございます」\\


　玄徳や関羽は、彼の言を聞いて大いに感じ、どうかしてこの人物を自分らの仲間へ留めおきたいと考えたが、張は、\\


「いやどうも私は臆病者で、とても戦争なさるあなた方の中にいる勇気はございません。なにかまた、お役に立つ時には出てきますから」といって、\ruby{倉皇}{そうこう}、何処ともなく立ち去ってしまった。\\


　千斤の鉄、百反の\ruby{織皮}{しょくひ}、五百両の金銀、思いがけない軍費を獲て、玄徳以下三人は、\\


「これぞ天のご援助」\\


　と、いやが上にも、心は奮い立った。\\


　早速、近郷の\ruby{鍛冶工}{かじこう}をよんできて、張飛は、一丈何尺という\ruby{蛇矛}{じゃぼこ}を\ruby{鍛}{う}ってくれと注文し、関羽は重さ何十斤という\ruby{偃月刀}{えんげつとう}を\ruby{鍛}{きた}えさせた。\\


　雑兵の鉄甲、\ruby{{\fontspec{jigmo}\char"76D4}}{かぶと}、槍、刀などもあわせて\ruby{誂}{あつら}え、それも日ならずしてできてきた。\\

\ruby{日月}{じつげつ}の\ruby{旗幟}{きし}。\\


飛龍の\ruby{幡}{ばん}。\\

\ruby{鞍}{くら}、\ruby{鏃}{やじり}。\\



　軍装はまず整った。\\


　その頃ようやく人数も二百人ばかりになった。\\


　もとより天下に臨むには足りない急仕立ての一小軍でしかなかったが、張飛の教練と、関羽の軍律と、劉玄徳の徳望とは、一卒にまでよく行きわたって、あたかも一箇の体のように、二百の兵は\ruby{挙手踏足}{きょしゅとうそく}、一音に動いた。\\


「では。――おっ母さん。行って参ります」\\


　劉玄徳は、ある日、武装して母にこう暇を告げた。\\


　兵馬は、粛々、彼の郷土から立って行った。劉玄徳の母は、それを桑の木の下からいつまでも見送っていた。泣くまいとしている眼が湯の泉のようになっていた。\\

\\

\subsection{転戦}

\\

\subsubsection{一}

\\


　それより前に、関羽は、玄徳の書をたずさえて、\ruby{幽州{\fontspec{jigmo}\char"6DBF}郡}{ゆうしゅうたくぐん}（河北省・{\fontspec{jigmo}\char"6DBF}県）の大守\ruby{劉焉}{りゅうえん}のもとへ使いしていた。\\


　太守劉焉は、何事かと、関羽を城館に入れて、\ruby{庁堂}{ちょうどう}で接見した。\\


　関羽は、礼をほどこして後、\\


「太守には今、士を四方に求めらるると聞く。果して然りや」\\


　と、訊ねた。\\


　関羽の威風は、堂々たるものであった。劉焉は、一見して、これ尋常人に非ずと思ったので、その\ruby{不遜}{ふそん}を\ruby{咎}{とが}めず、\\


「然り。諸所の駅路に高札を建てしめ、士を\ruby{募}{つの}ること急なり。\ruby{卿}{けい}もまた、\ruby{檄}{げき}に応じてきたれる偉丈夫なるか」と、いった。\\


　そこで関羽は、\\


「さん候。この国、黄賊の大軍に\ruby{攻蝕}{こうしょく}せらるること久しく、太守の軍、連年に\ruby{疲敗}{ひはい}し給い、各地の民倉は、挙げて賊の毒手にまかせ、百姓\ruby{蒼生}{そうせい}みな国主の無力と、賊の暴状に\ruby{哭}{な}かぬはなしと承る」\\


　あえて、媚びずおそれず、こう正直にいってからさらに重ねて、\\


「われら恩を久しく領下にうけて、この\ruby{秋}{とき}をむなしく\ruby{逸人}{いつじん}として\ruby{草廬}{そうろ}に\ruby{閑}{かん}を\ruby{偸}{ぬす}むをいさぎよしとせず、同志張飛その他二百余の有為の\ruby{輩}{ともがら}と団結して、劉玄徳を盟主と仰ぎ、太守の軍に入って、いささか報国の義をささげんとする者でござる。太守寛大、よくわれらの義心の兵を加え給うや否や」\\


　と、述べ、終りに、玄徳の手書を出して、一読を乞うた。\\


　\ruby{劉焉}{りゅうえん}は、聞くと、\\


「この\ruby{秋}{とき}にして、卿ら赤心の豪傑ら、劉焉の微力に援助せんとして訪ねらる、まさに、天祐のことともいうべきである。なんぞ、拒むの理があろうか。城門の\ruby{塵}{ちり}を掃き、客館に\ruby{旗飾}{きしょく}をほどこして、参会の日を待つであろう」\\


　といって、非常な歓びようであった。\\


「では、何月何日に、ご城下まで兵を\ruby{率}{ひき}いて参らん」と、約束して関羽は立帰ったのであるが、その折、はなしのついでに、義弟の張飛が、先ごろ、楼桑村の附近や\ruby{市}{いち}の関門などで、事の間違いから、太守の部下たる捕吏や役人などを殺傷したが、どうかその罪はゆるされたいと、一口ことわっておいたのである。\\


　そのせいか、あれっきり、\ruby{市}{いち}の関門からも、捕吏の人数はやって来なかった。いやそれのみか、あらかじめ、太守のほうから命令があったとみえ、劉玄徳以下の三傑に、二百余の郷兵が、突然、楼桑村から{\fontspec{jigmo}\char"6DBF}郡の府城へ向って出発する際には、関門のうえに小旗を立て、守備兵や役人は整列して、その行を鄭重に見送った。\\


　それと、眼をみはったのは、玄徳や張飛の顔を見知っている市の雑民たちで、\\


「やあ、先に行く大将は、\ruby{蓆売}{むしろ}りの劉さんじゃないか」\\


「そのそばに、馬にのって威張って行くのは、よく\ruby{猪}{いのこ}の肉を売りに出ていた呑んだくれの浪人者だぞ」\\


「なるほど。張だ、張だ」\\


「あの肉売りに、わしは酒代の貸しがあるんだが、弱ったなあ」\\


　などと群集のあいだから嘆声をもらして、見送っている酒売りもあった。\\


　義軍はやがて、{\fontspec{jigmo}\char"6DBF}郡の府に到着した。道々、風を慕って、日月の旗下に馳せ参じる者もあったりして、府城の大市へ着いた時は、総勢五百をかぞえられた。\\


　太守は、直ちに、玄徳らの三将を迎えて、その夜は、居館で歓迎の宴を張った。\\

\\

\subsubsection{二}

\\


　大将玄徳に会ってみるとまだ年も\ruby{二十歳台}{はたちだい}の青年であるが、\ruby{寡言沈厚}{かげんちんこう}のうちに、どこか大器の風さえうかがえるので、太守\ruby{劉焉}{りゅうえん}は、大いに好遇に努めた。\\


　なお、素姓を問えば、漢室の宗親にして、\ruby{中山靖王}{ちゅうざんせいおう}の\ruby{裔孫}{えいそん}とのことに、\\


「さもあらん」と、劉焉はうなずくことしきりでなおさら、親しみを改め、左右の関、張両将をあわせて、心から\ruby{敬}{うやま}いもした。\\


　折ふし。\\


　青州\ruby{大興山}{たいこうざん}の附近一帯（山東省\ruby{済南}{さいなん}の東）に跳梁している黄巾賊五万以上といわれる勢力に対して太守劉焉は、家臣の\ruby{校尉}{こうい}\ruby{鄒靖}{すうせい}を将として、大軍を附与し、にわかに、それへ馳け向わせた。\\


　関羽と、張飛は、それを知るとすぐ、玄徳へ向って、「人の歓待は、\ruby{冷}{さ}めやすいものでござる。歓宴長くとどまるべからずです。手はじめの出陣、進んでご加勢にお加わりなさい」と、すすめた。\\


　玄徳は、「自分もそう考えていたところだ。早速、太守へ進言しよう」と、劉焉に会って、その旨を申し出ると劉焉もよろこんで、\ruby{校尉}{こうい}\ruby{鄒靖}{すうせい}の先陣に参加することをゆるした。\\


　玄徳の軍五百余騎は、\ruby{初陣}{ういじん}とあって意気すでに天をのみ、日ならずして大興山の麓へ押しよせてみたところ、賊の五万は、嶮に\ruby{拠}{よ}って、利戦を策し、山の\ruby{襞}{ひだ}や谷あいへ\ruby{虱}{しらみ}のごとく長期の陣を備えていた。\\


　時、この地方の雨期をすぎて、すでに初夏の緑草豊かであった。\\


　合戦長きにわたらんか、賊は、地の利を得て、奇襲縦横にふるまい、諸州の\ruby{黄匪}{こうひ}、連絡をとって、いっせいに後路を断ち、征途の味方は重囲のうちに\ruby{殲滅}{せんめつ}の\ruby{厄}{やく}にあわんもはかりがたい。\\


　玄徳は、そう考えたので、\\


「いかに張飛、関羽。太守劉焉をはじめ、校尉鄒靖も、われらの手なみいかにと、その実力を見んとしておるに違いない。すでに、味方の先鋒たる以上、いたずらに、対峙して、味方に長陣の不利を招くべからずである。挺身、賊の陣近く斬入って、一気に戦いを決せんと思うがどうであろう」\\


　二人へ、計ると、「それこそ、同意」と、すぐ五百余騎を、鳥雲に備え立て、山麓まぢかへ迫ってからにわかに\ruby{鼓}{こ}を鳴らし\ruby{諸声}{もろごえ}あげて決戦を挑んだ。\\


　賊は、山の中腹から、鉄弓を射、\ruby{弩}{ど}をつるべ撃ちして、容易に動かなかったが、\\


「寄手は、たかのしれた小勢のうえに、国主の正規兵とはみえぬぞ、どこかそこらから狩り集めてきた烏合の雑軍。みなごろしにしてしまえ」\\


　賊の副将\ruby{{\fontspec{jigmo}\char"9127}茂}{とうも}という者、こう号令を下すや否や、\ruby{柵}{さく}を開いて、山上から逆落しに騎馬で馳けおりて来、\\


「やあやあ、\ruby{稗粕}{ひえかす}をなめて生きる、あわれな郷軍の百姓兵ども。官軍の名にまどわされて死骸の堤を築きに来りしか。愚かなる権力の楯につかわるるを止めよ。汝ら、槍をすて、馬を献じ、降を乞うなれば、わが将、\ruby{大方}{だいほう}\ruby{程遠志}{ていえんし}どのに申しあげて、\ruby{黄巾}{こうきん}をたまわり、肉食させて、世を楽しみ、その\ruby{痩骨}{やせぼね}を肥えさすであろう。否といわば、即座に包囲殲滅せん。耳あらば聞け、口あらば答えよ。――いかに、いかに！」と、呼ばわった。\\


　すると、寄手の陣頭より、おうと答えて、劉玄徳、左右に関羽、張飛をしたがえて、白馬を緑野の中央へすすめて来た。\\

\\

\subsubsection{三}

\\


「\ruby{推参}{すいさん}なり、\ruby{野鼠}{やそ}の将」\\


　玄徳は、賊将程遠志の前に駒を止めて、彼のうしろにひしめく黄巾賊の大軍へも轟けとばかりいった。\\


「天地ひらけて以来、まだ獣族の長く栄えたる例はなし。たとい、一時は人政を\ruby{紊}{みだ}し、暴力をもって権を奪うも、末路は野鼠の白骨と変るなからん。――醒めよ、われは、日月の\ruby{幡}{はた}を高くかかげ、暗黒の世に光明をもたらし、邪を\ruby{退}{しりぞ}け、\ruby{正}{せい}を明らかにするの義軍、いたずらに立ち向って、\ruby{生命}{いのち}をむだに落すな」\\


　聞くと、\ruby{程遠志}{ていえんし}は声をあげて、大笑し、\\


「白昼の大寝言、近ごろおもしろい。醒めよとは、うぬらのこと。いで」\\


　と、重さ八十斤と称する青龍刀をひッさげ、駒首おどらせて玄徳へかかってきた。\\


　玄徳はもとより武力の猛将ではない。泥土をあげて、\ruby{蹄}{ひづめ}を後ろへ返す。その間へ、待ちかまえていた張飛が、\\


「この下郎っ」\\


　おめきながら割って入り、先ごろ\ruby{鍛}{う}たせたばかりの丈余の\ruby{蛇矛}{じゃぼこ}――\ruby{牙形}{きばがた}の\ruby{大矛}{おおぼこ}を先につけた長柄を舞わして、賊将程遠志の\ruby{{\fontspec{jigmo}\char"76D4}}{かぶと}の鉢金から馬の背骨に至るまで斬り下げた。\\


「やあ、おのれよくも」\\


　賊の副将\ruby{{\fontspec{jigmo}\char"9127}茂}{とうも}は、乱れ立つ兵を励ましながら、逃げる玄徳を目がけて追いかけると、関羽が早くも騎馬をよせて、\\


「\ruby{豎子}{じゅし}っ、なんぞ死を急ぐ」\\


　虚空に鳴る\ruby{偃月刀}{えんげつとう}の一\ruby{揮}{き}、血けむり呼んで、人馬ともに、関羽の\ruby{葬}{ほうむ}るところとなった。\\


　賊の二将が打たれたので、残余の\ruby{鼠兵}{そへい}は、あわて乱れて、山谷のうちへ逃げこんでゆく。それを、追って打ち、包んでは\ruby{殲滅}{せんめつ}して賊の首を挙げること一万余。降人は容れて、部隊にゆるし、首級は村里の辻に\ruby{梟}{か}けならべて、\\


　――\ruby{天誅}{てんちゅう}はかくの如し。\\


　と、武威をしめした。\\


「\ruby{幸先}{さいさき}はいいぞ」\\


　張飛は、関羽にいった。\\


「なあ兄貴、このぶんなら、五十州や百州の賊軍ぐらいは、半歳のまに片づいてしまうだろう。天下はまたたく間に、俺たちの\ruby{旗幟}{きし}によって、日月照々だ。安民楽土の世となるにきまっている。愉快だな。――しかし、戦争がそう早くなくなるのがさびしいが」\\


「ばかをいえ」\\


　関羽は、首をふった。\\


「世の中は、そう簡単でないよ。いつも\ruby{戦}{いくさ}はこんな調子だと思うと、大まちがいだぞ」\\


　大興山を後にして、一同はやがて幽州へ凱旋の\ruby{轡}{くつわ}をならべた。\\


　太守\ruby{劉焉}{りゅうえん}は、五百人の楽人に勝利の譜を吹奏させ、城門に旗の列を植えて、自身、凱旋軍を出迎えた。\\


　ところへ。\\


　軍馬のやすむいとまもなく、青州の城下（山東省済南の東・黄河口）から早馬が来て、\\


「大変です。すぐ援軍のご出馬を乞う」と、ある。\\


「何事か」と、劉焉が、使いのもたらした\ruby{牒文}{ちょうぶん}をひらいてみると、\\



当地方ノ黄巾ノ\ruby{賊徒等}{ゾクトラ}県郡ニ\ruby{蜂起}{ホウキ}シテ雲集シ青州ノ城囲マレ終ンヌ\ruby{落焼}{ラクショウ}ノ運命スデニ急ナリタダ友軍ノ来援ヲ待ツ\\

青州\ruby{太守}{タイシュ}\ruby{{\fontspec{jigmo}\char"9F94}景}{キョウケイ}


　と、あった。\\


　玄徳は、また進んで、\\


「願わくば\ruby{行}{ゆ}いて\ruby{援}{たす}けん」\\


　と申し出たので、太守劉焉はよろこんで、\ruby{校尉}{こうい}\ruby{鄒靖}{すうせい}の五千余騎に加えて、玄徳の義軍にその先鋒を依嘱した。\\

\\

\subsubsection{四}

\\


　時はすでに夏だった。\\


　青州の野についてみると、賊数万の軍は、すべて黄の旗と、\ruby{八卦}{はっけ}の文を\ruby{証}{しるし}とした\ruby{幡}{はん}をかざして、その勢い、天日をも\ruby{侮}{あなど}っていた。\\


「なにほどのことがあろう」と、玄徳も、先頃の初陣で、難なく勝った手ごころから、五百余騎の先鋒で、当ってみたが、結果は大失敗だった。\\


　一敗地にまみれて、あやうく全滅をまぬがれ、三十里も退いた。\\


「これはだいぶ強い」\\


　玄徳は、関羽へ計った。\\


　関羽は、\\


「\ruby{寡}{か}をもって、衆を破るには、兵法によるしかありません」と一策を献じた。\\


　玄徳は、よく人の言を用いた。そこで、総大将の\ruby{鄒靖}{すうせい}の陣へ、使いを立て、謀事をしめしあわせて、作戦を立て直した。\\


　まず、総軍のうち、関羽は約千の兵をひっさげて、右翼となり、張飛も同数の兵力を持って、丘の陰にひそんだ。\\


　本軍の\ruby{鄒靖}{すうせい}と玄徳とは、正面からすすんで、敵の主勢力へ、総攻撃の態を示し、頃あいを計って、わざと、潮のごとく逃げ乱れた。\\


「追えや」\\


「討てや」\\


　と、図にのって、賊の大軍は、陣形もなく追撃してきた。\\


「よしっ」\\


　玄徳が、駒を返して、充分誘導してきた敵へ当り始めた時、丘陵の陰や、\ruby{曠野}{こうや}の\ruby{黍}{きび}の中から、夕立雲のように湧いて出た関羽、張飛の両軍が、敵の主勢力を、完全にふくろづつみにして、みなごろしにかかった。\\


　太陽は、血に煙った。\\


　草も馬の尾も、血のかからない物はなかった。\\


「それっ、今だ」\\


　逃げる賊軍を追って、そのまま味方は青州の城下まで迫った。\\


　青州の城兵は、\\


　――援軍来る！\\


　と知ると、城門をひらいて、討って出た。なだれを打って、逃げてきた賊軍は、城下に火を放ち、自分のつけた炎を墓場として、ほとんど、自滅するかのような敗亡を遂げてしまった。\\


　青州の太守\ruby{{\fontspec{jigmo}\char"9F94}景}{きょうけい}は、\\


「もし、\ruby{卿}{けい}らの来援がなければ、この城は、すでに今日は賊徒の享楽の宴会場になっていたであろう」\\


　と、人々を重く賞して、三日三晩は、夜も日も、歓呼の楽器と万歳の声にみちあふれていた。\\


　鄒靖は、軍を収めて、\\


「もはや、お\ruby{暇}{いとま}せん」\\


　と、幽州へ引揚げて行ったが、その際、劉玄徳は、鄒靖に向って、\\


「ずっと以前――私の少年の頃ですが、郷里の楼桑村に来て、しばらくかくれていた\ruby{盧植}{ろしょく}という人物がありました。私は、その盧植先生について、初めて文を学び、兵法を説き教えられたのです。その後先生はどうしたかと、時おり、思い出すのでしたが、近頃うわさに聞けば、盧植先生は官に仕えて、\ruby{中郎将}{ちゅうろうしょう}に任ぜられ、今では勅令をうけて、遠く\ruby{広宗}{こうそう}（山東省）の野に戦っていると聞きます。――しかしそこの賊徒は、\ruby{黄匪}{こうひ}の首領張角将軍直属の正規兵だということですから、さだめしご苦戦と察しられるので、これから行って、師弟の旧恩、いささかご加勢してあげたいと思うのです」と、心のうちをもらした。\\


　そして、自分はこれから、広宗の征野へ、旧師の軍を援けにおもむくから、幽州の城下へ帰ったら、どうか、その旨を、悪しからず太守へお伝えねがいたいと、伝言を頼んだ。\\


　もとより義軍であるから、鄒靖も引止めはしない。\\


「しからば、貴下の手勢のみ率いて、兵糧そのほかの\ruby{賄}{まかない}、心のままにし給え」\\


　と、武人らしく、あっさりいって別れた。\\

\\

\subsubsection{五}

\\


　\ruby{討匪}{とうひ}将軍の\ruby{印綬}{いんじゅ}をおびて、遠く\ruby{洛陽}{らくよう}の王府から、黄河口の広宗の\ruby{野}{や}に下り、五万の官軍を率いて軍務についていた中郎将\ruby{盧植}{ろしょく}は、\\


「なに。\ruby{劉備玄徳}{りゅうびげんとく}という者がわしを訪ねてきたと？　……はてな、劉、玄徳、誰だろう」\\


　しきりに首をひねっていたが、まだ思い出せない\ruby{容子}{ようす}だった。\\


　戦地といっても、さすが漢朝の征旗を奉じてきている軍の本営だけに、将軍の室は、大きな寺院の中央を占め、境内から四門の外郭一帯にかけて、駐屯している兵馬の勢威は物々しいものであった。\\


「はっ。――確かに、劉備玄徳と仰っしゃって、将軍にお目にかかりたいと申して来ました」\\


　外門から取次いできた一人の兵はそういって、盧将軍の前に、直立の姿勢をとっていた。\\


「一人か」\\


「いいえ、五百人も連れてであります」\\


「五百人」\\


　\ruby{唖然}{あぜん}とした顔つきで、\\


「じゃあ、その玄徳とやらは、そんなにも自分の手勢をつれて来たのか」\\


「さようです。関羽、張飛、という二名の部将を従えて、お若いようですが、立派な人物です」\\


「はてなあ？」\\


　なおさら、思い当らない容子であったが、取次ぎの兵が、\\


「申し残しました。その仁は、\ruby{{\fontspec{jigmo}\char"6DBF}}{たく}県楼桑村の者で、将軍がそこに隠遁されていた時代に、読み書きのお教えをうけたことがあるとかいっておりました」\\


「ああ！　では\ruby{蓆}{むしろ}売りの劉少年かもしれない。いや、そういえば、あれからもう十年以上も経っておるから、よい若人になっている年頃だろう」\\


　盧植は、にわかに、なつかしく思ったとみえ、すぐ通せと命令した。もちろん、連れている兵は外門にとめ、二人の部将は、内部の\ruby{廂}{ひさし}まで入ることを許してである。\\


　やがて玄徳は通った。\\


　盧植は、ひと目見て、\\


「おお、やはりお前だったか。変ったのう」と、驚いた目をした。\\


「先生にも、その後は、\ruby{赫々}{かっかく}と洛陽にご武名の聞え高く、蔭ながらよろこんでおりました」\\


　玄徳は、そういって、盧植の\ruby{沓}{くつ}の前に退がり、昔に変らぬ師礼をとった。\\


　そして彼は、自分の素志をのべた上、願わくば、旧師の征軍に加わって、朝旗のもとに報国の働きを尽したいといった。\\


「よく来てくれた。少年時代の小さな師恩を思い出して、わざわざ援軍に来てくれたとは、近頃うれしいことだ。その心もちはすでに朝臣であり、国を愛する士の持つところのものだ。わが軍に参加して、大いに勲功をたててくれ」\\


　玄徳は、参戦をゆるされて、約二ヵ月ほど、盧植の軍を援けていたが、実戦に当ってみると、賊のほうが、三倍も多い大軍を擁しているし、兵の強さも、比較にならないほど、賊のほうが優勢だった。\\


　そのため、官軍のほうが、かえって守勢になり、いたずらに、滞陣の月日ばかり長びいていたのだった。\\


「軍器は立派だし、服装も剣も華やかだが、洛陽の官兵は、どうも戦意がない。都に残している女房子供のことだの、うまい酒だの、そんなことばかり思い出しているらしい」\\


　張飛は、時々、そんな不平を鳴らして、\\


「長兄。こんな軍にまじっていると、われわれまでが、だらけてしまう。去って、ほかに大丈夫の戦う意義のある戦場を見つけましょう」\\


　と、玄徳へいったが、玄徳は、師を歓ばせておきながら、師へ酬いることもなく去る法はないといって、きかなかった。\\


　そのうちに、\ruby{盧植}{ろしょく}のほうから、折入って、軍機にわたる一つの相談がもちかけられた。\\

\\

\subsubsection{六}

\\


　盧植がいうには、\\


　――そもそもこの地方は、\ruby{嶮岨}{けんそ}が多くて、守る賊軍に利があり、一気に破ろうとすれば、多大に味方を損じるので、心ならずも、こうして長期戦を張って、長陣をしている\ruby{理}{わけ}であるが、折入って、貴下に頼みたいというのは、賊の総大将張角の弟で\ruby{張宝}{ちょうほう}・\ruby{張梁}{ちょうりょう}のふたりは目下、\ruby{潁川}{えいせん}（河南省・許昌）のほうで暴威をふるっている。\\


　その方面へは、やはり洛陽の朝命をうけて、\ruby{皇甫嵩}{こうほすう}・\ruby{朱雋}{しゅしゅん}の二将軍が、官軍を率いて討伐に向っている。\\


　ここでも勝敗決せず、官軍は苦戦しているが、わが広宗の地よりも、戦うに益が多い。ひとつ貴下の手勢をもって、急に援軍におもむいてもらえまいか。\\


　賊の張梁・張宝の二軍が敗れたりと聞えれば、自然、広宗の賊軍も、戦意を喪失し、退路を断たれることをおそれて、潰走し始めることと思う。\\


「玄徳殿。行ってもらえまいか」\\


　盧植の相談であった。\\


「承知しました」\\


　玄徳は、もとより義をもって、旧師を援けにきたので、その旧師の頼みを、すげなく\ruby{拒}{こば}む気にはなれなかった。\\


　即刻、軍旅の支度をした。\\


　手勢五百に、盧植からつけてくれた千余の兵を加え、総勢千五百ばかりで、\ruby{潁川}{えいせん}の地へ急いだ。\\


　陣地へ着くと、さっそく官軍の将、\ruby{朱雋}{しゅしゅん}に会って、盧植の牒文を示し、\\


「お手伝いに参った」とあいさつすると、\\


「ははあ。何処で\ruby{雇}{やと}われた雑軍だな」と、朱雋は、しごく冷淡な応対だった。\\


　そして、玄徳へ、\\


「まあ、せいぜい働き給え。軍功さえ立てれば、正規の官軍に編入されもするし、貴公らにも、戦後、何か地方の小吏ぐらいな役目は仰せつかるから」\\


　などともいった。\\


　張飛は、\\


「ばかにしておる」\\


　と怒ったが、玄徳や関羽でなだめて、前線の陣地へ出た。\\


　食糧でも、軍務でも、また応対でも、冷遇はするが、与えられた戦場は、もっとも強力な敵の正面で、官軍の兵が、手をやいているところだ。\\


　地勢を見るに、ここは広宗地方とちがって、いちめんの原野と湖沼だった。\\


　敵は、折からの、背丈の高い夏草や\ruby{野黍}{のきび}のあいだに、虫のようにかくれて、時々、猛烈な奇襲をしてきた。\\


「さらば。一策がある」\\


　玄徳は、関羽と張飛に、自分の考えを告げてみた。\\


「名案です。長兄は、そもそも、いつのまにそんなに、孫呉の兵を\ruby{会得}{えとく}しておられたんですか」\\


　と、二人とも感心した。\\


　その晩、二\ruby{更}{こう}の頃。\\


　一部の兵力を、迂回させて、敵のうしろに廻し、張飛、関羽らは、真っ暗な野を這って、敵陣へ近づいた。\\


　そして、用意の物に、一斉に火を点じると、\\


「わあっ」\\


　と、\ruby{鬨}{とき}の声をあげて、炎の波のように、攻めこんだ。\\


　かねて、兵一名に、十\ruby{把}{ぱ}ずつの\ruby{松明}{たいまつ}を負わせ、それに火をつけて、なだれこんだのである。\\


　寝ごみを衝かれ、不意を襲われて、右往左往、あわて廻る敵陣の中へ、投げ松明の光は、花火のように舞い飛んだ。\\


　草は燃え、兵舎は焼け、逃げくずれる賊兵の軍衣にも、火がついていないのはなかった。\\


　すると彼方から、一\ruby{彪}{ぴょう}の軍馬が、燃えさかる草の火を蹴って進んできた。見れば、全軍みな\ruby{紅}{くれない}の旗をさし、真っ先に立った一名の英雄も、\ruby{兜}{かぶと}、\ruby{鎧}{よろい}、剣装、馬鞍、すべて火よりも赤い姿をしていた。\\

\\

\subsubsection{七}

\\


「やよ、それに来る豪傑。貴軍はそも、敵か味方か」\\


　玄徳のそばから大音で、関羽が彼方へ向って云った。\\


　先でも、玄徳たちを、\\


「官軍か賊軍か？」と疑っていたように、ぴたと一軍の前進を停めて、\\


「これは洛陽より南下した五千騎の官軍である。汝らこそ、黄匪に非ずや」\\


　と、呶鳴り返してきた。\\


　聞くと、玄徳は左将軍関羽、右将軍張飛だけを両側に従えて、兵を後方に残したまま数百歩駒をすすめ、\\


「戦場とて、失礼をいたした。それがしは\ruby{{\fontspec{jigmo}\char"6DBF}県}{たくけん}楼桑村の\ruby{草莽}{そうもう}より起って、いささか奉公を志し、討賊の戦場に参加しておる義軍の将、劉備玄徳という者です。それにおいである豪傑は、そも何ぴとなりや。願わくばご尊名をうかがいたい」\\


　いうと、紅の旗、紅の鎧、紅の鞍にまたがっている人物は、玄徳の会釈を、馬上でうけながら微笑をたたえ、「ごていねいな挨拶。それへ参って申さん」と、\ruby{赤夜叉}{あかやしゃ}の如く、すべて赤く\ruby{鎧}{よろ}った旗本七騎につつまれて、玄徳の間近まで馬をすすめて来た。\\


　近々と、その人物を見れば。\\


　年はまだ若い。肉薄く色白く、\ruby{細眼長髯}{さいがんちょうぜん}、\ruby{胆量}{たんりょう}人にこえ、その眸には、智謀はかり知れないものが見えた。\\


　声静かに、名乗っていう。\\


「われは\ruby{沛国{\fontspec{jigmo}\char"8B59}郡}{はいこくしょうぐん}（安徽省・毫県）の生れで、\ruby{曹操}{そうそう}\ruby{字}{あざな}は\ruby{孟徳}{もうとく}、\ruby{小字}{こあざな}は\ruby{阿瞞}{あまん}、また\ruby{吉利}{きつり}ともいう者です。すなわち漢の\ruby{相国}{しょうこく}\ruby{曹参}{そうさん}より二十四代の\ruby{後胤}{こういん}にして、\ruby{大鴻臚}{たいこうろ}\ruby{曹崇}{そうすう}が嫡男なり。洛陽にあっては、\ruby{官騎}{かんき}\ruby{都尉}{とい}に封ぜられ、今、朝命によって、五千余騎にて馳せ来り、幸いにも、貴軍の火攻めの計に乗じて、逃ぐる賊を討ち、賊徒の首を討つことその数を知らないほどです。――ひとつお互いに両軍声をあわせて、天下の泰平を一日もはやく地上へ呼ぶため、凱歌をあげましょう」\\


「結構です。では、曹操閣下が\ruby{矛}{ほこ}をあげて、両軍へ発声の指揮をしてください」\\


　玄徳が謙遜していうと、\\


「いやそれは違う。こよいの勝ち\ruby{軍}{いくさ}はひとえに貴軍の謀略と働きにあるのですから、玄徳殿が音頭をとるべきです」と、曹操も譲りあう。\\


「では、一緒に、指揮の矛を揚げましょう」\\


「なるほど、それならば」\\


　と、曹操も従って、両将は両軍のあいだに\ruby{轡}{くつわ}をならべ、そして三度、\ruby{鬨}{とき}の声をあわせて野をゆるがした。\\


　野火は燃えひろがるばかりで賊徒らの住む尺地も余さなかった。賊の大軍は、ほとんど、秋風に舞う木の葉のように四散した。\\


「愉快ですな」\\


　曹操は、かえりみて云った。\\


　兵をまとめて、両軍引揚げの先頭に立ちながら、玄徳は、彼と駒を並べ、彼と親しく話すかなりな時間を得た。\\


　彼の最前の名乗りは、あながち\ruby{鬼面}{きめん}人を\ruby{脅}{おど}すものではなかった。玄徳は正直に、彼の人物に尊敬を払った。\ruby{晋文匡扶}{しんぶんきょうふ}の才なきを笑い、\ruby{趙高王莽}{ちょうこうおうもう}の\ruby{計策}{はかり}なきを\ruby{嘲}{あざけ}って時々、自らの才を誇る風はあるが、兵法は呉子孫子をそらんじ、学識は孔孟の遠き弟子をもって任じ、話せば話すほど、深みもあり広さもある人物と思われた。\\


　それにひきかえて、本軍の総大将\ruby{朱雋}{しゅしゅん}は、かえって玄徳の武功をよろこばないのみか、玄徳が戻ってくると、すぐこう命令した。\\


「せっかく、\ruby{潁川}{えいせん}にまとまっていた賊軍を四散させてしまったので、必ず彼らは、大興山の友軍や広宗の張角軍と合体して、\ruby{盧植}{ろしょく}将軍のほうを、今度はうんと悩ますにちがいない。――貴公はすぐ広宗へ引っ返して、再び、盧植軍に加勢してやり給え。今夜だけ、馬を休めたら、すぐ発足するがよかろう」\\

\\

\subsection{\ruby{檻車}{かんしゃ}}

\\

\subsubsection{一}

\\


　義はあっても、\ruby{官爵}{かんしゃく}はない。勇はあっても、官旗を持たない。そのために玄徳の軍は、どこまでも、私兵としか扱われなかった。\\


（よく戦ってくれた）と、恩賞の沙汰か、ねぎらいの言葉でもあるかと思いのほか、休むいとまもなく、（ここはもうよいから、広宗の地方へ転戦して、盧将軍を\ruby{援}{たす}けにゆけ）\\


　という\ruby{朱雋}{しゅしゅん}の命令には、玄徳は素直な\ruby{質}{たち}なので、承知して戻ったが、関羽も、張飛も、それを聞くと、\\


「え。すぐにここを立てというんですか」\\


　と、むっとした顔色だった。ことに張飛は、\\


「怪しからん沙汰だ。いかに官軍の大将だからといって、そんな命令を、おうけしてくる法があるものか。昨夜から悪戦苦闘してくれた部下にだって、気の毒で、そんなことがいえるものか」と、激昂し、「長兄は、大人しいもんだから、洛陽の都会人などの眼から見るとなめられやすいのだ。拙者がかけ合ってくる」\\


　と、剣をつかんで、朱雋の本営へ出かけそうにしたので、玄徳よりは、同じ不快をこらえている関羽が、\\


「まあ待て」と、極力おさえた。\\


「ここで、腹を立てたら、折角、官軍へ協力した意義も武功も、みな水泡に帰してしまう。都会人て奴は、元来、わがままで思い上がっているものだ。しかし、黙ってわれわれが国事に尽していれば、いつか誠意は天聴にも達するだろう。眼前の利慾に怒るのは小人の\ruby{業}{わざ}だ。われわれは、もっと高い理想に向って起つはずじゃないか」\\


「でも\ruby{癪}{しゃく}にさわる」\\


「感情に負けるな」\\


「無礼なやつだ」\\


「分った。分った。もうそれでいいだろう」\\


　ようやく\ruby{宥}{なだ}めて、\\


「劉兄。お腹も立ちましょうが、戦場も世の中の一部です。広い世の中としてみればこんなことはありがちでしょう。即刻、この地を引揚げましょう」\\


　ついでに関羽は、玄徳の憂鬱もそういって慰めた。\\


　玄徳はもとより、そう腹も立っていない。こらえるとか、堪忍とか、二人はいっているが、彼自身は、生来の性質が微温的にできているのか、実際、朱雋の命令にしてもそう無礼とも無理とも思えないし、怒るほどに、気色を害されてもいなかったのである。\\


　兵には、一睡させて、せめて食糧もゆっくりとらせて、夜半から玄徳は、そこの陣地を引払った。\\


　きのうは西に戦い。\\


　きょうは東へ。\\


　毎日、五百の手勢と、行軍をつづけていても、私兵のあじけなさを、しみじみ思わずにいられなかった。\\


　部落を通れば、土民までが馬鹿にする。――その土民らを賊の\ruby{虐圧}{ぎゃくあつ}と、悪政の下から救って、安心楽土の幸福な民としてやろうというこの軍の精神であるのに――そのみすぼらしい雑軍的な装備を見て、\\


「なんじゃ。官軍でもなし、黄巾賊でもないのが、ぞろぞろ通りよる」\\


　などと、陽なたに手をかざし合って、\ruby{嘲弄}{ちようろう}するような眼をあつめながら見物していた。\\


　けれど、先頭の玄徳、張飛、関羽の三人だけは、人目をひいた。威風が道を払った。土民らの中には土下座して拝する者もあった。\\


　拝されても、嘲弄されても、玄徳はいずれにせよ、気にかけなかった。自分が畑に働いていた頃の気持をもって、土民の気持を理解しているからだった。\\

\\

\subsubsection{二}

\\


　駒を並べてくる関羽と張飛とはまだ\ruby{朱雋}{しゅしゅん}の無礼を思い出して、時々、腹が立ってくるものとみえ、官軍の風紀や、洛陽の都人士の軽薄を、しきりに声を大にして\ruby{罵}{ののし}っていた。\\


「およそ嫌なものは、\ruby{官爵}{かんしゃく}を誇って、朝廷のご威光を、自分の偉さみたいに、思い上がっている奴だ。天下の\ruby{紊}{みだ}るるは、天下の紊れに非ず、官の\ruby{廃頽}{たいはい}によるというが、洛陽育ちの役人や将軍のうちには、あんなのが沢山いるだろうて」\\


　と、関羽がいえば、\\


「そうさ。俺はよッぽど、朱雋の面へ、ヘドを吐きかけてやろうと思ったよ」と張飛もいう。\\


「はははは。貴公のヘドをかけられたら、朱雋も驚いたろうな。しかし彼一人が官僚臭の鼻もちならぬ人間というわけではない。漢室の\ruby{廟堂}{びょうどう}そのものが腐敗しているのだ。彼は、その中に\ruby{棲息}{せいそく}している時代人だから、その悪弊を持っているに過ぎない」\\


「それゃあ分っているが、とにかく俺は、目前の事実を憎むよ」\\


「いくら\ruby{黄匪}{こうひ}を討伐しても、中央の悪風を粛正しなければ、ほんとのよい時代はやって来まいな」\\


「黄巾の賊はなお討つに易し。廟堂の\ruby{鼠臣}{そしん}はついに\ruby{趁}{お}うも難し――か」\\


「その通りだ」\\


「考えれば考えるほど、俺たちの理想は遠い――」\\


　道をながめ、空を仰ぎ、両雄は嘆じ合っていた。\\


　少し前へ立って、馬を進めていた玄徳は、二人の声高なはなしを先刻から後ろ耳で聞いていたが、その時、振りかえって、\\


「いやいや両人、そう一概にいってしまったものではない。洛陽の将軍のうちにも、立派な人物は乏しくない」と、いった。\\


　玄徳は、言葉をつづけて、\\


「たとえば先頃、野火の戦野で出会って挨拶を交わした――\ruby{赤備}{あかぞな}えの一軍の大将、\ruby{孟徳曹操}{もうとくそうそう}などという人物は、まだ若いが、\ruby{人品}{じんぴん}といい、言語態度といい、まことに見あげたものだった。叡智の才を、洛陽の文化と、武勇とに磨いて、一個の人格に飽和させているところ、彼など真に官軍の将軍といって恥かしからぬ者であろう。ああいう武将というものは、やはり郷軍や地方の\ruby{草莽}{そうもう}のなかには見当らないと思うな」と、\ruby{賞}{ほ}めたたえた。\\


　それには、張飛も関羽も、同感であったが、浪人の通有性として官軍とか官僚とかいうと、まずその人物の真価をみるより先に、その色や臭いを嫌悪してかかるので、玄徳にそういわれるまでは、特に、曹操に対しても、感服する気にはなれなかったのである。\\


「ヤ。旗が見える」\\


　そのうちに、彼らの部下は、こういって指さし合った。玄徳は、馬を止めて、\\


「なにが来るのだろうか」と、関羽をかえりみた。\\


　関羽は、手をかざして、道の前方数十町の先を、眺めていた。そこは山陰になって、山と山の間へ道がうねっているので、太陽の光もかげり、何やら一団の人間と旗とが、こっちへさして来るのは分るが官軍やら黄巾賊の兵やら――また、地方を浮浪している雑軍やら、見当がつかなかった。\\


　だが、次第に近づくに従って、ようやく\ruby{旗幟}{きし}がはっきり分った。関羽が、それと答えた時には、従う兵らも口々に云い交わしていた。\\


「朝旗をたてている」\\


「アア。官軍だ」\\


「三百人ばかりの官軍の隊」\\


「だが、おかしいぞ、熊でも捕まえて入れてくるのか、\ruby{檻車}{かんしゃ}をひいて来るじゃないか」\\

\\

\subsubsection{三}

\\


　大きな鉄格子の\ruby{檻}{おり}である。車がついているので\ruby{驢}{ろ}にひかせることができる。まわりには、槍や棒を持った官兵が、怖い目をしながら警固してくる。\\


　その前に百名。\\


　その後ろに約百名。\\


　檻車を真ん中にして、七\ruby{旒}{りゅう}の朝旗は山風にひるがえっていた。そして、檻車の中に、揺られてくるのは、熊でも\ruby{豹}{ひょう}でもなかった。膝を抱いて、天日に\ruby{面}{おもて}を\ruby{俯}{ふ}せている、あわれなる人間であった。\\


　ばらばらっと、先頭から、一名の隊将と、一隊の兵が、馳け抜けてきて、玄徳の一行を、頭から\ruby{咎}{とが}めた。\\


「こらっ、待てっ」というふうにである。\\


　張飛も、ぱっと、玄徳の前へ駒を躍らせて、万一をかばいながら、\\


「なんだっ、虫けら」と、いい返した。\\


　いわずともよい言葉であったが、\ruby{潁川}{えいせん}以来、とかく官兵の\ruby{空威}{からい}ばりに、\ruby{業腹}{ごうはら}の煮えていたところなので、つい口をついて出てしまったのである。\\


　石は石を打って、火を発した。\\


「なんだと、官旗に対して、虫けらといったな」\\


「礼を知るをもって\ruby{人倫}{じんりん}の始まりという。礼儀をわきまえん奴は、虫けらも同然だ」\\


「だまれ、われわれは、洛陽の勅使、\ruby{左豊卿}{さほうきょう}の直属の軍だ。旗を見よ。朝旗が見えんか」\\


「王城の直軍とあれば、なおさらのことである。俺たちも、武勇奉公を任じる軍人だ。私軍といえど、この旗に対し、こらっ待てとはなんだ。礼をもって問えば、こちらも礼をもって答えてやる。出直してこい」\\


　丈八の蛇矛を\ruby{斜}{しゃ}に構えて、かっとにらみつけた。\\


　官兵はちぢみ上がったものの、虚勢を張ったてまえ、退きもならず、\ruby{生唾}{なまつば}をのんでいた。玄徳は、眼じらせで、関羽にこの場を扱うように促した。\\


　関羽は、心得て、\\


「あいや、これは\ruby{潁川}{えいせん}の\ruby{朱雋}{しゅしゅん}・\ruby{皇甫嵩}{こうほすう}の両軍に参加して、これより広宗へ引っ返して参る\ruby{{\fontspec{jigmo}\char"6DBF}県}{たくけん}の劉玄徳の手勢でござる。ことばの行きちがい、この\ruby{漢}{おとこ}の短慮はゆるし給え。――ついてはまた、貴下の軍は、これより何処へ参らるるか。そして、あれなる\ruby{檻車}{かんしゃ}にある人間は、賊将の張角でも\ruby{生擒}{いけど}ってこられたのであるか」\\


　詫びるところは詫び、\ruby{糺}{ただ}すところは筋目をただして、質問した。\\


　官兵の隊将は、それに、ほっとした顔つきを見せた。張飛の暴言も薬になったとみえ、今度は丁寧に、\\


「いやいや、あれなる檻車に押しこめてきた罪人は、先頃まで、広宗の征野にあって、官軍一方の将として、洛陽より派遣せられていた\ruby{中郎将盧植}{ちゅうろうしょうろしょく}でござる。」［＃「でござる。」」は底本では「でござる。」］\\


「えっ、盧植将軍ですって」\\


　玄徳は、思わず、驚きの声を放った。\\


「されば、吾々には詳しいことも分らぬが、今度勅命にて下られた\ruby{左豊卿}{さほうきょう}が、各地の軍状を視察中、盧植の軍務ぶりに不届きありと奏されたため、急に盧植の官職を\ruby{褫奪}{ちだつ}され、これよりその身\jruby{がら}{・・・・・・}を、一囚人として、都へ差し立てて行く途中なので――」\\


　と語った。\\


　玄徳も、関羽も、張飛も、\\


「嘘のような……」と、茫然たる面を見あわせたまま、しばしいうことばを知らなかった。\\


　玄徳はやがて、\\


「実は、盧植将軍は、自分の旧師にあたるお人なので、ぜひともひと目、お別れをお告げ申したいが、なんとか許してもらえまいか」と切に頼んだ。\\

\\

\subsubsection{四}

\\


「ははあ。では、罪人盧植は、貴公の旧師にあたる者か。それは定めし、ひと目でも会いたかろうな」\\


　守護の隊将は、玄徳の切な願いを、\ruby{肯}{き}くともなく、肯かぬともなく、すこぶるあいまいに口を濁して、「許してもよいが、公の役目のてまえもあるしな」と、意味ありげに呟いた。\\


　関羽は、玄徳の袖をひいて、彼は\ruby{賄賂}{わいろ}を求めているにちがいない。貧しい軍費ではあるが、幾分かをさいて、彼に与えるしかありますまいといった。\\


　張飛は、それを小耳にはさむと、怪しからぬことである。そんなことをしては\ruby{癖}{くせ}になる。もしきかなければ、武力に訴えて、盧将軍の檻車へ迫り、ご対面なさるがよい。自分が引受けて、警固の奴らは近寄せぬからといったが、玄徳は、\\


「いやいや、かりそめにも、朝廷の旗を奉じている兵や役人へ向って、さような暴行はなすべきでない。といって、師弟の情、このまま盧将軍と相見ずに別れるにも忍びないから――」\\


　といって、なにがしかの銀を、軍費のうちから出させて、関羽の手からそっと、守護の隊将へ手渡し、\\


「ひとつ、あなたのお力で」\\


　と、折入っていうと、\ruby{賄賂}{わいろ}の効き目は、手のひらを返したようにきいて、隊将は立ち戻って、檻車を\ruby{停}{と}め、\\


「しばらく、休め」　と、自分の率いている官兵に号令した。\\


　そしてわざと、彼らは見て見ぬふりして、路傍に槍を組んで休憩していた。\\


　玄徳は、騎をおりて、その間に、檻車のそばへ馳け寄り、がんじょうな鉄格子へすがりついて、\\


「先生っ。先生っ。玄徳でございます。いったい、このお姿は、どうなされたことでござりますぞ」\\


　と、嘆いた。\\


　膝を曲げて、\ruby{暗澹}{あんたん}と、顔を埋めたまま、檻車の中に背をまるくしていた盧植は、その声に、はっと眼を向けたが、\\


「おうっ」\\


　と、それこそ、さながら野獣のように、鉄格子のそばへ、跳びついてきて、\\


「玄徳か……」と、舌をつらせて\ruby{顫}{わなな}いた。\\


「いい所で会った。玄徳、聞いてくれ」\\


　盧植は、無念な涙に、眼も顔もいっぱいに曇らせながらいう。\\


「実は、こうだ。――先頃、貴公がわしの陣を去って、\ruby{潁川}{えいせん}のほうへ立ってから間もなく、勅使左豊という者が、軍監として戦況の検分に来たが、世事に\ruby{疎}{うと}いわしは、陣中であるし、天子の使いとして、彼を迎えるに、あまりに真面目すぎて、他の将軍連のように、左豊に\ruby{献物}{けんもつ}を贈らなかった。……するとあつかましい左豊は、我に\ruby{賄賂}{まいない}をあたえよと、自分の口から求めてきたが、陣中にある金銀は、みなこれ官の公金にして、兵器戦備の\ruby{費}{つい}えにする物、ほかに私財とてはなし。ことに、軍中なれば、吏に贈る財物など、何であろうかと、わしはまた、真っ正直に断った」\\


「……なるほど」\\


「すると、左豊は、盧植はわれを恥かしめたりと、ひどく恨んで帰ったそうだが、間もなく、身に覚えない罪名のもとに、軍職を\ruby{褫奪}{ちだつ}されてこんな浅ましい姿をさらして、都へ差立てられる身とはなってしもうた。……今思えば、わしもあまり一徹であったが、洛陽の顕官どもが、私利私腹のみ肥やして、君も思わず、民をかえりみず、ただ一身の栄利に\ruby{汲々}{きゅうきゅう}としておる\ruby{状}{さま}は、想像のほかだ。実に嘆かわしい。こんなことでは、後漢の霊帝の御世も、おそらく長くはあるまい。……ああどうなりゆく世の中やら」\\


　と、盧植は、身の不幸を悲しむよりも、さすがに、より以上、上下乱脈の世相の果てを、\ruby{痛哭}{つうこく}するのであった。\\

\\

\subsubsection{五}

\\


　慰めようにも慰めることばもなく、鉄格子をへだてた\ruby{盧植}{ろしょく}の手を握りしめて、玄徳も共にただ悲嘆の涙にくれていたが、\\


「いや先生、ご胸中はお察しいたしますが、いかに世が末になっても、罪なき者が罰せられて、悪人や奸吏がほしいままに、\ruby{栄耀}{えいよう}を\ruby{全}{まっと}うすることはありません。日月も雲におおわれ、山容も、\ruby{烟霧}{えんむ}に真の\ruby{象}{かたち}を現さない時もあります。そのうちに、ご\ruby{冤罪}{えんざい}は\ruby{拭}{ぬぐ}われて、また聖代に祝しあう日もありましょう。どうか、時節をお待ちください。お体を大切に、恥をしのんで、じっとここは、ご辛抱ください」\\


　と励ました。\\


「ありがとう」と、盧植もわれにかえって、「思わぬ所で、思わぬ人に会ったため、つい心もゆるみ、不覚な涙を見せてしもうた。……わしなどはすでに老朽の身だが、頼むのは、貴公たち将来のある青年へだ。……どうか\ruby{億生}{おくしょう}の民草のために、頼むぞ劉備」\\


「やります。先生」\\


「ああしかし」\\


「何ですか」\\


「わしの如き、老年になっても、まだ\ruby{佞人}{ねいじん}の策におち、檻車に生き恥をさらされるような不覚をするのだ。\ruby{汝}{おこと}らはことに年も若いし、世の経験に浅い身だ。くれぐれも、平時の処世に細心でなければ危ないぞ。戦を覚悟の戦場よりも、心をゆるめがちの平時のほうが、どれほど危険が多いか知れない」\\


「ご訓誡、\ruby{肝}{きも}に銘じておきます」\\


「では、あまり長くなっても、また迷惑がかかるといけないから――」\\


　と、盧植が、早く立去れかしと、玄徳を眼で\ruby{急}{せ}き立てていると、それまで、檻車の横にたたずんでいた張飛が、突然、\\


「やあ長兄。罪もなき恩師が、獄府へ引かれて行くのを、このまま見過すという法があろうか。今のはなしを聞くにつけ、また先頃からの鬱憤もかさんでおる。もはや張飛の堪忍の緒はきれた。――守護の官兵どもを、みなごろしにして、檻車を奪い盧植様をお助けしようではないか」\\


　と、大声でいい放ち、一方の関羽をかえりみて、\\


「兄貴、どうだ」と、相談した。\\


　耳こすりや、眼まぜでしめし合わすのではない。天地へ向って呶鳴るのである。いくら背中を向けて見ぬ振りをしている官兵でも、それには総立ちになって、色めかざるをえない。しかし、張飛の眼中には、蠅が舞いだした程にもなく、\\


「なにを黙っておるのか。長兄らは、官兵が怖いのか。義を見て\ruby{為}{な}さざるは勇なきなり。よしっ、それでは、俺ひとりでやる。なんの、こんな虫籠のような檻車一つ」\\


　いきなり張飛は、その鉄格子に手をかけて、猛虎のように、ゆすぶりだした。\\


　いつもあまり大きな声を出さないし、めったに顔いろを変えない玄徳が、それを見ると、\\


「張飛！　何をするかッ」と、大喝して、「かりそめにも、朝命の\ruby{科人}{とがにん}へ、汝、一野夫の身として、何をなさんとするか。師弟の情は忍び難いが、なお、私情に過ぎない。いやしくも天子の命とあらば、地を噛んでも伏すべきである。世々の道に\ruby{反}{そむ}かずということは、そもそも、われらの軍律の第一則であった。\ruby{強}{た}って、乱暴を働くにおいては、天子の臣に代り、また、わが軍律に照らして、劉玄徳が、まず汝の首を\ruby{刎}{は}ねん。――いかに張飛、なおさわぐや」\\


　と、かの名剣の柄をにぎって、\ruby{眦}{まなじり}を紅に裂き、この人にしてこの血相があるかと疑われるばかりな声で叱りつけた。\\

\\

\subsubsection{六}

\\


　――檻車は遠く去った。\\


　叱られて、思いとどまった張飛は、後ろの山のほうを向いて、見ていなかった。\\


　玄徳は、立っていた。\\


「…………」\\


　黙然と、凝視して、遠くなり行く師の檻車を、暗涙の中に見送っていた。\\


「……さ。参りましょう」\\


　関羽は、\ruby{促}{うなが}して、駒を寄せた。\\


　玄徳は、黙々と、騎上の人になったが、盧植の運命の急変が、よほど精神にこたえたとみえ、\\


「……ああ」と、なお嘆息しては、振向いていた。\\


　張飛は、つまらない顔していた。彼にとっては、正しい義憤としてやったことが、計らずも玄徳の怒りを買い、義盟の血をすすり合ってから初めてのような叱られ方をした。\\


　官兵どもは、それを見て、いい気味だというような嘲笑を浴びせた。張飛たるもの、腐らずにいられなかった。\\


「いけねえや、どうも家の大将は、すこし安物の\ruby{孔子}{こうし}にかぶれている気味だて」\\


　舌打ちしながら、彼も黙りこんだまま、\ruby{悄気}{しょげ}かえった姿を、駒にまかせていた。\\


　山峡の道を過ぎて、二州のわかれ道へきた。\\


　関羽は、駒を止めて、\\


「玄徳様」と、呼びかけた。\\


「これから南へ行けば広宗。北へさしてゆけば、郷里\ruby{{\fontspec{jigmo}\char"6DBF}県}{たくけん}の方角へ近づきます。いずれを選びますか」\\


「もとより、盧植先生が\ruby{囚}{とら}われの身となって、洛陽へ送られてしまったからには、義をもってそこへ援軍にゆく意味ももうなくなった。ひとまず、{\fontspec{jigmo}\char"6DBF}県へ帰ろうよ」\\


「そうしますか」\\


「うム」\\


「それがしも、先刻からいろいろ考えていたのですが、どうも、残念ながら、一時郷里へ退くしかないであろう――と思っていたので」\\


「転戦、また転戦。――なんの功名ももたらさず、郷家に待つ母上にも、なんとなく、会わせる顔もないここちがするが……帰ろうよ、{\fontspec{jigmo}\char"6DBF}県へ」\\


「はっ。――では」\\


　と、関羽は、騎首をめぐらして、後からつづいて来る五百余の手兵へ、\\


「北へ、北へ！」\\


　と、指して歩行の号令をかけ、そしてまた黙々と、歩みつづけた。\\


「あア――、あ、あ」\\


　張飛は、大あくびして、\\


「いったい、なんのために、俺たちは戦ったんだい。ちっともわけが分らない。――こうなると一刻もはやく、{\fontspec{jigmo}\char"6DBF}県の城内へ帰って、市の酒屋で久しぶりに、\ruby{猪}{いのこ}の\ruby{股}{もも}でもかじりながら、うまい酒でも飲みたいものだ」と、いった。\\


　関羽は、苦い顔して、\\


「おいおい、兵隊のいうようなことをいうな。一方の将として」\\


「だって、俺は、ほんとのことをいっているんだ。嘘ではない」\\


「貴様からして、そんなことをいったら、軍紀がゆるむじゃないか」\\


「軍紀のゆるみだしたのは、俺のせいじゃない。官軍官軍と、なんでも、官軍とさえいえば、意気地なく恐がる人間のせいだろ」\\


　不平満々なのである。\\


　その不平な気もちは、玄徳にも分っていた。玄徳もまた、不平であったからだ。そしてひと頃の張り切っていた\ruby{壮志}{そうし}のゆるみをどうしようもなかった。彼は、\ruby{女々}{めめ}しく郷里の母を想い出し、また、思うともなくい\ruby{鴻芙蓉}{こうふよう}の麗しい眉や眼などを、人知れず胸の\ruby{奥所}{おくが}に描いたりして、なんとなく士気の\ruby{沮喪}{そそう}した軍旅の\ruby{虚無}{きょむ}と不平をなぐさめていた。\\


　すると、突然、山崩れでもしたように、一方の山岳で、\ruby{鬨}{とき}の声が聞えた。\\

\\

\subsubsection{七}

\\


「何事か」\\


　玄徳は聞き耳たてていたが、四山にこだまする\ruby{銅鑼}{どら}、\ruby{兵鼓}{へいこ}の響きに、\\


「張飛。物見せよ」と、すぐ命じた。\\


「心得た」\\


　と張飛は駒を飛ばして、山のほうへ向って行ったが、しばらくすると戻ってきて、\\


「広宗の方面から逃げくずれて来る官軍を、\ruby{黄巾}{こうきん}の\ruby{総帥}{そうすい}\ruby{張角}{ちょうかく}の軍が、\ruby{大賢良師}{たいけんりょうし}と書いた旗を進め、勢いに乗って、追撃してくるのでござる」と、報告した。\\


　玄徳は、驚いて、\\


「では、広宗の官軍は、総敗北となったのか。――罪なき盧植将軍を、檻車に囚えて、洛陽へ差し立てたりなどしたために、たちまち、官軍は統制を失って、賊にその虚をつかれたのであろう」\\


　と、嘆じた。\\


　張飛は、むしろ小気味よげに、\\


「いや、そればかりでなく、官軍の士風そのものが、長い平和になれ、気弱にながれ、思い上がっているからだ」と、関羽へいった。\\


　関羽は、それに答えず、\\


「長兄。どうしますか」\\


　と玄徳へ計った。\\


　玄徳は、ためらいなく、\\


「皇室を重んじ、秩序をみだす賊子を討ち、民の\ruby{安寧}{あんねい}を護らんとは、われわれの初めからの鉄則である。官の士風や軍紀をつかさどる者に、面白からぬ人物があるからというて、官軍そのものが潰滅するのを、\ruby{拱手傍観}{きょうしゅぼうかん}していてもよいものではない」\\


　と、即座に、援軍に馳せつけて、賊の追撃を、山路で中断した。そしてさんざんにこれを悩ましたり、また、奇策をめぐらして、張角大方師の本軍まで攪乱した上、勢いを挽回した官軍と合体して、五十里あまりも賊軍を追って引揚げた。\\


　広宗から敗走してきた官軍の大将は、\ruby{董卓}{とうたく}という将軍だった。\\


　からくも、総敗北を盛返して、ほっと一息つくと、将軍は、幕僚にたずねた。\\


「いったい、かの山嶮で、不意にわが軍へ加勢し、賊の後方を攪乱した軍隊は、いずれ味方には相違あるまいが、どこの部隊に属する将士か」\\


「さあ。どこの隊でしょう」\\


「汝らも知らんのか」\\


「誰もわきまえぬようです」\\


「しからば、その部将に会って、自身訊ねてみよう。これへ呼んでこい」\\


　幕僚は、直ちに、玄徳たちへ董卓の意をつたえた。\\


　玄徳は、\ruby{左将}{さしょう}関羽、\ruby{右将}{うしょう}張飛を従えて、董卓の面前へ進んだ。\\


　董卓は、椅子を与える前に、三名の姓名をたずねて、\\


「洛陽の王軍に、\ruby{卿}{けい}らのごとき勇将があることは、まだ\ruby{寡聞}{かぶん}にして聞かなかったが、いったい諸君は、なんという官職に就かれておるのか」と、身分を\ruby{糺}{ただ}した。\\


　玄徳は、無爵無官の身をむしろ誇るように、自分らは、正規の官軍ではなく、天下万民のために、大志を奮い起して立った一地方の義軍であると答えた。\\


「……ふうむ。すると、{\fontspec{jigmo}\char"6DBF}県の楼桑村から出た私兵か。つまり雑軍というわけだな」\\


　董卓の応対ぶりは、言葉つきからして違ってきた。露骨な軽蔑を鼻先に見せていうのだった。しかも、\\


「――ああそうか。じゃあ我が軍に\ruby{従}{つ}いて、大いに働くがよいさ。給料や手当は、いずれ沙汰させるからな」\\


　と同席するさえ、自分の\ruby{沽券}{こけん}にかかわるように、董卓はいうとすぐ\ruby{帷幕}{いばく}のうちへ隠れてしまった。\\

\\

\subsubsection{八}

\\


　官軍にとっては、大功を立てたのだ。董卓にとっては、生命の親だといってもよいのだ。\\


　然るに！\\


　何ぞ、\ruby{遇}{ぐう}するの、無礼。\\


　士を遇する道を知らぬにも程がある。\\


「…………」\\


　玄徳も、張飛と関羽も、董卓のうしろ姿を見送ったまま、茫然としていた。\\


「うぬっ」\\


　憤然と、張飛は、彼のかくれた\ruby{幕}{とばり}の奥へ、躍り入ろうとした。\\


　獅子のように、髪を立てて。\\


　そして剣を手に。\\


「あっ、何処へ行く」\\


　玄徳は、驚いて、張飛のうしろから組み止めながら、\\


「こらっ、また、わるい短慮を出すか」と、叱った。\\


「でも。でも」\\


　張飛は、怒りやまなかった。\\


「――ちッ、畜生っ。官位がなんだっ。官職がない者は、人間でないように思ってやがる。馬鹿野郎ッ。民力があっての官位だぞ。賊軍にさえ、蹴ちらされて、逃げまわって来やがったくせに」\\


「これッ、鎮まらんか」\\


「離してくれ」\\


「離さん。関羽関羽。なぜ見ているか、一緒に、張飛を止めてくれい」\\


「いや関羽、止めてくれるな。おれはもう、堪忍の\ruby{緒}{お}を切った。――功を立てて恩賞もないのは、まだ我慢もするが、なんだ、あの軽蔑したあいさつは。――人を雑軍かとぬかしおった。私兵かと、鼻であしらいやがった。――離してくれ、董卓の素ッ首を、この\ruby{蛇矛}{じゃぼこ}で一太刀にかッ飛ばして見せるから」\\


「待て。……まあ待て。……腹が立つのは、貴様ばかりではない。だが、小人の小人ぶりに、いちいち腹を立てていたひには、とても大事はなせぬぞ。天下、小人に満ちいる時だ」\\


　玄徳は、抱き止めたまま、声をしぼって\ruby{諭}{さと}した。\\


「しかし、なんであろうと、董卓は皇室の武臣である。朝臣を\ruby{弑逆}{しいぎゃく}すれば、理非にかかわらず、叛逆の賊子といわれねばならぬ。それに、董卓には、この大軍があるのだ。われわれも共に、ここで斬り死しなければならぬ。聞きわけてくれ張飛。われわれは、犬死するために、起ったのではあるまいが」\\


「……ち、ち、ちく生ッ」\\


　張飛は、床を、大きく\ruby{沓}{くつ}で踏み鳴らして、男泣きに、声をあげて泣いた。\\


「口惜しい」\\


　彼は、坐りこんで、まだ泣いていた。この忍耐をしなければ、世のために戦えないのか、義を\ruby{唱}{とな}えても、遂になすことはできないのかと考えると悲しくなってくるのだった。\\


「さ。外へ出よう」\\


　赤ン坊をあやすように、玄徳と関羽の二人して、彼を、左右から抱き起こした。\\


　そして、その夜、「こんな所に長居していると、いつまた、張飛が虫を起さないとも限らないから」と、董卓の陣を去って、手兵五百と共に、月下の曠野を、\ruby{蕭々}{しょうしょう}と、風を負って歩いた。\\


　わびしき雑軍。\\


　そして官職のない将僚。\\


　一軍の\ruby{漂泊}{さすらい}は、こうして再び続いた。夜ごとに、月は白く小さく、曠野は果てなくまた露が深かった。\\


　渡り鳥が、大陸をゆく。\\


　もう秋なのだ。\\


　いちどは郷里の{\fontspec{jigmo}\char"6DBF}県へ帰ろうとしたが、それも残念でならないし、あまりに無意義――という関羽の意見に、張飛も、将来は何事も我慢しようと同意したので、玄徳を先頭にしたこの渡り鳥にも似た一軍は、また、以前の\ruby{潁川}{えいせん}地方にある黄匪討伐軍本部――\ruby{朱雋}{しゅしゅん}の陣地へと志して行ったのであった。\\

\\

\subsection{\ruby{秋風陣}{しゅうふうじん}}

\\

\subsubsection{一}

\\


　潁川の地へ行きついてみると、そこにはすでに官軍の一部隊しか残っていなかった。大将軍の\ruby{朱雋}{しゅしゅん}も\ruby{皇甫嵩}{こうほすう}も、賊軍を追いせばめて、遠く河南の曲陽や\ruby{宛城}{えんじょう}方面へ移駐しているとのことであった。\\


「さしも\ruby{旺}{さかん}だった黄巾賊の勢力も、洛陽の派遣軍のために、次第に各地で討伐され、そろそろ自壊しはじめたようですな」\\


　関羽がいうと、\\


「つまらない事になった」\\


　と、張飛はしきりと、今のうちに功を立てねば、いつの時か風雲に乗ぜんと、\ruby{焦心}{あせ}るのであった。\\


「――義軍なんぞ小功を思わん。\ruby{義胆}{ぎたん}なんぞ風雲を要せん」\\


　劉玄徳は、独りいった。\\


　\ruby{雁}{かり}の列のように、漂泊の小軍隊はまた、南へ向って、旅をつづけた。\\


　黄河を渡った。\\


　兵たちは、馬に水を飼った。\\


　玄徳は、黄いろい大河に眼をやると、\ruby{憶}{おも}いを深くして、\\


「ああ、悠久なる\ruby{哉}{かな}」\\


　と、呟いた。\\


　四、五年前に見た黄河もこの通りだった。おそらく百年、千年の後も、黄河の水は、この通りにあるだろう。\\


　天地の悠久を思うと、人間の一瞬がはかなく感じられた。小功は思わないが、しきりと、生きている間の生甲斐と、意義ある仕事を残さんとする誓願が念じられてくる。\\


「この\ruby{畔}{ほとり}で、半日もじっと若い空想にふけっていたことがある。――洛陽船から茶を\ruby{購}{あがな}おうと思って」\\


　茶を思えば、同時に、母が憶われてくる。\\


　この秋、いかに\ruby{在}{お}わすか。足の冷えや、持病が出てはこぬだろうか。ご不自由はどうあろうか。\\


　いやいや母は、そんなことすら忘れて、ひたすら、子が大業をなす日を待っておられるであろう。それと共に、いかに聡明な母でも、実際の戦場の事情やら、また実地に当る軍人同士のあいだにも、常の社会と変らない難しい感情やら争いやらあって、なかなか武力と正義の信条一点張りで、世に出られないことなどは、お察しもつくまい。ご想像にも及ぶまい。\\


　だから以来、なんのよい便りもなく、月日をむなしく送っている子をお考えになると、\\


（\ruby{阿備}{あび}は、何をしているやら）\\


　と、さだめしふがいない者と、\ruby{焦}{じ}れッたく思っておいでになるに相違ない。\\


「そうだ。せめて、体だけは無事なことでも、お便りしておこうか」\\


　玄徳は、思いつめて、騎の鞍をおろし、その鞍に結びつけてある旅具の中から、\ruby{翰墨}{かんぼく}と筆を取りだして、母へ便りを書きはじめた。\\


　駒に水を飼って、休んでいた兵たちも、玄徳が\ruby{箋葉}{せんよう}に筆をとっているのを見ると、\\


「おれも」\\


「吾も」\\


　と、何か書きはじめた。\\


　誰にも、故郷がある。姉妹兄弟がある。玄徳は思いやって、\\


「故郷へ手紙をやりたい者は、わしの手もとへ持ってこい。親のある者は、親へ無事の消息をしたがよいぞ」と、いった。\\


　兵たちは、それぞれ紙片や木皮へ、何か書いて持ってきた。玄徳はそれを一\ruby{嚢}{のう}に納めて、実直な兵を一人撰抜し、\\


「おまえは、この手紙の\ruby{嚢}{ふくろ}をたずさえて、それぞれの郷里の家へ、郵送する役目に当れ」\\


　と、路費を与えて、すぐ立たせた。\\


　そして落日に染まった黄河を、騎と兵と荷駄とは、黒いかたまりになって、浅瀬は\ruby{徒渉}{としょう}し、深い所は\ruby{筏}{いかだ}に\ruby{棹}{さお}さして、対岸へ渡って行った。\\

\\

\subsubsection{二}

\\


　先頃から河南の地方に、何十万とむらがっている賊の大軍と戦っていた大将軍\ruby{朱雋}{しゅしゅん}は、思いのほか賊軍が手ごわいし、味方の死傷はおびただしいので、\\


「いかがはせん」と、内心\ruby{煩悶}{はんもん}して、苦戦の憂いを顔にきざんでいたところだった。\\


　そこへ、\\


「\ruby{潁川}{えいせん}から広宗へ向った玄徳の隊が、形勢の変化に、途中から引っ返してきて、ただ今、着陣いたしましたが」と、幕僚から知らせがあった。\\


　朱雋はそれを聞くと、\\


「やあ、それはよい所へ来た。すぐ通せ、失礼のないように」\\


　と、前とは、打って変って、鄭重に待遇した。陣中ながら、洛陽の美酒を開き、料理番に牛など裂かせて、\\


「長途、おつかれであろう」と、歓待した。\\


　正直な張飛は、前の不快もわすれて、すっかり感激してしまい、\\


「士は\ruby{己}{おのれ}を知る者の為に死す、である」\\


　などと酔った機嫌でいった。\\


　だが歓待の代償は義軍全体の生命に近いものを求められた。\\


　翌日。\\


「早速だが、豪傑にひとつ、打破っていただきたい方面がある」\\


　と、朱雋は、玄徳らの軍に、そこから約三十里ほど先の山地に陣取っている頑強な敵陣の突破を命じた。\\


　否む理由はないので、\\


「心得た」と、義軍は、朱雋の部下三千を加えて、そこの高地へ攻めて行った。\\


　やがて、山麓の野に近づくと天候が悪くなった。雨こそ降らないが、密雲低く垂れて、烈風は草を飛ばし、沼地の水は霧になって、兵馬の行くてを\ruby{晦}{くら}くした。\\


「やあ、これはまた、賊軍の大将の張宝が、妖気を起して、われらを皆ごろしにすると見えたるぞ。気をつけろ。樹の根や草につかまって、烈風に吹きとばされぬ用心をしたがいいぞ」\\


　朱雋からつけてよこした部隊から、誰いうとなく、こんな声が起って、恐怖はたちまち全軍を\ruby{蔽}{おお}った。\\


「ばかなっ」\\


　関羽は怒って、\\


「世に理のなき妖術などがあろうか。\ruby{武夫}{もののふ}たるものが、\ruby{幻妖}{げんよう}の術に怖れて、木の根にすがり、大地を這い、戦意を失うとは、何たるざまぞ。すすめや者ども、関羽の行く所には妖気も避けよう」\\


　と大声で鼓舞したが、\\


「妖術にはかなわぬ。あたら生命をわざわざおとすようなものだ」\\


　と、朱雋の兵は、なんといっても前進しないのである。\\


　聞けば、この高地へ向った官軍は、これまでにも何度攻めても、全滅になっているというのであった。黄巾賊の\ruby{大方師}{だいほうし}張角の弟にあたる張宝は、有名な妖術つかいで、それがこの高地の山谷の奥に陣取っているためであるという。\\


　そう聞くと張飛は、\\


「妖術とは、\ruby{外道}{げどう}魔物のする\ruby{業}{わざ}だ。天地\ruby{闢}{ひら}けて以来、まだかつて方術者が天下を取ったためしはあるまい。\ruby{怖}{お}じる心、おそれる\ruby{眼}{まなこ}、わななく魂を惑わす術を、妖術とはいうのだ。怖れるな、惑うな。――進まぬやつは、軍律に照らして斬り捨てるぞ」\\


　と、軍のうしろにまわって、手に\ruby{蛇矛}{じゃぼこ}を抜きはらい、督戦に努めた。\\


　朱雋の兵は、敵の妖術にも恐怖したが、張飛の蛇矛にはなお恐れて、やむなくわっと、黒風へ向って前進しだした。\\

\\

\subsubsection{三}

\\


　その日は、天候もよくなかったに違いないが、戦場の地勢もことに悪かった。寄手にとっては、甚だしく不利な地の利にいやでも置かれるように、そこの高地は自然にできている。\\


　\ruby{峨々}{がが}たる山が、道の両わきに、鉄門のように\ruby{聳}{そび}えている。そこを突破すれば、高地の沢から、山地一帯の敵へ肉薄できるのだが、そこまでが、近づけないのだった。\\


「鉄門峡まで行かぬうちに、いつも味方はみなごろしになる。豪傑、どうか無謀はやめて、引っ返し給え」\\


　と、朱雋の軍隊の者は、部将からして、\ruby{怯}{ひる}み上がっていうほどだから、兵卒が皆、恐怖して自由に動かないのも無理ではなかった。\\


　だが、張飛は、\\


「それは、いつもの寄手が弱いからだ。きょうは、われわれの義軍が先に立って進路を斬りひらく、武夫たる者は、戦場で死ぬのは、本望ではないか。死ねや、死ねや」と、督戦に声をからした。\\


　先鋒は、ゆるい\ruby{砂礫}{されき}の丘を這って、もう鉄門峡のまぢかまで、攻め上っていた。\ruby{朱雋}{しゅしゅん}軍も、張飛の蛇矛に斬り捨てられるよりはと、その後から、芋虫の群れが動くように這い上がった。\\


　すると、たちまち、一陣の風雷、天地を震動して木も砂礫も人も、中天へ吹きあげられるかとおぼえた時、一方の山峡の頂に、陣鼓を鳴らし、\ruby{銅鑼}{どら}を打ちとどろかせて、\\


　――わあっ。わあっ。\\


　と、烈風も圧するような\ruby{鬨}{とき}の声がきこえた。寄手は皆地へ伏し、眼をふさぎ、耳を忘れていたが、その声に振り仰ぐと、山峡の\ruby{絶巓}{ぜってん}はいくらか平盤な地になっているとみえて、そこに賊の一群が見え「\ruby{地公将軍}{ちこうしょうぐん}」と書いた旗や、八\ruby{卦}{け}の文を印した黄色の\ruby{幟}{のぼり}、\ruby{幡}{はた}など立て並べて、\\


「死神につかれた軍が、またも\ruby{黄泉}{よみじ}へ急いで来つるぞ。\ruby{冥途}{めいど}の\ruby{扉}{と}を開けてやれ」\\


　と、声を合わせて笑った。\\


　その中に一人、遠目にもわかる異相の巨漢があった。口に\ruby{魔符}{まふ}を噛み、髪をさばき、\ruby{印}{いん}をむすんでなにやら呪文を唱えている容子だったが、それと共に烈風は益{\fontspec{jigmo}\char"303B}つのって、\ruby{晦冥}{かいめい}な天地に、人の形や魔の形をした赤、青、黄などの紙片がまるで五彩の火のように降ってきた。\\


「やあ、魔軍が来た」\\


「賊将張宝が、\ruby{呪}{じゅ}を唱えて、天空から\ruby{羅刹}{らせつ}の援軍を呼び出したぞ」\\


　朱雋の兵は、わめき合うと、逃げ惑って、途も失い、ただ右往左往うろたえるのみだった。\\


　張飛の督戦も、もう\ruby{効}{き}かなかった。朱雋の兵があまり恐れるので、義軍の兵にも恐怖症がうつったようである。そして風魔と砂礫にぶつけられて、全軍、進むことも退くこともできなくなってしまった時、赤い\ruby{紙片}{かみきれ}や青い紙片の魔物や武者は、それ皆が、生ける夜叉か羅刹の軍のように見えて、寄手は完全に闘志を失ってしまった。\\


　事実。\\


　そうしている間に、無数の矢や岩石や火器は、うなりをあげ、煙をふいて、寄手の上に降ってきたのである。またたくうちに、全軍の半分以上は、動かないものになっていた。\\


「敗れた！　負けたっ」\\


　玄徳は、軍を率いてから初めて惨たる敗戦の味をいま知った。\\


　そう叫ぶと、\\


「関羽っ。張飛っ。はや兵を\ruby{退}{ひ}けっ――兵を退けっ」\\


　そして自分もまっしぐらに、駒首を逆落しに向けかえし、砂礫とともに山裾へ馳け下った。\\

\\

\subsubsection{四}

\\


　敗軍を収めて、約二十里の外へ退き、その夜、玄徳は関羽、張飛のふたりと共に、\ruby{帷幕}{いばく}のうちで軍議をこらした。\\


「残念だ、きょうまで、こんな敗北はしたことがないが」と、張飛がいう。\\


　関羽は、腕を\ruby{拱}{く}んでいたが、\\


「\ruby{朱雋}{しゅしゅん}の兵が、戦わぬうちから、あのように恐怖しているところを見ると、何か、あそこには不思議がある。張宝の幻術も、実際、ばかにはできぬかも知れぬ」と、呟いた。\\


「幻術の不思議は、わしには\ruby{解}{と}けている。それは、あの鉄門峡の地形にあるのだ。あの峡谷には、常に雲霧が立ちこめていて、その気流が、烈風となって、峡門から\ruby{麓}{ふもと}へいつも吹いているのだと思う」\\


　これは玄徳の説である。\\


「なるほど」と二人とも初めて、そうかと気づいた顔つきだった。\\


「だから少しでも天候の悪い日には、ほかの土地より何十倍も強い風が吹きまくる。この辺が、晴天の日でも、峡門には、黒雲がわだかまり、砂礫が飛び、煙雨が降り\ruby{荒}{すさ}んでいる」\\


「ははあ、大きに」\\


「好んで、それへ向ってゆくので、近づけばいつも、賊と戦う前に、天候と戦うようなものになる。張宝の地公将軍とやらは、奸智に\ruby{長}{た}けているとみえて、その自然の気象を、自己の妖術かの如く、巧みに使って、\ruby{藁}{わら}人形の武者や、紙の\ruby{魔形}{まぎょう}など降らせて、朱雋軍の愚かな恐怖をもてあそんでいたものであろう」\\


「さすがに、ご活眼です。いかにも、それに違いありません。けれど、あの山の賊軍を攻めるには、あの峡門から攻めかかるほかありますまい」\\


「ない。――それ故に、朱雋はわざと、われわれを、この攻め口へ当らせたのだ」\\


　玄徳は、沈痛にいった。\\


　関羽、張飛の二人も、良い策もない、唇をむすんで、陣の曠野へ眼をそらした。\\


　折から仲秋の月は、満目の曠野に露をきらめかせ、二十里外の彼方に黒々と見える臥牛のような山岳のあたりは、味方を悩ませた悪天候も嘘ごとのように、大気と月光の\ruby{下}{もと}に横たわっていた。\\


「いや、ある、ある」\\


　突然、張飛が、自問自答して云いだした。\\


「攻め口が、ほかにないとはいわさん。長兄、一策があるぞ」\\


「どうするのか」\\


「あの絶壁を\ruby{攀}{よ}じ登って、賊の予測しない所から不意に衝きくずせば、なんの造作もない」\\


「登れようか、あの断崖絶壁へ」\\


「登れそうに見える所から登ったのでは、奇襲にはならない。誰の眼にも、登れそうに見えない場所から登るのが、用兵の策というものであろう」\\


「張飛にしては、珍しい名言を吐いたものだ。その通りである。登れぬものときめてしまうのは、人間の観念で、その眼だけの観念を超えて、実際に懸命に当ってみれば案外やすやすと登れるような例はいくらでもあることだ」\\


　さらに、三名は、密議をねって、翌る日の作戦に備えた。\\


　\ruby{朱雋}{しゅしゆん}軍の兵、約半分の数に、おびただしい旗や\ruby{幟}{のぼり}を持たせ、また、\ruby{銅鑼}{どら}や\ruby{鼓}{こ}を打ち鳴らさせて、きのうのように峡門の正面から、強襲するような\ruby{態}{てい}を敵へ見せかけた。\\


　一方、張飛、関羽の両将に、幕下の\ruby{強者}{つわもの}と、朱雋軍の一部の兵を率きつれた玄徳は、峡門から十里ほど北方の絶壁へひそかに這いすすみ、惨澹たる苦心のもとに、山の一端へ\ruby{攀}{よ}じ登ることに成功した。\\


　そしてなお、士気を鼓舞するために、すべての兵が\ruby{山巓}{さんてん}の一端へ登りきると、そこで玄徳と関羽は、おごそかなる\ruby{破邪攘魔}{はじゃじょうま}の祈祷を天地へ向って捧げるの儀式を行った。\\

\\

\subsubsection{五}

\\


　敵を前にしながら、わざとそんな所で、おごそかな祈祷の儀式などしたのは、玄徳直属の義軍の中にも、張宝の幻術を内心怖れている兵がたくさんいるらしく見えたからであった。\\


　式が終ると、\\


「見よ」\\


　玄徳は空を指していった。\\


「きょうの一天には、風魔もない、迅雷もない、すでに、破邪の祈祷で、張宝の幻術は通力を失ったのだ」\\


　兵は答えるに、万雷のような\ruby{喊声}{かんせい}をもってした。\\


　関羽と張飛は、それと共に、\\


「それ、魔軍の\ruby{砦}{とりで}を踏みつぶせ」\\


　と軍を二手にわけて、峰づたいに張宝の本拠へ攻めよせた。\\


　地公将軍の\ruby{旗幟}{きし}を立てて、賊将の張宝は、例によって、鉄門峡の寄手を悩ましにでかけていた。\\


　すると、思わざる山中に、突然\ruby{鬨}{とき}の声があがった。彼は、味方を振返って、\\


「裏切り者が出たか」と、訊ねた。\\


　実際、そう考えたのは、彼だけではなかった。裏切り者裏切り者という声が、何処ともなく伝わった。\\


　張宝は、\\


「\ruby{不埓}{ふらち}な奴、何者か、成敗してくれん」\\


　と、そこの守りを、賊の一将にいいつけて、自身、わずかの部下を連れて、山谷の奥にある――ちょうど\ruby{螺}{ら}の穴のような渓谷を、\ruby{驢}{ろ}に鞭打って帰ってきた。\\


　するとかたわらの沢の密林から、一すじの矢が飛んできて、張宝のこめかみにぐざと立った。張宝はほとばしる黒血へ手をやって、わッと口を開きながら矢を抜いた。しかし\ruby{鏃}{やじり}はふかく頭蓋の中に止まって、矢柄だけしか抜けてこなかったくらいなので、とたんに、彼の巨躯は、鞍の上から真っ逆さまに落ちていた。\\


「賊将の張宝は射止めたるぞ。劉玄徳、ここに黄匪の大方張角の弟、地公将軍を討ち取ったり」\\


　次に、どこかで玄徳の大音声がきこえると、四方の山沢、みな鼓を鳴らし、奔激の渓流、\ruby{挙}{こぞ}って\ruby{鬨}{とき}をあげ、草木みな兵と\ruby{化}{な}ったかと思われた。玄徳の兵は、一斉に衝いて出で、あわてふためく張宝の部下をみなごろしにした。\\


　山谷の奥からも、同時に黒煙\ruby{濛々}{もうもう}とたち昇った。張飛か、関羽の手勢が、本拠の\ruby{砦}{とりで}に、火をかけたものらしい。\\


　上流から流れてくる渓水は、みるまに紅の奔流と化した。山吠え、谷叫び、火は山火事となって、三日三晩燃えとおした。\\


　首\ruby{馘}{き}る数一万余、黒焦げとなった賊兵の死骸幾千幾万なるを知らない。殲滅戦の続けらるること七日余り、玄徳は、赫々たる武勲を負って\ruby{朱雋}{しゅしゅん}の本営へ引揚げた。\\


　朱雋は、玄徳を見ると、\\


「やあ、\ruby{足下}{そっか}は実に運がいい。\ruby{戦}{いくさ}にも、運不運があるものでな」と、いった。\\


「ははあ、そうですか。ひと口に、武運ということもありますからね」\\


　玄徳は、なんの感情にも動かされないで、軽く笑った。\\


　朱雋は、さらにいう。\\


「自分のひきうけている野戦のほうは、まだいっこう勝敗がつかない。山谷の賊は、ふくろの鼠としやすいが、野陣の敵兵は、押せばどこまでも、逃げられるので弱るよ」\\


「ごもっともです」\\


　それにも、玄徳はただ、笑ってみせたのみであった。\\


　然るところ、ここに、先陣から伝令が来て、一つの異変を告げた。\\

\\

\subsubsection{六}

\\


　伝令の告げるには、\\


「先に戦没した賊将張宝の兄弟\ruby{張梁}{ちょうりょう}という者、天公将軍の名を称し、久しくこの曠野の陣後にあって、督軍しておりましたが、張宝すでに討たれぬと聞いて、にわかに大兵をひきまとめ、陽城へたて籠って、城壁を高くし、この冬を守って越えんとする策を取るかに見うけられます」\\


　とのことだった。\\


　朱雋は、聞くと、\\


「冬にかかっては、雪に凍え、食糧の運輸にも、困難になる。ことに\ruby{都聞}{みやこきこ}えもおもしろくない。今のうちに攻めおとせ」\\


　総攻撃の令を下した。\\


　大軍は陽城を囲み、攻めること急であった。しかし、賊城は要害堅固をきわめ、城内には多年積んだ食物が豊富なので、一月余も費やしたが、城壁の一角も奪れなかった。\\


「困った。困った」\\


　朱雋は本営で時折ため息をもらしたが、玄徳は聞えぬ顔をしていた。\\


　よせばいいに、そんな時、張飛が朱雋へいった。\\


「将軍。野戦では、押せば\ruby{退}{ひ}くしで、戦いにくいでしょうが、こんどは、敵も城の中ですから、袋の鼠を捕るようなものでしょう」\\


　朱雋は、まずい顔をした。\\


　そこへ遠方から使いが来て、新しい情報をもたらした。それもしかし朱雋の機嫌をよくさせるものではなかった。\\


　曲陽の方面には、朱雋と共に、討伐大将軍の任を負って下っていた\ruby{董卓}{とうたく}・\ruby{皇甫嵩}{こうほすう}の両軍が、賊の大方張角の大兵と戦っていた。使いはその方面のことを知らせに来たものだった。\\


　董卓と皇甫嵩のほうは、朱雋のいういわゆる武運がよかったのか、七度戦って七度勝つといった按配であった。ところへまた、黄賊の総帥張角が、陣中で病没したため、総攻撃に出て、一挙に賊軍を潰滅させ、降人を収めること十五万、辻に\ruby{梟}{か}くるところの賊首何千、さらに、張角を\ruby{埋}{い}けた\ruby{墳}{つか}をあばいてその首級を洛陽へ\ruby{上}{のぼ}せ、\\


（戦果かくの如し）と、報告した。\\


　大賢良師張角と称していた\ruby{首魁}{しゅかい}こそ、天下に満つる乱賊の首体である。張宝は先に討たれたりといっても、その弟にすぎず、張梁なおありといっても、これもその一肢体でしかない。\\


　朝廷の\ruby{御感}{ぎょかん}は斜めならず、\\


（征賊第一勲）\\


　として、\ruby{皇甫嵩}{こうほすう}を\ruby{車騎将軍}{しゃきしょうぐん}に任じ、\ruby{益州}{えきしゅう}の\ruby{牧}{ぼく}に封ぜられ、そのほか恩賞の令を受けた者がたくさんある。わけても、陣中常に赤い甲冑を着て通った武騎校尉\ruby{曹操}{そうそう}も、功によって、\ruby{済南}{さいなん}（山東省・黄河南岸）の\ruby{相}{しょう}に封じられたとのことであった。\\


　自分が逆境の中に、他人の栄達を聞いて、共によろこびを感じるほど、\ruby{朱雋}{しゅしゅん}は寛度でない。彼はなお、\ruby{焦心}{あせ}りだして、\\


「一刻もはやく、この城を攻め陥し、汝らも、朝廷の恩賞にあずかり、封土へ帰って、栄達の日を楽しまずや」と、幕僚をはげました。\\


　もちろん、玄徳らも、協力を惜しまなかった。攻撃に次ぐ攻撃をもって、城壁に当り、さしも頑強な賊軍をして、眠るまもない防戦に疲れさせた。\\


　城内の賊の中に、\ruby{厳政}{げんせい}という男があった。これは方針をかえる時だとさとったので、ひそかに朱雋に内通しておき、賊将張梁の首を斬って、\\


「願わくば、\ruby{悔悟}{かいご}の兵らに、王威の恩浴を垂れたまえ」と、軍門に降ってきた。\\


　陽城を\ruby{墜}{おと}した勢いで、\\


「さらに、与党を狩りつくせ」\\


　と、朱雋の軍六万は、\ruby{宛城}{えんじょう}（湖北省・荊門県附近）へ迫って行った。そこには、黄巾の残党、\ruby{孫仲}{そんちゅう}・\ruby{韓忠}{かんちゅう}・\ruby{趙弘}{ちょうこう}の三賊将がたて籠っていた。\\

\\

\subsubsection{七}

\\


「賊には\ruby{援}{たす}けもないし、城内の兵糧もいたずらに敗戦の兵を多く容れたから、またたく間に尽きるであろう」\\


　朱雋は、陣頭に立って、賊の宛城の運命を、かく\ruby{卜}{うらな}った。\\


　朱雋軍六万は、宛城の周囲をとりまいて、水も漏らさぬ布陣を詰めた。\\


　賊軍は、\\


「やぶれかぶれ」の策を選んだか、連日、城門をひらいて、戦を挑み、官兵賊兵、相互におびただしい死傷を毎日積んだ。\\


　然しいかんせん、城内の兵糧はもう乏しくて、賊は飢渇に瀕してきた。そこで賊将韓忠は遂に、降使を立てて、\\


「仁慈を垂れ給え」と、降伏を申し出た。\\


　朱雋は、怒って、\\


「\ruby{窮}{きゅう}すれば、\ruby{憐}{あわれ}みを乞い、勢いを得れば、暴魔の威をふるう、今日に至っては、仁慈もなにもない」\\


　と、降参の使者を斬って、なおも苛烈に攻撃を加えた。\\


　玄徳は彼に\ruby{諫}{いさ}めた。\\


「将軍、賢慮し給え。昔、漢の高祖の天下を\ruby{統}{す}べたまいしは、よく降人を容れてそれを用いたためといわれています」\\


　将軍は、\ruby{嘲笑}{あざわら}って、\\


「ばかをいい給え。それは時代による。あの頃は、\ruby{秦}{しん}の世が乱れて項羽のような\jruby{がさつ}{・・・・・・}者の私議暴論が横行して、天下に定まれる君主もなかった時勢だろ、ゆえに高祖は、\ruby{讐}{あだ}ある者でも、降参すれば、手なずけて用うことに腐心したのである。また、秦の乱世のそれと、今日の黄賊とは、その質がちがう。生きる利なく、窮地に墜ちたがゆえに、降を乞うてきた賊を、\ruby{愍}{あわ}れみをかけて、救けなどしたら、それはかえって\ruby{寇}{あだ}を長じさせ、世道人心に、悪業を奨励するようなものではないか。この際、断じて、賊の根を絶たねばいかん」\\


「いや、伺ってみると、たいへんごもっともです」\\


　玄徳は、彼の説に伏した。\\


「では、攻めて城内の賊を、殲滅するとしてもです。こう四方、一門も\ruby{遁}{のが}れる隙間なく囲んで攻めては、城兵は、死の\ruby{一途}{いちず}に結束し、恐ろしい最後の力を奮いだすにきまっています。味方の損害もおびただしいことになりましょう。一方の門だけは、逃げ口を与えておいて、三方からこれを攻めるべきではありますまいか」\\


「なるほど、その説はよろしい」\\


　朱雋は、直ちに、命令を変更して、急激に攻めたてた。\\


　\ruby{東南}{たつみ}の一門だけ開いて、三方から鼓をならし、火を放った。\\


　果たして、城内の賊は、乱れ立って一方へくずれた。\\


　朱雋は、騎を飛ばして、乱軍の中に、賊将の韓忠を見かけ、鉄弓で射とめた。\\


　韓忠の首を、槍に突き刺させて、従者に高く振り上げさせ、\\


「征賊大将軍朱雋、賊徒の将、韓忠をかく葬ったり。われと名乗る者やなおある」\\


　と、得意になって呶鳴った。\\


　すると、残る賊将の\ruby{趙弘}{ちょうこう}、\ruby{孫仲}{そんちゅう}のふたりは、\\


「あいつが朱雋か」と、火炎の中を、\ruby{黒驢}{こくろ}を飛ばして、名のりかけてきた。\\


　朱雋は、たまらじと、自軍のうちへ逃げこんだ。韓忠親分の\ruby{讐}{かたき}と怒りに燃えた賊兵は、朱雋を追って、朱雋の軍の真ん中を突破し、まったくの乱軍を呈した。\\


　賊の一に対して、官兵は十人も死んだ。朱雋につづいて、官軍はわれがちに十里も後ろへ退却した。\\


　賊軍は、気をもり返して、城壁の火を消し、再び四方の門を固くして、\\


「さあいつでも来い」と構えなおした。\\


　その日の\ruby{黄昏}{たそが}れ、多くの傷兵が、惨として夕月の野に横たわっている官軍の陣営へ、何処からきたか、一\ruby{彪}{ぴょう}の軍馬が馳けきたった。\\

\\

\subsubsection{八}

\\


「何者か」\\


　と、玄徳らは、やがて近づいて陣門に入るその軍馬を、幕舎の傍らから見ていた。\\


　総勢、約千五百の兵。\\


　隊伍は整然、歩武堂々。\\


「そもこの精鋭を\ruby{統}{す}べる将はいかなる人物か」を、それだけでも思わすに足るものだった。\\


　見てあれば。\\


　その隊伍の真っ先に、旗手、鼓手の兵を立て、続いてすぐ後から、一頭の\ruby{青驪}{せいり}にまたがって、威風あたりを払ってくる人がある。\\


　これなんその一軍の大将であろう。\ruby{広額}{こうがく}、\ruby{濶面}{かつめん}、唇は\ruby{丹}{たん}のようで、眉は\ruby{峨眉山}{がびさん}の半月のごとく高くして鋭い。\ruby{熊腰}{ゆうよう}にして\ruby{虎態}{こたい}、いわゆる威あって\ruby{猛}{たけ}からず、見るからに大人の風を備えている。\\


「誰かな？」\\


「誰なのやら」\\


　関羽も張飛も、見まもっていたが、ほどなく陣門の衛将が、名を\ruby{糺}{ただ}すに答える声が、遠くながら聞えてきた。\\


「これは\ruby{呉郡富春}{ごぐんふしゅん}（浙江省・富陽市）の産で、\ruby{孫堅}{そんけん}、\ruby{字}{あざな}は\ruby{文台}{ぶんだい}という者です。\ruby{古}{いにしえ}の孫子が末葉であります。官は\ruby{下{\fontspec{jigmo}\char"90B3}}{かひ}の\ruby{丞}{じょう}ですが、このたび王軍、黄巾の賊徒を諸州に討つと承って、手飼いの兵千五百を率い、いささか年来の恩沢にむくゆべく、官軍のお味方たらんとして馳せ参じた者であります。――\ruby{朱雋}{しゅしゅん}将軍へよろしくお取次を乞う」\\


　堂々たる態度であった。\\


　また、\ruby{音吐}{おんど}も朗々と聞えた。\\


「…………」\\


　関羽と張飛は、顔を見合わせた。先には、\ruby{潁川}{えいせん}の野で、\ruby{曹操}{そうそう}を見、今ここにまた、\ruby{孫堅}{そんけん}という一人物を見て、\\


「やはり世間はひろい。\ruby{秀}{ひい}でた人物がいないではない。ただ、世の平静なる時は、いないように見えるだけだ」と、感じたらしかった。\\


　同じ、その世間を、\\


「甘くはできないぞ」\\


　という気持を抱いたであろう。なにしろ、孫堅の入陣は、その卒伍までが、立派だった。\\


　孫堅の来援を聞いて、\\


「いや呉郡富春に、英傑ありと、かねてはなしに聞いていたが、よくぞ来てくれた」\\


　と、朱雋はななめならずよろこんで迎えた。\\


　きょうさんざんな敗軍の日ではあったし、朱雋は、大いに力を得て、翌日は、孫堅が\ruby{准泗}{わいし}の精鋭千五百をも加えて、\\


「一挙に」と、\ruby{宛城}{えんじょう}へ迫った。\\


　即ち、新手の孫堅には、南門の攻撃に当らせ、玄徳には北門を攻めさせ、自身は西門から攻めかかって、東門の一方は、前日の策のとおり、わざわざ道をひらいておいた。\\


「洛陽の将士に笑わるるなかれ」\\


　と、孫堅は、新手でもあるので、またたく間に、南門を衝き破り、彼自身も青毛の駒をおりて、濠を越え、単身、城壁へよじ登って、\\


「呉郡の孫堅を知らずや」\\


　と賊兵の中へ躍り入った。\\


　刀を舞わして孫堅が賊を斬ること二十余人、それに当って、噴血を浴びない者はなかった。\\


　賊将の\ruby{趙弘}{ちょうこう}は、\\


「ふがいなし、\ruby{彼奴}{きゃつ}、何ほどのことやあらん」\\


　\ruby{赫怒}{かくど}して孫堅に名のりかけ、烈戦二十余合、火をとばしたが、孫堅はあくまでつかれた色も見せず、たちまち趙弘を斬って捨てた。\\


　もう一名の賊将孫仲は、それを眺めて、かなわじと思ったか、敗走する味方の賊兵の中にまぎれこんで、早くも東門から逃げ走ってしまった。\\

\\

\subsubsection{九}

\\


　その時。\\


　ひゅっと、どこか天空で、弦を放たれた一矢の矢うなりがした。\\


　矢は、東門の望楼のほとりから、斜めに線を描いて、怒濤のように、われがちと敗走してゆく賊兵の中へ飛んだが、狙いあやまたず、今しも\ruby{金蘭橋}{きんらんきょう}の外門まで落ちて行った賊将孫仲の\ruby{頸}{うなじ}を射ぬき、孫仲は馬上からもんどり打って、それさえ眼に入らぬ賊兵の足にたちまち踏みつぶされたかに見えた。\\


「あの首、掻き取ってこい」\\


　玄徳は、部下に命じた。\\


　望楼のかたわらの壁上に鉄弓を持って立ち、目ぼしい賊を射ていたのは彼であった。\\


　一方、官軍の\ruby{朱雋}{しゅしゅん}も、孫堅も城中に攻め入って、首をとること数万級、各所の火災を鎮め、孫仲・趙弘・韓忠三賊将の首を城外に\ruby{梟}{か}け、市民に布告を発し、城頭の余燼まだ煙る空に、高々と、王旗をひるがえした。\\


「漢室万歳」\\


「洛陽軍万歳」\\


「朱雋大将軍万歳」\\


　南陽の諸郡もことごとく平定した。\\


　かの大賢良師張角が、戸ごとに貼らせた黄いろい\ruby{呪符}{じゅふ}もすべてはがされて、黄巾の兇徒は、まったく影をひそめ、万戸泰平を謳歌するかに思われた。\\


　しかし、天下の乱は、天下の草民から意味なく起るものではない。むしろその禍根は、民土の低きよりも、廟堂の高きにあった。川下よりも川上の水源にあった。政を奉ずる者より、政をつかさどる者にあった。地方よりも中央にあった。\\


　けれど腐れる者ほど自己の腐臭には気づかない。また、時流のうごきは眼に見えない。\\


　とまれ官軍は\ruby{旺}{さかん}だった。征賊大将軍は功なって、洛陽へ凱旋した。\\


　洛陽の城府は、挙げて、遠征の兵馬を迎え、市は五彩旗に染まり、夜は万燈にいろどられ、城内城下、七日七夜というもの酒の泉と音楽の狂いと、酔どれの歌などで沸くばかりであった。\\


　王城の府、洛陽は千万戸という。さすがに古い伝統の都だけに、物資は富み、文化は\ruby{絢爛}{けんらん}だった。佳人貴顕たちの往来は目を奪うばかり美しい。帝城は金壁にかこまれ、\ruby{瑠璃}{るり}の瓦を重ね、百官の\ruby{驢車}{ろしゃ}は、\ruby{翡翠門}{ひすいもん}に花のよどむような\ruby{雑鬧}{ざっとう}を呈している。天下のどこに一人の飢民でもあるか、今の時代を乱兆と悲しむいわれがあるのか、この\ruby{殷賑}{いんしん}に立って、\ruby{旺}{さかん}なる夕べの楽音を耳にし、\ruby{万斛}{ばんこく}の油が一夜にともされるという騒曲の灯の、宵早き有様を眺むれば、むしろ、世を憂え嘆く者のことばが不思議なくらいである。\\


　けれど。\\


　二十里の野外、そこに\ruby{連}{つら}なる外城の壁からもし一歩出てみるならば、秋は更けて、木も草も枯れ、いたずらに高き城壁に、\ruby{蔓}{つる}草の離々たる葉のみわずかに紅く、日暮れれば花々の闇一色、夜\ruby{暁}{あ}ければ\ruby{颯々}{さっさつ}の秋風ばかり\ruby{哭}{な}いて、所々の水辺に、寒げに啼く牛の仔と、灰色の空をかすめる\ruby{鴻}{こう}の影を時たまに仰ぐくらいなものであった。\\


　そこに。\\


　無口に\ruby{屯}{たむろ}している人間が、枯れ木や草をあつめて焚火をしながら、わずかに朝夕の霜の寒さをしのいでいた。\\


　玄徳たちの義軍であった。\\


　義軍は、外城の門の一つに立って、門番の役を命じられている。\\


　といえば、まだ体裁はよいが、正規の官軍でなし、官職のない将卒なので、三軍洛陽に凱旋の日も、ここに停められて、内城から先へは入れられないのであった。\\


　\ruby{鴻}{こう}が飛んでゆく。\\


　\ruby{野芙蓉}{のふよう}にゆらぐ秋風が白い。\\


「…………」\\


　玄徳も関羽も、この頃は、無口であった。\\


　あわれな卒伍は、まだ洛陽の温かい菜の味も知らない。\ruby{土龍}{もぐら}のように、鉄門の蔭に、かがまっていた。\\


　張飛も黙然と、水ばなをすすっては、時折、ひどく虚無に\ruby{囚}{とら}われたような顔をして、空行く鴻の影を見ていた。\\

\\

\subsection{十\ruby{常侍}{じょうじ}}

\\

\subsubsection{一}

\\


「\ruby{劉氏}{りゅううじ}。もし、劉氏ではありませんか」\\


　誰か呼びかける人があった。\\


　その日、劉玄徳は、\ruby{朱雋}{しゅしゅん}の官邸を訪ねることがあって、王城内の禁門の辺りを歩いていた。\\


　振向いてみると、それは\ruby{郎中張均}{ろうちゅうちょうきん}であった。張均は今、参内するところらしく、従者に\ruby{輿}{こし}をかつがせそれに乗っていたが、玄徳の姿を見かけたので、\\


「\ruby{沓}{くつ}を」と従者に命じて、輿から身をおろしていた。\\


「おう、どなたかと思うたら、張均閣下でいらっしゃいましたか」\\


　玄徳は、敬礼をほどこした。\\


　この人はかつて、\ruby{盧植}{ろしょく}をおとしいれた黄門左豊などと共に、監軍の勅使として、征野へ巡察に来たことがある。その折、玄徳とも知って、お互いに世事を談じ、\ruby{抱懐}{ほうかい}を話し合ったりしたこともある間なので、\\


「思いがけない所でお目にかかりましたな、ご健勝のていで、何よりに存じます」\\


　と、\ruby{久濶}{きゆうかつ}を\ruby{叙}{の}べた。\\


　郎中\ruby{張均}{ちょうきん}は、そういう玄徳の、従者も連れていない、しかも、かつて見た征衣のまま、この寒空を孤影悄然と歩いている様子をいぶかしげに打眺めて、\\


「貴公は今どこに何をしておられるのですか。少しお痩せになっているようにも見えるが」\\


　と、かえって玄徳の境遇を反問した。\\


　玄徳は、ありのままに、なにぶんにも自分には官職がないし、部下は私兵と見なされているので、凱旋の後も、外城より入るを許されず、また、忠誠の兵たちにも、この冬に向って、一枚の暖かい軍衣、一片の賞禄をもわけ与えることができないので、せめて外城の門衛に立っていても、霜をしのぐに足る暖衣と食糧とを恵まれんことを乞うために、きょう\ruby{朱雋}{しゅしゅん}将軍の官宅まで、願書をたずさえて出向いて来たところです、と話した。\\


「ほ……」\\


　張均は、驚いた顔して、\\


「では、足下はまだ、官職にもつかず、また、こんどの恩賞にもあずかっていないんですか」\\


　と、重ねて\ruby{糺}{ただ}した。\\


「はい、沙汰を待てとのことに、外城の門に\ruby{屯}{たむろ}しています。けれどもう冬は来るし、部下が\ruby{不愍}{ふびん}なので、お訴えに出てきたわけです」\\


「それは初めて知りました。\ruby{皇甫嵩}{こうほすう}将軍は、功によって、\ruby{益州}{えきしゅう}の太守に封ぜられ、朱雋は都へ凱旋するとただちに車騎将軍となり河南の\ruby{尹}{いん}に封ぜられている。あの孫堅さえ内縁あって、\ruby{別部司馬}{べつぶしば}に叙せられたほどだ。――いかに功がないといっても、貴君の功は孫堅以下ではない。いや或る意味では、こんどの掃匪征賊の戦で、最も苦戦に当って、忠誠をあらわした軍は、貴下の義軍であったといってもよいのに」\\


「…………」\\


　玄徳の面にも、\ruby{鬱々}{うつうつ}たるものがあった。ただ、彼は、朝廷の命なるがままに、思うまいとしているふうだった。そして部下の不愍を身の不遇以上にあわれと思いしめて噛んでいた唇の態であった。\\


「いや、よろしい」\\


　やがて張均はつよくいった。\\


「それも、これも、思い当ることがある。地方の騒賊を\ruby{掃}{はら}っても、\ruby{社稷}{しゃしょく}の\ruby{鼠巣}{そそう}を掃わなかったら、四海の平安を長く保つことはできぬ。賞罰の区々不公平な点ばかりでなく、嘆くべきことが実に多い。――貴君のことについては、特に帝へ\ruby{奏聞}{そうもん}しておこう。そのうちに明朗な恩浴をこうむることもあろうから、まあ気を腐らせずに待つがよい」\\


　郎中張均は、そう慰めて、玄徳とわかれ、やがて参内して、帝に拝謁した。\\

\\

\subsubsection{二}

\\


　めずらしく帝のお側には誰もいなかった。\\


　帝は、玉座からいわれた。\\


「張郎中。きょうは何か、\ruby{朕}{ちん}に、折入って懇願あるということだから、近臣はみな遠ざけておいたぞ。気がねなく思うことを申すがよい」\\


　張均は、階下に\ruby{拝跪}{はいき}して、\\


「帝のご聡明を信じて、臣張均は今日こそ、あえて、お気に入らぬことをも申しあげなければなりません。照々として、公明な御心をもて、暫時、お聴きくださいまし」\\


「なんじゃ」\\


「ほかでもありませんが、君側の十\ruby{常侍}{じょうじ}らのことについてです」\\


　十常侍ときくと、帝のお眸はすぐ横へ向いた。\\


　御気色がわるい――\\


　張均には分っていたが、ここを\ruby{冒}{おか}して真実の言をすすめるのが忠臣の道だと信じた。\\


「臣が多くを申しあげないでも、ご聡明な帝には、\ruby{疾}{と}くお気づきと存じますが、天下も今、ようやく平静に返ろうとして地方の乱賊も\ruby{終熄}{しゅうそく}したところです。この際、どうか君側の奸を\ruby{掃}{はら}い、ご粛正を上よりも示して、人民たちに暗天の憂えなからしめ、業に安んじ、ご徳政を謳歌するように、ご賢慮仰ぎたくぞんじまする」\\


「張郎中。なんできょうに限って、突然そんなことを云いだすのか」\\


「いや、十常侍らが政事を\ruby{紊}{みだ}して帝の御徳を\ruby{晦}{くろ}うし奉っている事はきょうのことではありません。私のみの憂いではありません。天下万民の怨みとするところです」\\


「怨み？」\\


「はい。たとえば、こんどの黄巾の乱でも、その賞罰には、十常侍らの私心が、いろいろ働いていると聞いています。\ruby{賄賂}{まいない}をうけた者には、功なき者へも官禄を与え、しからざる者は、罪なくても官を\ruby{貶}{おと}し、いやもう、ひどい沙汰です」\\


　帝の御気色は、いよいよ曇って見えた。けれど、帝は何もいわれなかった。\\


　十常侍というのは、十人の内官のことだった。民間の者は、彼らを\ruby{宦官}{かんがん}と称した。君側の権をにぎり\ruby{後宮}{こうきゅう}にも勢力があった。\\


　\ruby{議郎}{ぎろう}\ruby{張譲}{ちょうじょう}、\ruby{議郎}{ぎろう}\ruby{趙忠}{ちょうちゅう}、\ruby{議郎}{ぎろう}\ruby{段珪}{だんけい}、\ruby{議郎}{ぎろう}\ruby{夏輝}{かき}――などという十名が中心となって、\ruby{枢密}{すうみつ}に結束をつくっていた。議郎とは、参議という意味の役である。だからどんな枢密の政事にもあずかった。帝はまだお若くあられるし、そういう古池の\jruby{ぬし}{・・・・・・}みたいな老獪と\ruby{曲者}{くせもの}がそろっているので、彼らが遂行しようと思うことは、どんな悪政でもやって通した。\\


　霊帝はまだご\ruby{若年}{じゃくねん}なので、その悪弊に気づかれていても、いかんともする\ruby{術}{すべ}をご存じない。また、張均の\ruby{苦諫}{くかん}に感動されても、何というお答えもでなかった。ただ眼を宮中の\ruby{苑}{にわ}へそらしておられた。\\


「――遊ばしませ、ご断行なさいませ。今がその時です。陛下、ひとえに、ご賢慮をお決し下さいませ」\\


　張均は、口を\ruby{酸}{す}くし、われとわが忠誠の情熱に、\ruby{眦}{まなじり}に涙をたたえて諫言した。\\


　遂には、玉座に迫って、帝の\ruby{御衣}{ぎょい}にすがって、\ruby{泣訴}{きゅうそ}した。帝は、当惑そうに、\\


「では、張郎中、\ruby{朕}{ちん}に、どうせいというのか」と、問われた。\\


　ここぞと、張均は、\\


「十常侍らを獄に下して、その首を刎ね、南郊に\ruby{梟}{か}けて、諸人に罪文と共に示し給われば、人心おのずから平安となって、天下は」\\


　云いかけた時である。\\


「だまれっ。――まず汝の首より先に獄門に梟けん」\\


　と、\ruby{帳}{とばり}の蔭から怒った声がして、それと共に十常侍十名の者が躍り出した。みな\ruby{髪}{はつ}を逆立て、\ruby{眦}{まなじり}をあげながら、張均へ迫った。\\


　張均は、あッと驚きのあまり昏倒してしまった。\\


　手当されて、後に、典医から薬湯をもらったが、それを飲むと眠ったまま死んでしまった。\\

\\

\subsubsection{三}

\\


　張均は、その時、そんな死に方をしなくても、帝へ忠諫したことを十常侍に聴かれていたから、必ずや、後に命を\ruby{完}{まっと}うすることはできなかったろう。\\


　十常侍も、以来、\\


「油断しておると、とんでもない忠義ぶった奴が現れるぞ」\\


　と気がついたか、\ruby{誡}{いまし}め合って、帝の周囲はもとより、内外の政にわたって、大いに警戒しているふうであった。\\


　それもあるし、帝ご自身も、功ある者のうちに、恩賞にももれて不遇をかこち、不平を抑えている者がすくなくないのに気がつかれたか、特に、勲功の再調査と、第二期の恩賞の実施とを沙汰された。\\


　張均のことがあったので、十常侍も反対せず、むしろ自分らの善政ぶりを示すように、ほんの形ばかりな辞令を交付した。\\


　その中に、劉備玄徳の名もあった。\\


　それによって、玄徳は、\ruby{中山府}{ちゅうざんふ}（河北省・定県）の\ruby{安喜県}{あんきけん}の\ruby{尉}{い}という官職についた。\\


　\ruby{県尉}{けんい}といえば、片田舎の一警察署長といったような官職にすぎなかったが、帝命をもって叙せられたことであるから、それでも玄徳は、ふかく恩を謝して、関羽、張飛を従えて、即座に、任地へ出発した。\\


　もちろん、一官吏となったのであるから、多くの手兵をつれてゆくことは許されないし、必要もないので五百余の手兵は、これを王城の軍府に託して、編入してもらい、ほんの二十人ばかりの者を従者として連れて行ったに過ぎなかった。\\


　その冬は、任地でこえた。\\


　わずか四ヵ月ばかりしか経たないうちに、彼が役についてから、県中の政治は大いに\ruby{革}{あらた}まった。\\


　強盗悪逆の徒は、影をひそめ、良民は徳政に服して、平和な毎日を楽しんだ。\\


「張飛も関羽も、自己の器量に比べては、今の小吏のするような仕事は不服だろうが、しばらくは、現在に忠実であって貰いたい。時節はあせっても求め難い」\\


　玄徳は、時おり二人をそういって慰めた。それは彼自身を慰める言葉でもあった。\\


　その代り、県尉の任についてからも、玄徳は、彼らを下役のようには使わなかった。共に貧しきにおり、夜も床を同じゅうして寝た。\\


　するとやがて、河北の野に芽ぐみだした春とともに、\\


「天子の使いこの地に来る」\\


　と、伝えられた。\\


　勅使の使命は、\\


「このたび、黄巾の賊を平定したるに、軍功ありといつわりて、\ruby{政廟}{せいびょう}の内縁などたのみ、みだりに官爵をうけ或いは、功ありと自称して、州都に私威を振舞う者多く聞え、よくよく、正邪を\ruby{糺}{ただ}さるべし」\\


　という\ruby{詔}{みことのり}を奉じて\ruby{下向}{げこう}してきた者であった。\\


　そういう沙汰が、役所へ達しられてから間もなく、この安喜県へも、\ruby{督郵}{とくゆう}が下って来た。\\


　玄徳らは、さっそく関羽、張飛などを従えて、督郵の行列を道に出迎えた。\\


　何しろ、使いは、地方巡察の勅を奉じてきた大官であるから、玄徳たちは、地に坐して、最高の礼をとった。\\


　すると、馬上の督郵は、\\


「ここか安喜県とは。ひどい田舎だな。何、県城はないのか。役所はどこだ。県尉を呼べ。今夜の旅館はどこか、案内させて、ひとまずそこで休息しよう」\\


　と、いいながら、\ruby{傲然}{ごうぜん}と、そこらを見廻した。\\

\\

\subsubsection{四}

\\


　勅使督郵の人もなげな傲慢さを眺めて、\\


「いやに役目を鼻にかけるやつだ」と、関羽、張飛は、かたはらいたく思ったが、虫を抑えて、一行の車騎に従い、県の役館へはいった。\\


　やがて、玄徳は、衣服を正して、彼の前に、挨拶に出た。\\


　督郵は、左右に、随員の吏を侍立させ、さながら自身が帝王のような顔して、高座に構えこんでいた。\\


「おまえは何だ」\\


　知れきっているくせに、督郵は上から玄徳らを見下ろした。\\


「県尉玄徳です。はるばるのご下向、ご苦労にございました」\\


　\ruby{拝}{はい}をほどこすと、\\


「ああおまえが当地の県の尉か。途々、われわれ勅使の一行が参ると、うすぎたない住民どもが、車騎に近づいたり、指さしたりなど、はなはだ\ruby{猥雑}{わいざつ}な\ruby{態}{てい}で見物しておったが、かりそめにも、勅使を迎えるに、なんということだ。思うに平常の取締りも手ぬるいとみえる。もちっと王威を知らしめなければいかんよ」\\


「はい」\\


「旅館のほうの準備は整うておるかな」\\


「地方のこととて、諸事おもてなしはできませんが」\\


「われわれは、きれい好きで、飲食は贅沢である。田舎のことだから仕方がないが\ruby{卿}{けい}らが、勅使を遇するに、どういう心をもって歓待するか、その心もちを見ようと思う」\\


　意味ありげなことをいったが、玄徳には、よく解し得なかった。けれど、帝王の命をもって下ってきた勅使であるから、真心をもって、応接した。\\


　そして、ひとまず退がろうとすると、督郵はまた訊いた。\\


「\ruby{尉}{い}玄徳。いったい卿は、当所の出身の者か、他県から赴任してきたのか」\\


「されば、自分の郷家は\ruby{{\fontspec{jigmo}\char"6DBF}県}{たくけん}で、家系は、\ruby{中山靖王}{ちゅうざんせいおう}の\ruby{後胤}{こういん}であります。久しく土民の中にひそんでいましたが、この度ようやく、黄巾の乱に小功あって、当県の尉に叙せられた者であります」\\


　と、いうと、\\


「こらっ、黙れ」\\


　督郵は、突然、高座から叱るようにどなった。\\


「中山靖王の後胤であるとかいったな。\ruby{怪}{け}しからんことである。そもそも、このたび、帝がわれわれ臣下に命じて、各地を巡察せしめられたのは、そういう\ruby{大法螺}{おおぼら}をふいたり、軍功のある者だなどといつわって、自称豪傑や、自任官職の\ruby{輩}{やから}が横行する由を、お聞きになられたからである。汝の如き賤しき者が、天子の宗族などといつわって、愚民に臨んでおるのは、怪しからぬ不敬である。――ただちに帝へ奏聞し奉って、追っての沙汰をいたすであろうぞ。退がれっ」\\


「……はっ」\\


「退がれ」\\


「…………」\\


　玄徳は、唇をうごかしかけて、何かいわんとするふうだったが、益なしと考えたか、黙然と礼をして去った。\\


「いぶかしい人だ」\\


　彼は、督郵の随員に、そっと一室で面会を求めた。\\


　そして、何で勅使が、ご不興なのであろうかと、原因をきいてみた。\\


　随員の下吏は、\\


「それや、あんた知れきっているじゃありませんか、なぜ今日、督郵閣下の前に出る時、\ruby{賄賂}{まいない}の\ruby{金帛}{きんぱく}を、自分の姿ほども積んでお見せしなかったんです。そしてわれわれ随員にも、それ相当のことを、いちはやく袖の下からすることが肝腎ですよ。何よりの歓迎というもんですな。ですからいったでしょう督郵様も、いかに遇するか心を見ておるぞよってね」\\


　玄徳は、唖然として、私館へ帰って行った。\\

\\

\subsubsection{五}

\\


　私館へ帰っても、彼は、\ruby{怏々}{おうおう}と楽しまぬ顔いろであった。\\


「県の土民は、みな貧しい者ばかりだ。しかも一定の税は徴収して、中央へ送らなければならぬ。その上、なんで巡察の勅使や、大勢の随員に、彼らの満足するような\ruby{賄賂}{わいろ}を贈る余裕があろう。賄賂も土民の汗あぶらから出さねばならぬに、よくほかの県吏には、そんなことができるものだ」\\


　玄徳は、嘆息した。\\


　次の日になっても、玄徳のほうからなんの贈り物もこないので、督郵は、\\


「県吏をよべ」と、他の吏人を呼びつけ、\\


「尉玄徳は、\ruby{不埓}{ふらち}な\ruby{漢}{おとこ}である。天子の宗族などと僭称しておるのみか、ここの百姓どもから、いろいろと\ruby{怨嗟}{えんさ}の声を耳にする。すぐ帝へ奏聞して、ご処罰を仰ぐから、汝は、県吏を代表して、訴状をしたためろ」といった。\\


　玄徳の徳に服してこそはいるが、玄徳に何の落度も考えられない県の吏は、恐れわななくのみで、答えも知らなかった。\\


　すると、督郵も重ねて、\\


「訴状を書かんか、書かねば汝も同罪と見なすぞ」と、脅した。\\


　やむなく、県の吏は、ありもしない罪状を、督郵のいうままに並べて、訴状に書いた。督郵は、それを都へ急送し、帝の沙汰を待って、玄徳を厳罰に処せんと称した。\\


　この四、五日。\\


「どうも面白くねえ」\\


　張飛は、酒ばかりのんでいた。\\


　そう飲んでばかりいるのを、玄徳や関羽に知れると、意見されるし、また、この数日、玄徳の顔いろも、関羽の顔いろも、はなはだ憂鬱なので、彼はひとり、\\


「……どうも面白くねえ」をくり返して、どこで飲むのか、姿を見せず飲んでいた。\\


　その張飛が、\ruby{熟柿}{じゅくし}のような顔をして、\ruby{驢}{ろ}に乗って歩いていた。町中の者は、県の\ruby{吏人}{やくにん}なので、驢と行きちがうと、丁寧に礼をしたが、張飛は、驢の上から落ちそうな恰好して、居眠っていた。\\


「やい。どこまで行く気だ」\\


　眼をさますと、張飛は、乗っている驢にたずねた。驢は、てこてこと、軽い\ruby{蹄}{ひづめ}をただ運んでいた。\\


「おや、なんだ？」\\


　役所の門前をながめると、七、八十名の百姓や町の者が、土下座して、なにか\ruby{喚}{わめ}いたり、頭を地へすりつけたりしていた。\\


　張飛は、驢をおりて、\\


「みんな、どうしたんだ。おまえら、なにを役所へ泣訴しておるんだ」と、どなった。\\


　張飛のすがたを見ると、百姓たちは、声をそろえていった。\\


「旦那はまだなにもご存じないんですか。勅使さまは、県の吏人に、訴状を書かせて、都へさし送ったと申しますに」\\


「なんの訴状をだ」\\


「日頃、わしらが、お慕い申している、尉の玄徳さまが、百姓いじめなさるとか、\ruby{苛税}{かぜい}をしぼり取って、私腹を肥やしなすっているとか、何でも、二十ヵ条も罪をかき並べて、都へその訴状が差廻され、お沙汰が来次第に罰せられるとうわさに聞きましたで。……わしら、百姓どもは、玄徳さまを、親のように思っているので、皆の衆と打揃うて、勅使さまへおすがりにきたところ、\ruby{下吏}{したやく}たちに叩き出され、この通り、役所の門まで閉められてしもうたので、ぜひなくこうしているとこでござりまする」\\


　聞くと、張飛は、毛虫のような眉をあげて、閉めきってある役館の門をはったと睨みつけた。\\

\\

\subsection{\ruby{打風乱柳}{だふうらんりゅう}}

\\

\subsubsection{一}

\\


「おい」\\


　張飛はいった。大地に坐っている大勢の百姓町民へ向って、\\


「おまえ達は、退散しろ。これから俺がやることに、後で、かかり合いになるといけないぞ」\\


　しかし百姓たちは、泥酔しているらしい張飛が、何をやりだすのかと、そこを起っても、まだ附近から眺めていた。\\


　張飛は、門を打ち叩いて、\\


「番人どもっ、開けろ、開けなければ、ぶちこわすぞっ」と、どなり出した。\\


　役館の番卒は、「何者だっ」と、中から覗き合っていたが、\ruby{重棗}{ちょうそう}の如き\ruby{面}{おもて}に、\ruby{虎髯}{こぜん}を逆だて、怒れる形相に\ruby{抹{\fontspec{jigmo}\char"7843}}{まっしゅ}をそそいだ\ruby{巨漢}{おおおとこ}が、そこを揺りうごかしているので、\\


「誰だ、誰だ？」と、さわぎ立ち、県尉玄徳の部下だと聞くと、\ruby{督郵}{とくゆう}の家来たちは、\\


「開けてはならぬぞ」と、厳命した。そして人数をかためて、門の内へさらにまた、幾重にも\ruby{人墻}{ひとがき}を立ててひしめき合っていた。\\


　その気配に、張飛はいよいよ怒りを心頭に発して、\\


「よしっ、その分ならば！」\\


　門の柱へ両手をかけたと思うと、\ruby{地震}{ない}のようにみりみりとそれは揺れだして、あれよと人々の驚くうちに、すさまじい物音立てて内側へ仆れた。\\


　中にいた番卒や督郵の家来たちは、逃げおくれて、幾人かその下敷になった。張飛は、豹の如く、その上を躍り越えて、\\


「督郵はどこにいるかっ。督郵に会わんっ」と、\ruby{咆哮}{ほうこう}した。\\


　番卒たちは、それと見て、\\


「やるな」\\


「捕えろ」と、さえぎったが、\\


「えい、邪魔な」\\


　とばかり張飛は投げとばす、踏みつぶす、撲りたおす、あたかも一陣の旋風が、塵を巻いて\ruby{翔}{か}けるように、役館の奥へと躍りこんで行った。\\


　折から勅使督郵は、昼日中というのに、\ruby{帳}{とばり}を垂れて、この田舎町のひなびた\ruby{唄}{うた}い\ruby{女}{め}などを相手に酒をのんでいたところだった。\\


　\ruby{淫}{みだ}らな\ruby{胡弓}{こきゅう}の音を聞きつけて、張飛がその室をうかがうと、果たして正面の\ruby{榻}{とう}に美人を擁して酔いしれている高官がある。まぎれもない督郵だ。\\


　張飛は、帳を払って、\\


「やいっ\ruby{佞吏}{ねいり}、腐れ\ruby{吏人}{やくにん}。よくもわが義兄玄徳に汚名をぬりつけ、偽罪の訴状を作って都へ\ruby{上}{のぼ}せたな。先頃からの傲慢無礼といい、勅使たる身がこの態たらくといい、もはや堪忍はならぬ。天に代って、汝を懲らしめてやるからそう思え」\\


　眼は百錬の鏡にも似、髯はさかしまに立って、丹の如き\ruby{唇}{くち}を裂いた。\\


「――きゃっ」と、胡弓や琴をほうりだして\ruby{妓}{おんな}たちは\ruby{榻}{とう}の下へ逃げこんだ。\\


　督郵も、ちぢみ上がって、\\


「なんじゃ、待て、乱暴なことをするな」\\


　と、ふるえ声で、逃げかけるのを、張飛はとびかかって、\\


「どこへ行く」\\


　軽く一つ撲ったが、督郵は\ruby{顎}{あご}でもはずしたように、ぐわっと、歯をむいたままふん\ruby{反}{ぞ}った。\\


「じたばたするな」\\


　張飛は、その体を軽々と横に引っ抱えると、また疾風のように外へ出て行った。\\

\\

\subsubsection{二}

\\


　門外へ出てくると、\\


「犬にでも喰われろ」\\


　と、張飛は、引っ抱えてきた督郵のからだを、大地へたたきつけて罵った。\\


「汝のような腐敗した\ruby{佞吏}{ねいり}がいるから、天下が乱れるのだ。乱賊は打つも、佞吏を懲らす者はない。人のなし得ぬ正義をなし、人の抗し得ぬ権力に抗す。それを\ruby{旗幟}{きし}とする義軍の張飛を知らずや。やいっ」\\


　督郵の顔を踏んづけて、張飛がいうと、督郵は、手足をばたばたさせて、\\


「者どもっ。この\ruby{狼藉}{ろうぜき}を。――この乱暴者を、\ruby{搦}{から}め捕れ。誰かいないか」\\


　悲鳴に似た声でわめいた。\\


「やかましい」\\


　\ruby{髻}{もとどり}をつかんで引廻した上、張飛は、門前の巨きな柳の樹に目をつけて、\\


「そうだ、見せしめのために」\\


　と、督郵の両手を有りあう縄で縛りあげ、その縄尻を柳の枝に投げて、吊しあげた。\\


　柳から\ruby{生}{な}った人間のように、督郵の足は宙に浮いた。張飛は、彼が暴れても落ちないように縄の端を幹に巻いて、\\


「どうだ、やいっ」\\


　と、一本の柳の枝を折って、まずぴしりと一つ撲った。\\


「痛いっ」\\


「あたり前だ」と、また一つ打ち、\\


「悪吏の虐政に苦しむ人民の\ruby{傷}{いた}みはこんなものじゃないぞ。汝も、\ruby{廟鼠}{びょうそ}の一匹だろう。かの十\ruby{常侍}{じょうじ}などいう\ruby{佞臣}{ねいしん}の端くれだろう。その醜い面をさらせよ。その卑しい鼻の穴を天日に向けて\ruby{哭}{な}けっ。――こうか、こうか、こうしてやる」\\


　柳の枝は、すぐ粉々になった。\\


　また新しい柳の枝を折って撲りつけるのだった。三十、四十、五十、二百以上も打ちすえた。\\


　督郵は、見栄もなく、ひイひイひイと声をあげて、\\


「ゆるせ」と、泣き声だし、\\


「待て、待ってくれ。なんでもいう通りにするから」\\


　と、遂には、涙さえこぼして、あわれっぽく叫んだが、\\


「だめだ。その手は食わぬ」\\


　と、張飛は、乱打をやめなかった。\\


　その日も玄徳は、私宅に閉じこもって、\ruby{怏々}{おうおう}とすぐれない一日を過していたが、誰やらあわただしく門をたたく者があるので自身出てみると四、五名の百姓が、\\


「大変です。今、張飛さまが、お酒に酔って、役所の門をぶちこわし、勅使の高官を、柳の木に吊しあげて打ちすえております」\\


　と、告げて去った。玄徳は驚いて、そのまま馳けだして行った。\\


　折ふし居合せた関羽も、\\


「ちぇッ、張飛のやつ、また持病を起したか」\\


　と、舌打ちしながら、玄徳の後から馳けつけた。\\


　見ると、柳に吊されている督郵は、衣裳もやぶれ、\ruby{脛}{はぎ}は血を流し、顔面は紫いろにふくれていた。もう少し遅かったら、すんでのこと、撲り殺されていたであろう。\\


　仰天して、玄徳は、\\


「これっ、何をする」と、張飛の腕くびをつかんで叱りつけた。\\


　張飛は、大息つきながら、\\


「いや、止めないで下さい。民を害する逆賊とはこいつのことです。息のねを止めないでは俺の虫がおさまらん」\\


　と、玄徳のさえぎりなどは物ともせず、さらに、\ruby{柳鞭}{りゅうべん}をうならせて、督郵のからだを所きらわず打ちつづけた。\\

\\

\subsubsection{三}

\\


　悲鳴を放って、張飛の鞭にもがいていた督郵は、柳の梢から玄徳のすがたを見つけて、\\


「おお、それへ来たのは、県尉玄徳ではないか。公の部下の張飛が、酒に酔って、わしをかくの如く殺そうとしている。どうか早く止めてくれ。もしわしを助けてくれたなら、このまま、張飛の罪も不問にし、おん身には、帝に急使を立てて前の訴状をとどめ、代わるに充分な恩爵をもって\ruby{酬}{むく}ゆるであろう」と、叫んでまた、\\


「はやく助けてくれ」\\


　と何度も悲鳴をくり返した。\\


　そのいやしい言葉を聞くと、張飛の暴を制しかけていた玄徳も、かえって止める意志をさまたげられた。\\


　けれど、彼は、いかに\ruby{醜汚}{しゅうお}な人間であろうとも、勅命をうけて下った天子の使いである。玄徳は、叱咤して、\\


「止めぬかっ張飛」と、彼の手から柳の枝を奪い、その枝をもって張飛の肩を一つ打った。\\


　玄徳に打たれたことは初めてである。さすがの張飛も、はっと顔色を醒まして棒立ちになった。もちろん不平満々たる色をあらわしてではあったが。\\


　玄徳は、柳の幹の縄を解いて、督郵のからだを大地へ下ろしてやった。すると、それまで、是とも非ともいわず黙って見ていた関羽が、つと馳け寄ってきて、\\


「長兄。お待ちなさい」\\


「なぜ」\\


「そんな人間を助けてやったところで、所詮、むだなことです」\\


「何をいう。わしはこの人間から利を得るために助けようとするのではない。ただ、天子の御名を\ruby{畏}{おそ}るるのみだ」\\


「わかっております。しかしそういうお気持も、いったいどこに通じましょうか。前には、身命を\ruby{賭}{と}して、大功を立てておられながら、わずか一県の尉に封ぜられたのみか、今また、督郵のごとき腐敗した中央の吏に、最大の侮辱をうけ、黙っていれば、罪もなき罪におとし入れられようとしているではありませんか」\\


「……ぜひもない」\\


「ぜひもないことはありません。こんな不法は蹴とばすべきです。先頃からそれがしもつらつら思うに、\ruby{枳棘叢中}{ききょくそうちゅう}\ruby{鸞鳳}{らんほう}の\ruby{栖}{す}む所に非ず――と昔からいいます。\ruby{棘}{いばら}や\ruby{枳}{からたち}のようなトゲの木の中には良い\ruby{鳳}{とり}は自然栖んでいない――というのです。われわれは栖む所を誤りました。\ruby{如}{し}かずいちど身を退いて、別に遠大の計をはかり直そうではありませんか」\\


　関羽には、時々、\ruby{訓}{おし}えられることが多い。やはり学問においては、彼が一日の長を持っていた。\\


　玄徳はいつも聴くべき言はよく聴く人であったが、今も、彼の言をじっと聞いているうちに、大きくうなずいて、\\


「そうだ。……いいことをいってくれた。我れ栖む所を誤てり」\\


　と、胸にかけていた県尉の\ruby{印綬}{いんじゅ}を解いて、督郵にいった。\\


「卿は、民を害する賊吏、今その\ruby{首}{こうべ}を斬って、これに\ruby{梟}{か}けるはいと易いことながら、恥を思わぬ悲鳴を聞けば、畜類にも\ruby{不愍}{ふびん}は生じる。あわれ、犬猫と思うて助けてとらせる。――そしてこの印綬は、卿に託しておく。我れ今、官を捨てて去る。中央へよろしくこの\ruby{趣}{おもむき}を取次ぎたまえ」\\


　そして張飛、関羽のふたりをかえりみて、\\


「さ。行こう」\\


　と、風の如くそこを去った。\\


　\ruby{霏々}{ひひ}と散りしいた柳葉の地上に督郵は、まだ何か、苦しげに\ruby{喚}{わめ}いていたが、玄徳らの姿が遠くなるまで、前に懲りて、近づいていたわり助ける者もなかった。\\

\\

\subsection{\ruby{岳南}{がくなん}の\ruby{佳人}{かじん}}

\\

\subsubsection{一}

\\


　いっさんに馳けた玄徳らは、ひとまず私宅に帰って、私信や文書の\ruby{反故}{ほご}などみな焼きすて、その夜のうちに、この地を退去すべくあわただしい身支度にかかった。\\


　官を捨てて野に去ろうとなると、これは張飛も大賛成で、わずかの手兵や召使いを集め、\\


「ご主人には今度、にわかに、思うことがあって、県の\ruby{尉}{い}たる官職を辞め、しばらく野に下って、悠々自適なさることになった。しかし、実はおれが勅使督郵を半殺しの目にあわせたのが\ruby{因}{もと}だ。ついては、身の落着きの目あてのある者は、家に帰れ。あてのない者は、病人たりとも、捨てては行かぬ。苦楽を共にする気でご主人に従って参れ」と、いい渡した。\\


　貰う物を貰って、自由にどこかへ去る者もあり、どこまでも、玄徳様に従ってと、残る者もあった。\\


　かくて夜に入るのを待ち、手廻りの家財を\ruby{驢}{ろ}や車に積み、同勢二十人ばかりで、遂に、官地安喜県を後に、闇にまぎれて落ちて行った。\\


　――一方の督郵は。\\


　あの後、間もなく、下吏の者が寄ってきて、役所の中へ抱え入れ、手当を加えたが、五体の傷は火のように痛むし、大熱を発して、幾刻かは、まるで人事不省であった。\\


　だが、やがて少し落着くと、\\


「県尉の玄徳はどうしたっ」\\


　と、うわ\ruby{言}{ごと}みたいに呶鳴った。\\


　その玄徳は、官の印綬を解いて、あなたの首へかけると、捨てぜりふをいって馳け走りましたが、今宵、一族をつれて夜逃げしてしまったという噂です――と側の者が告げると、\\


「なに。逃げ落ちたと。――ではあの張飛という奴もか」\\


「そうです」\\


「おのれ、このまま、おめおめと無事に、逃がしてなろうか。――つ、つかいを、すぐ急使をやれっ」\\


「都へですか」\\


「ばかっ。都へなど、使いを立てていたひには間にあうものか。ここの\ruby{定州}{ていしゅう}（河北省・保定・正定の間）の太守へだ」\\


「はっ。――何としてやりますか」\\


「玄徳、常に民を\ruby{虐}{ぎゃく}し、こんど勅使の巡察に、その罪状の発覚を恐るるや、かえって勅使に暴行を加え、良民を\ruby{煽動}{せんどう}して乱をたくめど、その事、いちはやく官の知るところとなり、一族をつれて夜にまぎれ、無断官地を捨てて逃げ去る――と」\\


「はっ。わかりました」\\


「待て。それだけではいかん。すぐさま、\ruby{迅兵}{じんぺい}をさし向けて、玄徳らを召捕え、都へご\ruby{檻送}{かんそう}くださるべしと、促すのだ」\\


「心得ました」\\


　早馬は、定州の府へ飛んだ。\\


　定州の太守は、\\


「すわ、大事」と、勅使の名におそれ、また、督郵の\ruby{詭弁}{きべん}にも、うまく乗せられて、八方へ物見を走らせ、玄徳たちの落ちていった先を探させた。\\


　数日の後。\\


「何者とも知れず、安喜県のほうから\ruby{代州}{だいしゅう}（山西省・代県）のほうへ向って、\ruby{驢車}{ろしゃ}に家財を積み、十数名の従者をつれ、そのうち三名は、驢に乗った浪人風の人物で、北へ北へとさして行ったということでありますが」\\


　との報告があった。\\


「それこそ、玄徳であろう。からめ捕って、都へ差立てろ」\\


　定州の太守の命をうけて、即座に鉄甲の迅兵約二百、ふた手にわかれて、玄徳らの一行を追いかけた。\\

\\

\subsubsection{二}

\\


　北へ、北へ、車馬と落ち行く人々の影はいそいだ。\\


　幾度か、他州の兵に襲われ、幾度か追手の\ruby{詭計}{きけい}に墜ちかかり、百難をこえ、ようやくにして代州の五台山下までたどりついた。\\


「張飛。御身の指図で、ここまではやって来たが何か落着く先の\ruby{目的}{めあて}はあるのか。――ここはもう、五台山の麓だが」\\


　関羽もいうし、玄徳も、実は案じていたらしく、\\


「いったい、これからどこへ落着こうという考えか」と、共々に訊ねた。\\


「ご安心なさるがよい」\\


　張飛は大のみこみで云った。そして\ruby{岳麓}{がくろく}の平和そうな村へ行き着くと、\\


「しばらく、ご一同は、その辺に車馬を休めて待っていて下さい」\\


　と、一人でどこへか立ち去ったが、ほどなく立ち帰ってきて、\\


「\ruby{劉大人}{りゅうたいじん}が承知してくれました。もう大船に乗った気でおいでなさい」\\


　と告げた。\\


「劉大人とは、どこの何をしておる人物かね」\\


「この土地の大地主です。まあ大きな郷士といったような家柄と思えばまちがいありません。常に百人や五十人の食客は平気で邸においているんですから、われわれ二十人やそこらの者が厄介になっても、先は平気です。またこの地方の人望家でもありますから、しばらく身をかくまっておいてもらうには、なによりな場所でしょうが」\\


「それは願ってもないことだが、御身との間がらは、どういう仲なのだ」\\


「劉大人も、今こそ、こんな田舎にかくれて、岳南の隠士などと気どっていますが、以前は、拙者の旧主\ruby{鴻家}{こうけ}とは血縁もあって、軍糧兵馬の相談役もなされ、何かと、旧主鴻家とは、往来しておったのであります。――その頃、自分も鴻家の一家臣として、ご懇意をねがっていたので、鴻家が滅亡の後も、実は、拙者の飲み\ruby{代}{しろ}だの、遺臣の始末などにも、ずいぶんご厄介になったもので」\\


「そうか。その上にまた、同勢二十人も、食客をつれこんでは、劉大人も、眉をひそめておいでだろう」\\


「そんな事はありません。非常に浪人を愛する人ですし、玄徳様のご素姓と、われわれ義軍が、官地を捨てて去ったことなど、つぶさにおはなししたところ、苦労人ですから、非常によく分ってくれて、二年でも三年でもいるがいいというわけなんで」\\


　張飛のことばに、\\


「そういう人物の邸なら身を寄せてもよかろう」\\


　と、玄徳も安心して、彼の案内について行った。\\


　すると、\ruby{岳麓}{がくろく}の\ruby{疎林}{そりん}のほとりに、一廓の宏壮な土塀が見えた。玄徳らを\ruby{誘}{いざな}いながら、張飛が、\\


「あの邸です。どうです、まるで豪族の家のようでしょう」と、自分の住居ででもあるように誇って云った。\\


　玄徳がふと\ruby{驢}{ろ}を止めて見ていると、その邸の並びの\ruby{杏}{あんず}の並木道を今、\ruby{鄙}{ひな}には\ruby{稀}{まれ}な麗人が、白馬に乗って通ってゆくのが見えた。美人の驢の後からは、ひとりの童子が、琴を\ruby{担}{にな}って、眠そうに供をして行った。\\

\\

\subsubsection{三}

\\


「はて、どこかで見たような」\\


　玄徳はふとそんな気がした。\\


　遠目ではあったが、妙に印象づけられた。もっとも、\ruby{殺伐}{さつばつ}な戦場生活だの、\ruby{僻地}{へきち}から\ruby{曠野}{こうや}を\ruby{流浪}{るろう}してきた身なので、よけいに、彼方の女性が美しく見えたのかもしれない。\\


　麗人は、すぐ広い土塀に囲まれた、豪家の門のうちへ入ってしまった。\\


「そこが劉大人の邸だ」\\


　と、たった今、張飛に教えられたばかりなので、さては劉家の息女かなどと、玄徳はひとり想像していた。\\


　ほどなく、玄徳らの一行も、そこの門前に着いた。一同は車を停め、驢から降りて、\ruby{埃}{ほこり}まみれな旅の姿をかえりみた。\\


　ここの\ruby{主}{あるじ}は浪人を愛し、常に多くの食客を養っているという。どんな人物であろうか、玄徳や関羽は、会わないうちはいろいろに想像された。\\


　けれど、張飛に案内されて、\ruby{南苑}{なんえん}の客館に通ってみると、まったく世の風雲も知らぬげな\ruby{長閑}{のど}けさで、浪人を愛するよりは、むしろ風流を愛すことのはなはだしい気持の逸人ではないかと思われた。\\


　やがてのこと、\\


「はい、てまえが、当家の\ruby{主}{あるじ}の\ruby{劉恢}{りゅうかい}です。ようお越しなされました。お身の上は最前、張飛どのから聞きましたが、どうぞお気がねなく、一年でも二年でも遊んでいてください。その代りこんな田舎ですから、何もおかまいはできませんよ、豊かにあるのは、酒ぐらいなもので」\\


　こう主の\ruby{劉恢}{りゅうかい}が出てきてのあいさつに、張飛は、\\


「ありがたい。酒さえあれば何年だっていられますよ」\\


　と、もう贅沢をいう。\\


　玄徳はいんぎんに、\\


「何分」\\


　と、しばらくの\ruby{逗留}{とうりゅう}を頼み、関羽も姓名や郷地を名乗って、将来の\ruby{高誼}{こうぎ}を仰いだ。\\


　劉大人は、いかにも大人らしい\ruby{寡言}{かげん}な人で、やがて召使いをよび、三名の部屋として、この南苑の客館を提供し、何かの事などいいつけ、ほどなく奥へかくれてしまった。\\


「どうです、落着くでしょう」\\


　張飛は手がら顔にいう。\\


「落着きすぎるくらいだ」と、関羽は笑って、\\


「ぼろを出さぬようにしてくれよ」\\


　と、暗に張飛の酒ぐせをたしなめた。\\


　年を越えた。春になった。\\


　五台山下の部落は、まことに平和なものだった。ここには、劉恢が土豪として、\ruby{村長}{むらおさ}の役目をも兼ねているせいか、悪吏も棲まず、匪賊の害もなかった。\\


　しかし、張飛や関羽は、その余りにも無事なのにむしろ苦しんだ。酒にも平和にも\ruby{倦}{う}んだ。\\


　それとは違って、玄徳は近ごろひどく無口であった。常に物思わしいふうが見える。\\


「長兄も、この頃はようやく、ふたたび戦野が恋しくなってきているのではないかな。風雲児、とみに元気がないが」\\


　ある時、関羽がいうと、\\


「いやいや、戦野が恋しいのじゃないさ」\\


　と張飛は首を振った。\\


「では、郷里の母御のことでも案じておられるのかな？」\\


「それもあろうが、原因はもっとべつなほうにある。おれはそう\ruby{覚}{さと}っているが、わざと会わせないんだ」\\


「ふウむ。原因があるのか」\\


「ある」\\


　\ruby{苦々}{にがにが}しげに張飛はいった。その顔つきで思い出した。近頃、南苑に\ruby{梨花}{りか}が咲いて、夜は春月がそれに霞んでまたなく\ruby{麗}{うるわ}しい。時折その梨苑をさまよう月よりも美しい佳人が見かけられる。そうするといつのまにか、この客館から玄徳のすがたが見えなくなるのだった。\\

\\

\subsubsection{四}

\\


　張飛のはなしを聞いて関羽にも思い当るふしがあった。関羽はそれから特に玄徳の容子に注目していた。\\


　すると、それから数日後の宵であった。その夜は\ruby{朧月}{おぼろづき}が麗しかった。五台山の半身をぼかした夜霞が野にかけ銀を\ruby{刷}{は}いたような朧をひいていた。\\


「おや、いつのまにか」\\


　気がついて関羽はつぶやいた。三名して食卓を囲んでいたのである。張飛は例によっていつまでも酒をのんでいるし、自分も、杯をもって相手になっていたが、玄徳は室を去ったとみえて、彼の空席の卓には、皿や\ruby{酒盞}{しゅさん}しか残っていない。\\


「そうだ」\\


　こよいこそ彼の行動をつきとめてみよう。関羽はそう考えたので張飛にも黙って急に部屋から出て行った。\\


　そして南苑の白い梨花の\ruby{径}{こみち}を忍びながら歩いては見まわした。\\


　もう奥の内苑に近い。主の\ruby{劉恢}{りゅうかい}のいる棟やその家族らの住む棟の燈火は林泉をとおして彼方に見えるのであった。\\


「はて。これから先へゆく筈もないし」\\


　関羽がたたずんでいると、ほど近い木の間を、誰か、\ruby{楚々}{そそ}と通る人があった。見ると、劉恢の\ruby{姪}{めい}とかいうこの家の妙齢な麗人であった。\\


「……ははあ？」\\


　関羽は自分の予感があたってかえって肌寒いここちがした。物事の裏とか、人の秘密とかには、むしろ\ruby{面}{おもて}を横にして、無関心でいたい彼であったが、つい後から忍んで行った。\\


　劉恢の姪という佳人は、やがて\ruby{鮮}{あざ}やかに月の下に立った。辺りには木蔭もなく物の蔭もなく、ただ広い芝生に夜露が宝石をまいたように光っていた。\\


　すると梨の花の\ruby{径}{こみち}からまたひとりの人影が\ruby{忽然}{こつぜん}と立ち上がった。それは花の中に隠れていた若い男性であった。\\


「オ。玄徳さま」\\


「\ruby{芙蓉}{ふよう}どの」\\


　ふたりは顔を見あわせてニコと笑み交わした。芙蓉の歯が実に美しかった。\\


　相寄って、\\


「よく出られましたね」\\


　玄徳がいう。\\


「ええ」\\


　芙蓉はさしうつ向く。\\


　そして梨畑のほうへ、ふたりは背を\ruby{擁}{よう}し合いながら歩み出して、\\


「劉恢は、あれでとても、厳格な人ですからね。……食客や豪傑たちには、やさしい温情を示す人ですけれど、家庭の者には、おそろしくやかましい人なんです。……ですから、……、こうして\ruby{苑}{にわ}へ出てくるにも、ずいぶん苦心して来るんですの」\\


「そうでしょう。何しろ、われわれのような食客が常に何十人もいるそうですからね。私も、関羽だの張飛だのという腹心の者が、同じ室にいて、眼を光らしているので、彼らにかくれて出てくるのもなかなか容易ではありません」\\


「なぜでしょうね」\\


「何がですか」\\


「そんなにお互いに苦労しながらでも、夜になると、どうしてもここへ出てきたいのは」\\


「私もそうです。自分の気もちがふしぎでなりません」\\


「美しい月ですこと」\\


「夏や秋の冴えた頃よりも、今頃がいいですね。夢みているようで」\\


　梨の花から梨の花の径をさまよって、二人は飽くことを知らぬげであった。夢みようと意識しながら、あえて、夢を追っているふうであった。\\

\\

\subsubsection{五}

\\


　この家の\ruby{深窓}{しんそう}の\ruby{佳人}{かじん}と玄徳とが、いつのまにか、\ruby{春宵}{しゅんしょう}の秘語を楽しむ仲になっているのを目撃して、関羽は、非常なおどろきと\ruby{狼狽}{ろうばい}をおぼえた。\\


「ああ、平和は雄志を\ruby{蝕}{むしば}む」\\


　彼は、慨嘆した。\\


　見まじきものを見たように関羽はあわてて後苑の梨畑から馳け戻ってきた。そして客館の食卓の部屋をのぞくと、張飛はただ一人でまだそこに酒を飲んでいた。\\


「おい」\\


「やあ、何処へ行っていたのだ」\\


「まだ飲んでいるのか」\\


「飲むよりほかに為すことはないじゃないか。いかに\ruby{脾肉}{ひにく}を嘆じたところで、時利あらず、風雲招かず、\ruby{蛟龍}{こうりょう}も淵に\ruby{潜}{ひそ}んでいるしかない。どうだ、貴公も酒の淵に潜まんか」\\


「一杯もらおう、実は今、いっぺんに酒が醒めてしまったところだ」\\


「どうしたのか」\\


「……張飛」\\


「ウム」\\


「おれは、貴様のように、いたずらに現在の世態や時節の来ぬことを、そう悲観はしないつもりだが、今夜はがっかりしてしまった。――\ruby{野}{や}に隠れ淵に潜むとも、いつか蛟龍は風雲を\ruby{捉}{とら}えずにいないと信じていたが」\\


「ひどく失望の態だな」\\


「もう一杯くれ」\\


「めずらしく飲むじゃないか」\\


「飲んでから話すよ」\\


「なんだ」\\


「実は今、おれは、人の秘密を見てしまった」\\


「秘密を」\\


「されば。先頃から貴様が謎めいたことをいうので、こよい玄徳様が出て行った後からそっと\ruby{尾}{つ}けて行ってみたのだ。するとどうだろう……ああ、おれは語るに忍びん。あんな柔弱な人物だとは思わなかった」\\


「なにを見たのだ一体」\\


「あろうことかあるまいことか。当家の\ruby{深窓}{しんそう}に養われている\ruby{芙蓉娘}{ふようじょう}とかいう麗人と、逢引きをしているではないか。ふたりはいつのまにか恋愛におちておるのだ。われわれ義軍の盟主ともある者が、一女性に心をとらわれなどして何ができよう」\\


「そのことか」\\


「貴様は前から知っていたのか」\\


「うすうすは」\\


「なぜわしに告げないのだ」\\


「でも、できてしまっているものは仕方がないからな」\\


　張飛も腐った顔つきしてつぶやいた。その顔を頬杖にのせて、片手で独り酒を\ruby{酌}{つ}いであおりながら、\\


「英傑児も、あまり平和な温床に長く置くと\ruby{黴}{かび}が生えだして、ああいうことになるんだな」\\


「志を得ぬ\ruby{鬱勃}{うつぼつ}をそういうほうへ誤魔化しはじめると、人間ももうおしまいだな。……また、あの女も女ではないか。あれは\ruby{劉恢}{りゅうかい}の娘でもないし、いったい何だ」\\


「訊かれると面目ない」\\


「なぜ？　なぜ貴様が面目ないのか」\\


「……実はその、あの芙蓉娘は拙者の旧主\ruby{鴻家}{こうけ}のご息女なので、劉恢どのも鴻家とは浅からぬ関係があった人だから、主家鴻家の没落後、おれが芙蓉娘をこの家へ連れてきて、\ruby{匿}{かくま}っておいてくれるように頼んだお方なのだ」\\


「え。では貴様の旧主のご息女なのか」\\


「まだ義盟を結ばない数年前のはなしだが、その芙蓉娘と玄徳様とは、\ruby{黄匪}{こうひ}に追われて、お互いに危うい災難に見舞われていた頃、偶然、或る地方の古塔の下で、出会ったことがあるので、とっくに双方とも知り合っていた仲なのだ」\\


「え。そんなに古いのか」\\


　関羽が呆れ顔した時、室の外に誰かの\ruby{沓音}{くつおと}が聞えた。\\

\\

\subsubsection{六}

\\


　\ruby{主}{あるじ}の劉恢であった。\\


　劉恢は、室内の様子を見て、\\


「おさしつかえないですか」と、二人の許しをうけてから入ってきた。そしていうには、\\


「困ったことができました。数日の内に、洛陽の巡察使と定州の太守が、この地方へ巡遊に来る。そしてわしの邸がその宿舎に当てられることになった。当然、あなた方の潜伏していることが発覚する。一時どこかへ隠れ場所をお移しなさらぬと危ないが」\\


　という相談であった。\\


　折も折である。\\


　関羽、張飛も、一時は途方にくれたここちがしたが、むしろこれは、天が自分らの\ruby{懶惰}{らんだ}を\ruby{誡}{いまし}むるものであると思って、\\


「いや、ご当家にも、だいぶ長い間の逗留となりました。そういうことがなくても、このへんで一転機する必要がありましょう。いずれわれわれども三名で相談の上、ご返辞申しあげます」\\


「なんともお気の毒じゃが。……なお、落着く先にお心当りもなければ、わしの信じる人物で安心のなる所へご紹介もして上げますから」\\


　劉恢は、そういって、戻って行った。\\


　後で、二人は顔見あわせて、\\


「玄徳様と芙蓉娘の仲を、主もさとってきて、これはいかんと、急にあんな口実をいってきたのではあるまいか」\\


「さあ。どうとも知れぬ」\\


「しかし、いい\ruby{機}{しお}だ」\\


「そうだ。玄徳様のためには至極いいことだ」\\


　翌朝。二人はさっそく、「\ruby{云々}{しかじか}のわけですが」と、玄徳に主の旨を伝えて、善後策をはかった。\\


　すると玄徳は、一時は、はっとした顔色だったが、直ぐうつ向いた眼ざしをきっとあげて、\\


「立退こう。恩人の劉大人にご迷惑をかけてもならぬし、自分もいつまで安閑とここにいる気もなかったところだから」と、いった。\\


　そういう玄徳の面には、深く現在の自身を反省しているらしい容子が見えた。\\


　そこで関羽は、思いきって、こういってみた。\\


「――ですが、お名残り惜しくはありませんか、この家の深窓の佳人に」\\


　玄徳は微笑のうちにも、幾分か\ruby{羞恥}{しゅうち}の色をたたえながら、\\


「\ruby{否}{いな}とよ、恋は路傍の花」\\


　と、答えた。\\


　その一言に、\\


「さすがは」\\


　と関羽も、自分の取越し苦労を打消し、すっかり眉をひらいた。\\


「そういうお気持なら安心ですが、実は、われわれの盟主たりまた、大望を抱いている英傑児が、一女性のために、壮志を\ruby{蝕}{むしば}まれてしまうなどとは、残念至極だと、張飛と共に、ひそかに案じていたところなのです。――ではあなたは飽くまで、芙蓉娘と本気で恋などにおちているわけではありますまいな」\\


「いや」\\


　玄徳は、正直にいった。\\


「恋をささやいている間は、恥かしいが、わしは本気で恋をささやいているよ。女を\ruby{欺}{あざむ}けない、また、自分も欺けない。ただ、恋あるのみだ」\\


「え……？」\\


「だが両君。乞う、安んじてくれ給え。玄徳はそれだけが全部にはなりきれない。恋のささやきも一ときの間だ。すぐわれに返る。\ruby{中山靖王}{ちゅうざんせいおう}の\ruby{後裔}{こうえい}劉備玄徳というわれに返る。寒村の\ruby{田夫}{でんぷ}から身を起し、義旗をひるがえしてからすでに両三年、戦野の\ruby{屍}{しかばね}となりつるか、洛陽の府にさまよえるか、と故郷には今なお、わが子の我を待ち給う老母もいる。なんで大志を失おうや。……両君も、それは安心して可なりである。玄徳を信じてくれい」\\

\\

\subsection{\ruby{故園}{こえん}}

\\

\subsubsection{一}

\\


　その翌日である。玄徳たち三名は、にわかに五台山麓の地、\ruby{劉恢}{りゅうかい}の邸宅から一時身を去ることになった。\\


　別れにのぞんで、主の劉恢は、\ruby{落魄}{らくはく}の豪傑玄徳らのために別離の小宴をひらいて、さていうには、\\


「また、時をうかがって、この地へぜひ戻っておいでなさい。お連れになってきた二十名の兵や\ruby{下僕}{しもべ}たちは、それまでてまえの邸に預かっておきましょう。そして今度お見えになった時こそ、再起のご準備におかかりなさい。黄巾の乱は小康を得ても、洛陽の王府そのものに\ruby{自潰}{じかい}の\ruby{兆}{きざ}しがあらわれてきています。せっかく、自重自愛して、どうか国家のために尽してください」\\


「ありがとう」\\


　四人は起って乾杯した。\\


　劉恢のいうように、ここへくる時連れて来た二十名ばかりの一族郎党の身は、皆、劉家に託しておいて、関羽、張飛、玄徳、思い思いに別れて一時身をかくすこととなった。\\


　が――劉家の門を出る時は、三人一緒に出た。世間の眼もあるので、劉恢はわざと見送らなかった。けれど、邸内の楼台から三名の姿が遠くなるまで独り見送っている美人があった。いうまでもなく\ruby{芙蓉娘}{ふようじょう}であった。\\


　張飛は知っていた。\\


　しかし、わざと何もいわなかった。玄徳も黙々と歩いていた。\\


　もう五台山の影も後ろに遠く霞んでから、張飛がそっと玄徳へいった。\\


「きのうお言葉を伺って、もう自分らもあなたの心事を疑うような気もちは抱いておりません。むしろ大丈夫の多情多恨のおこころを推察しておりますよ。例えば、私が酒を愛するようなものですからな」\\


　彼は、酒と恋を、一つものに考えているのだ。\\


　その程度だから、玄徳の心に同情するといっても、およそ玄徳の感傷とははなはだ遠いものにちがいなかった。\\


「――だが、長兄」と、張飛はまた、玄徳の顔をさし覗いて云った。\\


「豪傑は色に触るべからずという法はない。あなただって一生涯独身でいられるわけもない。ほんとに芙蓉娘がお好きならこの張飛が話してどんなことにでもします。拙者にとっては、旧主のご息女ではあるし、ああいう頼りのないお身の上ですからむしろあなたに願っても生涯を見ていただきたいくらいなものですよ。けれど今はいけませんな。時でないでしょう。志を得た後のことにね」\\


「わかったよ」\\


　玄徳は、うなずいた。\\


　それから州道の道標の下まで来ると、\\


「じゃあ、わしはここから一人別れて、ひとまず郷里の\ruby{{\fontspec{jigmo}\char"6DBF}県}{たくけん}へ行くからね、いずれまた、一度この五台山下へ戻って来るが」と、いった。\\


　張飛も、関羽も、各{\fontspec{jigmo}\char"303B}そこから別れて、ひとまず思い思いに落ちてゆくつもりであったが、片時の間も離れたことのない三人なので、さすがに寂しげに、\\


「こんどはいつここで会おう」\\


「この秋」\\


　玄徳がいう。二人はうなずいて、\\


「ではあなたはこれから{\fontspec{jigmo}\char"6DBF}県の母御のもとへおいでになるつもりですか」\\


「うム。ご無事なお顔だけ拝したら、またすぐ風雲の\ruby{裡}{うち}へ帰ってくる。涼秋の八月、再び三人して、五台山の月を見よう」\\


「おさらば」\\


「気をつけて」\\


「お互いに」\\


　三名は三方の道へ、しばし別離の姿をかえりみ合った。\\

\\

\subsubsection{二}

\\


　関羽と張飛のふたりに別れてから、玄徳は姿を土民のふうに変えて、ただ一人、故郷の\ruby{{\fontspec{jigmo}\char"6DBF}県楼桑村}{たくけんろうそうそん}へ、そっと帰って行った。\\


「ああ、桑の木も変らずにある……」\\


　何年ぶりかで、わが家の門を見た玄徳は、そこに立つと一番先に、例の巨きな桑の大樹を、懐かしげに見上げていた。\\


　――かたん。\\


　――ことん、かたん。\\


　すると\ruby{蓆}{むしろ}を織る\ruby{機}{はた}の音が家の裏のほうで聞えた。玄徳は、はっと心を打たれた。ここ両三年は馬上に長槍をとって、忘れはてていたが、幼少から衣食してきた\ruby{生業}{なりわい}の\ruby{莚織}{むしろおり}の\ruby{機}{はた}は、今なお、この故郷の家では休んでいなかった。\\


　その機を、その\ruby{筬}{おさ}を、今も十年一日のごとく動かしている者は誰だろうか。\\


　問うまでもない、玄徳の母であった。征野に立った息子の後を、ひとり留守している老いたる母にちがいなかった。\\


「いかにお淋しいことであったろう。また、ご不自由なことであったろう」\\


　家にはいらぬうちに、玄徳はもう\ruby{瞼}{まぶた}を涙でいっぱいにしていた。思えば幾年の間、転戦また転戦、故郷の母に衣食の費を送るいとまさえなかった。便りすら幾度か数えるほどしかしていなかった。\\


　――すみません。\\


　彼はまず故園の荒れたる門に心から詫びて、そして機の音の聞える裏のほうへ馳けこんで行った。\\


　ああそこに、黙然と、\ruby{蓆}{むしろ}を織っている白髪の人。――玄徳は見るなり後ろから馳け寄って、母の足もとへ、\\


「母上っ」\\


　ひざまずいた。\\


「――母上。わたくしです。今帰って参りました」\\


「……？」\\


　老母は、驚いた顔して、機の手を休めた。そして、玄徳のすがたをじっと見て、\\


「……阿備か」\\


　と、いった。\\


「長い間、お便りもろくにせず、定めし何かとご不自由でございましたろう。陣中心にまかせず、転戦からまた転戦と、戦に暮れておりましたために」\\


　子の言葉をさえぎるように、\\


「阿備。……そしておまえはいったい、なにしに帰ってきたのですか」\\


「はい」\\


　玄徳は地に面を伏せて、\\


「まだ志も達せず、晴れて母上にお目にかかる時機でもありませんが、先頃から官地を去って、野に\ruby{潜}{ひそ}んでおりますゆえ、役人たちの目をぬすんで、そっとひと目、ご無事なお顔を見に戻って参りました」\\


　老母の眼は明らかにうるんでみえた。髪もわずかのうちに梨の花を盛ったように雪白になっていた。眼もとの肉もやつれてみえるし――\ruby{機}{はた}にかけている手は\ruby{藁}{わら}ゴミで荒れている。\\


　しかし、以前にかわらないものは、子に対してじっと向ける眸の大きな愛と\ruby{峻厳}{しゅんげん}な強さであった。こぼれ落ちそうな涙をもこらえて、老母は、静かにいうのだった。\\


「阿備……」\\


「はい」\\


「それだけで、そなたはこの家へ帰っておいでなのかえ」\\


「え。……ええ」\\


「それだけで」\\


「――母上」\\


　すがり寄る玄徳の手を、老母は、藁ゴミとともに\ruby{裳}{もすそ}から払って、たしなめるようにきつく云った。\\


「なんです。\ruby{嬰児}{あかご}のように。……それで、おまえは憂国の偉丈夫ですか。帰ってきたものはぜひもないが、長居はなりませんぞ。こよい一晩休んだら、すぐ出てゆくがよい」\\

\\

\subsubsection{三}

\\


　思いのほかな母の不機嫌な\ruby{気色}{けしき}なのである。それも、自分を励まして下さるためと、劉玄徳は、かえって大きな愛の下に泣きぬれてしまった。\\


　母は、その子を、大地に見ながら、なお叱っていった。\\


「まだおまえが郷土を出てから、わずか二年か三年ではないか。貧しい武器と、訓練もない郷兵を集めて、このひろい天下の騒乱の中へ打って出たおまえが、たった三年やそこらで、功を遂げ名をあげて戻ってこようなどと……そんな夢みたいなことを母は考えて待っておりはしない。……世の中というものはそんな単純ではありません」\\


「母上。……玄徳の\ruby{過}{あやま}りでございました。どこへ行っても、自分の正義は通らず、戦っても戦っても、なんのために戦ったのか、この頃、ふと失意のあまり疑いを抱いたりして」\\


「\ruby{戦}{いくさ}に勝つことは、強い豪傑ならば、誰でもすることです。そういう正しい道のさまたげにも、自分自身を時折に襲ってくる弱い心にも打ち克たなければ、\ruby{所詮}{しょせん}、大事はなし遂げられるものではあるまいが」\\


「……そうです」\\


「ようく、お分りであろう。……もうそなたも三十に近い男児。それくらいなことは」\\


「はい」\\


「そこらの豪傑たちが、乱世に乗じて、一州一郡を\ruby{伐取}{きりと}りするような小さい望みとは違うはずです。漢の宗室の末孫、中山靖王の\ruby{裔}{えい}であるおまえが、万民のために、剣をとって起ったのですよ」\\


「はい」\\


「千億の民の幸いを思いなさい。老先のないこの母ひとりなどが何であろう。そなたの心が――せっかく奮い起した大志が――この母ひとりのために\ruby{鈍}{にぶ}るものならば、母は、億民のために生命を縮めても、そなたを励ましたいと思うほどですよ」\\


「あ。母上」\\


　玄徳は驚いて、ほんとにそういう決心もしかねない母の\ruby{袂}{たもと}にすがって、\\


「悪うござりました。もう決して\ruby{女々}{めめ}しい心はもちません。あしたの朝は、夜の明けぬうちにここを去りますから、どうかただ一晩だけお側において下さいまし」\\


「…………」\\


　老母も、くずれるように、地へ膝をついた。そして、玄徳の体を、そっと抱いて、白髪の\ruby{鬢}{びん}をふるわせながらささやいた。\\


「阿備や……。だが、わたしはね、亡きお父さんの代りにもなっていうのだよ。今のは、お父さまのお声だよ。お叱りだよ。――あしたの朝は、近所の人の人目にかからないように、暗いうちに立っておくれね」\\


　そういうと、老母はいそいそと\ruby{母屋}{おもや}のほうへ立ち去った。\\


　間もなく、\ruby{厨}{くりや}のほうから、\ruby{夕餉}{ゆうげ}を\ruby{炊}{かし}ぐ煙が這ってきた。失意の子のために、母はなにか温かい物でも夕餉にと煮炊きしているらしいのであった。\\


　玄徳は、その間に、\ruby{蓆機}{むしろばた}へ寄って、織りのこして行った幾枚かの蓆を織りあげていた。\\


　手もとが暗くなってくる。白い夕星がもう上にあった。\\


　\ruby{機}{はた}を離れて、彼はひとり、裏の桃林を逍遥していた。はや晩春なので、桃の花はみな散り尽して黒い花の\ruby{蕋}{しべ}を梢に見るだけであった。\\


「ああ。故園は変らない――」\\


　玄徳は嘆じた。\\


　桃花はまた春に若やぐが、母の白髪が再び黒くかえる日はない。春秋は人の身のうえにのみ短い。しかも自分の思う望みは遠くまた大きく、いつの日、彼の母が心のそこからよろこんでくれる時がくるだろうか、考えると、いたずらに大きな嘆声が出るばかりであった。\\


「――阿備やあ。阿備やあ」\\


　もう暗い母屋のほうでは、母が\ruby{夕餉}{ゆうげ}のできたことを告げて呼んでいる。玄徳は、なんの悩みもなかった少年の頃を思い出して、少年のように遠くから高く答えながら馳けだした。\\

\\

\subsection{\ruby{乱兆}{らんちょう}}

\\

\subsubsection{一}

\\


　時は、\ruby{中平}{ちゅうへい}六年の夏だった。\\


　\ruby{洛陽宮}{らくようきゅう}のうちに、\ruby{霊帝}{れいてい}は重い\ruby{病}{やまい}にかかられた。\\


　帝は病の\ruby{篤}{あつ}きを知られたか、\\


「\ruby{何進}{かしん}をよべ」\\


　と、\ruby{病褥}{びょうじょく}から仰せ出された。\\


　大将軍何進は、すぐ参内した。何進はもと牛や豚を\ruby{屠殺}{とさつ}して業としている者であったが、彼の妹が、洛陽にも稀な美人であったので、貴人の娘となって宮廷に入り、帝の\ruby{胤}{たね}をやどして\ruby{弁皇子}{べんおうじ}を生んだ。そして皇后となってからは\ruby{何后}{かこう}といわれている。\\


　そのため兄の何進も、一躍要職につき、権を握る身となったのである。\\


　何進は、病帝をなぐさめて、\\


「ご安心なさいまし。たとえ如何なることがあっても、何進がおります。また、皇子がいらっしゃいます」といって退がった。\\


　しかし、帝の気色は、\ruby{慰}{なぐさ}まないようであった。\\


　帝には、なお、複雑な\ruby{憂悶}{ゆうもん}があったのである。何后のほかに、王美人という寵姫があって、その腹にも皇子の\ruby{協}{きょう}が生れた。\\


　何后は、それを知って、大いに嫉妬し、ひそかに\ruby{鴆毒}{ちんどく}を盛って、王美人を殺してしまった。そして、\ruby{生}{な}さぬ仲の皇子協を、霊帝のおっ母さんにあたる\ruby{董}{とう}太后の手へあずけてしまったのである。\\


　ところが、董太后は、預けられた協皇子が可愛くてたまらなかった。帝もまた、何后の生んだ\ruby{弁}{べん}よりも、協に\ruby{不愍}{ふびん}を感じて偏愛されていた。\\


　で、十\ruby{常侍}{じょうじ}の\ruby{蹇碩}{けんせき}などが、時々そっと帝の病褥へ来てささやいた。\\


「もし、協皇子を、皇太子に立てたいという思し召ならば、まず何后の兄何進から先に\ruby{誅罰}{ちゅうばつ}なさらなければなりません。何進を殺すことが、\ruby{後患}{こうかん}をたつ\ruby{所以}{ゆえん}です」\\


「……ウム」\\


　帝は蒼白い顔でうなずかれた。\\


　自己の病は篤い。いつとも知れない命数。\\


　帝は決意すると急がれた。\\


　にわかに、何進の邸へ向って、\\


「急ぎ、参内せよ」と、勅令があった。\\


　何進は、変に思った。\\


「はてな。きのう参内したばかりなのに？」\\


　急に帝の病状でも変ったのかと考えて、家臣に探らせてみるとそうでもない。のみならず、十常侍の\ruby{蹇碩}{けんせき}らが、なにか\ruby{謀}{はか}っている\ruby{経緯}{いきさつ}がうすうす分ったので、\\


「\ruby{小癪}{こしゃく}な輩。そんな\ruby{策}{て}に乗る何進ではない」\\


　と、参内しないかわりに、廟堂の諸大臣を私館へ招いて、\\


「こういう事実がある。実に怪しからぬ陰謀だ。さなきだに天下みな、十常侍の輩を恨んで、機あらば、彼らの肉を\ruby{啖}{くら}わんとまで\ruby{怨嗟}{えんさ}している。おれもこの機会に、\ruby{宦官}{かんがん}どもをみな殺しにしようと思うが、諸公のご意見はどうだ」と、会議の席にはかった。\\


「…………」\\


　誰も皆、黙ってしまった。ただびっくりした眼ばかりであった。すると、座隅の一席からひとりの\ruby{白皙}{はくせき}の美丈夫が起立して、\\


「至極けっこうでしょう。しかし十常侍とその与党の勢力というものは、宮中においては、想像のほかと承ります。将軍、威あり実力ありといえども、うっかり手を焼くと、ご自身、\ruby{滅族}{めつぞく}の\ruby{禍}{わざわ}いを求めることになりはしませんか」と忠言を吐いた。\\


　見るとそれは、\ruby{典軍}{てんぐん}の校尉\ruby{曹操}{そうそう}であった。何進の眼から見ればまことに微々たる一将校でしかない。何進は苦い顔して、\\


「だまれっ。貴様のような若輩の一武人に、朝廷の内事が分ってたまるものか、ひかえろ」\\


　と、一言に叱りつけた。\\


　ために、座中白け渡って見えた時、折も折、霊帝がたった今\ruby{崩御}{ほうぎょ}されたという報らせが入った。\\

\\

\subsubsection{二}

\\


　何進は、その報らせを手にすると、会議の席へ戻ってきて、諸大臣以下一同に向い、\\


「ただ今、重大なる報らせがあったが、まだ公の発表ではないから、そのつもりで聞いて欲しい」と、前提し、厳粛なる口調で、次のように述べた。\\


「天子、ご\ruby{不例}{ふれい}久しきにわたっておったが、今日ついに、\ruby{嘉徳殿}{かとくでん}において、崩御あそばされた」\\


「…………」\\


　何進がそういい終っても、ややしばらくの間、会議の席は\ruby{寂}{せき}として、声を発する者もなかった。\\


　諸大臣の面上には、はっとしたような色が流れた。予期していたことながら、\\


　――どうなることか？\\


　と、この先の政治的な変動やら一身の\ruby{去就}{きょしゅう}に、\ruby{暗澹}{あんたん}たる動揺がかくしきれなかった。\\


　しかも場合が場合である。\\


　何進が、十常侍をみな殺しにせんと息まいてこの席に計り、十常侍らは、何進を\ruby{謀}{はか}って、亡き者にしようと、暗躍しているという折も折であった。\\


　そも、何の\ruby{兆}{きざ}しか。\\


　人々が一瞬自失したかのように、暗澹たる\ruby{危惧}{きぐ}の底に沈んで、\\


　――ああ、漢朝四百年の天下も今日から崩れ始める\ruby{兆}{きざ}しか。\\


　と、いうような予感に襲われたのも、決してむりではない。\\


　しばし、黙祷のうちに、人々は亡き霊帝をめぐる近年の宮廷の浅ましい限りの女人と権謀の争いやら、数々の悪政の頽廃を胸によびかえして、今さらのように、深い嘆息をもらし合った。\\


　　　　　　×　　　　　×　　　　　×\\


　霊帝は不幸なお方だった。\\


　何も知らなかった。十常侍たちの見せる「\ruby{偽飾}{ぎしょく}」ばかりを信じられて、世の中の「真実」というものは、何ひとつご存じなく死んでしまわれた。\\


　十常侍の一派にとっては、霊帝は即ち「\ruby{盲帝}{もうてい}」であった。\ruby{傀儡}{かいらい}にすぎなかった。玉座は彼らが暴政をふるい魔術をつかう恰好な壇上であり\ruby{帳}{とばり}であった。\\


　その悪政を数えたてればきりもないが、まず近年のことでは、黄巾の乱後、恩賞を与えた将軍や勲功者へ、裏からひそかに人をやって、\\


「公らの軍功を奏上して、公らはそれぞれ莫大な封禄の恩典にあずかりたるに、それを奏した十常侍に、なんの沙汰もせぬのは、非礼ではないか」\\


　などと\ruby{賄賂}{わいろ}の\jruby{なぞ}{・・・・・・}をかけたりした。\\


　恐れて、すぐ\ruby{賂}{まいない}を送った者もあるが、\ruby{皇甫嵩}{こうほすう}と、\ruby{朱雋}{しゅしゆん}の二将軍などは、\\


「何をばかな」\\


　と、一蹴したので、十常侍たちはこもごもに、天子に\ruby{讒}{ざん}したので、帝はたちまち、朱雋、皇甫嵩のふたりの官職を剥いで、それに代るに、\ruby{趙忠}{ちょうちゅう}を車騎将軍に任命した。\\


　また、\ruby{張譲}{ちょうじょう}その他の内官十三人を列侯に封じ、\ruby{司空張温}{しくうちょううん}を大尉に昇せたりしたので、そういう機運に乗った者は、十常侍に媚びおもねって、さらに彼らの勢力を増長させた。\\


　たまたま、\ruby{忠諫}{ちゅうかん}をすすめ、真実をいう良臣は、みな獄に下されて、斬られたり毒殺されたりした。\\


　従って宮廷の\ruby{紊}{みだ}れは、あざむかず、民間に反映して、地方にふたたび黄巾賊の残党やら、新しい\ruby{謀叛人}{むほんにん}が\ruby{蜂起}{ほうき}して、洛陽城下に天下の危機が聞えてきた。\\


　この動乱と風雲の再発に、人の運命も波浪にもてあそばれる如く転変をきわめたが、たまたま、幸いしたのは、前年来、不遇の地におわれて、代州の\ruby{劉恢}{りゅうかい}の情けにようやく身をかくしていた\ruby{劉備玄徳}{りゅうびげんとく}であった。\\

\\

\subsubsection{三}

\\


　\ruby{黄匪}{こうひ}の乱がやんでからまた間もなく、近年各地に蜂起した賊では、\ruby{漁陽}{ぎょよう}（河北省）を騒がした\ruby{張挙}{ちょうきょ}、\ruby{張純}{ちょうじゅん}の謀叛。\ruby{長沙}{ちょうさ}、\ruby{江夏}{こうか}（湖北省・麻城県附近）あたりの兵匪の乱などが最も大きなものだった。\\


「天下は泰平です。みな帝威に伏して、何事もありません」\\


　十常侍の輩は、口をあわせて、いつもそんなふうにしか、奏上していなかった。\\


　だが。\\


　長沙の乱へは、孫堅を向わせて、平定に努めていた。\\


　また\ruby{劉焉}{りゅうえん}を益州の\ruby{牧}{ぼく}に封じ、\ruby{劉虞}{りゅうぐ}を幽州に封じて、\ruby{四川}{しせん}や漁陽方面の賊を討伐させていた。\\


　その頃。\\


　故郷の\ruby{{\fontspec{jigmo}\char"6DBF}県}{たくけん}から再び戻って、\ruby{代州}{だいしゅう}の\ruby{劉恢}{りゅうかい}の邸に身を寄せていた玄徳は、\ruby{主}{あるじ}劉恢から（時節は来た。これをたずさえて、幽州の\ruby{劉虞}{りゅうぐ}を訪ねてゆき給え。\ruby{虞}{ぐ}は自分の親友だから、君の人物を見ればきっと重用するだろう）\\


　といわれて、一通の紹介状をもらった。\\


　玄徳は恩を謝して、直ちに、関羽張飛などの一族をつれ、劉虞の所へ行った。劉虞はちょうど、中央の命令で、漁陽に起った乱賊を\ruby{誅伐}{ちゅうばつ}にゆく出陣の折であったから、大いによろこんで、\\


（よし。君らの一身はひきうけた）と、自分の軍隊に編入して、戦場へつれて行った。\\


　四川、漁陽の乱も、ようやく一時の平定を見たので、その後、劉虞は朝廷へ表をたてまつって、玄徳の勲功あることを大いにたたえた。\\


　同時に、廟堂の\ruby{公孫{\fontspec{jigmo}\char"74DA}}{こうそんさん}も、\\


（玄徳なる者は、前々黄賊の大乱の折にも抜群の功労があったものです）と、\ruby{上聞}{じょうぶん}に達したので、朝廷でも捨ておかれず、\ruby{詔}{みことのり}を下して、彼を\ruby{平原県}{へいげんけん}（山東省・平原）の令に封じた。\\


　で、玄徳は、即時、一族を率いて任地の平原へさし下った。行ってみると、ここは地味\ruby{豊饒}{ほうじょう}で\ruby{銭粮}{せんろう}の蓄えも官倉に満ちているので、\\


（天、我に兵馬を養わしむ）と、みな非常に元気づいた。そこで玄徳以下、張飛や関羽たちも、ようやくここに\ruby{酬}{むく}いられて、前進一歩の地をしめ、大いに武を練り兵を講じ、駿馬に\ruby{燕麦}{えんばく}を飼って、平原の一角から時雲の去来をにらんでいた。\\


　――果たせるかな。\\


　一雲去れば一風生じ、征野に賊を\ruby{掃}{はら}い去れば、宮中の\ruby{瑠璃殿裡}{るりでんり}に\ruby{冠帯}{かんたい}の\ruby{魔魅}{まみ}や\ruby{金釵}{きんさい}の百鬼は跳梁して、内外いよいよ多事の折から、一夜の黒風に霊帝は崩ぜられてしまった。\\


　\ruby{紛乱}{ふんらん}はいよいよ紛乱を見るであろう。漢室四百年の\ruby{末期相}{まっきそう}はようやくここに\ruby{瓦崩}{がほう}のひびきをたてたのである。――いかになりゆく世の末やらん、と霊帝崩御の由を知るとともに、人々みな色を失って、呆然、足もとの大地が\ruby{九仞}{きゅうじん}の底へめりこむような顔をしたのも、あながち、平常の心がけなき者とばかり\ruby{嗤}{わら}えもしないことであった。\\


　　　　　　×　　　　　×　　　　　×\\


　会議の席も、\ruby{寂}{せき}としてしまい、\ruby{咳声}{しわぶき}をする者すらなかったが、そこへまた、あわただしく、\\


「将軍。お耳を」と、室外にちらと影を見せた者があった。\\


　\ruby{何進}{かしん}に通じている禁門の武官\ruby{潘隠}{はんいん}であった。\\


「オ、潘隠か。なんだ」\\


　何進はすぐ会議の席をはずし、外廊で何かひそひそ潘隠のささやきを聞いていた。\\

\\

\subsubsection{四}

\\


　潘隠が告げていうには、\\


「十常侍の輩は例によって、帝の崩御と同時に、謀議をこらし、帝の死を隠しておいて、まずあなたを宮中に召し、後の禍いを除いてから\ruby{喪}{も}を発し、協皇子を立てて御位を継がしめようという\ruby{魂胆}{こんたん}に密議は一決を見たようであります。――きっと今に、宮中から帝の名をもって、将軍に参内せよと、使いがやってくるにちがいありません」\\


　何進は聞いて、\\


「\ruby{獣}{けだもの}めら、よしっ、それならそれで俺にも考えがある」\\


　\ruby{憤怒}{ふんど}して、会議の壇に戻り、潘隠の密報を諸大臣や、並いる文武官に公然とぶちまけて発表した。\\


　ところへ案の定、宮中からお召しという使者が来邸して、\\


「天子、今ご気息も危うし。\ruby{枕頭}{ちんとう}に公を召して、漢室の後事を託せんと\ruby{宣}{のたま}わる。いそぎ参内あるべし」と、うやうやしくいった。\\


「狸め」\\


　何進は、潘隠へ向って、\\


「こいつを血祭にしろ」と命じるや否や、再び、会衆の前に立って、\\


「もう俺の堪忍はやぶれた。\ruby{断乎}{だんこ}として俺は欲することをやるぞ！」と呶鳴った。\\


　すると、先に忠言して何進に一喝された典軍の校尉\ruby{曹操}{そうそう}が、ふたたび沈黙を破って、\\


「将軍将軍。今日ついに断を下して計をなさんとするならば、まず、天子の位を正してしかる後に賊を討つことをなし給え」と叫んだ。\\


　何進も、今度は前のように、だまれとはいわなかった。大きくうなずいて、\\


「誰か我がために、新帝を正して、\ruby{宮闕}{きゅうけつ}の謀賊どもを討ち尽さん者やある」\\


　\ruby{爛}{らん}たる眼をして、衆席を見まわすと、時に、彼の声に応じて、\\


「\ruby{司隷校尉}{しれいこうい}\ruby{袁紹}{えんしょう}ありっ！」と名乗って起った者がある。\\


　人々の\ruby{首}{こうべ}は、一斉にそのほうへ振向いた。見ればその人は、\ruby{貌相}{ぼうそう}\ruby{魁偉}{かいい}胸ひろく\ruby{双肩}{そうけん}威風をたたえ、武芸抜群の勇将とは見られた。\\


　これなん、漢の司徒\ruby{袁安}{えんあん}が孫、\ruby{袁逢}{えんほう}が子、\ruby{袁紹}{えんしょう}であった。袁紹\ruby{字}{あざな}は\ruby{本初}{ほんしょ}といい、\ruby{汝南}{じょなん}\ruby{汝陽}{じょよう}（河南省・\ruby{淮河}{わいが}上流の北岸）の名門で門下に多数の吏事武将を輩出し、彼も現在は漢室の司隷校尉の職にあった。\\


　袁紹は、昂然とのべた。\\


「願わくば自分に精兵五千を授け給え。直ちに禁門に入って、新帝を\ruby{擁立}{ようりつ}し奉り、多年禁廷に巣くう内官どもをことごとく\ruby{誅滅}{ちゅうめつ}して見せましょう」\\


　何進はよろこんで、\\


「行けっ」と、号令した。\\


　この一声に洛陽の王府は一転戦雲の天と修羅の地になったのである。\\


　袁紹は、たちまち鉄甲に身を\ruby{鎧}{よろ}い、\ruby{御林}{ぎょりん}の近衛兵五千をひっさげて、\ruby{内裏}{だいり}まで押通った。王城の八門、市中の衛門のこらず閉じて戒厳令を布き、入るも出ずるも味方以外は断乎として一人も通すなと命じた。\\


　その間に。\\


　何進もまた、車騎将軍たる武装をして\ruby{何{\fontspec{jigmo}\char"9852}}{かぎょう}、\ruby{荀攸}{じゅんゆう}、\ruby{鄭泰}{ていたい}などの一族や大臣三十余名を\ruby{伴}{ともな}い、陸続と宮門に入り、霊帝の\ruby{柩}{ひつぎ}のまえに、彼が支持する\ruby{弁太子}{べんたいし}を立たせて、即座に、新帝ご即位を宣言し、自分の発声で、百官に万歳を唱えさせた。\\

\\

\subsection{\ruby{舞刀飛首}{ぶとうひしゅ}}

\\

\subsubsection{一}

\\


　百官の拝礼が終って、\\


「新帝万歳」の声が、喪の\ruby{禁苑}{きんえん}をゆるがすと共に、\ruby{御林軍}{ぎょりんぐん}（近衛兵）を指揮する\ruby{袁紹}{えんしょう}は、\\


「次には、陰謀の首魁\ruby{蹇碩}{けんせき}を血まつりにあげん」\\


　と、剣を抜いて宣言した。\\


　そしてみずから宮中を捜しまわって、蹇碩のすがたを見つけ、\\


「おのれっ」と、何処までもと追いかけた。\\


　蹇碩はふるえ上がって、懸命に逃げまわったが、度を失って御苑の花壇の陰へ這いこんでいたところを、何者かに尻から槍で突き殺されてしまった。\\


　彼を突き殺したのは、同じ仲間の十常侍\ruby{郭勝}{かくしょう}だともいわれているし、そこらにまで、乱入していた一兵士だともいわれているが、いずれにせよ、それすら分らない程、もう\ruby{宮闕}{きゅうけつ}の内外は大混乱を呈して、人々の眼も血ばしり、気も\ruby{逆上}{あが}っていたにちがいなかった。\\


　\ruby{袁紹}{えんしょう}は、さらに気負って、何進の前へ行き、\\


「将軍、なんで無言のままこの混乱を見ているんですか。時は今ですぞ、宮廷の\ruby{癌}{がん}、\ruby{社稷}{しゃしょく}の\ruby{鼠賊}{そぞく}ども、十常侍の輩を一匹残らず殺してしまわなければいけません。この機を逸したら、再び\ruby{臍}{ほぞ}を噛むような日がやってきますぞ」と、進言した。\\


「ウむ。……むむ」\\


　何進はうなずいていた。\\


　けれど顔色は蒼白で、日頃の元気も見えない。元来、小心な何進、一時は憤怒にかられて、この大事をあえて求めたが、一瞬のまに禁門の内外はこの世ながらの修羅地獄と化し、自分を殺そうと謀った\ruby{蹇碩}{けんせき}も殺されたと聞いたので、一時の怒りもさめて、むしろ自分のつけた火の果てなくひろがりそうな光景に、呆然と戦慄をおぼえているらしい容子であった。\\


　その間に。\\


　一方十常侍の面々は、\\


「すわ、大変」と、狼狽して、\ruby{張譲}{ちょうじょう}を始め、おのおの生きた心地もなく、内宮へ逃げこんで、窮余の一策とばかり、何進の妹にして皇后の位置にある\ruby{何后}{かこう}の\ruby{裙下}{くんか}にひざまずいて、百拝、\ruby{憐愍}{れんびん}を乞うた。\\


「よい、よい。安心せい」\\


　何后はすぐ、兄の何進を呼びにやった。\\


　そして何進をなだめた。\\


「私たち兄妹が、\ruby{微賤}{びせん}の身から今日の\ruby{富貴}{ふうき}となったのも、そのはじめは十常侍たちの内官の推薦があったからではありませんか」\\


　何進は、妹にそういわれると、むかし牛の屠殺をしていた頃の貧しい自分の姿が思い出された。\\


「なに、俺は、俺を殺そうと謀った蹇碩の奴さえ\ruby{誅戮}{ちゅうりく}すればいいのだ」\\


　内宮を出ると、何進は、右往左往する味方や宮内官たちを、鎮撫する気でいった。\\


「蹇碩は、すでに誅罰した。彼は我を害さんとしたから斬ったのである。我に害意なき者には、我また害意なし。安心して鎮まれ！」\\


　すると、それを聞いて、\\


「将軍、何をばかなことをいうんですか」\\


　と、\ruby{袁紹}{えんしょう}は血刀を持ったまま彼の前へきて、その\ruby{軽忽}{けいこつ}を責めた。\\


「この大事を挙げながら、そんな手ぬるい宣言を将軍の口から発しては困ります。今にして、\ruby{宮闕}{きゅうけつ}の\ruby{癌}{がん}を除き、根を刈り尽しておかなければ、後日かならず後悔なさいますぞ」\\


「いや、そういうな。宮門の火の手が、洛陽一面の火の手になり、洛陽の火の手が、天下を\ruby{燎原}{りょうげん}の火としてしまったら取返しがつかんじゃないか」\\


　何進の優柔不断は、とうとう袁紹の言を容れなかった。\\

\\

\subsubsection{二}

\\


　一時、禁門の兵乱は、治まったかに見えた。\\


　その後。\\


　\ruby{何后}{かこう}、何進の一族は、\\


「邪魔ものは\ruby{董太后}{とうたいこう}である」\\


　と、悪策をめぐらして、太后を\ruby{河間}{かかん}（河北省・河間県）という片田舎へ\ruby{遷}{うつ}してしまった。\\


　故霊帝の母公たる董太后も、今は彼らの勢力に拒む力もなかった。これというのも、前帝の\ruby{寵妃}{ちょうひ}だった王美人の生んだ協皇子を愛するのあまり、何后、何進らの一族から睨まれた結果と――ぜひなき運命の\ruby{輦}{くるま}のうちに涙にくれながら都離れた地方へ送られて行った。\\


　けれど、何后も何進も、それでもまだ不安を覚えて、ひそかに後から刺客をやって、董太后を殺してしまった。\\


　わずかの間に董太后はふたたび洛陽の帝城に還ってきたが、それは\ruby{柩}{ひつぎ}の中に冷たい\ruby{空骸}{むくろ}となって戻られたのであった。\\


　京師では大葬が\ruby{執行}{とりおこな}われた。\\


　けれど、何進は、\\


「病中――」と称して、宮中へも世間へも顔を出さなかった。\\


　彼は怒りっぽい。\\


　しかも、小心であった。\\


　彼は自己や一門の栄華のために大悪もあえてする。けれど小心な彼は半面でまた、ひどく世間に気がねし、自らも責めている。\\


　要するに何進は、下賤から人臣の上に立ったが、大なる野望家にもなりきれず、ほんとの悪人にもなりきれず、位階冠帯は重きに過ぎて、\ruby{右顧左眄}{うこさべん}、気ばかり病んでいるつまらない人物だった。\\


　貝殻が人の跫音に貝のフタをしているように、門から出ないので、或る日、袁紹は何進の邸を訪ねて、\\


「どうしました将軍」と、見舞った。\\


「どうもせんよ」\\


「お元気がないじゃないですか」\\


「そんなことはない」\\


「――ところで、聞きましたか」\\


「何を？　……じゃね」\\


「董太后のお生命をちぢめた者は何進なりと、また、例の\ruby{宦官}{かんがん}どもが、しきりと流言を放っているのを」\\


「……ふウむ」\\


「だから私がいわない事ではありません。今からでも遅くないでしょう。あくまでも、\ruby{彼奴}{きゃつ}らは\ruby{癌}{がん}ですよ。根こそぎ切ってしまわなければ、どう懲らしても、日が経てばすぐ芽を生やし、根を張って、増長わがまま、陰謀暗躍、手がつけられない物になるんです」\\


「……む、む」\\


「ご決断なさい」\\


「考えておこう」\\


　煮え切らない顔つきである。\\


　袁紹は舌打ちして帰った。\\


　\ruby{奴僕}{ぬぼく}の中に、\ruby{宦官}{かんがん}たちのまわし者が住みこんでいる。\\


「袁紹が来てこうこうだ」とすぐ密報する。\\


　諜報をうけて、\\


「また、大変だ」と、宦官らはあわてた。――だが、危険になると、消火栓のような便利な手がある。何進の妹の何后へすがって泣訴することであった。\\


「いいよ」\\


　何后は、彼らからあやされている\ruby{簾中}{れんちゅう}の人形だったが、兄へは権威を持っていた。\\


「何進をおよび」\\


　また、始まった。\\


「兄さん、あなたは、悪い部下にそそのかされて、またこの平和な宮中を乱脈に騒がすようなことを考えなどなさりはしないでしょうね。禁裡の内務を宦官がつかさどるのは、漢の宮中の伝統で、それを憎んだり殺したりするのは、宗廟に対して非礼ではありませんか」\\


　釘を刺すと、何進は、\\


「おれはなにもそんなことを考えておりはせぬが……」\\


　と、あいまいに答えたのみで退出してしまった。\\

\\

\subsubsection{三}

\\


　宮門から退出してくると、\\


「将軍。どうでした」\\


　と、彼の乗物の蔭に待っていた武将が、参内の\ruby{吉左右}{きっそう}を小声でたずねた。\\


「ア。……\ruby{袁紹}{えんしょう}か」\\


「何太后に召されたと聞いたので、案じていたところです。何か、\ruby{宦官}{かんがん}の問題で、ご内談があったのでしょう」\\


「……ム。あったにはあったが」\\


「ご決意を告げましたか」\\


「いや、こちらから云いださないうちに、太后から、\ruby{憐愍}{れんびん}の取りなしがあったので」\\


「いけません」\\


　袁紹は、断乎としていった。\\


「そこが、将軍の弱点です。宦官どもは、一面にあなたを陥し入れるように、陰謀や悪宣伝を放って、露顕しかかると、太后の\ruby{裳}{も}やお袖にすがって、泣き声で訴えます。――お気の弱い太后と、太后のいうことには\ruby{反}{そむ}かないあなたの急所を、彼らはみこんでやっている仕事ですからな」\\


「なるほど……」\\


　そういわれると、何進も、気づくところがあった。\\


「今です。今のうちです。今日をおいて、いつの日かありましょう。よろしく、四方の英雄に\ruby{檄}{げき}を飛ばし、もって\ruby{万代}{ばんだい}の計を、一挙に定められるべきです」\\


　彼の熱弁には、何進もうごかされるのである。なるほどと思い――それもそうだと思い、いつのまにか、\\


「よしっ、やろう。実はおれもそれくらいのことは考えていたのだ」と、いってしまった。\\


　二人の密談を、乗物のおいてある樹蔭の近くで聞いていた者がある。典軍の校尉曹操であった。\\


　曹操は、独りせせら笑って、\\


「ばかな煽動をする奴もあればあるものだ。\ruby{癌}{がん}は体じゅうにできている物じゃない。一個の元兇を抜けばいいのだ。宦官のうちの首謀者をつまんで牢へぶちこめば、刑吏の手でも事は片づくのに、諸方の英雄へ檄を飛ばしたりなどしたら、漢室の\ruby{紊乱}{びんらん}はたちまち諸州の野望家のうかがい知るところとなり、争覇の分脈は、諸国の群雄と、複雑な糸をひいて、天下はたちまち大乱になろう」\\


　それから、彼はまた、何進の\ruby{輦}{くるま}について歩きながら、\\


「……失敗するにきまっている。さあ、その先は、どんなふうに風雲が\ruby{旋}{めぐ}るか」\\


　と、独りごとにいっていた。\\


　けれど、曹操は、もう自分の考えを、何進に直言はしなかった。その点、袁紹の如く真っ正直な熱弁家でもないし、何進のような小胆者とも違う彼であった。\\


　彼は今、天下に多い野望家とつぶやいたが、彼自身もその一人ではなかろうか。\ruby{白皙秀眉}{はくせきしゅうび}、\ruby{丹唇}{たんしん}をむすんで、\ruby{唯々}{いい}として何進の警固についてはいるが、どうもその輦の中にある上官よりも典軍の一将校たる彼のほうが、もっと底の深い、もっと肚も黒い、そしてもっと\ruby{器}{うつわ}も大きな曲者ではなかろうかと見られた。\\


　　　　　　×　　　　　×　　　　　×\\


　ここに、\ruby{西涼}{せいりょう}（\ruby{甘粛省}{かんしゅくしょう}・蘭州）の地にある\ruby{董卓}{とうたく}は、前に黄巾賊の討伐の際、その司令官ぶりは至って\ruby{香}{かんば}しくなく、乱後、朝廷から罪を問われるところだったが、内官の十常侍一派をたくみに買収したので、不問に終ったのみか、かえって顕官の地位を占めて、今では西涼の\ruby{刺史}{しし}、兵二十万の軍力をさえ擁していた。\\


　その董卓の手へ、\\


「洛陽からです」\\


　と或る日、一片の\ruby{檄}{げき}が、密使の手から届けられた。\\

\\

\subsubsection{四}

\\


　洛陽にある\ruby{何進}{かしん}は、先ごろ来、檄を諸州の英雄に飛ばして、\\



天下の府、\ruby{枢廟}{すうびょう}の\ruby{弊}{へい}や今きわまる。よろしく公明の\ruby{旌旗}{せいき}を林集し、正大の雲会を遂げ、もって、\ruby{昭々}{しょうしょう}日月の下に万代の革政を諸公と共に正さん。\\



　といったような意味を伝え、その反響いかにと待っていたところ、やがて諸国から続々と、\\


「\ruby{上洛参会}{じょうらくさんかい}」\\


　とか、或いは、\\


「提兵援助」\\


　などという答文をたずさえた使者が日夜早馬で先触れして来て、彼の館門を叩いた。\\


「西涼の\ruby{董卓}{とうたく}も、兵をさげてやって来るようですが」\\


　――\ruby{御史}{ぎょし}の\ruby{鄭泰}{ていたい}なる者が、何進の前に来て云った。\\


「檄文は、董卓へもお出しになったんですか？」\\


「む。……出した」\\


「彼は、\ruby{豺狼}{さいろう}のような男だとよく人はいいます。京師へ豺狼を引入れたら人を喰いちらしはしませんかな」\\


　\ruby{鄭泰}{ていたい}が憂えると、\\


「わしも同感だ」\\


　と、室の一隅で、参謀の幕将たちと、一面の地形図をひらいていた一老将が、歩を何進のほうへ移してきながら云った。\\


　見ると、中郎将\ruby{盧植}{ろしょく}である。\\


　彼は黄匪討伐の征野から\ruby{讒}{ざん}せられて、\ruby{檻車}{かんしゃ}で都へ送られ、一度は軍の裁廷で罪を宣せられたが、後、彼を陥し入れた\ruby{左豊}{さほう}の失脚とともに、\ruby{免}{ゆる}されて再び中郎将の原職に復していたのである。\\


「おそらく董卓は、檄文を見て時こそ来れりとよろこんだに違いない。政廟の革正をよろこぶのでなく、乱をよろこび、自己の野望を乗ずべき時としてです。――わしも董卓の人物はよく知っておるが、あんな\ruby{漢}{おとこ}をもし禁廷に入れたら、どんな禍患を生じるやも計り知れん」\\


　盧植は、わざと、鄭泰のほうへ向って話しかけた。暗に何進を\ruby{諫}{いさ}めたのである。だが何進は、用いなかった。\\


「そう諸君のように、疑心をもっては、天下の英雄を操縦はできんよ」\\


「――ですが」\\


　\ruby{鄭泰}{ていたい}がなお、苦言を呈しかけると何進はすこし不機嫌に、\\


「まだまだ、君たちは、大事を共に謀るに足りんなあ」と、いった。\\


　鄭泰も、盧植も、\\


「……そうですか」\\


　と、後のことばを胸にのんで退がってしまった。そしてこの両者をはじめ、心ある朝臣たちも、こんなことを伝え聞いて、そろそろ何進の人間に\ruby{見限}{みき}りをつけだして離れてしまった。\\


「董卓どのの兵馬は、もう\ruby{{\fontspec{jigmo}\char"6FA0}池}{べんち}（河南省・洛陽西方）まで来ているそうです」\\


　何進は、部下から聞いて、\\


「なぜすぐにやって来んのか。迎えをやれ」と、しばしば使いを出した。\\


　けれど、董卓は、\\


「長途を来たので、兵馬にも少し休養させてから」\\


　とか、軍備を整えてとか、何度催促されても、それ以上動いて来なかった。何進の催促を\ruby{馬耳東風}{ばじとうふう}に、\ruby{豺狼}{さいろう}の眼をかがやかしつつ、ひそかに、\ruby{眈々}{たんたん}と洛内の気配をうかがっているのであった。\\

\\

\subsubsection{五}

\\


　一方。宮城内の十常侍らも、何進が諸国へ\ruby{檄}{げき}をとばしたり、檄に応じて董卓などが、{\fontspec{jigmo}\char"6FA0}池附近にまできて駐軍しているなどの大事を、知らないでいる筈はない。\\


「さてこそ」と、彼らはあわてながらも対策を講ずるに急だった。そこで張譲らはひそかに手配にかかり、\ruby{刀斧}{とうふ}鉄弓をたずさえた禁中の兵を、嘉徳門や長楽宮の内門にまでみっしり伏せておいて、何太后をだまし何進を召すの親書を書かせた。\\


　宮門を出た使者は平和時のように、わざと\ruby{美車金鞍}{びしゃきんあん}をかがやかせ、なにも知らぬ顔して、書を何進の館門へとどけた。\\


「いけません」\\


　何進の側臣たちは、即座に十常侍らの\ruby{陥穽}{かんせい}を\ruby{看破}{みやぶ}って諫めた。\\


「太后の\ruby{御詔}{ごしょう}とて、この際、信用はできません。危ない限りです。一歩もご門外に出ることはなさらぬほうが賢明です」\\


　こういわれると、それに対して自分にない器量をも見せたいのが何進の病であった。\\


「なにをいう。宮中の病廃を正し、政権の正大を期し、やがては天下に臨まんとするこの何進である。十常侍らの\ruby{輩}{ともがら}が我に何かせん。彼らごとき\ruby{廟鼠輩}{びょうそはい}を怖れて、何進門を閉ざせりと聞えたら天下の英雄どもも、かえって予を見くびるであろう」\\


　変にその日は強がった。\\


　すぐ車騎の用意を命じ、その代り鉄甲の精兵五百に、物々しく護衛させて、参内に出向いた。果たせるかな、\ruby{青鎖門}{せいさもん}まで来ると、\\


「兵馬は禁門に入ることならん。門外にて待ちませい」\\


　と隔てられ、何進は、数名の従者だけつれて入った。それでも彼は\ruby{傲然}{ごうぜん}、胸をそらし、威風を示して歩いて行ったが、嘉徳門のあたりまでかかると、\\


「豚殺し待てっ」\\


　と、物陰から呶鳴られて、あっとたじろぐ間に、前後左右、十常侍一味の軍士たちに取巻かれていた。\\


　躍りでた\ruby{張譲}{ちょうじょう}は、\\


「何進っ、汝は元来、洛陽の裏町に、豚を屠殺して、からくも生きていた貧賤ではなかったか。それを、今日の栄位まで昇ったのは、そもそも誰のおかげと思うか。われわれが陰に陽に、汝の妹を天子に\ruby{薦}{すす}め奉り、汝をも推挙したおかげであるぞ。――この恩知らずめ！」と、面罵した。\\


　何進は、真ッ蒼になって、\\


「しまった！」\\


　と口走ったが、時すでに遅しである。諸所の宮門はみな閉ざされ、逃げまわるにも\ruby{刀斧}{とうふ}鉄槍、身を囲んで、一尺の隙もなかった。\\


「――わッっ。だっ！」\\


　何進はなにか絶叫した。空へでも飛び上がってしまう気であったか、躍り上がって、体を三度ほどぐるぐるまわした。張譲は、跳びかかって、\\


「下郎っ。思い知ったか」\\


　と、真二つに斬りさげた。\\


　青鎮門外ではわいわいと騒がしい声が起っていた。なにかしら宮門の中におかしな空気を感じだしたものとみえ、\\


「何将軍はまだ退出になりませんか」\\


「将軍に急用ができましたから、早くお車に召されたいと告げて下さい」\\


　などと喚いて動揺しているのであった。\\


　すると、城門の\ruby{墻壁}{しょうへき}の上から、武装の宮兵が一名首を出して、\\


「やかましいッ。鎮まれ。汝らの主人何進は、\ruby{謀叛}{むほん}のかどによって査問に付せられ、ただ今、かくの如く罪に服して処置は終った。これを車にのせて立帰れっ」\\


　なにか\ruby{蹴鞠}{けまり}ほどな黒い物がそこからほうられてきたので、外にいた面々は、急いで拾い上げてみると、唇を噛んだ蒼い何進の生首であった。\\

\\

\subsection{\ruby{蛍}{ほたる}の\ruby{彷徨}{さまよ}い}

\\

\subsubsection{一}

\\


　何進の幕将で中軍の校尉\ruby{袁紹}{えんしょう}は、何進の首を抱いて、\\


「おのれ」と、青鎖門を睨んだ。\\


　同じ何進の部下、\ruby{呉匡}{ごきょう}も、\\


「おぼえていろ」と、怒髪を逆だて、宮門に火を放つと五百の精兵を駆って、なだれこんだ。\\


「十常侍をみなごろしにしろ」\\


「\ruby{宦官}{かんがん}どもを焼きつくせ」\\


　華麗な宮殿は、たちまち土足の暴兵に占領された。炎と、黒煙と、悲鳴と矢うなりの\ruby{旋風}{つむじかぜ}であった。\\


「\ruby{汝}{うぬ}もかっ」\\


「おのれもかっ」\\


　宦官と見た者は、見つかり次第に殺された。宮中深く棲んでいた十常侍の輩なので、兵はどれが誰だかよく分らないが、\ruby{髯}{ひげ}のない男だの、俳優のように\jruby{にやけ}{・・・・・・}て美装している内官は、みんなそれと見なして首を刎ねたり突き殺したりした。\\


　十常侍\ruby{趙忠}{ちょうちゅう}や\ruby{郭勝}{かくしょう}などという連中も、\ruby{西宮翠花門}{せいきゅうすいかもん}まで逃げ転んできたが、鉄弓に射止められて、虫の息で這っているところを、ずたずたに斬りきざまれ、手足は翠花楼の大屋根にいる\ruby{鴉}{からす}へ投げられ、首は西苑の湖中へ跳ねとばされた。\\


　天日も\ruby{晦}{くら}く、地は燃ゆる。\\


　女人たちの棲む後宮の悲鳴は、雲にこだまし地底まで届くようだった。\\


　その中を、十常侍一派の張譲、\ruby{段珪}{だんけい}のふたりは、新帝と何太后と、新帝の弟にあたる協皇子――帝が即位してからは、\ruby{陳留王}{ちんりゅうおう}といわれている――の三人を黒煙のうちから救け出して、北宮\ruby{翡翠門}{ひすいもん}からいち早く逃げ出す準備をしていた。\\


　ところへ。\\


　\ruby{戈}{ほこ}を引っさげ、身を鎧い、悍馬に泡を噛ませてきた一老将がある。宮門に変ありと、火の手を見るとともに馳せつけてきた中郎将盧植であった。\\


「待てっ毒賊。帝を擁し、太后をとって、\ruby{何地}{いずち}へゆかんとするかっ」\\


　大喝して、馬上から降りるまに張譲たちは、新帝と陳留王の車馬に鞭打って逃げてしまった。\\


　ただ何太后だけは、盧植の手にひき留められた。\\


　折ふし、宮中各所の火災を、懸命に部下を指揮して消し止めていた校尉曹操に出会ったので、ふたりは、\\


「新帝のご帰還あるまで、しばし、大権をお執りくだされたい」\\


　と請い、一方諸方に兵を派して、新帝と陳留王の後を追わせた。\\


　洛陽の巷にも火が降っていた。兵乱は今にも全市に及ぶであろうと、家財商品を負って避難する民衆で混乱は極まっている。その中を――張譲らの馬と、新帝、皇弟を乗せた\ruby{輦}{くるま}は、逃げまどう老父を\ruby{轢}{ひ}き、幼子を蹴とばして、躍るが如く、城門の郊外遠くまで逃げ落ちてきた。\\


　けれど、輦の車輪はこわれ、張譲らの馬も傷ついたり、ぬかるみへ脚を入れたりして、みな\ruby{徒歩}{かち}にならなければならなかった。\\


「――ああ」\\


　帝は、時々、よろめいた。\\


　そして大きく嘆息された。\\


　かえりみれば、洛陽の空は、夜になってまだ赤かった。\\


「もう少しのご辛抱です」\\


　張譲らは、帝を離すまいとした。帝を擁することが自分らの強味だからである。\\


　草原の果てに、\ruby{北{\fontspec{jigmo}\char"9099}山}{ほくぼうざん}が見えた。夜は暗い。もう三\ruby{更}{こう}に近いであろう。すると一隊の人馬がおって来た。張譲は観念した。追手と直感したからである。\\


「もうだめだっ」\\


　無念を叫びながら、張譲は、自ら河に飛込んで自殺してしまった。帝と、帝の弟の\ruby{陳留王}{ちんりゅうおう}とは、河原の草の\ruby{裡}{うち}へ抱き合って、しばし近づく兵馬に耳をすましておられた。\\

\\

\subsubsection{二}

\\


　やがて河を越えて驟雨のように馳け去って行ったのは、河南の\ruby{中部掾史}{ちゅうぶえんし}、\ruby{閔貢}{びんこう}の兵馬であったが、なにも気づかず、またたくまに闇に消え去ってしまった。\\


「…………」\\


　しゅく、しゅく……と新帝は草むらの中で泣き声をもらした。\\


　皇弟陳留王は、わりあいにしっかりした声で、\\


「ああ\ruby{飢餓}{きが}をお覚えになりましたね。ごもっともです。私も、今朝から水一滴のんでいませんし、馴れない道を、夢中で歩いてきたので、身を起そうとしてもただ身がふるえるばかりです」と、慰めて――「けれど、この河原の草の中で、このまま夜を明かすこともできません。ことに、ひどい夜露、お体にもさわります。――歩けるだけ歩いてみましょう。どこか民家でもあるかもしれません」\\


「…………」\\


　帝は微かにうなずいた。\\


　二人は、衣の\ruby{袂}{たもと}と袂とを結び合わせ、「迷わないように」と、闇を歩いた。\\


　\ruby{茨}{いばら}か、\ruby{野棗}{のなつめ}か、とげばかりが脚を刺した。帝も陳留王も生れて初めて、こうした世のあることを知ったので、生きた気もちもなかった。\\


「ああ、蛍が……」\\


　陳留王はさけんだ。\\


　大きな蛍の群れが、風のまにまに一かたまりになって、眼のまえをふわふわ飛んでゆく、蛍の光でも非常に心づよくなった。\\


　夜が明けかけた――\\


　もう歩けない。\\


　新帝はよろめいたまま起き上がらなかった。陳留王も、\\


「ああ」と、腰をついてしまった。\\


　\ruby{昏々}{こんこん}と、しばらくは何もしらなかった。誰かそのうちに起す者がある。\\


「どこから来た？」と、訊ねるのである。\\


　見まわすと、古びた荘院の土塀が近くにある。そこの\ruby{主}{あるじ}かもしれない。\\


「いったい、そなた達は、\ruby{何人}{なにびと}のお子か」\\


　と、重ねて問う。\\


　陳留王は、まだしっかりした声を持っていた。帝を指さして、\\


「先頃、ご即位されたばかりの新帝陛下です。十常侍の乱で、宮門から遁れてきたが、侍臣たちはみなちりぢりになり、ようやく、私がお供をしてこれまで来たのです」と、いった。\\


　主は、仰天して、\\


「そして、あなたは」と、眼をまろくした。\\


「わしは、帝の弟、陳留王という者である」\\


「げっ、では真の？　……」\\


　主は、驚きあわてた様で、帝を扶けて、荘院のうちへ迎え入れた。古びた田舎\ruby{邸}{やしき}である。\\


「申しおくれました。自分儀は、先朝にお仕え申していた\ruby{司徒}{しと}\ruby{崔烈}{さいれつ}の弟で、\ruby{崔毅}{さいき}という者であります。十常侍の徒輩が、あまりにも賢を追い邪を容れて、目をおおうばかりな暴状に、官吏がいやになって、\ruby{野}{や}に隠れていた者でございます」\\


　主は改めて礼をほどこした。\\


　その夜明け頃――\\


　河へ投身して死んだ張譲を見捨てて、\ruby{段珪}{だんけい}はひとり野道を逃げ惑うてきたが、途中、\ruby{閔貢}{びんこう}の隊に見つかって、天子の行方を訊かれたが、知らないと答えると、\\


「不忠者め」\\


　と、閔貢は、馬上から一\ruby{颯}{さつ}に斬ってしまった。そしてその首を、鞍に結びつけ兵へ向って、\\


「なにせい、この地方に来られたに違いない」と、捜査の手分けを命じ、自身もただ一騎馳け、\ruby{彼方此方}{あなたこなた}と、\ruby{血眼}{ちまなこ}で尋ねあるいていた。\\

\\

\subsubsection{三}

\\


　\ruby{崔毅}{さいき}の家をかこむ木立の空に、炊煙があがっていた。\\


　帝と陳留王のふたりを\ruby{匿}{かく}しておいた\ruby{茅屋}{あばらや}の板戸を開いて、崔毅は、\\


「田舎です、なにもありませんが、飢えをおしのぎ遊ばすだけと\ruby{思}{おぼ}し\ruby{召}{め}して、この\ruby{粥}{かゆ}など一時召上がっていてください」と、食事を捧げた。\\


　帝も、皇弟も、浅ましきばかりがつがつと粥をすすられた。\\


　崔毅は涙を催して、\\


「安心して、お眠りください。外はてまえが見張っておりますから」と、告げて退がった。\\


　荒れた\ruby{傾}{かし}いだ荘院の門に立ったまま、崔毅は半日も立っていた。\\


　すると、\ruby{戛々}{かつかつ}と、馬蹄の音が木立の下を踏んでくる。\\


「誰か？」\\


　どきっとしながらも、何くわぬ顔して、\ruby{箒}{ほうき}の手をうごかしていた。\\


「おいおい、家の主、なにか喰う物はないか。湯なと一杯恵んでくれい」\\


　声に振向くと、それは馬上の\ruby{閔貢}{びんこう}であった。\\


　崔毅は、彼の馬の鞍に結いつけてある生々しい首級を見て、\\


「おやすいことです。――ですが豪傑、その首は一体、誰の首です」\\


　閔貢は問われると、\\


「知らずや、これは十常侍張譲などと共に、久しく廟堂に巣くって、天下の害をなした段珪という男だ」\\


「えっ、ではあなたはどなたですか」\\


「河南の\ruby{掾史閔貢}{えんしびんこう}という者だが、昨夜来、帝のお行方が知れないので、ほうぼうお捜し申しておるのだ」\\


「ああ、では！」\\


　崔毅は、手をあげて、奥のほうへ転んで行った。\\


　閔貢は怪しんで、馬をつなぎ、後から駈けて行った。\\


「お味方の豪傑が、お迎えにやって来ましたよ」\\


　崔毅の声に、藁の上で眠っていた帝と陳留王は、夢かとばかり狂喜した。そしてなお、閔貢の拝座するすがたを見ると、うれし泣きに抱き合って号泣された。\\



帝も帝におわさず\\


王また王に非ず\\


千乗万騎走るなる\\

\ruby{北{\fontspec{jigmo}\char"9099}}{ほくぼう}の草野、\ruby{夏}{なつ}\ruby{茫々}{ぼうぼう}\\



　――思いあわせればこの夏の初め頃から、洛陽の童女のなかにこんな歌が\ruby{流行}{はや}っていた。天に口なく、無心の童歌をして、今日のことを予言していたものだろうか。\\


「天下一日も帝なかるべからずです。さあ、一刻も早く、都へご還幸なされませ」\\


　閔貢のことばに、崔毅は、自分の\ruby{厩}{うまや}から、一匹の\ruby{痩馬}{そうば}を曳いてきて、帝に献上した。\\


　閔貢は、自分の馬に、陳留王を乗せて、二騎の口輪をつかみ、門を出て、諸所へ散らかっている兵をよび集めた。\\


　二、三里ほど来ると、\\


「おお、帝はご無事でおわしたか」\\


　校尉\ruby{袁紹}{えんしょう}が馳せ出会う。\\


　また、司徒\ruby{王允}{おういん}、太尉\ruby{楊彪}{ようひょう}、\ruby{左軍校尉}{さぐんこうい}\ruby{淳于瓊}{じゅんうけい}、右軍の\ruby{趙萌}{ちょうぼう}、同じく\ruby{後軍校尉}{ごぐんこうい}\ruby{鮑信}{ほうしん}などがめいめい数百騎をひきいて来合せ、帝にまみえて、みな\ruby{哭}{な}いた。\\


「還御を盛んにし、洛陽の市民にも安心させん」\\


　と、段珪の首を、早馬で先へ送り、洛陽の市街に\ruby{曝}{さら}し首として、同時に、帝のご無事と還幸を布告した。\\


　かくて帝の\ruby{御駕}{ぎょが}は、郊外の近くまでさしかかって来た。するとたちまち彼方の丘の陰から\ruby{旺}{さかん}なる兵気馬塵が立ち昇り、一隊の旌旗、天をおおって見えたので、\\


「や、や？」とばかり、随身の将卒百官、みな色を失って立ちすくんだ。\\

\\

\subsubsection{四}

\\


「敵か？」\\


「そも、\ruby{何人}{なにびと}の軍ぞ」\\


　帝をはじめ、茫然、疑い怖れているばかりだったが、時に\ruby{袁紹}{えんしょう}あって、\ruby{鹵簿}{ろぼ}の前へ馬をすすめ、\\


「それへ来るは、何者の軍隊か。帝いま、皇城に還り給う。道をふさぐは不敬ではないか」\\


　と、大喝した。\\


　すると、\\


「おうっ。吾なり」\\


　と吠えるが如き答が、正面へきた軍の真ん中に轟き聞えた。\\


　千\ruby{翻}{ぽん}の旗、\ruby{錦繍}{きんしゅう}の\ruby{幡旗}{はんき}、さっと隊を開いたかと見れば駿馬は\ruby{龍爪}{りゅうそう}を掻いて、堂々たる鞍上の一偉夫を、袁紹の前へと馳け寄せてきた。\\


　これなん先頃から洛陽郊外の\ruby{{\fontspec{jigmo}\char"6FA0}池}{べんち}に兵馬を\ruby{駐}{と}めたまま、何進が再三召し呼んでも動かなかった\ruby{惑星}{わくせい}の人――\ruby{西涼}{せいりょう}の\ruby{刺史}{しし}\ruby{董卓}{とうたく}であった。\\


　董卓、\ruby{字}{あざな}は\ruby{仲穎}{ちゅうえい}、\ruby{隴西臨{\fontspec{jigmo}\char"6D2E}}{ろうせいりんとう}（甘粛省岷県）の生れである。身長八尺、腰の太さ十囲という。肉脂豊重、眼細く、\ruby{豺智}{さいち}の光り針がごとく人を刺す。\\


　\ruby{袁紹}{えんしょう}が、\\


「何者だっ」\\


　と、咎めたが、部将などは眼中にないといった態度で、\\


「天子はいずこに\ruby{在}{おわ}すか」\\


　と、\ruby{鹵簿}{ろぼ}の間近まで寄ってくる様子なのだ。帝は、戦慄されて、お答えもなし得ないし、百官も皆、怖れわななき、さすがの袁紹さえも、その容態の立派さに、呆っ気にとられて\ruby{阻}{はば}めもできなかった。\\


　すると、帝の御駕のすぐうしろから、\\


「ひかえろッ」\\


　\ruby{涼}{すず}やかに叱った者がある。\\


　凜たる音声に、董卓も思わず駒をすこし\ruby{退}{ひ}いて、\\


「何。控えろと。――そういう者は誰だっ」と眼をみはった。\\


「おまえこそ、名をいえ」\\


　こういって馬を前へ出してきたのは、皇弟の陳留王であった。帝よりも年下の紅顔の少年なのである。\\


「……あっ。皇弟の陳留王でいらっしゃいますな」\\


　董卓も、気がついてあわてて、馬上で礼儀をした。\\


　陳留王は、あくまで頭を高く、\\


「そうだ。そちは誰だ」\\


「西涼の刺史董卓です」\\


「その董卓が、何しに来たか。――\ruby{聖駕}{せいが}をお迎えに参ったのか、それとも奪い取ろうという気で来たか」\\


「はっ……」\\


「いずれだ！」\\


「お迎えに参ったのでござる」\\


「お迎えに参りながら、天子のこれにましますに、下馬せぬ無礼者があるかっ、なぜ、馬をおりん！」\\


　身なりは小さいが、王の声は実に峻烈であった。威厳に打たれたか、董卓は二言もなく、あわてて馬からとびおりて、道のかたわらに退き、謹んで帝の車駕を拝した。\\


　陳留王は、それを見ると、帝に代って、\\


「大儀であった」\\


　と、董卓へ言葉を下した。\\


　\ruby{鹵簿}{ろぼ}は難なく、洛陽へさして進んだ。心ひそかに舌を巻いたのは董卓であった。天性備わる陳留王の威風にふかく胆を奪われて、\\


「これは、今の帝を廃して、陳留王を御位に立てたほうが……？」\\


　と、いう大野望が、早くもこの時、彼の胸には芽を\ruby{兆}{きざ}していた。\\

\\

\subsection{\ruby{呂布}{りょふ}}

\\

\subsubsection{一}

\\


　洛陽の\ruby{余燼}{よじん}も、ようやく\ruby{熄}{や}んだ。\\


　帝と皇弟の車駕も、かくて無事に宮門へ還幸になった。\\


　\ruby{何太后}{かたいこう}は、帝を迎えると、\\


「おお」\\


　と、共に相擁したまま、しばらくは\ruby{嗚咽}{おえつ}にむせんでいた。\\


　そして太后はすぐ、\\


「\ruby{玉璽}{ぎょくじ}を――」\\


　と、帝のお手にそれを戻そうとして求めたが、いつのまにか紛失していた。\\


　伝国の玉璽が見えなくなったことは漢室として大問題である。だがそれだけに、絶対に秘密にしていたが、いつか洩れたとみえてひそかに聞く者は、\\


「ああ。またそんな\ruby{亡兆}{ぼうちょう}がありましたか」と、眉をひそめた。\\


　\ruby{董卓}{とうたく}はその後、\ruby{{\fontspec{jigmo}\char"6FA0}池}{べんち}の兵陣を、すぐ城外まで移してきて、自身は毎日、千騎の鉄兵をひきつれて市街王城をわが物顔に横行していた。\\


「寄るな」\\


「咎められるな」\\


　人民は\ruby{恟々}{きょうきょう}と、道をひらいて避けた。\\


　その頃、\ruby{并州}{へいしゅう}の丁原、\ruby{河内}{かだい}の太守\ruby{王匡}{おうきょう}、東郡の\ruby{喬瑁}{きょうぼう}などと諸将がおくればせに先の詔書に依って上洛して来たが、董卓軍の有様を見て皆、なすことを知らなかった。\\


　後軍の校尉\ruby{鮑信}{ほうしん}は、ある時、\ruby{袁紹}{えんしょう}に向ってそっとささやいた。\\


「どうかしなければいかんでしょう。あいつらの\ruby{沓}{くつ}は、\ruby{内裏}{だいり}も街もいっしょくたに濶歩しておる」\\


「なんのことだ」\\


「知れきったことでしょう。\ruby{董卓}{とうたく}とその周りの連中ですよ」\\


「だまっていたまえ」\\


「なぜです。私は、安からぬ思いがしてなりませんが」\\


「でも、この頃ようやく、宮廷も少しお静かになりかけたところだからな」\\


　鮑信はまた、同じような憂えを、司徒の\ruby{王允}{おういん}にもらした。けれど司法官たる王允でも、董卓のような大物となるとどうしようもなかった。\\


　網をたずさえた\ruby{漁夫}{りょうし}が、鯨をながめて嘆じるように、\\


「ううむ。まったくだ。同感だ。だが、どうしようもないじゃないか」\\


　\ruby{疎髯}{そぜん}をつまんで、とがった顎を引っ張りながら、そううそぶくだけだった。\\


「やんぬる\ruby{哉}{かな}――」\\


　鮑信は、嫌になって、自分の手勢だけを\ruby{引具}{ひきぐ}し、泰山の閑地へ逃避してしまった。\\


　去る者は去り、\ruby{媚}{こ}ぶる者は媚びて董卓の勢力につき、彼の勢いは日増しに\ruby{旺}{さかん}になるばかりだった。\\


　董卓の性格は、その軍に、彼の態度に、ようやく露骨にあらわれてきた。\\


「\ruby{李儒}{りじゅ}」\\


「はい」\\


「断行しようと思うがどうだろう。もういいだろう」\\


　董卓は、\ruby{股肱}{ここう}の李儒に計った。それは、かねて彼の腹中にあった画策で、現在の天子を廃し、彼の見こんだ陳留王を位につけて、宮廷を私しようという大野望であった。\\


　李儒は、よろしいでしょうと云った。時機は今です、早くおやりなさいともつけ加えた。これも彼に劣らぬ暴逆家だ。しかし董卓は気にいった。\\


　翌日。温明園で大宴会がひらかれた。招きの主人名はいうまでもなく董卓である。ゆえに、その威を怖れて欠席した者はほとんどなかったというてよい。文武の百官はみな集まった。\\


「みなお揃いになりました」\\


　侍臣から知らせると、董卓は容態をつくろって、\ruby{轅門}{えんもん}の前でゆらりと駒をおり、宝石をちりばめた剣を\ruby{佩}{は}いて悠々と席へついた。\\

\\

\subsubsection{二}

\\


　美酒玉杯、数巡して、\\


「今日の宴に列せられた諸公にむかって、予は一言提議したい」\\


　董卓は起って、おもむろにこう発言した。\\


　なにをいうのかと、一同は静まり返った。董卓はその肥満した体をぐっとそらすと、\\


「予は思う。天子は\ruby{天稟}{てんぴん}の玉質であらねばならぬ。万民の\ruby{景仰}{けいぎょう}をあつめるに足るお方であらねばならぬ。\ruby{宗廟社稷}{そうびょうしゃしょく}を護りかためて揺ぎなき仁徳を兼ね備えておわさねばならぬ。しかるに、不幸にも新帝は\ruby{薄志懦弱}{はくしだじゃく}である。漢室のため、われわれ臣民の常に憂うるところである」\\


　大問題だ。\\


　聞く者みな色を\ruby{醒}{さ}ました。\\


　董卓は、\ruby{寂}{せき}としてしまった百官の頭上を見まわして、左の\ruby{拳}{こぶし}を、剣帯に当てがい、右の手をつよく振った。\\


「ここにおいて、予はあえていおう。憂うるなかれ諸卿と。幸いにも、皇弟\ruby{陳留王}{ちんりゅうおう}こそは、学を好み、聡明におわし、天質\ruby{玲瓏}{れいろう}、まことに天子の\ruby{器}{うつわ}といってよい。今や天下多事、よろしくこの際ただ今の天子に代うるに、陳留王をもってし、帝座の廃立を決行したいと考えるが、いかがあろうか。異論あるものは立って意見を述べ給え」\\


　驚くべき大事を、彼は宣言同様にいいだしたのである。広い大宴席に\ruby{咳声}{せき}ひとつ聞えなかった。気をのまれた形でもあろう。董卓は、俺に反対する者などあるわけもない――といったように、自信のみちた眼で眺めまわした。\\


　すると、百官の席のうちから、突として誰か立つ音がした。一斉に人々の首は彼のほうを見た。\\


　\ruby{并州}{へいしゅう}の刺史\ruby{丁原}{ていげん}である。\\


「吾輩は起立した、反対の表示である」\\


　董卓は\jruby{くわっ}{・・・・・・}と睨めて、\\


「木像を見ようとは思わない。反対なら反対の意見を吐け」\\


「天子の座は、天子の御意にあるものである。臣下の私議するものではない」\\


「私議はせん。故におれは公論に\ruby{糺}{ただ}しておるのじゃっ」\\


「先帝の正統なる\ruby{御嫡子}{おんちゃくし}たる今の帝に、なんの\ruby{瑕瑾}{かきん}やあらん、咎めやあらん。こんな所で、帝位の廃立を議するとは何事だ。おそらく、\ruby{纂奪}{さんだつ}を\ruby{企}{たくら}む者でなくば、そんな暴言は吐けまい」\\


　皮肉ると、董卓は、\\


「だまれっ、われに\ruby{反}{そむ}く者は死あるのみだぞ」\\


　\ruby{繍袍}{しゅうほう}の袖をはねて、\ruby{佩剣}{はいけん}の柄に手をかけた。\\


「なにをする気か」\\


　丁原は、\jruby{びく}{・・・・・・}ともしなかった。\\


　それも道理、彼のうしろには、一個の偉丈夫が儼然と立っていて、\\


（丁原に指でもさしてみろ）といわんばかり恐ろしい顔していた。\\


　\ruby{爛々}{らんらん}たるその\ruby{眸}{ひとみ}、\ruby{凜々}{りんりん}たる威風、見るからに\ruby{猛豹}{もうひょう}の気がある。\\


　董卓の股肱として、常に秘書のごとく側へついている\ruby{李儒}{りじゅ}は、あわてて主人の袖を引っぱった。\\


「きょうは折角の\ruby{御宴}{ぎょえん}です。かたくるしい国政向きのことなどは、席を改めて、他日になすっては如何です。とかく酒気のあるところでは、論議はまとまりません」\\


「……む、うむ」\\


　董卓も、気づいたので、不承不承、剣の柄から手をさげた。しかしどうも、丁原のうしろに立っている男が気になってたまらなかった。\\

\\

\subsubsection{三}

\\


　――けれど、董卓の野望は、丁原に反対されたぐらいで、決してしぼみはしなかった。\\


　大饗宴の席は一時、そんなことで白け渡ったが、酒杯の交歓ひとしきりあると、董卓はまた起って、\\


「最前、予の述べたところ、おそらく諸君の意中であり、天下の公論と思うがどうだろう」\\


　と、重ねて\ruby{糺}{ただ}した。\\


　すると、席にあった中郎将\ruby{盧植}{ろしょく}が、率直に、彼を意見した。\\


「もうお止めなさい。あまり我意を押しつけようとなさると、天子の廃立に名分をかりて、董公ご自身が、\ruby{簒奪}{さんだつ}の\ruby{肚}{はら}があるのではないかと人が疑います。昔、\ruby{殷}{いん}の\ruby{太甲}{たいこう}\ruby{無道}{むどう}でありしため、\ruby{伊尹}{いいん}これを\ruby{桐宮}{とうきゅう}に放ち、漢の\ruby{昌邑}{しょうゆう}が王位に登って――」\\


　なにか、故事をひいて、学者らしく諫言しかけると、董卓は、\\


「だまれっ、だまれっ――貴様も血祭りに首を出したいのか」\\


　と激怒して、周囲の武将をかえりみ、\\


「彼を斬れっ。斬っちまえ。斬らんかっ」と指さし震えた。\\


　けれど、李儒は、押止め、\\


「いけません」と、いった。\\


「盧植は海内の学者です。中郎将としてよりも、\ruby{大儒}{たいじゅ}として名が知られています。それを董卓が殺したと天下へ聞えることは、あなたの不徳になります。ご損です」\\


「では、追っ払えっ」\\


　董卓は、またつづけざまに怒号した。\\


「官職を引っ\ruby{剥}{ぱ}いでだぞ。――盧植を官に置こうという者はおれの相手だ」\\


　もう、誰も止めなかった。\\


　盧植は、官を逐われた。この日から先、彼は世を見限って、\ruby{上谷}{じょうこく}の\ruby{閑野}{かんや}にかくれてしまった。\\


　それは、さておき、饗宴もこんなふうで、殺伐な散会となってしまった。帝位廃立の議は、またの日にしてと、百官は逃げ腰に閉会の乾杯を\ruby{強}{し}いてあげた。\\


　司徒\ruby{王允}{おういん}などは、真っ先にこそこそ帰った。董卓はなお、丁原の反対に根をもって、\ruby{轅門}{えんもん}に待ちうけて、彼を斬って捨てんと、剣を按じていた。\\


　ところが。\\


　最前から轅門の外に、黒馬に踏みまたがって、手に\ruby{方天戟}{ほうてんげき}をひっさげ、しきりと帰る客を物色したり、門内をうかがったりしている風貌非凡な若者がある。\\


　ちらと、董卓の眼にとまったので、彼は\ruby{李儒}{りじゅ}を呼んで訊ねた。李は外をのぞいて、\\


「あれですよ、最前、丁原のうしろに突っ立っていた男は」\\


「あれか。はてな、身なりが違うが」\\


「武装して出直して来たんでしょう。怖ろしい奴です。丁原の養子で、\ruby{呂布}{りょふ}という人間です。\ruby{五原郡}{ごげんぐん}（内蒙古・五原市）の生れで、\ruby{字}{あざな}は\ruby{奉先}{ほうせん}、弓馬の達者で天下無双と聞えています。あんな奴にかまったら\ruby{大事}{おおごと}ですよ。避けるに\ruby{如}{し}くはなし。見ぬふりをしているに限ります」\\


　聞いていた董卓は、にわかに恐れを覚え、あわてて園内の一亭へ隠れこんでしまった。\\


　重ね重ね彼は呂布のために丁原を討ち損じたので、呂布の姿を、夢の中にまで大きく見た。どうも忘れ得なかった。\\


　するとその翌日。\\


　こともにわかに、丁原が兵を率いて、董卓の陣を急に襲ってきた。彼は聞くや否や、大いに怒って、たちまち身を鎧い、陣頭へ出て見ていると、たしかに昨日の呂布、黄金の\ruby{{\fontspec{jigmo}\char"76D4}}{かぶと}をいただき、百\ruby{花戦袍}{かせんぽう}を着、\ruby{唐猊}{からしし}の鎧に、\ruby{獅蛮}{しばん}の\ruby{宝帯}{ほうたい}をかけ、方天戟をさげて、縦横無尽に馬上から斬りまくっている有様に――董卓は敵ながら見とれてしまい、また内心ふかく怖れおののいた。\\

\\

\subsection{\ruby{赤兎馬}{せきとば}}

\\

\subsubsection{一}

\\


　その日の戦いは、\ruby{董卓}{とうたく}の大敗に帰してしまった。\\


　\ruby{呂布}{りょふ}の勇猛には、それに当る者もなかった。\ruby{丁原}{ていげん}も、十方に馬を躍らせて、董卓軍を蹴ちらし、大将董卓のすがたを乱軍の中に見かけると、\\


「\ruby{簒逆}{さんぎゃく}の賊、これにありしか」と、馳け迫って、\\


「漢の天下、内官の\ruby{弊悪}{へいあく}にみだれ、万民みな塗炭の苦しみをうく。しかるに、汝は涼州の一\ruby{刺史}{しし}、国家に一寸の功もなく、ただ\ruby{乱隙}{らんげき}をうかがって、野望を遂げんとし、みだりに帝位の廃立を議するなど、身のほど知らずな逆賊というべきである。いでその\ruby{素頭}{すこうべ}を刎ねて、\ruby{巷}{ちまた}に\ruby{梟}{か}け、洛陽の民の祭に供せん」\\


　と討ってかかった。\\


　董卓は、一言もなく、敵の優勢に怖れ、自身の恥ずる心にひるんで、あわてて味方の楯の内へ逃げこんでしまった。\\


　そんなわけで董卓の軍は、その日、士気のあがらないことおびただしく、董卓も腐りきった態で、遠く陣を退いてしまった。\\


　夜――\\


　本陣の燈下に、彼は諸将を呼んで嘆息した。\\


「敵の丁原はともかく、養子の呂布がいるうちは勝ち目がない。呂布さえおれの配下にすれば、天下はわが\ruby{掌}{たなごころ}のものだが――」\\


　すると、諸将のうちから、\\


「将軍。嘆ずるには及びません」と、いった者がある。\\


　人々がかえりみると、\ruby{虎賁中郎将}{こほんちゅうろうしょう}の\ruby{李粛}{りしゅく}であった。\\


「李粛か。なんの策がある？」\\


「あります。私に、将軍の愛馬\ruby{赤兎}{せきと}と一\ruby{嚢}{ふくろ}の金銀珠玉をお託しください」\\


「それをどうするのか」\\


「幸いにも、私は、呂布と同郷の生れです。彼は勇猛ですが賢才ではありません。以上の二品に、私の持っている三寸\ruby{不爛}{ふらん}の舌をもって、呂布を訪れ、将軍のお望みを、きっとかなえてみせましょう」\\


「ふム。成功するかな？」\\


「まず、おまかせ下さい」\\


　でもまだ迷っている顔つきで、董卓は、側にいる\ruby{李儒}{りじゅ}の意見をきいた。\\


「どうしよう。李粛はあのように申すが」\\


　すると李儒は、\\


「天下を得るために、なんで一匹の馬をお惜しみになるんです」と、いった。\\


「なるほど」\\


　董卓は大きくうなずいて、李粛の献策を容れることにし、秘蔵の名馬\ruby{赤兎}{せきと}と、一\ruby{嚢}{ふくろ}の金銀珠玉とを彼に託した。\\


　赤兎は稀代の名馬で、一日よく千里を走るといわれ、馬体は真っ赤で、風をついて\ruby{奔馳}{ほんち}する時は、その\ruby{鬣}{たてがみ}が炎の流るるように見え、将軍の赤兎といえば、知らない者はないくらいだった。\\


　李粛は、二人の従者にその名馬をひかせ、金銀珠玉をたずさえて、その翌晩、ひそかに呂布の陣営を訪問した。\\


　呂布は彼を見ると、\\


「やあ、貴公か」と、手を打ってよろこび、「君と予とは、同郷の友だがその後お互いに消息も聞かない。いったい今はどうしているのか」と、帳中へ迎え入れた。\\


　李粛も、\ruby{久濶}{きゅうかつ}を\ruby{叙}{じょ}して、\\


「自分は漢朝に仕えて、今では\ruby{虎賁}{こほん}中郎将の職を奉じている。君も、\ruby{社稷}{しゃしょく}を扶けて大いに国事に尽していると聞いて、実は今夜、祝いに来たわけだ」と、いった。\\

\\

\subsubsection{二}

\\


　その時、呂布はふと耳をそばだてて、李粛へ訊いた。\\


「今、陣外にいなないたのは、君の乗馬か、啼き声だけでもわかるが、素晴らしい名馬を持っているじゃないか」\\


「いや、外につないであるのは、自分の乗用ではない。\ruby{足下}{そっか}に進上するために、わざわざ従者に曳かせて来たのだ。気に入るかどうか、見てくれ給え」と、外へ誘った。\\


　呂布は、\ruby{赤兎馬}{せきとば}を一見すると、\\


「これは稀代の逸駿だ」と驚嘆して、\\


「こんな贈り物を受けても、おれはなにも酬いるものがないが」\\


　と、陣中ながら酒宴をもうけて歓待に努める容子は、心の底からよろこんでいるふうだった。\\


　酒、たけなわの頃を計って、\\


「だが呂布君。折角、君に贈った馬だが、赤兎馬のことは、足下の父がよく知っておるから、必ず君の手からとり上げてしまうだろう。それが残念だな」\\


　李粛がいうと、\\


「は……何をいうのか、君はだいぶ酔ってきたな」\\


「どうして」\\


「吾輩の父は、もう世を去ってこの世に\ruby{亡}{な}い人じゃないか。なんでおれの馬を奪おう」\\


「いやいや。わしのいうのは足下の実父ではない。養父の\ruby{丁原}{ていげん}のことだ」\\


「あ。養父のことか」\\


「思えば、足下ほどな武勇才略を備えながら、\ruby{墻}{かき}の\ruby{内}{うち}の羊みたいに飼われているのは、実に惜しいものだ」\\


「けれど、父亡き後、久しく丁原の邸に養われてきた身だから、今さら、どうにもならん」\\


「ならん？　……そうかなあ」\\


「おれだって、若いし、大いに雄才を伸ばしてみたい気もするが」\\


「そこだ、呂布君。\ruby{良禽}{りょうきん}は木を選んで\ruby{棲}{す}むという。日月は\ruby{遷}{うつ}りやすし。空しく青春の時を過すのは愚かではないか」\\


「む、む。……では李君。貴公のみるところでは、今の朝臣の中で、英雄とゆるしてよい人は、一体誰だと思うか」\\


　李粛は一言のもとに、\\


「それやあ、\ruby{董卓}{とうたく}将軍さ」といった。\\


「賢を敬い、士に篤く、寛仁徳望を兼備している英傑といえば董卓をおいては、ほかに人物はない。必ずや将来大業をなす人はまずあの将軍だろうな」\\


「そうかなあ。……やはり」\\


「足下はどう思う」\\


「いや、実はこの呂布も、日頃そう考えているが、何しろ丁原と仲が悪いし、それに縁もないので――」\\


　聞きもあえず李粛は、たずさえてきた金銀珠玉をそれに取りだして、\\


「これこそ、その董卓公から、貴公へ礼物として送られた物だ。実は、予はその使いとして来たわけだ」\\


「えっ。これを」\\


「赤兎馬もご自身の愛馬で、一城とも取換えられぬ――といっておられるほど秘蔵していた馬だが、ご辺の武勇を慕って、どうか上げてくれというお言葉じゃ」\\


「ああ。それまでにこの呂布を愛し給うか。何をもって、おれは知己の篤い志に酬いたらいいのか」\\


「いや、それはやすいことだ。耳を貸し給え」と、李粛はすり寄った。\\


　陣帳風暗く、夜は\ruby{更}{ふ}けかけていた。兵はみな\ruby{睡}{ねむ}りに落ち、時おり、馴れぬ\ruby{厩}{うまや}につながれた赤兎馬が、\ruby{静寂}{しじま}を破って、\ruby{蹄}{ひづめ}の音をさせているだけだった。\\

\\

\subsubsection{三}

\\


「……よしっ」\\


　呂布は大きくうなずいた。\\


　何事かを、その耳へささやいた李粛は、彼の怪しくかがやく眼を見つめながら、そばを離れて、\\


「善は急げという。ご決心がついたら直ぐやり給え。予は、ここで酒を酌んで、\ruby{吉左右}{きっそう}を待っていよう」と、煽動した。\\


　呂布は、直ちに出て行った。\\


　そして営の中軍へ入って、丁原の幕中をうかがった。\\


　丁原は、\ruby{燈火}{ともしび}をかかげて、書物を見ていたが、何者か入ってきた様子に、\\


「誰だっ」と、振向いた。\\


　血相の変った呂布が剣を抜いて突っ立っているので、\ruby{愕然}{がくぜん}と立ち、\\


「呂布ではないか。何事だ、その血相は」\\


「何事でもない。大丈夫たるものなんで汝がごとき\ruby{凡爺}{ぼんや}の子となって\ruby{朽}{く}ちん」\\


「ばッ、ばかっ。もう一度いってみい」\\


「何を」\\


　呂布は、躍りかかるや否や、一刀のもとに、丁原を斬り伏せ、その首を落した。\\


　黒血は燈火を消し、夜は惨として暗澹であった。\\


　呂布は、狂える如く、中軍に立って、\\


「丁原を斬った。丁原は不仁なるゆえに、これを斬った。志ある者はわれにつけ。不服な者は、我を去れっ」と、大呼して馳けた。\\


　中軍は騒ぎ立った。去る者、従う者、混乱を極めたが、半ばは、ぜひなく呂布についてとどまった。\\


　この騒ぎが揚ると、\\


「大事成れり」と、李粛は手を打っていた。\\


　やがて直ちに、呂布を伴い、\ruby{董卓}{とうたく}の陣へ帰ってきて、事の次第を報告すると、\\


「でかしたり李粛」と、董卓のよろこびもまた、非常なものであった。\\


　翌日、特に、呂布のために盛宴をひらいて、董卓自身が出迎えるというほどの歓待ぶりであった。\\


　呂布は、贈られたところの赤兎馬にまたがって来たが、鞍をおりて、\\


「士はおのれを知る者の為に死すといいます。今、暗きを捨てて明らかなるに仕う日に会い、こんなうれしいことはありません」と、\ruby{拝跪}{はいき}していった。\\


　董卓もまた、\\


「今、大業の途に、足下のごとき俊猛をわが軍に迎えて、\ruby{旱苗}{かんびょう}に雨を見るような気がする」\\


　と、手をとって、酒宴の席へ迎え入れた。\\


　呂布は、有頂天になった。\\


　しかもまた、黄金の\ruby{甲}{よろい}と\ruby{錦袍}{きんぽう}とをその日の引出物として貰った。恐るべき毒にまわされて、呂布は有頂天に酔った。好漢、惜しむらくは眼前の慾望にくらんで、遂に、青雲の大志を踏み誤ってしまった。\\


　　　　　　　×　　　　　×　　　　　×\\


　呂布は、\ruby{檻}{おり}に入った。\\


　董卓はもう怖ろしい者あるを知らない。その威勢は、旭日のように\ruby{旺}{さかん}だった。\\


　自分は、前将軍を領し、弟の\ruby{董旻}{とうびん}を、左将軍に任じ、呂布を\ruby{騎都尉}{きとい}中郎将の\ruby{都亭侯}{とていこう}に封じた。\\


　思うことができないことはない。\\


　――が、まだ一つ、残っている問題がある。帝位の廃立である。李儒はまた、側にあって、しきりにその実現を彼にすすめた。\\


「よろしい。今度は断行しよう」\\


　董卓は、省中に大饗宴を催して再び百官を一堂に招いた。\\

\\

\subsubsection{四}

\\


　洛陽の都会人は、宴楽が好きである。わけて朝廷の百官は皆、舞楽をたしなみ、酒を愛し、長夜にわたるも辞さない酔客が多かった。\\


（――今日は、この間の饗宴の時よりも、だいぶ\ruby{和}{なご}やかに浮いているな）\\


　董卓は、大会場の空気を見まわして、そう察していた。\\


　時分は好し――と、\\


「諸卿！」\\


　彼は、卓から起って、一場の挨拶を試みた。\\


　初めの演舌は、至極、主人側としてのお座\jruby{なり}{・・・・・・}なものであったから、人々はみな一斉に\ruby{酒盞}{しゅさん}をあげて、\\


「謝す。謝す」と、声を和し、拍手の音も、しばし鳴りもやまなかった。\\


　董卓は、その沸騰ぶりを、自分への人気と見て、\\


「さて。――いつぞやは遂に諸公のご明判を仰いで議決するまでに至らなかったが、きょうはこの盛会と吉日を\ruby{卜}{ぼく}して、過日、未解決におわった大問題をぜひ一決して、さらに\ruby{盞}{さん}を重ねたいと思うのであるが、諸公のお考えは如何であるか」\\


　と、現皇帝の廃位と陳留王の\ruby{即位推戴}{そくいすいたい}のことを、突然、いいだした。\\


　熱湯が\ruby{冷}{さ}めたように、饗宴の席は、一時に\jruby{しん}{・・・・・・}としてしまった。\\


「…………」\\


「…………」\\


　誰も彼も、この重大問題となると\ruby{唖}{おし}のように黙ってしまった。\\


　すると、一つの席から、\\


「否！　否！」と叫んだ者がある。\\


　中軍の校尉\ruby{袁紹}{えんしょう}であった。\\


　袁紹は、敢然、反対の口火を切っていった。\\


「\ruby{借問}{しゃもん}する！　董将軍。――あなたは何がために、好んで平地に波瀾を招くか。一度ならず二度までも、現皇帝を廃して、陳留王をして御位にかわらしめんなどと、陰謀めいたことを提議されるのか」\\


　董卓は、剣に手をかけて、\\


「だまれっ。陰謀とは何か」\\


「廃帝の議をひそかに計るのが陰謀でなくてなんだ」\\


　袁紹も負けずに呶鳴った。\\


　董卓はまッ青になって、\\


「いつ密議したか。朝廷の百官を前において自分は信ずるところをいっておるのだ」\\


「この宴は私席である。朝議を議するならば、なぜ帝の玉座の前で、なお多くの重臣や、太后のご出座をも仰いでせんか」\\


「えいっ、やかましいっ。私席で嫌なら、汝よりまず去れ」\\


「去らん。おれは、陰謀の宴に頑張って、誰が賛成するか、監視してやる」\\


「いったな、貴様はこの董卓の剣は切れないと思っておるのか」\\


「暴言だっ。――諸君っ、今の声を、なんと聞くか」\\


「天下の権は、予の自由だ。予の説に不満な輩は、袁紹と共に、席を出て行けっ」\\


「ああ。妖雷声をなす、天日も\ruby{真}{ま}っ\ruby{晦}{くら}だ」\\


「世まい言を申しておると、一刀両断だぞ。去れっ、去れっ、異端者め」\\


「誰がおるか、こんな所に」\\


　袁紹は、身をふるわせながら、席を蹴って飛び出した。\\


　その夜のうち、彼は、官へ辞表を出して、遠く\ruby{冀州}{きしゅう}の地へ\ruby{奔}{はし}ってしまった。\\

\\

\subsubsection{五}

\\


　席を蹴って、袁紹が出て行ってしまうと、董卓は、やにわに、客席の一方を強くさして、\\


「\ruby{太傅}{たいふ}\ruby{袁隗}{えんかい}！　袁隗をこれへ引っ張ってこい」\\


　と、左右の武士に命じた。\\


　袁隗はまッ青な顔をして、董卓の前へ引きずられて来た。彼は、袁紹の伯父にあたる者だった。\\


「こら、汝の\ruby{甥}{おい}が、予を恥かしめた上、無礼を極めて出て行った態は、その眼でしかと見ていたであろうが。――ここで汝の首を斬ることを予は知っているが、その前に、ひと言訊いてつかわす。この世と\ruby{冥途}{めいど}の辻に立ったと心得て、肚をすえて返答をせい」\\


「はっ……はいっ」\\


「汝は、この董卓が宣言した帝位廃立をどう思う？　賛同するか、それとも、甥の奴と同じ考えか」\\


「尊命の如し――であります」\\


「尊命の如しとは！」\\


「あなたのご宣言が正しいと存じます」\\


「よしっ。しからばその首をつなぎ止めてやろう。ほかの者はどうだ。我すでに大事を宣せり。\ruby{背}{そむ}く者は、軍法をもって問わん」\\


　剣をあげて、雷の如くいった。\\


　並いる百官も、\ruby{慴伏}{しょうふく}して、もう誰ひとり反対をさけぶ者もなかった。\\


　董卓は、かくて、威圧的に百官に宣誓させて、また、\\


「\ruby{侍中}{じちゅう}\ruby{周{\fontspec{jigmo}\char"6BD6}}{しゅうひ}！　校尉\ruby{伍瓊}{ごけい}！　議郎\ruby{何{\fontspec{jigmo}\char"9852}}{かぎょう}！　――」\\


　と、いちいち役名と名を呼びあげて、その起立を見ながら厳命を発した。\\


「我に\ruby{背}{そむ}いた\ruby{袁紹}{えんしょう}は、必ずや夜のうちに、本国冀州へさして逃げて帰る心にちがいない。彼にも兵力があるから油断はするな。すぐ精兵を率いて追い討ちに打って取れ」\\


「はっ」\\


　三将のうち、二人は命を奉じて、すぐ去りかけたが、侍中周{\fontspec{jigmo}\char"6BD6}のみは、\\


「あいや、おそれながら、仰せはご短慮かと存じます。上策とは思われません」\\


「周{\fontspec{jigmo}\char"6BD6}っ。汝も背く者か」\\


「いえ、袁紹の首一つをとるために、大乱の生じるのを怖れるからです。彼は平常、恩徳を布き、門下には\ruby{吏人}{やくにん}も多く、国には財があります。袁紹\ruby{叛旗}{はんき}を立てたりと聞えれば、山東の国々ことごとく騒いで、それらが、一時にものをいいますぞ」\\


「ぜひもない。予に背く者は討つあるのみだ」\\


「ですが、元来、袁紹という人物は、思慮はあるようでも、決断のない男です。それに天下の大勢を知らず、ただ憤怒に駆られてこの席を出たものの、あれは一種の恐怖です。なんであなたの覇業を妨げるほどな害をなし得ましょうや。むしろ喰らわすに利をもってし、彼を一郡の太守に封じ、そっとしておくに限ります」\\


「そうかなあ？」\\


　座右をかえりみて呟くと、\ruby{蔡{\fontspec{jigmo}\char"9095}}{さいよう}も大きに道理であると、それに賛意を表した。\\


「では、袁紹を追い討ちにするのは、見あわせとしよう」\\


「それがいいです、上策と申すものです」\\


　口々からでる\ruby{讃礼}{さんらい}の声を聞くと、董卓はにわかに気が変って、\\


「使いを立てて、袁紹を\ruby{渤海郡}{ぼっかいぐん}の太守に任命すると伝えろ」\\


　と、厳命を変更した。\\


　その後。\\


　九月\ruby{朔日}{ついたち}のことである。\\


　董卓は、帝を嘉徳殿に請じて、その日、文武の百官に、\\


　――今日出仕せぬ者は、斬首に処せん。\\


　という布告を発した。そして殿上に抜剣して、玉座をも\jruby{しり}{・・・・・・}目に、\\


「\ruby{李儒}{りじゅ}、宣文を読め」\\


　と\ruby{股肱}{ここう}の彼にいいつけた。\\

\\

\subsubsection{六}

\\


　予定の計画である。李儒は、はっと答えるなり、用意の宣言文をひらいて、\\


「\ruby{策文}{さくもん}っ――」\\


　と高らかに読み始めた。\\



孝霊皇帝\\

\ruby{眉寿}{ビジュ}ノ\ruby{祚}{サイワイ}ヲ\ruby{究}{キワ}メズ\\


早ク臣子ヲ\ruby{棄給}{ステタマ}ウ\\


皇帝\ruby{承}{ウ}ケツイデ\\


海内側望ス\\


而シテ天資\ruby{軽佻}{ケイチョウ}\\


威儀ツツシマズシテ\ruby{慢惰}{マンダ}\\


凶徳スデニアラワレ\\


神器ヲ\ruby{損}{ソコナ}イ\ruby{辱}{ハズカ}シメ宗廟ケガル\\

\ruby{太后}{タイコウ}マタ\ruby{教}{オシ}エニ母儀ナク\\

\ruby{政治}{マツリゴト}\ruby{統}{スベ}テ荒乱\\


衆論ココニ起ル\ruby{大革}{タイカク}ノ道\\



　李儒は、さらに声を大にして読みつづけていた。\\


　百官の\ruby{面}{おもて}は色を失い、玉座の帝はおおのき\ruby{慄}{ふる}え、嘉徳殿上\ruby{寂}{せき}として墓場のようになってしまった。\\


　すると突然、\\


「ああ、ああ……」\\


　と、\ruby{嗚咽}{おえつ}して泣く声が流れた。帝の側にいた\ruby{何太后}{かたいこう}であった。\\


　太后は涙にむせぶの余り、ついに椅子から坐りくずれ、帝のすそにすがりついて、\\


「誰がなんといっても、あなたは漢の皇帝です。うごいてはいけませんよ。玉座から降ってはなりませんよ」\\


　と、いった。\\


　董卓は、剣を片手に、\\


「今、李儒が読み上げた通り、帝は\ruby{闇愚}{あんぐ}にして威儀なく、太后は教えにくらく母儀の\ruby{賢}{けん}がない。――依って今日より、現帝を\ruby{弘農王}{こうのうおう}とし、何太后は永安宮に押しこめ、代るに\ruby{陳留王}{ちんりゅうおう}をもって、われらの皇帝として\ruby{奉戴}{ほうたい}する」\\


　いいながら、帝を玉座から引き降ろして、その\ruby{璽綬}{じじゅ}を解き、北面して臣下の列の中へ無理に立たせた。\\


　そして、泣き狂う何太后をも、即座にその\ruby{后衣}{こうい}を\ruby{剥}{は}いで、\ruby{平衣}{へいい}とさせ、後列へしりぞけたので、群臣も思わず眼をおおうた。\\


　時に。\\


　ただ一人、大音をあげて、\\


「待てっ逆臣っ。汝董卓、そも誰から大権をうけて、天を\ruby{欺}{あざむ}き、聖明の天子を、\ruby{強}{し}いてひそかに廃せんとするか。――\ruby{如}{し}かず！　汝と共に刺しちがえて死のう」\\


　いうや否、群臣のうちから騒ぎだして、董卓を目がけて短剣を突きかけてきた者があった。\\


　\ruby{尚書}{しょうしょ}\ruby{丁管}{ていかん}という若い純真な宮内官であった。\\


　董卓は、おどろいて身をかわしながら、醜い声をあげて救けを呼んだ。\\


　刹那――\\


「うぬっ、何するかっ」\\


　横から跳びついた\ruby{李儒}{りじゅ}が、抜打ちに丁管の首を斬った。同時に、武士らの刃もいちどに丁管の五体に集まり、殿上はこの若い一義人の鮮血で彩られた。\\


　さはあれ、ここに。\\


　董卓は遂にその目的を達し、陳留王を立てて天子の位につけ奉り、百官もまた彼の暴威に怖れて、万歳を唱和した。\\


　そして、新しき皇帝を\ruby{献帝}{けんてい}と申上げることになった。\\


　だが、献帝はまだ年少である。何事も董卓の意のままだった。\\


　即位の式がすむと、董卓は自分を\ruby{相国}{しょうこく}に封じ、\ruby{楊彪}{ようひょう}を司徒とし、\ruby{黄{\fontspec{jigmo}\char"742C}}{こうえん}を太尉に、\ruby{荀爽}{じゅんそう}を司空に、\ruby{韓馥}{かんふく}を\ruby{冀州}{きしゅう}の牧に、\ruby{張資}{ちょうし}を南陽の太守に――といったように、地方官の任命も\ruby{輦下}{れんか}の朝臣の登用も、みな自分の腹心をもって当て、自分は相国として、宮中にも\ruby{沓}{くつ}をはき、剣を\ruby{佩}{は}いて、その肥大した体躯をそらしてわが物顔に殿上に横行していた。\\


　同時に。\\


　年号も\ruby{初平}{しょへい}元年と改められた。\\

\\

\subsection{\ruby{春園走獣}{しゅんえんそうじゅう}}

\\

\subsubsection{一}

\\


　まだ若い廃帝は、明け暮れ泣いてばかりいる母の\ruby{何太后}{かたいこう}と共に、永安宮の幽居に深く閉じこめられたまま、春をむなしく、月にも花にも、ただ悲しみを誘わるるばかりだった。\\


　\ruby{董卓}{とうたく}は、そこの衛兵に、\\


「監視を怠るな」と厳命しておいた。\\


　見張りの衛兵は、春の\ruby{日永}{ひなが}を、あくびしていたが、ふと\ruby{幽楼}{ゆうろう}の上から、哀しげな\ruby{詩}{うた}の声が聞えてきたので、聞くともなく耳を澄ましていると、\\



春は来ぬ\\


けむる\ruby{嫩草}{わかくさ}に\\

\ruby{{\fontspec{jigmo}\char"88CA}々}{じょうじょう}たり\\


双燕は飛ぶ\\


ながむれば都の水\\


遠く一すじ青し\\

\ruby{碧雲}{へきうん}深きところ\\


これみなわが旧宮殿\\

\ruby{堤上}{ていじょう}、義人はなきや\\


忠と義とによって\\


誰か、晴らさん\\


わが心中の怨みを――\\



　衛兵は、聞くと、その詩を覚え書にかいて、\\


「\ruby{相国}{しょうこく}。廃帝の弘農王が、こんな詩を作って歌っていました」\\


　と、密告した。董卓は、それを見ると、\\


「\ruby{李儒}{りじゅ}はいないか」\\


　と呼び立てた。そして、その詩を李儒に示して、\\


「これを見ろ、幽宮におりながら、こんな悲歌を作っている。生かしておいては必ずや後の害になろう。何太后も廃帝も、おまえの処分にまかせる。殺して来い」と、いいつけた。\\


「承知しました」\\


　李儒はもとより暴獣の爪のような男だ。情けもあらばこそ、すぐ十人ばかりの屈強な兵を連れて、永安宮へ馳せつけた。\\


「どこにおるか、王は」\\


　彼はずかずか楼上へ登って行った。折ふし弘農王と何太后とは、楼の上で春の憂いに沈んでおられ、突然、李儒のすがたを見たのでぎょっとした\ruby{容子}{ようす}だった。\\


　李儒は笑って、\\


「なにもびっくりなさる事はありません。この春日を慰め奉れ、と相国から酒をお贈り申しにきたのです。これは延寿酒といって、百歳の\ruby{齢}{よわい}を延ぶる美酒です。さあ一\ruby{盞}{さん}おあがりなさい」\\


　携えてきた一壺の酒を取り出して杯を\ruby{強}{し}いると、廃帝は、眉をひそめて、\\


「それは毒酒であろう」と、涙をたたえた。\\


　太后も顔を振って、\\


「\ruby{相国}{しょうこく}がわたし達へ、延寿酒を贈られるわけはない。李儒、これが毒酒でないなら、そなたがまず先に飲んでお見せなさい」といった。\\


　李儒は、眼を怒らして、\\


「なに、飲まぬと。――それならば、この二品をお受けなさるか」\\


　と、\ruby{練絹}{ねりぎぬ}の縄と短刀とを、突きつけた。\\


「……おお。我に死ねとか」\\


「いずれでも好きなほうを選ぶがよい」\\


　李儒は冷然と毒づいた。\\


　弘農王は、涙の中に、\\



ああ、天道は\ruby{易}{かわ}れり\\


人の道もあらじ\\

\ruby{万乗}{ばんじょう}の\ruby{位}{くらい}をすてて\\


われ何ぞ安からん\\


臣に迫られて\ruby{命}{めい}はせまる\\


ただ\ruby{潸々}{さんさん}、涙あるのみ\\



　と、悲歌をうたってそれへ泣きもだえた。\\


　太后は、\jruby{はった}{・・・・・・}と李儒を睨めつけて、\\


「国賊！　\ruby{匹夫}{ひっぷ}！　おまえ達の滅亡も、決して長い先ではありませぬぞ。――ああ兄の\ruby{何進}{かしん}が愚かなため、こんな獣どもを都へ呼び入れてしまったのだ」\\


　\ruby{罵}{ののし}り狂うのを、李儒はやかましいとばかり、その襟がみをつかみ寄せて、高楼の欄から投げ落してしまった。\\

\\

\subsubsection{二}

\\


「どうしたか」\\


　董卓は美酒を飲みながら、李儒の\ruby{吉左右}{きっそう}を待っていた。\\


　やがて李儒は、\ruby{袍}{ほう}を血まみれに汚して戻ってきたが、いきなり提げていた二つの首を突きだして、\\


「\ruby{相国}{しょうこく}、ご命令通り致してきました」と、いった。\\


　弘農王の首と、何太后の首であった。\\


　二つとも首は眼をふさいでいたが、その眼がかっと開いて、今にも飛びつきそうに、董卓には見えた。\\


　さすがに眉をひそめて、\\


「そんな物、見せんでもいい。城外へ埋めてしまえ」\\


　それから彼は、日夜、大酒をあおって、禁中の宮内官といい、後宮の女官といい、気に入らぬ者は立ちどころに殺し、夜は天子の\ruby{床}{しょう}に横たわって春眠をむさぼった。\\


　或る日。\\


　彼は陽城を出て、四頭立ての\ruby{驢車}{ろしゃ}に美人を大勢のせ、酔うた彼は、\ruby{馭者}{ぎょしゃ}の真似をしながら、城外の梅林の花ざかりを逍遥していた。\\


　ところが、ちょうど村社の祭日だったので、なにも知らない農民の男女が晴れ着を飾って帰ってきた。\\


　\ruby{董相国}{とうしょうこく}は、それを見かけ、\\


「農民のくせに、この晴日を、田へも出ずに、着飾って歩くなど、不届きな怠け者だ。天下の百姓の見せしめに召捕えろ」と、驢車の上で、急に怒りだした。\\


　突然、相国の従兵に追われて、若い男女は悲鳴をあげて逃げ散った。そのうち逃げ遅れた者を兵が\ruby{拉}{らっ}して来ると、\\


「\ruby{牛裂}{うしざ}きにしろ」\\


　と、相国は\ruby{威猛高}{いたけだか}に命じた。\\


　手脚に縄を縛りつけて、二頭の\ruby{奔牛}{ほんぎゅう}にしばりつけ、東西へ向けて鞭打つのである。手脚を裂かれた人間の血は、梅園の大地を\ruby{紅}{くれない}に汚した。\\


「いや、花見よりも、よほど面白かった」\\


　驢車は\ruby{黄昏}{たそがれ}に陽城へ向って帰還しかけた。\\


　するとある\ruby{巷}{ちまた}の角から、\\


「逆賊ッ」と、\ruby{喚}{おめ}いて、不意に驢車へ飛びついて来た\ruby{漢}{おとこ}がある。\\


　美姫たちは、悲鳴をあげ、驢は狂い合って、\ruby{端}{はし}なくも、大混乱をよび起した。\\


「何するか、\ruby{下司}{げす}っ」\\


　肥大な体躯の持主である相国は、身うごきは敏速を欠くが、力はおそろしく強かった。\\


　\ruby{精悍}{せいかん}な刺客の男は、驢車へ足を踏みかけて、短剣を引抜き、相国の大きな腹を目がけて勢いよく突ッかけて行ったのであったが、董相国にその剣を叩き落され、しっかと、抱きすくめられてしまったので、どうすることもできなかった。\\


「\ruby{曲者}{しれもの}め。誰に頼まれた」\\


「残念だ」\\


「名を申せ」\\


「…………」\\


「誰か、叛逆を企む奴らの与党だろう。さあ、誰に頼まれたか」\\


　すると、苦しげに、刺客はさけんだ。\\


「叛逆とは、臣下が君にそむくことだ。おれは貴様などの臣下であった覚えはない。――おれは朝廷の臣、越騎校尉の\ruby{伍俘}{ごふ}だっ」\\


「斬れッ、こいつを」\\


　驢車から蹴落すとともに、董卓の武士たちは伍俘の全身に無数の刃と槍を加えて、\ruby{塩辛}{しおから}のようにしてしまった。\\


　　　　　　×　　　　　×　　　　　×\\


　都を落ちて、遠く\ruby{渤海郡}{ぼっかいぐん}（河北省）の太守に封じられた\ruby{袁紹}{えんしょう}はその後、洛陽の情勢を聞くにつけ、\ruby{鬱勃}{うつぼつ}としていたが、遂に矢も楯もたまらなくなって、在京の同志で三公の重職にある司徒\ruby{王允}{おういん}へ、ひそかに書を飛ばし、激越な辞句で奮起を促してきた。\\


　だが、王允は、その書簡を手にしてからも、日夜心で苦しむだけで、董相国を討つ計はなにも持たなかった。\\

\\

\subsubsection{三}

\\


　日々、朝廷に上がって、政務にたずさわっていても、\ruby{王允}{おういん}はそんなわけで、少しも勤めに気がのらなかった。心中ひとり\ruby{怏々}{おうおう}と\ruby{悶}{もだ}えを抱いていた。\\


　ところがある日、董相国の息のかかった高官は誰も見えず、皆、前朝廷の旧臣ばかりが一室にいあわせたので、（これぞ、天の与え）とひそかによろこんで、急に座中へ向って誘いかけた。\\


「実は、今日は、此方の誕生日なのじゃが、どうでしょう、\ruby{竹裏館}{ちくりかん}の\ruby{別業}{べつぎょう}のほうへ、諸卿お揃いで\ruby{駕}{が}を\ruby{枉}{ま}げてくれませんか」\\


「ぜひ伺って、公の\ruby{寿}{ことぶき}を祝しましょう」\\


　誰も、差支えをいわなかった。\\


　董卓系の人間をのぞいて、水入らずに話したい気持が、期せずして、誰にも鬱していたからであった。\\


　別業の竹裏館へ、王允は先へ帰ってひそかに宴席の支度をしていた。やがて宵から忍びやかに前朝廷の公卿たちが集まった。\\


　時を得ぬ不遇な人々の密会なので、初めからなんとなく、座中はしめっぽい。その上にまた、酒のすすみだした頃、王允は、冷たい杯を見入って、ほろりと涙をこぼした。\\


　見とがめた客の一人が、\\


「王公。せっかく、およろこびの誕生の宴だというのに、なんで落涙されるのですか」といった。\\


　王允は、長大息をして、\\


「されば、自分の福寿も、今日の有様では、祝う気持にもなれんのじゃ。――不肖、前朝以来、三公の一座を占め、\ruby{政}{まつりごと}にあずかりなから［＃「あずかりなから」はママ］、董卓の勢いはどうすることもできんのじゃ。耳に万民の\ruby{怨嗟}{えんさ}を聞き、眼に漢室の衰亡を見ながら、なんでわが\ruby{寿筵}{じゅえん}に酔えようか」\\


　といって、指で\ruby{瞼}{まぶた}を拭った。\\


　聞くと一座の者も皆、\\


「ああ――」と、大息して、「こんな世に生れ合わせなければよかった。昔、漢の高祖三尺の剣をひっさげて白蛇を斬り、天下を鎮め給うてより王統ここに四百年、なんぞはからん、この末世に生れ合わせようとは」\\


「まったく、われわれの運も悪いものだ。こんな時勢に巡り合ったのは」\\


「――というて、少し大きな声でもして、董相国やその一類の\ruby{誹謗}{ひぼう}をなせば、この首の無事は\ruby{保}{たも}てないし」\\


　などと各{\fontspec{jigmo}\char"303B}、涙やら愚痴やらこぼして\ruby{燭}{しょく}もめいるばかりであったが、その時、末座のほうから突然、\\


「わはははは。あはッはははは」\\


　手を叩いて、誰か笑う者があった。公卿たちは、びっくりして、末席を振返った。見るとそこに若年の一朝臣が、独りで杯をあげ、白面に紅潮をみなぎらせて、人々が泣いたり愚痴るのを、さっきからおかしげに眺めていた。\\


　王允は、その無礼をとがめ、\\


「誰かと思えば、そちは校尉\ruby{曹操}{そうそう}ではないか。なんで笑うか」\\


　すると、曹操はなお笑って、\\


「いや、すみません。しかしこれが笑わずにおられましょうか。朝廷の諸大臣たる方々が、夜は泣いて暁に至り、昼は悲しんで暮れに及び、寄るとさわると泣いてばかりいらっしゃる。これでは天下万民もみな泣き暮しになるわけですな。おまけに、誕生祝いというのに、わさわさ集まって、また泣き上戸の泣き競べとは――。わはははは。失礼ですが、どうもおかしくって、笑いが止まりませんよ。あははは、あははは」\\


「やかましいっ。汝はそもそも、相国\ruby{曹参}{そうさん}が\ruby{後胤}{こういん}で、四百年来、代々漢室の大恩をうけて来ながら、今の朝廷の有様が、悲しくないのか。われわれの憂いが、そんなにおかしいのか。返答によってはゆるさんぞ」\\

\\

\subsubsection{四}

\\


「これは意外なお怒りを――」と、曹操はやや真面目に改まって、\\


「それがしとて何も理のないことを笑ったわけではありません。時の\ruby{大臣}{おとど}ともあろう方々が、\ruby{女童}{おんなわらべ}の如く、日夜めそめそ悲嘆しておらるるのみで\ruby{董卓}{とうたく}を\ruby{誅伏}{ちゅうふく}する\ruby{計}{はかりごと}といったら何もありはしない。――そんな意気地なしなら、時勢を慨嘆したりなどせずに、美人の腰掛けになって胡弓でも聴きながら感涙を流していたらよかろうに――と思ったのでつい笑ってしまった次第です」\\


　と臆面もなくいった。\\


　曹操の皮肉に\ruby{王允}{おういん}をはじめ公卿たちもむっと色をなして、座は白け渡ったが、\\


「しからば何か、そちはそのような広言を吐くからには、董卓を殺す計でもあるというのか。その自信があっての大言か」\\


　王允が再び\ruby{急}{せ}きこんでなじったので、人々は、彼の返答いかにと、\ruby{固唾}{かたず}をのんで、曹操の白い面に\ruby{眸}{ひとみ}をあつめた。\\


「なくてどうしましょう！」\\


　毅然として彼は眉をあげ、\\


「不才ながら小生におまかしあれば、董卓が首を斬って、洛陽の門に\ruby{梟}{か}けてごらんに入れん」\\


　と明言した。\\


　王允は、彼の自信ありげな言葉に、かえって喜色をあらわし、\\


「曹校尉、もし今の言に偽りがないならば、まことに天が義人を地上にくだして、万民の苦しみを助け給うものだ。そも、君にいかなる計やある。願わくば聞かしてもらいたいが」\\


「されば、それがしが常に董相国に近づいて、表面、\ruby{媚}{こ}びへつらって仕えているのは、何を隠そう、隙もあれば彼をひと思いに刺し殺そうと内心誓っているからです」\\


「えっ。……では君には\ruby{疾}{と}くよりそれまでの決心を持っていたのか」\\


「さもなくて、何の大笑大言を諸卿に呈しましょう」\\


「ああ、天下になおこの義人あったか」\\


　王允はことごとく感じて、人々もまたほっと喜色をみなぎらした。\\


　すると曹操は、「時に、王公に小生から、一つのご無心がありますが」といいだした。\\


「何か、遠慮なくいうてみい」\\


「ほかではありませんが、王家には昔より七宝をちりばめた稀代の名刀が伝来されておる由、常々、承っておりますが、董卓を刺すために、願わくばその名刀を、小生にお貸し下さいませんか」\\


「それは、目的さえ必ず仕遂げてくれるならば……」\\


「その儀は、きっとやりのけて見せます。董相国も近頃では、それがしを\ruby{寵愛}{ちょうあい}して、まったく腹心の者同様にみていますから、近づいて一断に斬殺することは、なんの造作もありません」\\


「うム。それさえ首尾よく参るものなら、天下の大幸というべきだ。なんで家宝の名刀一つをそのために惜しもうや」\\


　と、王允はすぐ家臣に命じて、秘蔵の七宝剣を取りだし、手ずからそれを曹操に授け、かつ云った。\\


「しかし、もし仕損じて、事\ruby{顕}{あらわ}れたら一大事だぞ、充分心して\ruby{行}{おこな}えよ」\\


「乞う、\ruby{安}{やす}んじて下さい」\\


　曹操は剣を受け、その夜の酒宴も終ったので、颯爽として帰途についた。七宝の利剣は燦として夜光の珠の帯の如く、彼の腰間にかがやいていた。\\

\\

\subsection{\ruby{白面郎}{はくめんろう}「\ruby{曹操}{そうそう}」}

\\

\subsubsection{一}

\\


　\ruby{曹操}{そうそう}はまだ若い人だ。にわかに、彼の存在は近ごろ大きなものとなったが、その\ruby{年歯風采}{ねんしふうさい}はなお、白面の一青年でしかない。\\


　年二十で、初めて洛陽の北都尉に任じられてから、数年のうちにその才幹は認められ、朝廷の少壮武官に列して、禁中紛乱、時局多事の中を、よく失脚もせず、いよいよその地歩を\ruby{占}{し}めて、新旧勢力の大官中に伍し、いつのまにか若年ながら\ruby{錚々}{そうそう}たる朝臣の一員となっているところ、早くも凡物でない\ruby{圭角}{けいかく}は現れていた。\\


　竹裏館の秘密会で、\ruby{王允}{おういん}もいったとおり、彼の家柄は、元来名門であって、高祖覇業を立てて以来の――漢の\ruby{丞相}{じょうしょう}\ruby{曹参}{そうさん}が\ruby{末孫}{ばっそん}だといわれている。\\


　生れは\ruby{沛国{\fontspec{jigmo}\char"8B59}郡}{はいこくしょうぐん}（安徽省・毫県）の産であるが、その父\ruby{曹嵩}{そうすう}は、宮内官たりし職を辞して、早くから野に下り、今では\ruby{陳留}{ちんりゅう}（河南省・開封の東南）に住んでいて、老齢だがなお健在であった。\\


　その父曹嵩も、\\


「この子は\ruby{鳳眼}{ほうがん}だ」\\


　といって、幼少の時から、大勢の子のうちでも、特に曹操を可愛がっていた。\\


　鳳眼というのは\ruby{鳳凰}{ほうおう}の眼のように細くてしかも光があるという意味であった。\\


　少年の頃になると、色は白く、髪は\ruby{漆黒}{しっこく}で、\ruby{丹唇明眸}{たんしんめいぼう}、中肉の美少年ではあり、しかも学舎の教師も、里人も、「\ruby{恐}{こわ}いようなお\ruby{児}{こ}だ」と、その鬼才に怖れた。\\


　こんなこともあった。\\


　少年の曹操は、学問など一を聞いて十を知るで、書物などにかじりついている日はちっとも見えない。\ruby{游猟}{ゆうりょう}が好きで弓を持って\ruby{獣}{けもの}を追ったり、早熟で不良を集めて村娘を\ruby{誘拐}{かどわか}したり、そんなことばかりやっていた。\\


「困った奴だ」\\


　叔父なる人が、将来を案じて、彼の父へひそかに忠告した。\\


「あまり可愛がり過ぎるからいけない。親の目には、子の良い才ばかり見えて、\ruby{奸才}{かんさい}は見えないからな」\\


　父の曹嵩も、ちらちら良くないことを耳にしていた折なので、早速曹操を呼びつけて、厳しく叱り、一晩中お談義を聞かせた。\\


　翌る日、叔父がやって来た。\\


　すると曹操は、ふいに門前に卒倒して、\ruby{癲癇}{てんかん}の発作に襲われたみたいな苦悶をした。\\


　仮病とは知らず、正直な叔父は驚きあわてて奥の父親へ告げた。\\


　父の曹嵩も、可愛い曹操のことなので、顔色を変えて飛びだして来た。――ところが曹操は門前に遊んでいて、いつもと何も変わったところは見えない。\\


「曹操、曹操」\\


「なんです、お父さん」\\


「なんともないのか。今、叔父御が駆けこんで来て、お前が\ruby{癲癇}{てんかん}を起してひッくり返っている、大変だぞ、直ぐ行ってみろ、といわれて仰天して見にきたのだが」\\


「ヘエ……。どうしてそんな嘘ッぱちを叔父さんは知らせたんでしょう。私はこの通り何でもありませんのに」\\


「変な人だな」\\


「まったく、叔父さんは変な人ですよ。嘘をいって、人が驚いたり困ったりするのを見るのが趣味らしいんです。村の人もいっていますね。――坊っちゃんは、あの叔父さんに何か憎まれてやしませんかッて。なんでも、わたしの事を放蕩息子だの、困り者だの、また癲癇持ちだのって、方々へ行って、しゃべりちらしているらしいんですよ」\\


　曹操は、\jruby{けろり}{・・・・・・}とした顔で、そういった。彼の父は、そのことがあってからというもの、何事があっても、叔父の言葉は信じなくなってしまった。\\


「甘いもンだな。親父は」\\


　曹操はいい気になって、いよいよ機謀縦横に\ruby{悪戯}{わるさ}をしたり、\ruby{放埓}{ほうらつ}な日を送って育った。\\

\\

\subsubsection{二}

\\


　二十歳まで、これという職業にもつかず、家産はあるし、名門の子だし、叔父の予言どおり困り息子で通ってきた曹操だった。\\


　しかし、人の憎みも多いかわり、一面任侠の\ruby{風}{ふう}もあるので、\\


「気の\ruby{利}{き}いた人だ」\\


　とか、また、\\


「曹操は話せるよ。いざという時は頼みになるからね」\\


　と、彼を取り巻く一種の人気といったようなものもあった。\\


　そういう友達の中でも、\ruby{橋玄}{きょうげん}とか、\ruby{何{\fontspec{jigmo}\char"9852}}{かぎょう}とかいう人々は、むしろ彼の縦横な策略の才を\ruby{異}{い}なりとして、\\


「今に、天下は乱れるだろう。一朝、乱麻となったが最後、これを収拾するのは、よほどな人物でなければできん。或いは後に、天下を安んずべき人間は、ああいったふうな\ruby{漢}{おとこ}かも知れんな」\\


　と、青年たちの集まった場所で、真面目にいったこともある。\\


　その橋玄が、ある折、曹操へ向っていった。\\


「君は、まだ無名だが、僕は君を有為の青年と見ているのだ。折があったら、\ruby{許子将}{きょししょう}という人と交わるがいい」\\


「子将とは、どんな人物かね」\\


　曹操が問うと、\\


「非常に人物の鑑識に\ruby{長}{た}けている。学者でもあるし」\\


「つまり人相\ruby{観}{み}だね」\\


「あんないい加減なものじゃない。もっと\ruby{炯眼}{けいがん}な人物批評家だよ」\\


「おもしろい。一度訪うてみよう」\\


　曹操は一日、その許子将を訪れた。座中、弟子や客らしいのが大勢いた。曹操は名乗って、彼の\ruby{忌憚}{きたん}ない「曹操評」を聞かしてもらおうと思ったが、子将は、冷たい眼で一\ruby{眄}{べん}したのみで、\ruby{卑}{いや}しんでろくに答えてくれない。\\


「ふふん……」\\


　曹操も、持前の皮肉がつい鼻先へ出て、こう\ruby{揶揄}{やゆ}した。\\


「――先生、池の魚は毎度\ruby{鑑}{み}ておいでらしいが、まだ大海の巨鯨は、この部屋で鑑たことがありませんね」\\


　すると、許子将は、学究らしい薄べったくて、黒ずんだ唇から、抜けた歯をあらわして、\\


「\ruby{豎子}{じゅし}、何をいう！　お前なんぞは、治世の能臣、乱世の\ruby{姦雄}{かんゆう}だ」\\


　と、初めて答えた。\\


　聞くと、曹操は、\\


「乱世の姦雄だと。――結構だ」\\


　彼は、満足して去った。\\


　間もなく。\\


　年二十で、初めて\ruby{北都尉}{ほくとい}の職についた。\\


　任は皇宮の警吏である。彼は就任早々、\ruby{掟}{おきて}を厳守し、犯す者は高官でも、ビシビシ罰した。時めく十常侍の\ruby{蹇碩}{けんせき}の身寄りの者でも、禁を破って、夜、帯刀で禁門の附近を歩いていたというので、曹操に棒で殴りつけられたことがあったりした程である。\\


「あの\ruby{弱冠}{じゃっかん}の警吏は、犯すと\ruby{仮借}{かしゃく}しないぞ」\\


　彼の名はかえって高まった。\\


　わずかな間に、騎都尉に昇進し、そして黄巾賊の乱が地方に起ると共に、征討軍に編入され、\ruby{潁川}{えいせん}その他の地方に転戦して、いつも紅の旗、紅の鞍、紅の鎧という人目立つ備え立てで征野を疾駆していたことは、かつて、張梁、張宝の賊軍を\ruby{潁川}{えいせん}の草原に火攻めにした折、――そこで行き会った\ruby{劉玄徳}{りゅうげんとく}とその\ruby{旗下}{きか}の関羽、張飛たちも、\\


（そも、何者？）\\


　と、目を見はったことのあるとおりである。\\


　そうした彼。\\


　そうした人となりの\ruby{驍騎}{ぎょうき}校尉曹操であった。\\


　\ruby{王允}{おういん}の家に伝わる七宝の名刀を譲りうけて、\ruby{董相国}{とうしょうこく}を刺すと誓って帰った曹操は、その夜、剣を抱いて床に横たわり、果たしてどんな夢を描いていたろうか。\\

\\

\subsubsection{三}

\\


　その翌日である。\\


　曹操は、いつものように\ruby{丞相府}{じょうしょうふ}へ出仕した。\\


「相国はどちらにおいでか」\\


　と、小役人に訊ねると、\\


「ただ今、小閣へ入られて、書院でご休息になっている」\\


　とのことなので、彼は直ちにそこへ行って、挨拶をした。董相国は、\ruby{牀}{しょう}の上に身を投げだして、茶をのんでいる様子。側には、\ruby{屹}{きっ}と、\ruby{呂布}{りょふ}が侍立していた。\\


「出仕が遅いじゃないか」\\


　曹操の顔を見るや否や、董卓はそういって咎めた。\\


　実際、陽はすでに\ruby{三竿}{さんかん}、丞相府の各庁でも、みなひと仕事すまして\ruby{午}{ひる}の休息をしている時分だった。\\


「恐れいります。なにぶん、私の持ち馬は痩せおとろえた老馬で道が遅いものですから」\\


「良い馬を持たぬのか」\\


「はい。薄給の身ですから、良馬は望んでもなかなか\ruby{購}{あがな}えません」\\


「呂布」と、董卓は振り向いて、\\


「わしの\ruby{厩}{うまや}から、どれか手ごろなのを一頭選んできて、曹操につかわせ」\\


「はっ」\\


　呂布は、閣の外へ出て行った。\\


　曹操は、彼が去ったので、\\


　――しめた！\\


　と、心は躍りはやったが、董卓とても、武勇はあり大力の持主である。\\


（仕損じては――）\\


　となお、大事をとって、彼の\ruby{隙}{すき}をうかがっていると、董卓はひどく肥満しているので、少し長くその体を\ruby{牀}{しょう}に正していると、すぐくたびれてしまうらしい。\\


　ごろりと、背を向けて、牀の上へ横になった。\\


（今だ！　天の与え）\\


　曹操は、心にさけびながら、七宝剣の柄に手をかけ、さっと抜くなり刃を背へまわして、牀の下へ近づきかけた。\\


　すると、名刀の\ruby{光鋩}{こうぼう}が、董卓の側なる壁の鏡に、\ruby{陽炎}{かげろう}の如くピカリと映った。\\


　むくりと、起き上がって、\\


「曹操、今の光は何だ？」\\


　と、鋭い眼をそそいだ。\\


　曹操は、刃を納めるいとまもなく、ぎょッとしたが、さあらぬ顔して、\\


「はっ、近頃それがしが、稀代の名刀を手に入れましたので、お気に召したら、献上したいと思って、\ruby{佩}{は}いて参りました。尊覧に入れる前に、そっと拭っておりましたので、その光鋩が室にみちたのでございましょう」\\


　と騒ぐ色もなく、剣を差出した。\\


「ふウむ。……どれ見せい」\\


　手に取って見ているところへ、呂布が戻ってきた。\\


　董卓は、気に入ったらしく、\\


「なるほど、名剣だ。どうだこの刀は」\\


　と、呂布へ見せた。\\


　曹操は、すかさず、\\


「\ruby{鞘}{さや}はこれです。七宝の\ruby{篏飾}{かんしょく}、なんと見事ではありませんか」\\


　と、呂布のほうへ、鞘をも渡した。\\


　呂布は無言のまま、\ruby{刃}{やいば}を鞘におさめて手に預かった。そして、\\


「馬を見給え」と促すと、曹操は、\\


「はっ、有難く拝領いたします」\\


　と、急いで庭上へ出て、呂布がひいて来た駿馬の\ruby{鬣}{たてがみ}をなでながら、\\


「あ。これは\ruby{逸物}{いちもつ}らしい。願わくば相国の\ruby{御前}{おんまえ}で、ひと当て試し乗りに乗ってみたいものですな」\\


　という言葉に、董卓もつい、図に乗せられて、\\


「よかろう。試してみい」\\


　とゆるすと、曹操はハッとばかり鞍へ飛び移り、にわかにひと鞭あてるや否や、丞相府の門外へ馳けだして、それなり帰ってこなかった。\\

\\

\subsubsection{四}

\\


「まだ戻らんか」\\


　董卓は、不審を起して、\\


「試し乗りだといいながら、いったい何処まで馳けて行ったのだ――曹操のやつは」\\


　と、何度も呟いた。\\


　呂布は初めて、口を開いた。\\


「丞相、彼はおそらく、もう此処に帰りますまい」\\


「どうして？」\\


「最前、あなたへ名刀を献じた時の挙動からして、どうも\ruby{腑}{ふ}に落ちない点があります」\\


「ム。あの時の彼奴の素振りは、わしも少し変だと思ったが」\\


「お馬を賜わり、これ幸いと、風を喰らって逃げ去ったのかも知れませんぞ」\\


「――とすれば、捨ておけん\ruby{曲者}{くせもの}だが。\ruby{李儒}{りじゅ}を呼べ。とにかく、李儒を！」\\


　と、急に甲高くいって、巨きな躯を牀からおろした。\\


　李儒は来て、つぶさに仔細を聞くと、\\


「それは、しまったことをした。\ruby{豹}{ひょう}を\ruby{檻}{おり}から出したも同じです。彼の妻子は都の外にありますから、てッきり相国のお命を狙っていたに違いありません」\\


「憎ッくい奴め。李儒、どうしたものだろう」\\


「一刻も早く、お召しといって、彼の住居へ人をやってごらんなさい。二心なければ参りましょうが、おそらくもうその家にもおりますまい」\\


　念のためと、直ちに、使い番の兵六、七騎をやってみたが、果たして李儒の言葉どおりであった。\\


　そしてなお、使い番から告げることには――\\


「つい今しがた、その曹操は、黄毛の駿馬にまたがって、飛ぶが如く東門を乗打ちして行ったので、番兵がまた馬でそれを追いかけ、ようやく城外へ出る関門でとらえて詰問したところ、曹操がいうには――我は丞相の急命を帯びてにわかに使いに立つなり。汝ら、我をはばめて大事の急用を遅滞さすからには、後に董相国よりいかなるお咎めがあらんも知れぬぞ――とのことなので、誰も疑う者なく、曹操はそのまま\ruby{鞭}{むち}を上げて関門を越え、行方のほども相知れぬ由にござります」\\


　とのことであった。\\


「さてこそ」と、董卓は、怒気のみなぎった顔に、朱をそそいで云った。\\


「小才のきく奴と、日頃、恩をほどこして、目をかけてやった予の寵愛につけ上がり、予にそむくとは八ツ裂きにしても飽きたらん匹夫だ。李儒っ――」\\


「はっ」\\


「彼の人相服装を画かせ、諸国へ写しを配布して、厳重に布令をまわせ」\\


「承知しました」\\


「もし、曹操を\ruby{生擒}{いけど}ってきた者あらば、\ruby{万戸侯}{ばんここう}に封じ、その首を丞相府に献じくる者には、千金の賞を与えるであろうと」\\


「すぐ手配しましょう」\\


　李儒が退がりかけると、\\


「待て。それから」と早口に、董卓はなお、言葉をつけ加えた。\\


「この細工は、思うに、白面郎の曹操一人だけの仕事ではなかろう。きっとほかにも、同謀の与類があるに相違ない」\\


「もちろんでしょう」\\


「なおもって、重大事だ。曹操への手配や追手にばかり気を取られずに一方、都下の与類を\ruby{虱}{しらみ}つぶしに詮議して、引っ捕えたら拷問にかけろ」\\


「はっ、その辺も、抜かりなく急速に手を廻しましょう」\\


　李儒は大股に去って、\ruby{捕囚庁}{ほしゅうちょう}の\ruby{吏人}{やくにん}を呼びあつめ、物々しい活動の指令を発していった。\\

\\

\\

\\

\end{document}
